\pdfoutput=1
\documentclass[11pt]{amsart}

\usepackage[margin=1.1in]{geometry}
\usepackage{amsmath,amssymb,amsthm,mathtools,mathrsfs}
\usepackage{graphicx}
\usepackage{array}
\usepackage{longtable}
\usepackage{enumitem}
\usepackage{microtype}
\usepackage{aliascnt}
\usepackage{tikz}
\usetikzlibrary{arrows.meta,positioning,shapes.geometric}
\usepackage{hyperref}
\usepackage[nameinlink,capitalize]{cleveref}

\hypersetup{
  colorlinks=true,
  linkcolor=blue,
  citecolor=blue,
  urlcolor=blue,
  pdftitle={Local Exclusion of Residual Type I Navier-Stokes Seregin Limits},
  pdfauthor={Guillem Duran Ballester},
  pdfsubject={Navier-Stokes Type I singularities and Seregin limits},
  pdfkeywords={Navier-Stokes equations, Type I singularities, ancient solutions, Seregin limits, Liouville theorems, local concentration}
}

\numberwithin{equation}{section}

\theoremstyle{plain}
\newtheorem{theorem}{Theorem}[section]

\newaliascnt{proposition}{theorem}
\newtheorem{proposition}[proposition]{Proposition}
\aliascntresetthe{proposition}

\newaliascnt{lemma}{theorem}
\newtheorem{lemma}[lemma]{Lemma}
\aliascntresetthe{lemma}

\newaliascnt{corollary}{theorem}
\newtheorem{corollary}[corollary]{Corollary}
\aliascntresetthe{corollary}

\theoremstyle{definition}
\newaliascnt{definition}{theorem}
\newtheorem{definition}[definition]{Definition}
\aliascntresetthe{definition}

\theoremstyle{remark}
\newaliascnt{remark}{theorem}
\newtheorem{remark}[remark]{Remark}
\aliascntresetthe{remark}

\crefname{theorem}{Theorem}{Theorems}
\Crefname{theorem}{Theorem}{Theorems}
\crefname{lemma}{Lemma}{Lemmas}
\Crefname{lemma}{Lemma}{Lemmas}
\crefname{proposition}{Proposition}{Propositions}
\Crefname{proposition}{Proposition}{Propositions}
\crefname{corollary}{Corollary}{Corollaries}
\Crefname{corollary}{Corollary}{Corollaries}
\crefname{definition}{Definition}{Definitions}
\Crefname{definition}{Definition}{Definitions}
\crefname{remark}{Remark}{Remarks}
\Crefname{remark}{Remark}{Remarks}
\crefalias{theorem}{theorem}
\crefalias{lemma}{lemma}
\crefalias{proposition}{proposition}
\crefalias{corollary}{corollary}
\crefalias{definition}{definition}
\crefalias{remark}{remark}

\newcommand{\R}{\mathbb R}
\newcommand{\N}{\mathbb N}
\newcommand{\eps}{\varepsilon}
\newcommand{\loc}{\mathrm{loc}}
\newcommand{\CKN}{\mathrm{CKN}}
\newcommand{\NS}{\mathrm{NS}}
\newcommand{\RNSE}{\mathrm{RNSE}}
\newcommand{\divergence}{\operatorname{div}}
\newcommand{\dist}{\operatorname{dist}}
\newcommand{\supp}{\operatorname{supp}}
\newcommand{\esssup}{\operatorname*{ess\,sup}}
\newcommand{\calS}{\mathcal S}
\newcommand{\calR}{\mathcal R}
\newcommand{\calX}{\mathcal X}
\newcommand{\calU}{\mathcal U}
\newcommand{\calE}{\mathcal E}
\newcommand{\Vel}{\mathcal V}
\newcommand{\calA}{\mathcal A}
\newcommand{\calM}{\mathcal M}
\newcommand{\calF}{\mathcal F}
\newcommand{\calC}{\mathcal C}
\newcommand{\calG}{\mathcal G}
\newcommand{\calK}{\mathcal K}
\newcommand{\calH}{\mathcal H}
\newcolumntype{L}[1]{>{\raggedright\arraybackslash}p{#1}}
\newcommand{\Cax}{\mathcal C_{\mathrm{ax}}}
\newcommand{\Crot}{\mathcal C_{\mathrm{rot}}}
\newcommand{\Cstat}{\mathcal C_{\mathrm{stat}}^{\mathrm{ret}}}
\newcommand{\Cactive}{\mathcal C_{\mathrm{conc}}}
\newcommand{\Caff}{\mathcal C_{\mathrm{aff}}}
\newcommand{\Cgen}{\mathcal C_{\mathrm{gen}}}
\newcommand{\Ccritdiff}{\mathcal C_{\mathrm{critdiff}}}
\newcommand{\Ccohcrit}{\mathcal C_{\mathrm{cohcrit}}}
\newcommand{\Chomcrit}{\mathcal C_{\mathrm{homcrit}}}
\newcommand{\CDSStail}{\mathcal C_{\mathrm{DSS\text{-}tail}}}
\newcommand{\Capercrit}{\mathcal C_{\mathrm{apercrit}}}
\newcommand{\Tcrit}{\mathfrak T_M}
\newcommand{\Tcoh}{\mathfrak T_{\mathrm{coh}}}
\newcommand{\Tyoung}{\mathfrak T_{\mathrm{Young}}}
\newcommand{\Tlogdiff}{\mathfrak T_{\log\mathrm{diff}}}
\newcommand{\Clogper}{\mathcal C_{\mathrm{logper}}}
\newcommand{\Capcrit}{\mathcal C_{\mathrm{ap}}}
\newcommand{\Cyoung}{\mathcal C_{\mathrm{Young}}}
\newcommand{\Clogdiff}{\mathcal C_{\log\mathrm{diff}}}
\newcommand{\MesoscopicReductionTheorem}{\Cref{thm:first-bad-mesoscopic-reduction}}
\newcommand{\oscsp}{\operatorname{osc}_{3,\mathrm{sp}}}
\newcommand{\actthree}{\operatorname{osc}_{3,\mathrm{ct}}}
\newcommand{\bP}{\mathbb P}
% SetupLabel: render each imported setup paper result by final companion-paper
% theorem/lemma/proposition number.  The stable source keys remain internal.
\makeatletter
\newcommand{\SetupNum}[2]{\expandafter\gdef\csname SetupNum@#1\endcsname{#2}}
\SetupNum{p0:thm:singular-positive}{Theorem~5.7}
\SetupNum{p0:prop:terminal-concentration}{Proposition~5.9}
\SetupNum{p0:thm:local-weak-serrin-typeI}{Theorem~8.3}
\SetupNum{p0:cor:paper0-local-typeI-admissible}{Corollary~8.4}
\SetupNum{p0:lem:local-seregin-extraction}{Lemma~8.6}
\SetupNum{p1:thm:extraction}{Theorem~9.1}
\SetupNum{p1:prop:ancient-limit}{Proposition~11.12}
\SetupNum{p1:prop:nonzero}{Proposition~12.2}
\SetupNum{p1:lem:pressure}{Lemma~11.2}
\SetupNum{p1:lem:gauge-compatibility}{Lemma~11.3}
\SetupNum{p1:lem:pressure-atlas}{Lemma~11.4}
\SetupNum{p1:lem:mild-representation}{Lemma~15.3}
\SetupNum{p2:lem:mild-limit}{Lemma~20.4}
\SetupNum{p2:lem:physical-mild}{Lemma~25.4}
\SetupNum{p1:thm:classical-classes}{Theorem~15.8}
\SetupNum{p1:thm:known-structure-decay-liouville}{Theorem~15.10}
\SetupNum{p2:thm:no-tight-normalized-collection}{Theorem~26.1}
\SetupNum{p2:thm:tight-liouville}{Theorem~25.8}
\SetupNum{p2:thm:AB-endpoint}{Theorem~25.5}
\SetupNum{p1:hyp:no-remainder}{Hypothesis~16.1}
\SetupNum{p1:thm:final-assembly}{Theorem~33.1}
\newcommand{\SetupLabel}[1]{%
   \@ifundefined{SetupNum@#1}%
      {\textup{the corresponding setup result}}%
      {\csname SetupNum@#1\endcsname\ of the setup paper}%
}
\makeatother

\title[Residual Type I Seregin limit exclusion]
{Local Exclusion of Residual Type I Navier--Stokes Seregin Limits}

\author{Guillem Duran Ballester}
\date{}

\subjclass[2020]{35Q30, 35B44, 35B40, 76D05}
\keywords{Navier--Stokes equations, Type I singularities, ancient solutions, Seregin limits, Liouville theorems, local concentration}

\begin{document}
\raggedbottom

\begin{abstract}
This paper proves the residual-class hypothesis left open by the companion
setup paper for local pointwise Type-I Navier--Stokes blow-up.  Starting from
the setup Seregin extraction, every remaining bounded centered ancient profile
with retained local velocity \(L^3\)-concentration is assigned to an ordered
residual decomposition.  Axisymmetric bounded-circulation profiles are excluded
by a centered angular-circulation absorption argument.  Rotational and
stationary-hull profiles are reduced to a terminal concentration theorem.  The
generic terminal branch is closed by a local concentration-set decomposition,
diffuse-defect and critical-tail compactifications, a minimal mesoscopic-scale
rigidity theorem, exclusion of separated retained profiles, and an endpoint
sequence-\(L^3\) Liouville argument for bounded mild ancient solutions.  No
global critical-norm estimate or global Liouville theorem for bounded centered
profiles is used.
\end{abstract}

\maketitle
\tableofcontents

\section{Introduction and main theorem}
\label{sec:introduction}

The incompressible Navier--Stokes equations on
\(\R^3\times(0,T)\) are written as
\begin{equation}\label{eq:NS-physical}
   \partial_tu-\Delta u+(u\cdot\nabla)u+\nabla p=0,
   \qquad \nabla\cdot u=0 .
\end{equation}
The residual analysis in this paper begins after the standard Type I blow-up
construction of Seregin has been performed \cite{Seregin2012}.  Thus the basic solution is a smooth bounded
ancient solution of the centered equation
\begin{equation}\label{eq:centered-main}
   \partial_\tau V-\Delta V+\frac12 y\cdot\nabla V+\frac12 V
   +(V\cdot\nabla)V+\nabla P=0,
   \qquad \nabla\cdot V=0 ,
\end{equation}
on \(\R^3_y\times\R_\tau\).  The equation is obtained from the physical
ancient representative \((U,Q)\) by
\begin{equation}\label{eq:centered-transform}
   y=\frac{x}{\sqrt{-t}},\qquad
   \tau=-\log(-t),\qquad
   V(y,\tau)=\sqrt{-t}\,U(x,t),
   \qquad
   P(y,\tau)=(-t)Q(x,t),
\end{equation}
where \(t<0\) and pressures are always understood modulo addition of a
function of time.  Conversely,
\begin{equation}\label{eq:physical-pullback-main}
   U(x,t)=(-t)^{-1/2}
   V\!\left(\frac{x}{\sqrt{-t}},-\log(-t)\right),
   \qquad t<0 .
\end{equation}

The purpose of this paper is to present the refined residual decomposition as a
single standard PDE argument, using only the reductions already established in
the companion setup paper.  In this two-paper presentation, the current paper
depends only on the imported setup results stated below and on cited
Navier--Stokes regularity and Liouville literature
\cite{KNSS2009,Seregin2012}.  The
rotational-flux estimates, terminal concentration-set decompositions,
exclusions of local compactness failure and diffuse exterior concentration, and
finite separated-profile exclusions are proved in this paper from local
estimates rather than assumed as global estimates.

This paper does not reconstruct the Seregin extraction or the endpoint mildness
theory.  Those are imported from the setup paper.  The sole new task is to
prove that the residual alternative left open there is empty, namely
\[
   \calR^\#(\calS;I,J)=\varnothing .
\]

\subsection{Proof architecture}
\label{subsec:proof-architecture}
The retained normalized residual class \(\calR^\#(\calS;I,J)\) of
\Cref{def:coarse-residual} is the complement, inside the normalized Seregin
collection \(\calS\), of the alternatives already removed in the setup paper.
After the previously treated small, stationary \(L^3\), uniformly
\(L^3\)-tight, and closed structured/decay alternatives have been removed
(\Cref{hyp:previous-nonresidual}), \(\calR^\#(\calS;I,J)\) is divided in
\Cref{def:refined-decomposition} into four analytic classes -- axisymmetric
bounded-circulation profiles \(\Cax\), rotational relative equilibria
\(\Crot\), degenerate stationary-hull profiles with retained local velocity
concentration \(\Cstat\), and the generic terminal residual \(\Cgen\) -- and
the terminal generic case is itself refined into the affine/parasitic quotient
stratum \(\Caff\), the log-diffuse and Young/defect critical-tail alternatives
\(\Clogdiff,\Cyoung\), the coherent homogeneous, log-periodic, and aperiodic
critical-tail alternatives \(\Chomcrit,\Clogper,\Capcrit\), and the residual
remainder \(\Cgen\).

The closure of these classes is organized as follows.
\begin{enumerate}[label=\textup{(\arabic*)},leftmargin=2.2em]
   \item \(\Cax\) is excluded directly by a centered angular-circulation
   equation and an axis-absorption Liouville theorem
   (\Cref{thm:R1-exclusion}).
   \item \(\Crot\) is reduced by a local Coriolis-flux identity to either
   lower classes already excluded by the setup paper or the terminal stratification assembly
   (\Cref{prop:R2-reduction-to-terminal}); its closure is obtained only
   after the terminal assembly has excluded retained velocity concentration
   profiles.
   \item \(\Cstat\) is reduced, by scale reselection for time-translation
   hulls, to a stationary profile in a degenerate stationary family which
   occurs as a Seregin limit and still carries local velocity concentration,
   and is then excluded by the terminal stratification assembly theorem
   (\Cref{prop:R3S-R3-terminal-reduction}).
   \item \(\Cgen\) and the critical-tail subclasses are closed by covariant
   tail normalization, concentration-set extraction, diffuse-defect
   compactness, critical-tail compactification, recurrence for chains of
   concentration profiles, exclusion of finite separated concentration
   profiles, indecomposable-profile sequence-\(L^3\) completeness,
   parasitic-free mildness inheritance, and the Albritton--Barker endpoint
   Liouville theorem.
\end{enumerate}
The paper distinguishes statements inherited from the setup paper, local
compactness estimates and reduction results proved here, and the minimal
mesoscopic-scale reduction theorem
(\Cref{thm:first-bad-mesoscopic-reduction}) which excludes the
critical-tail boundary defects.  The affine stratum and the coherent
critical-tail alternatives are closed here; the log-diffuse and Young
alternatives are both reduced to the same minimal mesoscopic-scale theorem,
which is then proved by pressure compactness and a bounded ancient variance
Liouville argument.  No global Coriolis--Leray Liouville theorem for bounded
profiles or global critical-norm estimate is used in the residual closure.
A short reader's guide is given in \Cref{subsec:reading-guide} below.

\begin{definition}[Type I convention used in the chain]
\label{def:pointwise-typeI-convention}
In this paper and in the final two-paper chain, a terminal singular point
\(z_*=(x_*,T)\) is called \emph{local pointwise Type~I} if there are
\(\rho>0\) and \(M<\infty\) such that
\[
   \esssup_{T-\rho^2<t<T}
   \sqrt{T-t}\,\|u(t)\|_{L^\infty(B_\rho(x_*))}\le M .
\]
This is the local pointwise Type~I hypothesis used in the setup paper and in
\Cref{sec:local-typeI-implication}.  For suitable weak solutions this envelope lies
in the endpoint weak Serrin class \(L^{2,\infty}_tL^\infty_x\).  The
Albritton--Barker weak-Serrin local Type~I estimate, recorded below as
\Cref{thm:AB-weak-serrin-typeI-estimate}, upgrades it on smaller cylinders to
the scale-invariant local energy bounds needed for terminal velocity
no-escape.  Thus the terminal no-escape condition used by the setup paper is
automatic for the local pointwise Type~I singularities considered here.
The paper does not use the phrase ``Type~I'' to mean an unrelated global
\(L^\infty\), local CKN, weak-\(L^3\), or Serrin scale condition unless that
condition is explicitly stated.
\end{definition}

\begin{remark}[Translation table for standard Type I hypotheses]
\label{rem:typeI-translation-table}
The global pointwise condition
\[
   \sup_{t<T}\sqrt{T-t}\,
      \|u(t)\|_{L^\infty(\mathbb R^3)}<\infty
\]
implies \Cref{def:pointwise-typeI-convention} at every terminal point.  Any
stronger local criterion that yields the displayed \(L^\infty\) envelope on a
neighborhood of \(x_*\) enters the same argument, provided the solution is
suitable in the corresponding cylinder.  The endpoint weak Serrin estimate then
gives the local scale-invariant energy bounds and terminal no-escape.
Purely CKN-type, weak-\(L^3\), or other critical criteria are not silently
identified with this hypothesis; they require their own implication into the
local pointwise envelope, or directly into the local Type~I estimates, before the
present exclusion applies.
\end{remark}

\subsection{Imported setup results and remaining hypothesis}
\label{subsec:imported-setup-results}

\begin{theorem}[Imported setup results]
\label{thm:imported-setup-results}
The companion setup paper proves the following facts, in the exact theorem
numbers listed below.
\begin{enumerate}[label=\textup{(\arabic*)},leftmargin=2.2em]
   \item Positive terminal local velocity concentration at every local pointwise Type~I
   singular point.
   \begin{quote}\small
   \SetupLabel{p0:thm:singular-positive}\\
   \SetupLabel{p0:prop:terminal-concentration}
   \end{quote}
   \item Local pointwise Type~I control gives an admissible Seregin extraction
   sequence with pressure gauges, no-escape, and retained unit-scale velocity
   concentration.
   \begin{quote}\small
   \SetupLabel{p0:thm:local-weak-serrin-typeI}\\
   \SetupLabel{p0:cor:paper0-local-typeI-admissible}\\
   \SetupLabel{p0:lem:local-seregin-extraction}
   \end{quote}
   \item Every normalized Seregin limit produced by this extraction is a
   bounded ancient suitable weak solution, admits the pressure atlas of the
   setup paper, and has the mild ancient representative required by the
   endpoint Liouville theorem.
   \begin{quote}\small
   \SetupLabel{p1:thm:extraction}\\
   \SetupLabel{p1:prop:ancient-limit}\\
   \SetupLabel{p1:prop:nonzero}\\
   \SetupLabel{p1:lem:pressure}\\
   \SetupLabel{p1:lem:gauge-compatibility}\\
   \SetupLabel{p1:lem:pressure-atlas}\\
   \SetupLabel{p1:lem:mild-representation}\\
   \SetupLabel{p2:lem:mild-limit}\\
   \SetupLabel{p2:lem:physical-mild}
   \end{quote}
   \item The previously classified small, stationary \(L^3\), uniformly
   tight, and known structured/decay alternatives are empty.
   \begin{quote}\small
   \SetupLabel{p1:thm:classical-classes}\\
   \SetupLabel{p1:thm:known-structure-decay-liouville}\\
   \SetupLabel{p2:thm:no-tight-normalized-collection}\\
   \SetupLabel{p2:thm:tight-liouville}
   \end{quote}
   \item The endpoint and mildness results include the Albritton--Barker input
   and the setup paper proves local pointwise Type~I exclusion conditionally on
   the residual-class hypothesis.
   \begin{quote}\small
   \SetupLabel{p2:thm:AB-endpoint}\\
   \SetupLabel{p1:hyp:no-remainder}\\
   \SetupLabel{p1:thm:final-assembly}
   \end{quote}
\end{enumerate}
The purpose of the present paper is to prove the residual closure
\(\calR^\#(\calS;I,J)=\varnothing\).  This verifies
\SetupLabel{p1:hyp:no-remainder}.
\end{theorem}

\begin{remark}[Use of imported labels]
\label{rem:imported-label-convention}
Below, the phrase ``by the imported setup results'' means that the terminal
profile is bounded, locally suitable, pressure-normalized by the setup pressure
atlas, parasitic-free modulo the affine class, and mild after finite terminal
shift, with the exact imported labels listed in
\Cref{thm:imported-setup-results}.  Paper~IV proves only the residual closure
\(\calR^\#(\calS;I,J)=\varnothing\).
\end{remark}

\begin{enumerate}[label=\textup{(\arabic*)},leftmargin=2.2em]
   \item Local concentration and terminal no-escape entry: positive local
   concentration depends on CKN theory and the setup paper; terminal no-escape follows
   from the endpoint weak Serrin local Type~I bounds in
   \Cref{thm:AB-weak-serrin-typeI-estimate,thm:local-typeI-terminal-energy-estimates}.
   \item Seregin extraction and parasitic-free normalization: follows from
   the setup paper and the mildness and compactness arguments in this paper.
   \item Small and stationary \(L^3\) strata: handled in the setup paper.
   \item Uniformly tight stratum: handled in the setup paper.
   \item Closed structured/decay strata: handled in the setup paper.
   \item Axisymmetric and rotational residual strata: proved in this paper.
   \item Terminal concentration-set decomposition, recurrence for concentration chains, and
   diffuse-defect compactness: proved in this paper.
   \item Critical-tail compactification and realized critical-tail
   decomposition: proved in this paper.
   \item Residual indecomposable-profile closure: proved in this paper, with
   the final sequence-\(L^3\) endpoint supplied by Albritton--Barker.
\end{enumerate}

\begin{figure}[tbp]
\centering
\resizebox{0.78\textwidth}{!}{%
\begin{tikzpicture}[
   node distance=1.10cm,
   box/.style={draw, rounded corners=2pt, align=center, text width=4.8cm,
      inner sep=4pt},
   widebox/.style={draw, rounded corners=2pt, align=center, text width=7.8cm,
      inner sep=4pt},
   arrow/.style={-{Latex[length=2mm]}, thick},
   lab/.style={align=left, font=\scriptsize, text width=4.7cm}
]
\node[box] (typeI) {local pointwise Type~I singularity};
\node[box, below=of typeI] (admissible) {admissible Type~I compactness setting};
\node[widebox, below=of admissible] (seregin)
   {nonzero bounded centered Seregin profile\\with retained velocity concentration};
\node[box, below=of seregin] (split)
   {inherited lower classes\\or\\\(\calR(\calS;I,J)\)};
\node[widebox, below=of split] (decomp)
   {\(\Cax,\ \Crot,\ \Cstat,\ \Caff,\ \Clogdiff,\ \Cyoung\)\\
    \(\Chomcrit,\ \Clogper,\ \Capcrit,\ \Cgen\)};
\node[box, below=of decomp] (contradiction) {contradiction};

\draw[arrow] (typeI) -- node[lab, right]
   {endpoint weak Serrin bounds\\and terminal no-escape} (admissible);
\draw[arrow] (admissible) -- node[lab, right] {Seregin extraction} (seregin);
\draw[arrow] (seregin) -- (split);
\draw[arrow] (split) -- node[lab, right] {complete residual decomposition} (decomp);
\draw[arrow] (decomp) -- node[lab, right] {case closures in this paper}
   (contradiction);
\end{tikzpicture}}
\caption{Decision diagram for the Type~I profile decomposition.  The first
arrow uses the local pointwise Type~I to terminal no-escape implication supplied by
the endpoint weak Serrin local Type~I estimate.  The inherited lower classes in
the split box are already closed in the setup paper.  The residual alternatives in the boxed
list are excluded by the case-exclusion theorems cited in the proof of
\Cref{thm:refined-residual-closure}.}
\label{fig:typeI-state-stratification-tree}
\end{figure}

\begin{theorem}[Admissible Type I Seregin extraction theorem]
\label{hyp:base-seregin-hypotheses}
An admissible finite-energy Type I singular point produces a normalized Seregin
ancient limit \((V,P)\) satisfying \eqref{eq:centered-main}, with
\[
   \|V\|_{L^\infty(\R^3\times\R)}\le M<\infty,
\]
local suitability and local pressure-gauge compactness, and a compact
parabolic cylinder \(K\Subset\R^3\times\R\) for which
\begin{equation}\label{eq:positive-velocity-main}
   \iint_K |V|^3\,dy\,d\tau
   \ge \eta_0>0 .
\end{equation}
The collection of all normalized Seregin limits generated by the original
solution is
denoted by \(\calS\).
\end{theorem}

\begin{proof}
This is the base admissible Seregin extraction theorem included in
\Cref{thm:imported-setup-results}: velocity nonvanishing, local suitability,
and local pressure compactness are imported setup conclusions.  The present
paper uses this inherited reduction and does not reprove the Type I blow-up
compactness theorem.
\end{proof}

\begin{definition}[Earlier alternatives inherited from the setup paper]
\label{hyp:previous-nonresidual}
Within a fixed normalized Seregin collection \(\calS\), the small-amplitude,
stationary \(L^3\), uniformly \(L^3\)-tight, and closed structured/decay
classes are the non-residual alternatives already established in the setup
paper.  The small-amplitude and stationary \(L^3\) alternatives are closed
there; the uniformly tight class is excluded there; and the closed structured/decay
class is precisely the part proved empty there.  Any structured profile not
belonging to that closed class is residual by definition and is handled by
the present paper.  The present paper does not reprove those earlier
alternatives and does not add any hypothesis to their statements.
\end{definition}

\begin{definition}[Coarse residual class]
\label{def:coarse-residual}
Let \(I,J\) be the index sets used in the base Type I residual
criterion.  The pre-quotient coarse residual class
\[
   \calR_{\rm pre}(\calS;I,J)
\]
is the complement in \(\calS\) of the already removed alternatives in
\Cref{hyp:previous-nonresidual}.  The retained normalized residual class
\[
   \calR^\#(\calS;I,J)
\]
is obtained from \(\calR_{\rm pre}(\calS;I,J)\) by passing to the canonical
affine/parasitic representative and by removing any representative assigned to
the affine/parasitic lower stratum.  Unless the word unreduced is used explicitly,
\(\calR(\calS;I,J)\) denotes this retained normalized residual class.  Thus an
``empty affine/parasitic stratum'' below always means that no retained
normalized residual representative remains in that quotient stratum.
\end{definition}

\begin{definition}[Refined residual decomposition]
\label{def:refined-decomposition}
The refined residual decomposition is an ordered definition, not a separate
compactness theorem.  Starting with the pre-quotient coarse residual class
\(\calR_{\rm pre}(\calS;I,J)\), one removes the following alternatives in
order; the affine/parasitic alternative is then recorded as a quotient stratum
with no retained normalized representative in \(\calR^\#(\calS;I,J)\).  After
the first nine alternatives have been removed, whatever remains is by definition
the generic terminal residual \(\Cgen\):
\[
\begin{gathered}
   \Cax,\qquad \Crot,\qquad \Cstat,\qquad \Caff,\\
   \Clogdiff,\qquad \Cyoung,\qquad
   \Chomcrit,\qquad \Clogper,\qquad \Capcrit,\qquad \Cgen .
\end{gathered}
\]
Here \(\Cax\) is the axisymmetric bounded-circulation class, \(\Crot\) is the
rotational relative-equilibrium class, \(\Cstat\) is the degenerate
stationary-hull class with retained local velocity concentration, \(\Caff\) is the
affine/parasitic lower stratum selected by affine-normalized oscillatory entry,
\(\Clogdiff\) and \(\Cyoung\) are respectively the log-diffuse and Young/defect
critical-tail alternatives, \(\Chomcrit,\Clogper,\Capcrit\) are the coherent
homogeneous, log-periodic, and aperiodic critical-tail alternatives, and \(\Cgen\)
is the generic terminal residual class.  The inclusion below is part of this
ordered definition, written for the retained normalized residual class; the
appearance of \(\Caff\) records the quotient alternative before retained
normalization, and \(\Caff\cap\calR^\#(\calS;I,J)=\varnothing\) is proved in
\Cref{thm:oscillatory-entry-normalization}.  The substantive content is the
closure of each non-quotient alternative, especially \(\Cgen\).  Equivalently,
\begin{equation}\label{eq:complete-residual-decomposition}
   \calR(\calS;I,J)
   \subset
   \Cax\cup\Crot\cup\Cstat\cup\Caff
   \cup\Clogdiff\cup\Cyoung\cup\Chomcrit\cup\Clogper\cup\Capcrit
   \cup\Cgen .
\end{equation}
\end{definition}

\begin{definition}[Profiles realized from the original blow-up sequence]
\label{def:sequence-realized-residual-object}
Fix a normalized Seregin collection \(\calS\) and a retained residual candidate
\((V,P)\in\calR(\calS;I,J)\).  An auxiliary profile or limiting measure used in the residual proof is
\emph{realized from the original blow-up sequence} if it is obtained from \((V,P)\), or from the terminal
sequence which generated \((V,P)\), by a finite or diagonal sequence of the
following operations:
\begin{enumerate}[label=\textup{(\roman*)}]
   \item exact centered symmetries \(T_a\), \(\Theta_s\), and the covariant
   observer maps \(\calR_z\);
   \item unreduced parabolic recenterings before passing to centered variables, with
   the local pressure gauges supplied by Seregin compactness;
   \item retained tail limits or concentrating successors with
   \(\Vel_1\ge\eta_*\) after the fixed admissible pressure gauges;
   \item normalized diffuse or critical-tail defect limits obtained by
   restricting the associated defect measure to a positive-measure window and
   renormalizing it;
   \item minimal mesoscopic blow-downs selected from such retained or
   normalized defect windows.
\end{enumerate}
When the shorter phrase ``realized from the original sequence'' is used below,
it refers to this definition and to the original singularity-generated
candidate \((V,P)\), not merely to the immediately preceding normalized profile.
\end{definition}

\begin{definition}[Retained descendants and normalized support profiles]
\label{def:retained-vs-support-profile-descendant}
A \emph{retained descendant} is a descendant with a genuine inherited positive
unit velocity concentration bound
\[
   \liminf_n\Vel_1(\cdot;0)\ge\eta_*>0 .
\]
A \emph{normalized support profile}, also called a \emph{support profile
descendant}, is a normalized profile in the support of a defect measure or
normalized window measure associated with the original blow-up sequence.  It may arise from
unnormalized mass tending to zero and therefore need not carry the original
velocity threshold.  Log-window and Young/defect normalizations are called
retained only when the displayed unit velocity concentration bound is explicitly proved;
otherwise they are normalized support profiles.
\end{definition}

\begingroup
\footnotesize
\setlength{\tabcolsep}{2pt}
\begin{longtable}{@{}L{0.20\textwidth}L{0.17\textwidth}L{0.14\textwidth}L{0.15\textwidth}L{0.17\textwidth}L{0.11\textwidth}@{}}
\caption{Stability of the equation, gauges, and concentration bounds under limiting operations.}
\label{tab:ancestor-heredity-table}\\
\hline
\textbf{Operation} & \textbf{Equation} & \textbf{Boundedness} &
\textbf{Pressure gauge} & \textbf{Mildness} & \textbf{Threshold}\\
\hline
\endfirsthead
\hline
\textbf{Operation} & \textbf{Equation} & \textbf{Boundedness} &
\textbf{Pressure gauge} & \textbf{Mildness} & \textbf{Threshold}\\
\hline
\endhead
\hline
\endfoot
Centered time shift & yes & yes after finite shift & shifted gauge &
setup mildness & yes\\
Spatial recentering & covariant & yes & translated gauge & yes & yes\\
Parabolic rescaling & yes & normalized scale & scaled gauge & yes & yes\\
Tail recentering & quotient covariant & locally yes & tail gauge & setup
closure & support profile or retained\\
Log-window normalization & normalized defect equation & normalized & quotient
gauge & not automatic & support profile\\
Affine pressure quotient & modulo \(\Caff\) & yes & affine removed & yes
outside \(\Caff\) & specified\\
Critical-tail blow-down & limiting equation & locally yes & scaled/quotient
gauge & after bounded-origin realization & support profile\\
\hline
\end{longtable}
\endgroup

\begin{theorem}[Stability under realized limiting operations]
\label{thm:ancestor-realization-inheritance}
Every descendant, tail limit, concentrating successor, log-window normalization,
critical-tail profile, and minimal-scale blow-down used in the proof of
\Cref{thm:refined-residual-closure} is realized from the original blow-up sequence in the
sense of \Cref{def:sequence-realized-residual-object}.  More precisely, each
such limit satisfies one of the following alternatives.
\begin{enumerate}[label=\textup{(\alph*)}]
   \item It is a retained Seregin descendant of the original candidate, with
   the fixed descendant lower bound \(\Vel_1\ge\eta_*\).
   \item It is a normalized defect profile whose unnormalized defect
   measure has positive total mass on the selected window.  If every normalized
   window descendant is excluded, then the original defect is excluded.
   \item It lies in a previously closed lower stratum, and the retained
   descendant heredity theorem applies.
   \item It generates a retained descendant or a normalized defect descendant
   satisfying \textup{(a)}, \textup{(b)}, or \textup{(c)}.
\end{enumerate}
Consequently a contradiction obtained for a realized limiting profile is a
contradiction for the original retained residual candidate, not merely for an
auxiliary compactification limit.
The theorem is used only through the concrete support-transfer statements
proved for the individual compactifications below.
\end{theorem}

\begin{proof}
It suffices to verify the assertion for the operations listed in
\Cref{def:sequence-realized-residual-object}.  Exact centered
symmetries preserve the centered equation, local suitability, pressure gauges,
and retained local velocity concentration; this is recorded in
\Cref{lem:covariant-translation-invariance,thm:covariant-observer-calculus}.
Unreduced terminal recenterings before centering are the recenterings in the Seregin
extraction theorem of the setup paper and are compact after the gauges in
\Cref{prop:local-typeI-terminal-compactness}.  Retained tail limits and concentrating
successors are realized from the original sequence by diagonal lifting:
\Cref{lem:diagonal-lifting-descendants,thm:descendant-heredity} show that the
limit remains in the normalized compactness class and keeps the fixed retained velocity
threshold through its velocity part, up to the once-for-all threshold loss in
\Cref{def:threshold-hierarchy}.

Normalized defect profiles are obtained by restricting a positive defect
measure to a window and dividing by its mass.  The mass of the chosen window may
depend on the index and may tend to zero; this does not detach the normalized
profile from the original sequence.  It records a point in the support of the
defect compactification.  The renormalized-window heredity lemma
\Cref{lem:renormalized-log-window-heredity} below states explicitly that if all
such support profiles are assigned to closed lower alternatives, then the
defect measure itself has no surviving mass outside those alternatives.

Minimal-scale blow-downs are selected from retained or normalized defect windows;
their centers and scales therefore come with the same defect measure or
retained velocity lower bound.  The pressure/gauge normalization for those blow-downs
is treated in
\Cref{lem:first-bad-quotient-formulation,prop:normalized-pressure-representation,lem:first-bad-pressure-compactness}.
Thus every contradiction later in the proof is applied either to a retained
descendant of the original profile or to a normalized support profile whose
closure removes the corresponding defect.  This proves the stated
principle.
\end{proof}

\begin{lemma}[Support transfer in compact metrizable defect spaces]
\label{lem:compact-defect-support-transfer}
Let \(X\) be a compact metrizable defect space and let
\[
   \lambda_n=\frac{\mu_n\lfloor E_n}{\mu_n(E_n)}
   \qquad(\mu_n(E_n)>0)
\]
be normalized realized defect measures converging weakly to
\(\lambda\in\mathcal P(X)\).  Let \(A\subset X\) be closed.  If every
normalized support profile in \(\operatorname{supp}\lambda\) belongs to \(A\),
then the original normalized defect has no surviving mass in \(X\setminus A\).
Equivalently, if a positive amount of the normalized defect survives outside
\(A\), then there exists a support profile in \(X\setminus A\).
\end{lemma}

\begin{proof}
If positive normalized defect survives outside \(A\), then
\(\lambda(X\setminus A)>0\).  Since \(X\) is compact metrizable and \(A\) is
closed, \(X\setminus A\) is open.  By inner regularity choose a compact set
\(K\subset X\setminus A\) with \(\lambda(K)>0\).  Then
\(K\cap\operatorname{supp}\lambda\ne\varnothing\), and every point of this
intersection is a normalized support profile outside \(A\), contradicting the
support hypothesis.  Hence \(\lambda\) is supported on \(A\), which is exactly
the assertion that no normalized defect mass survives outside the closed
excluded set.
\end{proof}

\begin{lemma}[Transfer of contradictions to the original blow-up sequence]
\label{lem:ancestor-contradiction-transfer}
Let \((V,P)\in\calR(\calS;I,J)\) be a retained residual candidate, and let
\(\mathcal A\) be a closed union of lower alternatives which is stable under
the descendant operations used in this paper.  Suppose that every
limit realized from \((V,P)\) by the relevant construction
has one of the following two properties:
\begin{enumerate}[label=\textup{(\roman*)}]
   \item it is a retained Seregin descendant lying in \(\mathcal A\);
   \item it is a normalized support profile of a defect measure associated with the original sequence and lies
   in \(\mathcal A\).
\end{enumerate}
If \(\mathcal A\) has already been excluded in the ordered residual
decomposition, then the original candidate \((V,P)\) cannot exist.
\end{lemma}

\begin{proof}
In case \textup{(i)}, retained descendant heredity
\Cref{thm:descendant-heredity} identifies the descendant as an admissible
successor of the original residual candidate with the velocity threshold fixed
in \Cref{def:threshold-hierarchy}.  Since \(\mathcal A\) is closed under the
allowed limiting operations and is excluded, the retained descendant cannot
occur.

In case \textup{(ii)}, the normalized support profile is obtained by restricting a
positive defect measure to a window and renormalizing.  If all such
normalized support profiles lie in the closed excluded set \(\mathcal A\), then
regularity of the compactified defect measure implies that the defect
measure has no mass outside \(\mathcal A\); this is the support-transfer
mechanism of \Cref{lem:compact-defect-support-transfer}, instantiated below
for the individual compactifications.  Hence the defect has no surviving
support in the residual class after \(\mathcal A\) is removed.  In both cases
the construction cannot produce a retained residual profile from \((V,P)\),
contradicting the defining retained velocity concentration of the residual candidate.
\end{proof}

\begin{theorem}[Refined residual closure for the complete ordered decomposition]
\label{thm:refined-residual-closure}
Relative to the inherited alternatives in
\Cref{hyp:base-seregin-hypotheses,hyp:previous-nonresidual}, one has
\[
\begin{gathered}
   \Cax=\Crot=\Cstat=\Cgen=\varnothing,\\
   \Clogdiff=\Cyoung=\Chomcrit=\Clogper=\Capcrit=\varnothing .
\end{gathered}
\]
Moreover,
\[
   \Caff\cap\calR^\#(\calS;I,J)=\varnothing ,
\]
that is, no retained normalized residual representative remains in the
affine/parasitic quotient stratum.  Consequently
\[
   \calR^\#(\calS;I,J)=\varnothing .
\]
\end{theorem}

\begin{proof}
\textup{Deferred.}  The proof is given in
\Cref{sec:final-assembly}, after the terminal decomposition theorem and the
exclusion of all reduced cases.  The dependency order is recorded in
\Cref{tab:paperIV-dependency-order}.
\end{proof}

\begin{longtable}{@{}L{0.24\textwidth}L{0.45\textwidth}L{0.23\textwidth}@{}}
\caption{Dependency table for the residual closure proof.}
\label{tab:paperIV-dependency-order}\\
\hline
\textbf{Result} & \textbf{Uses} & \textbf{Does not use}\\
\hline
\endfirsthead
\hline
\textbf{Result} & \textbf{Uses} & \textbf{Does not use}\\
\hline
\endhead
\hline
\endfoot
Imported setup results &
Setup paper only: local concentration, local Type~I implication, normalized
Seregin extraction, pressure atlas, lower classes, mildness, and conditional
final assembly. &
Paper~IV residual closure.\\
R1 axisymmetric closure &
Centered circulation equation, axis regularity, killed-strip comparison, and
setup lower classes. &
Terminal stratification assembly theorem.\\
R2 reduction &
Rotational flux identities and finite unit-cylinder velocity concentration
cover; closure is obtained only after terminal stratification. &
R2 closure or terminal stratification conclusion.\\
R2 closure &
R2 reduction, terminal stratification assembly, finite separated-family exclusion,
and no atomic active profile. &
Standalone rotational Liouville theorem.\\
R3 reduction &
Hull realization, centered scale reselection, and retained velocity heredity;
closure is obtained only after terminal stratification. &
R3 closure or terminal stratification conclusion.\\
R3 closure &
R3 reduction, terminal stratification assembly, sequence-\(L^3\), and endpoint
atomic contradiction. &
Premature stationary-hull exclusion.\\
Concentration-set extraction &
Local compactness, raw atlas defect measures, compact observer windows, and pressure
atlas. &
Endpoint Albritton--Barker theorem.\\
Diffuse-defect compactification &
Concentration-set residual measure, observer compactification, and amenability of
\(\mathbb R^3\rtimes\mathbb R\). &
Generic closure.\\
Affine routing &
Affine normalization dichotomy and retained oscillation functional. &
Critical-tail closure.\\
Coherent critical-tail closure &
Bounded-origin realization and log-hull minimality. &
Log-diffuse exclusion, Young exclusion, hidden-scale exclusion.\\
Minimal mesoscopic reduction &
Local quotient blow-down, corrected scaled pressure compactness, and variance
Liouville theorem. &
Log-diffuse exclusion.\\
No hidden-scale variance &
Minimal mesoscopic reduction and coherent closure. &
Log-diffuse exclusion.\\
Young branch exclusion &
No hidden-scale variance for the non-Dirac subcase and defect-coordinate closure
for pressure/viscous subcases by definition of \(\Cyoung\). &
Log-diffuse exclusion.\\
Log-diffuse branch exclusion &
Log-window support profiles, no hidden-scale variance, and coherent closure. &
Young branch exclusion except through the hidden-scale theorem.\\
Complete critical-tail exclusion &
Regenerated concentration routing, affine routing, log-diffuse exclusion, Young
exclusion, and coherent closure. &
Generic final assembly.\\
Finite separated-profile exclusion &
Diagonal lifting, finite directed graph recurrence, no atomic profile, and no
infinite retained velocity concentration chain. &
Terminal assembly theorem and global \(L^3\) bound.\\
Atomic sequence-\(L^3\) &
Indecomposability and absence of diffuse, tail, separated, pressure-loss, and
critical-tail alternatives. &
Albritton--Barker endpoint conclusion.\\
Endpoint atomic contradiction &
Bounded mild finite terminal shifts, sequence-\(L^3\), and retained nonzero
velocity concentration. &
Residual decomposition.\\
Final assembly &
All previous rows and the ordered decomposition. &
Itself.\\
\hline
\end{longtable}

\begin{corollary}[Residual input from the setup paper]
\label{cor:paperI-residual-input}
For every normalized Seregin collection \(\calS\) generated by an
admissible finite-energy Type I singularity in the setup paper, the setup
remaining class is empty.
\end{corollary}

\begin{proof}[Proof deferred]
This is proved in \Cref{sec:final-assembly}, after
\Cref{thm:refined-residual-closure} has been established and inserted into the
setup paper's residual hypothesis.
\end{proof}

\begin{corollary}[Local pointwise Type I exclusion after residual closure]
\label{cor:typeI-exclusion-main}
No finite-energy local pointwise Type~I singularity can generate a normalized
Seregin collection \(\calS\).
\end{corollary}

\begin{proof}[Proof deferred]
This is the final two-paper consequence recorded in
\Cref{cor:two-paper-local-typeI-exclusion}; it is stated here only to identify
the consequence of residual closure.
\end{proof}

\subsection{How to read the proof}
\label{subsec:reading-guide}

A reader interested only in the residual closure can navigate the paper as
follows.  The proof is read top-down by class.

\begin{enumerate}[label=\textup{(\arabic*)},leftmargin=2.2em]
   \item \emph{Setup interface and Type-I implication}
   (\Cref{subsec:imported-setup-results,sec:local-typeI-implication}). Records the
   setup results used here as imported hypotheses and verifies that local pointwise
   Type-I singularities enter the admissible Seregin compactness setting.
   \item \emph{Residual decomposition} (\Cref{def:refined-decomposition}).
   The ordered list of alternatives is
   \[
      \Cax,\ \Crot,\ \Cstat,\ \Caff,\ \Clogdiff,\quad
      \Cyoung,\ \Chomcrit,\ \Clogper,\ \Capcrit,\ \Cgen .
   \]
   The inherited list of allowed limiting operations is recorded in
   \Cref{thm:ancestor-realization-inheritance}.
   \item \emph{Class \(\Cax\)}. Centered angular-circulation equation
   (\Cref{lem:centered-circulation-equation}), one-dimensional killed radial
   absorption, axis-absorption Liouville theorem, closure
   (\Cref{thm:R1-exclusion}).
   \item \emph{Class \(\Crot\)}. Coriolis Bernoulli identity, local flux
   dichotomy, reduction proposition
   (\Cref{prop:R2-reduction-to-terminal}).  Closure deferred until the terminal
   stratification assembly.
   \item \emph{Class \(\Cstat\)}. Stationary phase space, hull realization,
   centered scale reselection, reduction
   (\Cref{prop:R3-stationary-hull-reduction,prop:R3S-R3-terminal-reduction}).
   Closure deferred until the terminal stratification assembly.
   \item \emph{Generic terminal residual}. Tail geometry and covariant observer
   calculus, asymptotic hull, recurrent-core rigidity, terminal concentration
   profiles, separated-profile theory, Albritton--Barker endpoint, generic
   exhaustion (\Cref{thm:terminal-stratification,thm:R4-final-closure}).
   \item \emph{Final assembly} (\Cref{sec:final-assembly}). Combines all
   class closures and the deferred lower-class arguments.
\end{enumerate}

A reader who only wants to see the central technical core can read in this
order: the asymptotic hull and recurrent-core rigidity, the terminal
concentration-set extraction, the diffuse-defect compactification, the
critical-tail compactification, and the minimal mesoscopic-scale reduction
theorem (\Cref{thm:first-bad-mesoscopic-reduction}).  Everything else is
either tracking for the residual decomposition or reduction of the four
analytic classes into this terminal part of the argument.

\subsection{Local Type I singularities enter the terminal concentration setting}
\label{sec:local-typeI-implication}

The purpose of this subsection is to make explicit the local reduction that is
implicit in the two-paper proof architecture here.  There are two logically
different issues.  The local pointwise Type~I envelope gives Seregin compactness
on every compact negative-time cylinder
\[
   B_R\times[-S,-\sigma],\qquad \sigma>0,
\]
but this compactness is strictly away from the terminal slice \(t=0\).  It does
not alone imply that the positive \(L^3\)-mass in
\(B_1\times(-1,0)\) survives away from the terminal layer
\((- \sigma,0)\).  The latter is the terminal velocity no-escape
condition used in the setup paper.

The implication below first records the admissible form used by the unified setup
paper.  It then uses the endpoint weak Serrin local Type~I estimate of
Albritton--Barker to derive the required scale-invariant \(L^2\) bound on the
terminal interval.  Interpolation with the Type~I envelope gives terminal
velocity no-escape.

For this subsection, if \(z_*=(x_*,T)\) and \(r_n\downarrow0\), write
\[
   u_n(x,t):=r_nu(x_*+r_nx,T+r_n^2t),
   \qquad
   p_n(x,t):=r_n^2p(x_*+r_nx,T+r_n^2t).
\]
Also write \(Q_R:=B_R(0)\times(-R^2,0)\) and
\[
   Q_1(0,0)=B_1(0)\times(-1,0).
\]

\begin{definition}[Admissible local Type I concentration sequence]
\label{def:admissible-local-typeI-sequence}
Let \(z_*=(x_*,T)\) satisfy the local Type~I bound
\eqref{eq:local-typeI-implication-bound}, and let \(r_n\downarrow0\).  The
corresponding rescaled velocities \(u_n\) form an admissible local Type~I
concentration sequence if there are \(\eta_v>0\) such that
\begin{equation}\label{eq:admissible-local-typeI-positive-mass}
   \iint_{B_1\times(-1,0)} |u_n|^3\,dx\,dt\ge\eta_v
   \qquad\hbox{for every }n,
\end{equation}
and the terminal velocity no-escape condition holds:
\begin{equation}\label{eq:admissible-local-typeI-no-escape}
   \lim_{\sigma\downarrow0}
   \sup_n
   \iint_{B_1\times(-\sigma,0)} |u_n|^3\,dx\,dt=0 .
\end{equation}
Equivalently, after lowering the mass threshold, there are
\(\sigma_0\in(0,1)\) and \(\eta_1>0\) such that
\begin{equation}\label{eq:admissible-local-typeI-away-mass}
   \iint_{B_1\times(-1,-\sigma_0)} |u_n|^3\,dx\,dt\ge\eta_1
   \qquad\hbox{for all }n .
\end{equation}
\end{definition}

\begin{lemma}[No-escape leaves positive mass away from the terminal layer]
\label{lem:local-no-escape-away-mass}
If \(u_n\) is an admissible local Type~I concentration sequence in the sense of
\Cref{def:admissible-local-typeI-sequence}, then
\eqref{eq:admissible-local-typeI-away-mass} holds for some
\(\sigma_0\in(0,1)\) and some \(\eta_1>0\).
\end{lemma}

\begin{proof}
Choose \(\sigma_0\in(0,1)\) so small that
\[
   \sup_n
   \iint_{B_1\times(-\sigma_0,0)} |u_n|^3\,dx\,dt
   \le \frac12\eta_v ,
\]
which is possible by \eqref{eq:admissible-local-typeI-no-escape}.  Subtracting
this terminal-layer contribution from
\eqref{eq:admissible-local-typeI-positive-mass} gives
\[
   \iint_{B_1\times(-1,-\sigma_0)} |u_n|^3\,dx\,dt
   \ge \frac12\eta_v .
\]
Take \(\eta_1=\eta_v/2\).
\end{proof}

\begin{proposition}[Local Type I rescalings have Seregin compactness]
\label{prop:local-typeI-terminal-compactness}
Let \((u,p)\) be a suitable Leray--Hopf solution on
\(\R^3\times(0,T)\), let \(z_*=(x_*,T)\) be a singular point, and assume that
there exist \(\rho>0\) and \(M<\infty\) such that
\begin{equation}\label{eq:local-typeI-implication-bound}
   \esssup_{T-\rho^2<t<T}
   \sqrt{T-t}\,\|u(t)\|_{L^\infty(B_\rho(x_*))}\le M .
\end{equation}
Let \(r_n\downarrow0\), and let \((u_n,p_n)\) be the rescaled sequence above.
Then the following hold.
\begin{enumerate}[label=\textup{(\roman*)}]
\item For every compact cylinder
\[
   K=B_R\times[-S,-\sigma]\Subset\R^3\times(-\infty,0),
   \qquad R,S<\infty,\quad \sigma>0,
\]
there exists \(n_0(K)\) such that
\begin{equation}\label{eq:local-typeI-implication-unifLinfty}
   \|u_n\|_{L^\infty(K)}\le M\sigma^{-1/2},
   \qquad n\ge n_0(K).
\end{equation}
\item On every compact cylinder \(K\Subset\R^3\times(-\infty,0)\) there are
time-dependent pressure gauges \(a_{n,K}(t)\) and a constant \(C_K<\infty\)
such that, for \(n\ge n_0(K)\),
\begin{equation}\label{eq:local-typeI-implication-local-bounds}
   \|u_n\|_{L^\infty_tL^2_x(K)}
   +\|\nabla u_n\|_{L^2(K)}
   +\|p_n-a_{n,K}(t)\|_{L^{3/2}(K)}
   \le C_K .
\end{equation}
The constant \(C_K\) may depend on the fixed Leray energy level of the original
solution, but no concentration or no-escape assumption enters it.
\item After subtracting these local pressure gauges, the family
\((u_n,p_n)\) is locally precompact in the standard Seregin topology on
\(\R^3\times(-\infty,0)\).  Along every subsequence there is a further
subsequence and an ancient suitable weak solution \((U,Q)\) such that, on
every compact \(K\Subset\R^3\times(-\infty,0)\),
\[
   u_n\to U\quad\hbox{strongly in }L^q(K),\qquad 1\le q<\infty,
\]
and
\[
   p_n-a_{n,K}(t)\rightharpoonup Q
   \quad\hbox{weakly in }L^{3/2}(K).
\]
The limit is bounded and satisfies
\[
   |U(x,t)|\le M(-t)^{-1/2},\qquad t<0.
\]
\item More generally, let \(\mathfrak c=(z_n)\), \(z_n=(y_n,\tau_n)\), be a
parabolic recentering sequence for \((u_n,p_n)\) with \(\tau_n\le0\).  Assume
that \(\mathfrak c\) remains in the local Type~I window in the following
precise sense: for every \(R,S<\infty\), all sufficiently large \(n\) satisfy
\[
   r_n(|y_n|+R)<\rho,
   \qquad
   T+r_n^2(\tau_n-S)>T-\rho^2 .
\]
If \(\mathfrak c\) carries retained unit velocity concentration, then after
passing to a subsequence it has no loss of local compactness in the sense of
\Cref{def:terminal-measures}, and its recentered fields converge locally to a
bounded suitable terminal profile.  Recenterings which do not satisfy this
local-window condition are assigned by the compact-window-to-expanding-region
protocol to the exterior, diffuse, pressure/noncompact, or critical-tail
terminal alternatives rather than to the compact local profile class.
\end{enumerate}
\end{proposition}

\begin{proof}
Fix
\[
   K=B_R\times[-S,-\sigma]\Subset\R^3\times(-\infty,0).
\]
For \(n\) sufficiently large,
\[
   r_nR<\rho/2,\qquad r_n^2S<\rho^2/2.
\]
Hence, for a.e. \((x,t)\in K\),
\[
   x_*+r_nx\in B_\rho(x_*),
   \qquad
   T+r_n^2t\in(T-\rho^2,T).
\]
The local Type~I bound \eqref{eq:local-typeI-implication-bound} gives
\[
   |u_n(x,t)|
   =r_n|u(x_*+r_nx,T+r_n^2t)|
   \le r_n M\,(T-(T+r_n^2t))^{-1/2}
   =M(-t)^{-1/2}
   \le M\sigma^{-1/2}.
\]
This proves \eqref{eq:local-typeI-implication-unifLinfty}.  In particular,
\[
   \esssup_{-S<t<-\sigma}\int_{B_R}|u_n(x,t)|^2\,dx
   \le |B_R|M^2\sigma^{-1},
\]
and
\[
   \iint_K |u_n|^3\,dx\,dt
   \le |K|M^3\sigma^{-3/2}.
\]

We next record the pressure estimate, where the nonlocal pressure contribution
is localized in the local Type~I compactness argument.  We use the whole-space
Leray pressure representative, so at each rescaled time
\[
   p_n(\cdot,t)=\mathcal R_i\mathcal R_j
      \bigl(u_{n,i}u_{n,j}\bigr)(\cdot,t)
\]
modulo a function of time.  Fix a cutoff
\(\chi_R\in C_c^\infty(B_{8R})\) with \(\chi_R\equiv1\) on \(B_{4R}\), and
write in \(B_{4R}\)
\[
   p_n=q_n+h_n,\qquad
   q_n=\mathcal R_i\mathcal R_j(\chi_R u_{n,i}u_{n,j}).
\]
The local term satisfies
\[
   \|q_n\|_{L^{3/2}(B_{4R}\times[-S,-\sigma])}
   \le C_{R,S,\sigma,M}.
\]
For the harmonic part choose the gauge \(a_{n,K}(t)=h_n(0,t)\).  The standard
kernel difference estimate for the Riesz pressure kernel is as follows.  If
\[
   K_{ij}(x)=c_3\,\mathrm{p.v.}\,
      \frac{3x_ix_j-\delta_{ij}|x|^2}{|x|^5}
\]
is the Calder\'on--Zygmund pressure kernel, then for \(|x|\le2R\) and
\(|z|\ge4R\),
\[
   |K_{ij}(x-z)-K_{ij}(-z)|\le C_R|z|^{-4}.
\]
Decomposing the exterior region into dyadic annuli gives, for \(x\in B_{2R}\),
\[
   |h_n(x,t)-a_{n,K}(t)|
   \le
   C_R\sum_{m\ge2}2^{-4m}
   \int_{2^mR<|z|<2^{m+1}R}|u_n(z,t)|^2\,dz .
\]
For annuli whose physical image remains in \(B_\rho(x_*)\), namely while
\(2^{m+1}Rr_n\le\rho\), the local Type~I bound and \(-t\ge\sigma\) on \(K\)
give
\[
   \int_{2^mR<|z|<2^{m+1}R}|u_n(z,t)|^2\,dz
   \le C M^2\sigma^{-1}(2^mR)^3 .
\]
Multiplication by \(2^{-4m}\) leaves a summable \(2^{-m}\) factor, so the
corresponding part of the series is bounded by a geometric sum depending only
on \(R,M,\sigma\).  For the remaining outer annuli one uses only the finite
Leray energy of the original solution:
\[
   \int_{\R^3}|u_n(z,t)|^2\,dz
   =r_n^{-1}\int_{\R^3}|u(y,T+r_n^2t)|^2\,dy .
\]
If \(m_n\) is the first index for which \(2^{m_n}Rr_n\gtrsim\rho\), then
\[
   \sum_{m\ge m_n}2^{-4m}\int_{\R^3}|u_n(z,t)|^2\,dz
   \le
   C_{\rho,R}\,r_n^3\|u\|_{L^\infty(0,T;L^2)}^2
   \le C_{\rho,R,u}.
\]
Thus \(h_n-a_{n,K}\) is uniformly bounded on
\(B_{2R}\times[-S,-\sigma]\).  Since this cylinder has finite measure,
\[
   \|h_n-a_{n,K}\|_{L^{3/2}(B_{2R}\times[-S,-\sigma])}
   \le C_K .
\]
Together with the Calder\'on--Zygmund bound for \(q_n\), this gives
\[
   \|p_n-a_{n,K}(t)\|_{L^{3/2}(B_{2R}\times[-S,-\sigma])}
   \le C_K .
\]

Apply the local energy inequality for \((u_n,p_n)\) with a cutoff equal to one
on \(K\) and supported in a slightly larger compact cylinder
\(K'\Subset\R^3\times(-\infty,0)\).  The already obtained local \(L^2\),
\(L^3\), and pressure-gauge bounds control every term on the right-hand side,
and therefore
\[
   \iint_K |\nabla u_n|^2\,dx\,dt\le C_K .
\]
This proves \eqref{eq:local-typeI-implication-local-bounds}.

The Navier--Stokes equation now bounds \(\partial_tu_n\) locally in a negative
Sobolev space, for instance in \(L^{3/2}_tH^{-2}_x(K)\).  Aubin--Lions
compactness and the uniform \(L^\infty\) bound give strong local convergence of
the velocity in every finite \(L^q\), after passing to a subsequence.  The
pressure bound gives weak \(L^{3/2}\) convergence after gauges.  Passing to the
limit in the distributional equations and in the local energy inequality gives
an ancient suitable weak solution \((U,Q)\).  The pointwise Type~I bound passes
to the limit, so \(U\) is bounded on every compact negative-time cylinder.
This proves (iii).

Finally, let \(\mathfrak c=(y_n,\tau_n)\) be as in (iv), and set
\[
   u_n^{\mathfrak c}(y,s)=u_n(y+y_n,s+\tau_n),
   \qquad
   p_n^{\mathfrak c}(y,s)=p_n(y+y_n,s+\tau_n).
\]
The local-window condition, together with \(\tau_n\le0\), is what is
needed to repeat the preceding estimates on each fixed backward compact cylinder
in the \((y,s)\)-variables: the
physical cylinders remain in \(B_\rho(x_*)\times(T-\rho^2,T)\), the local
Type~I envelope gives the same \(L^\infty\) bound on compact windows, and the
pressure decomposition has the same local and finite-energy tail estimates.
Thus the two failure modes in \Cref{def:terminal-measures}, blow-up of the
local dissipation measure or failure of local pressure-gauge control, cannot
occur for such a retained local recentering sequence after subsequence
extraction.  The recentered fields therefore converge locally to a bounded
suitable terminal profile.  If the local-window condition fails, then the
recentering leaves every bounded terminal window attached to the singular
point, or its critical mass persists only through expanding regions; those are
precisely the exterior and diffuse terminal alternatives rather than loss of
local compactness inside the local Type~I implication.
\end{proof}

\subsection{Pressure-projected terminal energy estimate}
\label{subsec:pressure-projected-terminal-estimate}

The preceding compactness result is stated away from the terminal
slice.  The next estimate supplies the terminal-layer bound required for
velocity no-escape.  The local pressure projection removes the nonlocal pressure
flux from the local energy inequality; the remaining estimate is an interior
Caccioppoli propagation bound on balls contained in the local Type~I window.

Let \(v=u^{(r)}\), \(\pi=p^{(r)}\), and
\(\Lambda=\rho/r\).  After translating \(z_*\) to \((0,0)\), the local Type~I
bound gives
\[
   |v(x,s)|\le M(-s)^{-1/2}
   \qquad (x,s)\in B_\Lambda\times(-4,0),
\]
for all sufficiently small \(r\).  For \(a\in\mathbb R^3\), \(L\ge1\), and
\(-4<s<t<0\), write
\[
   A_L(a,s):=\int_{B_L(a)}|v(x,s)|^2\,dx,\qquad
   D_L(a;s,t):=\int_s^t\int_{B_L(a)}|\nabla v|^2\,dx\,d\sigma .
\]

\begin{lemma}[Local pressure-projected Caccioppoli inequality]
\label{lem:local-pressure-projected-caccioppoli}
There is a universal constant \(C\) with the following property.  Let
\((v,\pi)\) be a suitable weak solution in
\(B_{8L}(a)\times(s_-,0)\), \(L\ge1\), and assume
\[
   |v(x,s)|\le M(-s)^{-1/2}
   \qquad\hbox{on }B_{8L}(a)\times(s_-,0).
\]
If \(s_0\in(s_-,0)\) is a Lebesgue time for \(v\) and \(\nabla v\), then for
a.e. \(t\in(s_0,0)\),
\begin{equation}\label{eq:pressure-projected-caccioppoli}
\begin{aligned}
   \int_{B_L(a)}|w(x,t)|^2\,dx+D_L(a;s_0,t)
   &\le
   C A_{4L}(a,s_0)                                      \\
   &\quad
   +C\int_{s_0}^{t}
      \left(L^{-2}+ML^{-1}(-s)^{-1/2}\right)
      A_{4L}(a,s)\,ds .
\end{aligned}
\end{equation}
The estimate is entirely local on \(B_{8L}(a)\); no pressure tail from outside
\(B_{8L}(a)\) appears.  Here \(w=v+\nabla p_h\) is the local projected velocity
defined in the proof below.
\end{lemma}

\begin{proof}
The proof uses the standard local pressure-projection argument in the scale \(L\)
normalization.  Let \(B=B_{8L}(a)\).  Denote by \(E_B^*\) the Stokes pressure
projection on \(B\): for \(F\in W^{-1,q}(B)\), \(1<q<\infty\),
\(\nabla E_B^*(F)\) is the gradient part of the stationary Stokes solution in
\(B\) with zero boundary velocity.  It satisfies the scale-invariant estimates
\[
   \|\nabla E_B^*(\nabla\cdot F)\|_{L^q(B)}
   \le C_q\|F\|_{L^q(B)},\qquad
   \|\nabla E_B^*(\Delta w)\|_{L^2(B)}
   \le C\|\nabla w\|_{L^2(B)}.
\]
Set
\[
   \nabla p_h=-\nabla E_B^*(v),\qquad
   w:=v+\nabla p_h .
\]
The harmonic pressure \(p_h\) is harmonic in \(B\), and the projection estimates
give
\[
   \|w\|_{L^2(B)}+\|\nabla p_h\|_{L^2(B)}
   \le C\|v\|_{L^2(B)},\qquad
   \|\nabla w\|_{L^2(B_{6L}(a))}
   \le C\|\nabla v\|_{L^2(B)} .
\]
The projected local energy inequality for \(w\), obtained by applying the
Stokes projection to the distributional Navier--Stokes equations and testing
with \(\zeta^2 w\), reads
\begin{align}
   &\int |w(t)|^2\zeta^2
      +\int_{s_0}^{t}\int |\nabla v|^2\zeta^2              \notag\\
   &\qquad\le
      C\int |w(s_0)|^2\zeta^2
      +C\int_{s_0}^{t}\int |w|^2
        \left(|\partial_s\zeta|+|\Delta\zeta|+|\nabla\zeta|^2\right)
                                                                  \notag\\
   &\qquad\quad
      +C\int_{s_0}^{t}\int |v|\,|w|^2\,|\nabla\zeta|
      +C\int_{s_0}^{t}\int |v|\,|\nabla p_h|^2\,|\nabla\zeta| .
\label{eq:projected-local-energy}
\end{align}
Here \(\zeta\in C_c^\infty(B_{4L}(a))\), \(\zeta\equiv1\) on \(B_L(a)\), and
\[
   |\nabla\zeta|\le CL^{-1},\qquad
   |\Delta\zeta|+|\partial_s\zeta|\le CL^{-2}.
\]
The pressure-projection estimates allow \(w\) and \(\nabla p_h\) to be bounded
by \(v\) in \(L^2(B_{4L}(a))\).  Therefore the cutoff terms in
\eqref{eq:projected-local-energy} are bounded by
\[
   C L^{-2}\int_{s_0}^{t}A_{4L}(a,s)\,ds .
\]
The two convective terms are bounded, using the Type~I envelope, by
\[
   C L^{-1}\int_{s_0}^{t}
      \|v(s)\|_{L^\infty(B_{4L}(a))}A_{4L}(a,s)\,ds
   \le
   CML^{-1}\int_{s_0}^{t}(-s)^{-1/2}A_{4L}(a,s)\,ds .
\]
The initial projected energy is controlled by
\[
   \int |w(s_0)|^2\zeta^2\,dx
   \le C A_{4L}(a,s_0).
\]
Combining the preceding bounds gives
\eqref{eq:pressure-projected-caccioppoli}.
\end{proof}

\begin{definition}[Interior terminal energy estimate]
\label{def:interior-shielding-estimate}
The local pointwise Type~I window at \(z_*\) is said to satisfy the interior
terminal energy estimate if, for every fixed \(R<\infty\), there are
\[
   r_R>0,\qquad B_R<\infty,
\]
such that every rescaling \(v=u^{(r)}\), \(0<r<r_R\), with
\(\Lambda=\rho/r\), satisfies
\begin{equation}\label{eq:interior-shielding-estimate}
   \esssup_{-1<s<0}\int_{B_{8R}}|v(x,s)|^2\,dx
   +\int_{-1}^{0}\int_{B_{8R}}|\nabla v|^2\,dx\,ds
   \le B_R .
\end{equation}
Equivalently, the dependence on \(A_{4L}\) in
\Cref{lem:local-pressure-projected-caccioppoli} can be stopped inside a fixed
multiple of the target ball, uniformly as the boundary
\(\partial B_{\rho/r}\) of the local Type~I window recedes to infinity.
\end{definition}

\begin{proposition}[Pressure-projected reduction to an interior energy bound]
\label{prop:pressure-projection-reduces-to-shielding}
Assume the local pointwise Type~I bound \eqref{eq:local-typeI-implication-bound} at
\(z_*\).  If the interior energy estimate
\eqref{eq:interior-shielding-estimate} holds with \(R=1\), then the velocity
no-escape condition \eqref{eq:admissible-local-typeI-no-escape} holds for every
positive local Type~I concentration sequence at \(z_*\).  Consequently such a
sequence is admissible in the sense of
\Cref{def:admissible-local-typeI-sequence}.
\end{proposition}

\begin{proof}
Let \(u_n=u^{(r_n)}\) be any positive local Type~I concentration sequence.  By
\eqref{eq:interior-shielding-estimate} with \(R=1\), after discarding finitely
many indices,
\[
   \esssup_{-1<s<0}\int_{B_1}|u_n(x,s)|^2\,dx\le B_1 .
\]
The local Type~I envelope gives
\[
   \|u_n(s)\|_{L^\infty(B_1)}\le M(-s)^{-1/2}.
\]
Therefore, for \(0<\sigma<1\),
\[
\begin{aligned}
   \int_{-\sigma}^{0}\int_{B_1}|u_n|^3\,dx\,ds
   &\le
   \int_{-\sigma}^{0}
      \|u_n(s)\|_{L^\infty(B_1)}
      \int_{B_1}|u_n(x,s)|^2\,dx\,ds    \\
   &\le
   2MB_1\sigma^{1/2}.
\end{aligned}
\]
The finitely many discarded indices have absolutely continuous
\(|u_n|^3\)-integrals on \(B_1\times(-1,0)\), so taking the supremum over the
whole sequence and then letting \(\sigma\downarrow0\) gives
\eqref{eq:admissible-local-typeI-no-escape}.  The positive \(Q_1\) mass
hypothesis is the other part of
\Cref{def:admissible-local-typeI-sequence}.
\end{proof}

\begin{remark}[Pressure-projected reduction]
\label{rem:narrowed-final-implication}
\Cref{lem:local-pressure-projected-caccioppoli} shows, independently of the
Albritton--Barker endpoint input used below, that terminal no-escape reduces
to an interior terminal energy estimate: after local pressure projection, the local
energy inequality controls the projected velocity and the dissipation up to the
larger-scale quantity \(A_{4L}\).  The estimate
\eqref{eq:interior-shielding-estimate} prevents
terminal-layer concentration entering from the boundary of the local Type~I
window.  In the present paper this bound is supplied
instead by the endpoint weak Serrin local Type~I bounds in
\Cref{thm:AB-weak-serrin-typeI-estimate}.
\end{remark}

\begin{theorem}[Setup local Type-I implication]
\label{thm:AB-weak-serrin-typeI-estimate}
Let \((v,q)\) be a suitable weak solution in a parabolic cylinder
\[
   Q_\rho(z_0)=B_\rho(x_0)\times(t_0-\rho^2,t_0).
\]
Assume the local pointwise Type~I envelope
\[
   \esssup_{t_0-\rho^2<t<t_0}
   (t_0-t)^{1/2}\|v(t)\|_{L^\infty(B_\rho(x_0))}\le M<\infty .
\]
Then the setup paper proves that the terminal local energy, pressure-gauge,
and no-escape bounds stated in
\Cref{prop:local-typeI-terminal-compactness} hold on every smaller cylinder
\(Q_{\theta\rho}(z_0)\) with \(0<\theta<1\), with constants depending only on
\(\rho,\theta,M\) and on the suitability constants of \((v,q)\) on
\(Q_\rho(z_0)\).  In particular every rescaling sequence centered at a local
pointwise Type~I terminal singularity satisfies
\Cref{def:admissible-local-typeI-sequence}.
\end{theorem}

\begin{proof}
This theorem is imported verbatim from the companion setup paper; see
\Cref{thm:imported-setup-results}\,(2).  Paper~IV does not reprove the
weak-Serrin-to-local-Type~I estimate or the terminal no-escape estimate.  Its
only role here is to place first-generation terminal sequences inside the
admissible local-window class and to supply the local energy, pressure, and
no-escape bounds used by
\Cref{def:first-generation-local-window-condition}.
\end{proof}

\begin{remark}[Albritton--Barker as the underlying input]
\label{rem:AB-underlying-input}
The setup proof of
\Cref{thm:AB-weak-serrin-typeI-estimate} relies on the Albritton--Barker
endpoint weak-Serrin estimate, which gives finite scale-invariant
\(A,C,D,E\)-quantities on every \(Q_{\theta\rho}(z_0)\) under the local
pointwise Type~I envelope.  The constants depend on local suitable-solution
background quantities on \(Q_\rho(z_0)\), with no circularity because the
input envelope and the suitable-solution data are part of the hypothesis
data of the cylinder \(Q_\rho(z_0)\), while the conclusion is restricted to
\(Q_{\theta\rho}(z_0)\).  The detailed weak-Serrin discussion belongs to the
setup paper; it is recalled here only to identify the source of the constants
used in \Cref{prop:local-typeI-terminal-compactness}.
\end{remark}

\begin{theorem}[Local pointwise Type I implies terminal no-escape]
\label{thm:local-typeI-terminal-energy-estimates}
Let \((u,p)\) be a suitable Leray--Hopf solution on
\(\mathbb R^3\times(0,T)\), let \(z_*=(x_*,T)\), and assume that
\[
  u\in L^\infty(0,T;L^2(\mathbb R^3))
  \cap L^2(0,T;\dot H^1(\mathbb R^3)).
\]
Assume that \(z_*\) satisfies the local pointwise Type~I bound: there are
\(\rho>0\) and \(M<\infty\) such that
\begin{equation}\label{eq:local-typeI-terminal-noescape-assumption}
  \esssup_{T-\rho^2<t<T}
  \sqrt{T-t}\,\|u(t)\|_{L^\infty(B_\rho(x_*))}
  \le M .
\end{equation}
For \(0<r<\rho\), set
\[
  u^{(r)}(x,s)
  =
  r\,u(x_*+rx,T+r^2s),
  \qquad
  p^{(r)}(x,s)
  =
  r^2p(x_*+rx,T+r^2s).
\]
Then for every \(0<R<\infty\) there are constants
\[
  r_R\in(0,\rho/R),\qquad B_R<\infty,
\]
and time-dependent pressure gauges \(a_{r,R}(s)\), depending on
\(R,M,\rho,T\) and the local suitable-solution background quantities of
\((u,p)\) in \(Q_\rho(z_*)\), such that the following terminal local energy
estimate holds for every \(0<r<r_R\):
\begin{equation}\label{eq:local-terminal-energy-estimates}
\begin{aligned}
  &\esssup_{-1<s<0}\int_{B_R}|u^{(r)}(x,s)|^2\,dx
  +\int_{-1}^{0}\int_{B_R}|\nabla u^{(r)}|^2\,dx\,ds  \\
  &\qquad
  +\int_{-1}^{0}\int_{B_R}
      |p^{(r)}(x,s)-a_{r,R}(s)|^{3/2}\,dx\,ds
  \le B_R .
\end{aligned}
\end{equation}
Consequently, for every \(0<\sigma<1\) and every \(0<r<r_R\),
\begin{equation}\label{eq:local-terminal-noescape-bound}
  \int_{-\sigma}^{0}\int_{B_R}|u^{(r)}(x,s)|^3\,dx\,ds
  \le
  2MB_R\sigma^{1/2},
\end{equation}
and in particular
\begin{equation}\label{eq:local-terminal-noescape}
  \lim_{\sigma\downarrow0}
  \sup_{0<r<r_R}
  \int_{-\sigma}^{0}\int_{B_R}|u^{(r)}(x,s)|^3\,dx\,ds
  =0 .
\end{equation}
\end{theorem}

\begin{proof}
Fix once and for all \(0<\theta<1\), for instance \(\theta=1/2\), and put
\(L_R:=\max\{R,1\}\).  By decreasing \(r_R\), we may assume
\[
   B_{L_Rr}(x_*)\times(T-L_R^2r^2,T)\subset Q_{\theta\rho}(z_*)
   \qquad (0<r<r_R).
\]
The pointwise Type~I bound implies the endpoint weak Serrin estimate on
\(Q_\rho(z_*)\).  Indeed, after writing
\[
   f(t)=\|u(t)\|_{L^\infty(B_\rho(x_*))},
\]
one has
\[
   |\{t\in(T-\rho^2,T): f(t)>\alpha\}|
   \le M^2\alpha^{-2},
   \qquad \alpha>0,
\]
and therefore \(u\in L^{2,\infty}_tL^\infty_x(Q_\rho(z_*))\).  Applying
\Cref{thm:AB-weak-serrin-typeI-estimate} gives a finite constant \(K_\theta\)
such that \(A,C,D,E\le K_\theta\) on all cylinders contained in
\(Q_{\theta\rho}(z_*)\).

Taking the cylinder \(Q_{L_Rr}(z_*)\) gives
\[
   (L_Rr)^{-1}
   \esssup_{T-L_R^2r^2<t<T}
      \int_{B_{L_Rr}(x_*)}|u(x,t)|^2\,dx
   \le K_\theta ,
\]
and similarly the \(D\)- and \(E\)-bounds give
\[
   (L_Rr)^{-1}\int_{T-L_R^2r^2}^{T}
      \int_{B_{L_Rr}(x_*)}|\nabla u|^2\,dx\,dt
   \le K_\theta ,
\]
\[
   (L_Rr)^{-2}\int_{T-L_R^2r^2}^{T}
      \int_{B_{L_Rr}(x_*)}|p-a_{r,R}^{\rm phys}(t)|^{3/2}\,dx\,dt
   \le K_\theta
\]
for suitable time-dependent pressure gauges \(a_{r,R}^{\rm phys}(t)\).  Scaling
these three estimates and restricting to
\[
   B_R\times(-1,0)\subset B_{L_R}\times(-L_R^2,0)
\]
gives \eqref{eq:local-terminal-energy-estimates} with
\[
   a_{r,R}(s):=r^2a_{r,R}^{\rm phys}(T+r^2s),
\]
after enlarging the constant \(B_R\) by a harmless factor depending on \(R\)
and \(K_\theta\).

After translating \((x_*,T)\) to \((0,0)\), the rescaled fields also satisfy
\[
   |u^{(r)}(x,s)|\le M(-s)^{-1/2}
   \qquad (x,s)\in B_R\times(-1,0)
\]
for all \(r<r_R\), after decreasing \(r_R\) if necessary.  The terminal local
energy estimate established above gives
\[
   \esssup_{-1<s<0}\int_{B_R}|u^{(r)}(x,s)|^2\,dx\le B_R .
\]

For \(0<\sigma<1\), the local Type~I bound above and the kinetic bound give
\[
\begin{aligned}
  \int_{-\sigma}^{0}\int_{B_R}|u^{(r)}|^3\,dx\,ds
  &\le
  \int_{-\sigma}^{0}
     \|u^{(r)}(s)\|_{L^\infty(B_R)}
     \int_{B_R}|u^{(r)}(x,s)|^2\,dx\,ds  \\
  &\le
  MB_R\int_{-\sigma}^{0}(-s)^{-1/2}\,ds
  =
  2MB_R\sigma^{1/2}.
\end{aligned}
\]
This proves \eqref{eq:local-terminal-noescape-bound} and
\eqref{eq:local-terminal-noescape}.
\end{proof}

\begin{corollary}[No diffuse terminal-layer mass]
\label{cor:no-diffuse-terminal-layer-mass}
Under the hypotheses of \Cref{thm:local-typeI-terminal-energy-estimates}, for
every fixed \(R<\infty\) there do not exist \(r_n\downarrow0\),
\(\sigma_n\downarrow0\), and \(\varepsilon_0>0\) such that
\[
  r_n^{-2}
  \int_{T-\sigma_n r_n^2}^{T}
  \int_{B_{Rr_n}(x_*)}|u(x,t)|^3\,dx\,dt
  \ge \varepsilon_0
  \qquad\hbox{for all }n .
\]
\end{corollary}

\begin{proof}
For \(u_n=u^{(r_n)}\), the left hand side is
\[
   \int_{-\sigma_n}^{0}\int_{B_R}|u_n(x,s)|^3\,dx\,ds .
\]
By \eqref{eq:local-terminal-noescape-bound}, this is bounded by
\[
   2MB_R\sigma_n^{1/2}\to0,
\]
which contradicts any fixed positive lower bound.
\end{proof}

\begin{corollary}[Local pointwise Type I concentration sequences are admissible]
\label{cor:local-typeI-implies-admissible}
Let \((u,p)\) and \(z_*=(x_*,T)\) satisfy the hypotheses of
\Cref{thm:local-typeI-terminal-energy-estimates}.  Let \(r_n\downarrow0\) be any
velocity concentration sequence satisfying
\[
   \iint_{B_1\times(-1,0)} |u_n|^3\,dx\,dt\ge\eta_v>0,
   \qquad
   u_n(x,t)=r_nu(x_*+r_nx,T+r_n^2t).
\]
Then \(u_n\) is an admissible local Type~I concentration sequence in the sense
of \Cref{def:admissible-local-typeI-sequence}.
\end{corollary}

\begin{proof}
Apply \Cref{thm:local-typeI-terminal-energy-estimates} with \(R=1\).  After
discarding finitely many \(n\), all \(r_n<r_1\), and
\[
   \lim_{\sigma\downarrow0}
   \sup_n
   \int_{-\sigma}^{0}\int_{B_1}|u_n|^3\,dx\,dt=0 .
\]
The finitely many discarded indices do not affect the limit as
\(\sigma\downarrow0\), because each fixed suitable weak solution has
\(|u_n|^3\in L^1(B_1\times(-1,0))\).  The positive mass hypothesis is the other
condition in \Cref{def:admissible-local-typeI-sequence}.
\end{proof}

\begin{lemma}[Local Type I concentration enters the terminal dichotomy]
\label{lem:local-typeI-into-terminal-dichotomy}
Under the hypotheses of
\Cref{prop:local-typeI-terminal-compactness}, let \(r_n\downarrow0\) be a
velocity concentration sequence at \(z_*\) satisfying
\[
   \iint_{B_1\times(-1,0)} |u_n|^3\,dx\,dt\ge\eta_v>0 .
\]
Then \(r_n\) is admissible by \Cref{cor:local-typeI-implies-admissible}.  Let
\(\sigma_0,\eta_1\) be given by
\Cref{lem:local-no-escape-away-mass}.
For the associated terminal extraction sequence \((V_N,P_N)\) generated by a
Seregin limit of this sequence, choose the small-data threshold
\(\eps_{\rm sd}\le\eta_1/2\).  Then the origin extraction produces a bounded
ancient Seregin profile with retained local velocity concentration away from the
terminal layer.  In particular the retained velocity concentration generated
from the origin recentering has no loss of local compactness.  Any retained
local recentering which stays in the local Type~I window has the same compactness
property by \Cref{prop:local-typeI-terminal-compactness}; any concentration
leaving that window is assigned to the exterior or diffuse terminal remainder.
The paired terminal decomposition of \Cref{thm:terminal-exhaustion-main},
together with the mixed-profile exclusions and concentration-chain exclusions
\Cref{thm:no-mixed-compact-diffuse,thm:compact-active-descendant}, reduces the
retained possibilities to a single retained concentrating terminal profile or a
finite separated family of retained concentrating terminal profiles.
\end{lemma}

\begin{proof}
By \Cref{lem:local-no-escape-away-mass},
\[
   \iint_{B_1\times(-1,-\sigma_0)} |u_n|^3\,dx\,dt\ge\eta_1
   \qquad\hbox{for all }n .
\]
The cylinder \(B_1\times[-1,-\sigma_0]\) is compactly contained in
\(\mathbb R^3\times(-\infty,0)\).  Therefore
\Cref{prop:local-typeI-terminal-compactness} gives, after passing to a
subsequence, strong local \(L^3\) convergence on this cylinder:
\[
   u_n\to U
   \qquad\hbox{in }L^3(B_1\times(-1,-\sigma_0)).
\]
Consequently
\[
   \iint_{B_1\times(-1,-\sigma_0)} |U|^3\,dx\,dt\ge\eta_1 .
\]
Cover \(B_1\times(-1,-\sigma_0)\) by finitely many unit terminal cylinders
strictly away from \(t=0\).  At least one such cylinder carries a fixed positive
velocity lower bound, and since the raw atlas defect measure dominates the velocity
part, choosing \(\eps_{\rm sd}\le\eta_1/2\) after adjusting by this finite cover
makes the corresponding origin-frame recentering raw unit-window active at
threshold \(\eps_{\rm sd}\).
This argument is precisely where no-escape is used: without
\eqref{eq:admissible-local-typeI-no-escape}, the lower bound could remain
entirely in \(B_1\times(-\sigma,0)\) for every fixed \(\sigma>0\).

The distinguished recentering satisfies the local-window condition in
\Cref{prop:local-typeI-terminal-compactness}, so it is not in the
local-compactness-failure class.  Other retained local recenterings inside the
same Type~I window have the same compactness.  Recenterings leaving the window
are exterior or diffuse terminal remainders.  The paired decomposition
\Cref{thm:terminal-exhaustion-main} records these remainders alongside the
concentration set rather than treating them as mutually exclusive alternatives.  The
sole-source radiative cases are excluded by
\Cref{thm:terminal-inactive-exclusion}, mixed compact--diffuse or
escaping noncompact coexistence is handled by
\Cref{thm:no-mixed-compact-diffuse,lem:escaping-retained-recentering-routing}, and infinite chains of retained concentrating descendants are excluded
by \Cref{thm:compact-active-descendant}.  Hence
\Cref{cor:persistent-forces-active} leaves only a single retained concentrating
terminal profile or a finite separated family of retained concentrating terminal
profiles.
\end{proof}

\begin{theorem}[Local Type I singularities enter the terminal concentration setting]
\label{thm:local-typeI-entry-or-exclusion}
Let \((u,p)\) be a suitable Leray--Hopf solution on
\(\R^3\times(0,T)\), let \(z_*=(x_*,T)\) be a singular point, and assume that
\(z_*\) satisfies the local Type~I bound
\eqref{eq:local-typeI-implication-bound}.  Then there exists a bounded centered
Type~I Seregin profile with retained local velocity concentration generated from
\((u,p,z_*)\).  Consequently the local terminal concentration setting of
\Cref{thm:terminal-stratification} applies after the residual
closure theorem is proved.
\end{theorem}

\begin{proof}
By the local velocity concentration theorem in the setup paper, choose
\(r_n\downarrow0\) and \(\eta_v>0\) with
\[
   \iint_{B_1\times(-1,0)} |u_n|^3\,dx\,dt\ge\eta_v,
   \qquad
   u_n(x,t)=r_nu(x_*+r_nx,T+r_n^2t).
\]
\Cref{thm:local-typeI-terminal-energy-estimates,cor:local-typeI-implies-admissible}
show that this concentration sequence satisfies the terminal
velocity no-escape condition required by the setup paper.
Apply \Cref{lem:local-typeI-into-terminal-dichotomy} to this sequence.  The
local Type~I rescalings therefore produce at least one bounded terminal profile
carrying retained local velocity concentration.  After passing to the centered variables, the
bound \(|U(x,t)|\le M(-t)^{-1/2}\) becomes the global centered Type~I bound
\[
   |V(y,\tau)|\le M,
\]
so each retained profile generated by this implication is a bounded centered Type~I
Seregin profile.

If every such retained profile belongs to one of the previously closed lower
classes, then the contradiction is supplied by the setup paper:
small-amplitude and stationary \(L^3\) profiles are excluded there, the
uniformly tight class is excluded there, and the closed structured/decay class
is excluded there.  Otherwise choose a retained profile outside those lower
classes.  It is then a refined residual candidate, so
\Cref{thm:terminal-stratification} applies.  In that residual
setting, a finite separated family of retained concentrating terminal profiles
is excluded by \Cref{thm:no-finite-separated-profile-family}, infinite chains of
retained concentrating descendants are excluded by \Cref{thm:compact-active-descendant}, and mixed
compact--diffuse tails are excluded by \Cref{thm:no-mixed-compact-diffuse}.  The
remaining single retained profile is the terminal case closed in
\Cref{sec:final-assembly}.  Thus the present theorem supplies the entry into
the residual decomposition; the contradiction is the final two-paper corollary
proved after the residual closure.
\end{proof}

\begin{corollary}[Local pointwise Type~I exclusion after setup plus residual closure]
\label{cor:local-typeI-unconditional}
Let \((u,p)\) be a suitable Leray--Hopf solution on
\(\R^3\times(0,T)\).  Combining the residual closure
\Cref{thm:refined-residual-closure} of this paper with
the setup paper, the following holds: if \(z_*=(x_*,T)\) is a singular point
satisfying the local Type~I bound \eqref{eq:local-typeI-implication-bound}, then the
completed case closures give a contradiction.  Equivalently, in the local
Type~I/local Type~II dichotomy of the setup paper, every finite-energy
singular point belongs to the local Type~II alternative.  The conclusion is
the joint output of the two-paper chain; it is not asserted as a hidden
hypothesis inside Paper~IV alone.
\end{corollary}

\begin{proof}
The entry assertion is \Cref{thm:local-typeI-entry-or-exclusion}.  The
contradiction is deferred to the final assembly, where the residual closure is
proved and then inserted into the setup paper's conditional final theorem.  The
equivalence is the local Type~I/local Type~II dichotomy of the setup paper.
\end{proof}

\section{Ordered residual profile decomposition}
\label{sec:mermaid-decision-tree}

The following diagram records the ordered partition used in the
paper.  Each diamond is a stratification criterion.  Each terminal branch is a
possible singularity or candidate blow-up profile type selected by the criteria.
The local Type~II terminal alternative is the remaining regime after the local
pointwise Type~I case is excluded; every Type~I terminal alternative is
excluded either by the inherited alternatives or by the case-closure
theorems proved below.

\begin{figure}[tbp]
\centering
\resizebox{0.72\textwidth}{!}{%
\begin{tikzpicture}[
   node distance=0.72cm and 1.20cm,
   process/.style={draw, rounded corners=2pt, align=center, text width=3.35cm,
      inner sep=3pt},
   decision/.style={draw, diamond, aspect=2.15, align=center, text width=2.95cm,
      inner sep=1.3pt},
   terminalnode/.style={draw, rounded corners=2pt, align=center, text width=3.00cm,
      inner sep=3pt},
   arrow/.style={-{Latex[length=1.9mm]}, thick},
   elab/.style={font=\scriptsize, inner sep=1pt, fill=white}
]
\node[process] (S) {Finite-energy singular point \(z_*\)};
\node[decision, below=of S] (QI) {Local pointwise Type~I bound?};
\node[terminalnode, right=1.55cm of QI] (TII)
   {Local Type~II singularity\\(not addressed here)};
\node[process, below=of QI] (WS)
   {Endpoint weak-Serrin bounds\\terminal no-escape};
\node[process, below=of WS] (SE)
   {Admissible Seregin profile\\retained local velocity concentration};
\node[decision, below=of SE] (P13) {Inherited alternatives?};
\node[process, below=of P13] (R) {Coarse residual class\\\(\calR(\calS;I,J)\)};

\node[terminalnode, right=1.55cm of P13, yshift=1.95cm] (Lsmall)
   {Small Type~I profile};
\node[terminalnode, right=1.55cm of P13, yshift=0.65cm] (LstatL3)
   {Stationary \(L^3\)\\Type~I profile};
\node[terminalnode, right=1.55cm of P13, yshift=-0.65cm] (Ltight)
   {Tight Type~I profile};
\node[terminalnode, right=1.55cm of P13, yshift=-1.95cm] (Lstruct)
   {Structured Type~I profile};

\draw[arrow] (S) -- (QI);
\draw[arrow] (QI) -- node[elab, above] {fails} (TII);
\draw[arrow] (QI) -- node[elab, left] {holds} (WS);
\draw[arrow] (WS) -- (SE);
\draw[arrow] (SE) -- (P13);
\draw[arrow] (P13) -- node[elab, above] {small amplitude} (Lsmall);
\draw[arrow] (P13) -- node[elab, above] {stationary \(L^3\)} (LstatL3);
\draw[arrow] (P13) -- node[elab, above] {uniformly \(L^3\)-tight} (Ltight);
\draw[arrow] (P13) -- node[elab, above] {structured/decay} (Lstruct);
\draw[arrow] (P13) -- node[elab, left] {none} (R);
\end{tikzpicture}}
\caption{Entry diagram.  Local pointwise Type~I enters the
admissible Seregin compactness setting through endpoint weak Serrin estimates and terminal
no-escape.  If the bound fails, the point belongs to the local Type~II
alternative, which is not addressed here.}
\label{fig:rendered-entry-tree}
\end{figure}

\begin{figure}[tbp]
\centering
\resizebox{0.98\textwidth}{!}{%
\begin{tikzpicture}[
   node distance=0.50cm and 1.08cm,
   process/.style={draw, rounded corners=2pt, align=center, text width=3.10cm,
      inner sep=2.5pt},
   decision/.style={draw, diamond, aspect=2.15, align=center, text width=2.70cm,
      inner sep=1pt},
   terminalnode/.style={draw, rounded corners=2pt, align=center, text width=2.75cm,
      inner sep=2.5pt},
   arrow/.style={-{Latex[length=1.8mm]}, thick},
   elab/.style={font=\scriptsize, inner sep=1pt, fill=white}
]
\node[process] (R) {Coarse residual class\\\(\calR(\calS;I,J)\)};
\node[decision, below=of R] (Qax) {Axisymmetric with\\bounded circulation?};
\node[terminalnode, right=1.35cm of Qax] (Lax) {\(\Cax\) residual profile};
\node[decision, below=of Qax] (Qrot) {Rotational relative\\equilibrium?};
\node[terminalnode, right=1.35cm of Qrot] (Lrot) {\(\Crot\) residual profile};
\node[decision, below=of Qrot] (Qst)
   {Degenerate stationary hull\\with retained velocity concentration?};
\node[terminalnode, right=1.35cm of Qst] (Lst) {\(\Cstat\) residual profile};
\node[decision, below=of Qst] (Qaff)
   {Affine/parasitic after\\oscillatory normalization?};
\node[terminalnode, right=1.35cm of Qaff] (Laff) {\(\Caff\) lower stratum};
\node[decision, below=of Qaff] (Qtail) {Critical-tail\\boundary defect?};
\node[decision, below=of Qtail] (Qgen) {No lower stratum remains?};
\node[terminalnode, right=1.35cm of Qgen] (Lgen)
   {\(\Cgen\) generic terminal residual};

\node[terminalnode, right=1.35cm of Qtail, yshift=1.25cm] (Llog)
   {\(\Clogdiff\) critical tail};
\node[terminalnode, right=1.35cm of Qtail] (Lyng) {\(\Cyoung\) tail};
\node[decision, right=1.35cm of Qtail, yshift=-1.25cm] (Qcoh)
   {Coherent\\logarithmic type?};
\node[terminalnode, right=1.35cm of Qcoh, yshift=1.05cm] (Lhom)
   {\(\Chomcrit\) coherent tail};
\node[terminalnode, right=1.35cm of Qcoh] (Lper)
   {\(\Clogper\) coherent tail};
\node[terminalnode, right=1.35cm of Qcoh, yshift=-1.05cm] (Lap)
   {\(\Capcrit\) coherent tail};

\draw[arrow] (R) -- (Qax);
\draw[arrow] (Qax) -- node[elab, above] {yes} (Lax);
\draw[arrow] (Qax) -- node[elab, left] {no} (Qrot);
\draw[arrow] (Qrot) -- node[elab, above] {yes} (Lrot);
\draw[arrow] (Qrot) -- node[elab, left] {no} (Qst);
\draw[arrow] (Qst) -- node[elab, above] {yes} (Lst);
\draw[arrow] (Qst) -- node[elab, left] {no} (Qaff);
\draw[arrow] (Qaff) -- node[elab, above] {yes} (Laff);
\draw[arrow] (Qaff) -- node[elab, left] {no} (Qtail);
\draw[arrow] (Qtail) -- node[elab, above] {log escape} (Llog);
\draw[arrow] (Qtail) -- node[elab, above] {Young variance} (Lyng);
\draw[arrow] (Qtail) -- node[elab, above] {coherent tail} (Qcoh);
\draw[arrow] (Qcoh) -- node[elab, above] {homogeneous} (Lhom);
\draw[arrow] (Qcoh) -- node[elab, above] {log-periodic} (Lper);
\draw[arrow] (Qcoh) -- node[elab, above] {aperiodic} (Lap);
\draw[arrow] (Qtail) -- node[elab, left] {none} (Qgen);
\draw[arrow] (Qgen) -- node[elab, above] {yes} (Lgen);
\draw[arrow] (Qgen.west) -- ++(-0.95,0) |- node[elab, near end, left] {no}
   (R.west);
\end{tikzpicture}}
\caption{Residual stratification diagram.  The final return arrow only records
the ordered nature of the decomposition: if a lower stratum has not been
removed, the profile is assigned to that terminal alternative; after all lower
alternatives are removed,
the remaining residual class is \(\Cgen\).}
\label{fig:rendered-residual-tree}
\end{figure}

The loop from the last criterion back to the residual class is only a graphical
way to indicate the ordered nature of the decomposition: if a lower stratum has
not actually been removed, it is assigned to its corresponding terminal
alternative rather than
to \(\Cgen\).  After the ordered removals are performed, \(\Cgen\) is
the remaining residual class by definition.
For the table, write
\[
   \mathfrak T_\rho(u;z_*):=
   \esssup_{T-\rho^2<t<T}
   \sqrt{T-t}\,\|u(t)\|_{L^\infty(B_\rho(x_*))}.
\]

\begingroup
\footnotesize
\setlength{\tabcolsep}{4pt}
\renewcommand{\arraystretch}{1.25}
\begin{longtable}{@{}L{1.20in}L{4.65in}@{}}
\caption{Classification and closure table for the residual decomposition diagram.  Each row gives
the criterion selecting the terminal alternative, the mechanism that excludes
it when it is a Type~I alternative, and the statement used at that point.
For the inherited alternatives, the displayed condition is a representative
form of the inherited definition; the precise definitions and closures are those
of the companion setup paper, as summarized in
\Cref{thm:imported-setup-results}.}
\label{tab:decision-tree-classification}\\
\hline
\textbf{Terminal alternative} & \textbf{Criterion, exclusion mechanism, and reference}\\
\hline
\endfirsthead
\hline
\textbf{Terminal alternative} & \textbf{Criterion, exclusion mechanism, and reference}\\
\hline
\endhead
\hline
\endfoot
Local Type~II singularity &
\textbf{Criterion.}  For every \(\rho>0\),
\(\mathfrak T_\rho(u;z_*)=\infty\).
\textbf{Role.}  This is the complementary alternative left by the local
Type~I/Type~II dichotomy after local pointwise Type~I is excluded.
\textbf{Link.}  \Cref{cor:local-typeI-unconditional}.\\

Small Type~I profile &
\textbf{Criterion.}  Inherited small-amplitude class; in the setup normalization,
\(\|V\|_{L^\infty(\R^3\times\R)}\le\eps_{\rm sm}\).
\textbf{Exclusion.}  Small-data/regularity closure recorded in the setup
paper.
\textbf{Link.}  \Cref{hyp:previous-nonresidual}.\\

Stationary \(L^3\) Type~I profile &
\textbf{Criterion.}  Inherited stationary class; representative condition
\(\partial_\tau V=0\) and \(V\in L^3(\R^3)\).
\textbf{Exclusion.}  Stationary \(L^3\) Liouville/regularity closure recorded
in the setup paper.
\textbf{Link.}  \Cref{hyp:previous-nonresidual}.\\

Tight Type~I profile &
\textbf{Criterion.}  Inherited tight class; representative condition
\(\lim_{R\to\infty}\sup_\tau\int_{|y|>R}|V(y,\tau)|^3\,dy=0\).
\textbf{Exclusion.}  Endpoint \(L^3\) rigidity for uniformly tight ancient
dynamics.
\textbf{Link.}  \Cref{hyp:previous-nonresidual}.\\

Structured Type~I profile &
\textbf{Criterion.}  Inherited closed structured/decay class,
\(V\in\mathcal C_{\rm struct/decay}^{\rm setup}\).
\textbf{Exclusion.}  Closed structured/decay alternatives recorded in the
setup paper.
\textbf{Link.}  \Cref{hyp:previous-nonresidual}.\\

\(\Cax\) residual profile &
\textbf{Criterion.}  Axisymmetric residual profile with bounded centered
circulation \(\Gamma(\rho,Z,\tau)=\rho V_\theta\) and
\(\|\Gamma\|_{L^\infty}<\infty\).
\textbf{Exclusion.}  Centered circulation equation plus axis-absorption
Liouville theorem; the profile becomes no-swirl and falls into an earlier
structured class.
\textbf{Link.}  \Cref{thm:R1-exclusion}.\\

\(\Crot\) residual profile &
\textbf{Criterion.}  Rotational relative equilibrium
\(V(y,\tau)=Q(\tau)W(Q(\tau)^Ty)\), \(Q(\tau)=e^{\tau\Omega}\).
\textbf{Reduction/closure.}  Localized Coriolis-flux reduction sends the profile
into lower classes already excluded by the setup paper or the terminal velocity-concentration
decomposition; closure is only after terminal stratification assembly.
\textbf{Link.}  \Cref{cor:R2-closure}.\\

\(\Cstat\) residual profile &
\textbf{Criterion.}  Degenerate stationary-hull retained profile: along
time translations \(T_{s_k}V\to W\), where \(W\) is stationary and
\(\Vel_1(W;0)\ge\eta_*\).
\textbf{Exclusion.}  Scale reselection reduces to a stationary retained
concentration profile, then terminal concentration excludes it.
\textbf{Link.}  \Cref{cor:R3-closure}.\\

\(\Caff\) affine/parasitic routing row &
\textbf{Criterion.}  The affine-normalized oscillation vanishes:
\(\Phi(S_cU)=0\) for an admissible affine symmetry \(S_c\).
\textbf{Routing.}  Oscillatory entry normalization assigns such profiles to the
affine/parasitic quotient row before retention.
\textbf{Link.}  \Cref{thm:oscillatory-entry-normalization}.\\

\(\Clogdiff\) critical tail &
\textbf{Criterion.}  Logarithmic escape component in the realized
critical-tail compactification: \(\Lambda\ne0\).
\textbf{Exclusion.}  Log-window descendant selection plus minimal
mesoscopic-scale reduction and hidden-scale compactness.
\textbf{Link.}  \Cref{cor:log-diffuse-branch-exclusion}.\\

\(\Cyoung\) tail &
\textbf{Criterion.}  Tight annular critical-tail profile with non-Dirac Young
measure: \(\Lambda=0\) and
\(\mathcal Y_{x,s}\ne\delta_{\bar W(x,s)}\) on a set of positive measure.
\textbf{Exclusion.}  No hidden-scale variance gives strong
\(L^3_{\loc}\) compactness, hence a Dirac Young measure.
\textbf{Link.}  \Cref{cor:young-branch-exclusion}.\\

\(\Chomcrit\) coherent tail &
\textbf{Criterion.}  Coherent \((-1)\)-homogeneous tail:
\(W(r\omega,s)=r^{-1}A(\omega,s)\).
\textbf{Exclusion.}  Bounded-origin realization gives
\(W\in L^\infty_{\loc}\) near \(r=0\), forcing \(A=0\).
\textbf{Link.}  \Cref{cor:coherent-critical-tail-branch-closure}.\\

\(\Clogper\) coherent tail &
\textbf{Criterion.}  Nonzero log-periodic coherent tail:
\(Z(\rho+L,\omega,s)=Z(\rho,\omega,s)\), \(L\ne0\), where
\(Z=e^\rho W(e^\rho\omega,s)\).
\textbf{Exclusion.}  Bounded-origin realization gives
\(Z(\rho,\cdot,\cdot)\to0\) as \(\rho\to-\infty\); periodicity forces
\(Z\equiv0\).
\textbf{Link.}  \Cref{cor:coherent-critical-tail-branch-closure}.\\

\(\Capcrit\) coherent tail &
\textbf{Criterion.}  Nonzero aperiodic minimal logarithmic hull:
\(\mathcal H(Z)\) is minimal, nonperiodic, and nonzero.
\textbf{Exclusion.}  Bounded-origin realization puts \(0\) in the compact
log-translation hull; minimality forces the hull to be \(\{0\}\).
\textbf{Link.}  \Cref{cor:coherent-critical-tail-branch-closure}.\\

\(\Cgen\) generic terminal residual &
\textbf{Criterion.}  Ordered complement
\[
   V\in\calR(\calS;I,J)
   \setminus\bigcup_{\mathcal B\in\mathfrak B_{\rm lower}}\mathcal B ,
\]
where \(\mathfrak B_{\rm lower}\) is the list of all preceding Type~I
alternatives.
\textbf{Exclusion.}  Paired concentration-set decomposition, finite
separated-family exclusion, terminal indecomposability, sequence-\(L^3\),
parasitic-free mildness, and the Albritton--Barker endpoint Liouville theorem.
\textbf{Link.}  \Cref{thm:R4-final-closure}.\\
\end{longtable}
\endgroup

\section{Shared local notation}
\label{sec:shared-notation}

For \(z_0=(y_0,\tau_0)\) and \(r>0\), write
\[
   Q_r(z_0):=B_r(y_0)\times(\tau_0-r^2,\tau_0).
\]

\begin{definition}[Scale-invariant CKN concentration and defect measures]
\label{def:local-critical-quantity}
For a centered profile \((V,P)\), define the velocity concentration
\begin{equation}\label{eq:velocity-activity-main}
   \Vel_r(V;z_0)
   :=
   r^{-2}\iint_{Q_r(z_0)} |V|^3\,dy\,d\tau,
\end{equation}
and the quotient pressure concentration
\begin{equation}\label{eq:pressure-activity-main}
   D_r(P;z_0)
   :=
   r^{-2}
   \inf_{c\in L^{3/2}(\tau_0-r^2,\tau_0)}
   \iint_{Q_r(z_0)}
      |P(y,\tau)-c(\tau)|^{3/2}\,dy\,d\tau .
\end{equation}
The scale-invariant local CKN quantity is
\begin{equation}\label{eq:local-energy-main}
   \calE_r(V,P;z_0):=\Vel_r(V;z_0)+D_r(P;z_0).
\end{equation}
Thus \(\calE_r\) is parabolically scale-invariant under the Navier--Stokes
scaling.

If an explicit pressure atlas \(\mathfrak a\) has been chosen, with local
representative \(a_{z_0,r}(\tau)\), we also write
\begin{equation}\label{eq:atlas-dependent-activity-main}
   \calE_r^{\mathfrak a}(V,P;z_0)
   :=
   r^{-2}
   \iint_{Q_r(z_0)}
      \left(|V|^3+|P-a_{z_0,r}(\tau)|^{3/2}\right)\,dy\,d\tau .
\end{equation}
Then
\[
   \calE_r(V,P;z_0)\le \calE_r^{\mathfrak a}(V,P;z_0).
\]
Thresholds and regularity statements use the quotient quantity \(\calE_r\);
compactness statements use a fixed setup pressure atlas when an explicit
representative is needed.

For a terminal sequence \((V_n,P_n)\) with chosen pressure gauges \(a_n(\tau)\),
the atlas-dependent Radon defect measure is
\begin{equation}\label{eq:raw-defect-measure-main}
   d\mu_n
   :=
   \left(|V_n|^3+|P_n-a_n(\tau)|^{3/2}\right)\,dy\,d\tau .
\end{equation}
The measure \(\mu_n\) is not denoted by \(\calE_r\).  On unit cylinders the
factor \(1^{-2}\) is invisible, so unit-scale atlas mass and
unit-scale CKN concentration are compatible, with the atlas quantity dominating the
quotient quantity.
The velocity quantity \(\Vel_r\) is the part used to certify nonzero retained
limits under weak pressure convergence; the pressure-inclusive quantity
\(\calE_r\) is used for CKN regularity and pressure-normalized smallness.
\end{definition}

\begin{remark}[Two-scale dictionary]
\label{rem:two-scale-dictionary}
We use two different normalizations.  The quantity \(\calE_r\) is parabolic and
CKN-critical, hence carries the factor \(r^{-2}\).  By contrast, the tail
quantities \(\mathcal A_n(z,r)\) are spatial/logarithmic averaged oscillations
on fixed centered time windows; their natural normalization is \(r^{-3}\).
These two quantities are never identified.
\end{remark}

\begin{definition}[Velocity retention at a level and unit CKN activity]
\label{def:velocity-retention-unit-activity}
Fix \(0<\eta_*,\eps_*<\infty\).  For \(\eta>0\), a terminal profile
\((U,\Pi)\) has retained local velocity concentration at level \(\eta\) if,
after the terminal normalization under consideration,
\[
   \Vel_1(U;0)\ge \eta .
\]
It has positive retained local velocity concentration if this holds for some
\(\eta>0\).  When no level is displayed, the setup-level retained descendant
threshold is \(\eta_*\).
It has unit CKN activity if
\[
   \calE_1(U,\Pi;0)\ge \eps_* .
\]
Only retained velocity concentration is used to pass nontriviality to limits
from approximating sequences.  Pressure-inclusive CKN activity is used for
local regularity, pressure compactness, and threshold smallness, and any
pressure-only lower bound is treated through the pressure-gauge alternatives
below.
\end{definition}

\begin{definition}[Threshold hierarchy]
\label{def:threshold-hierarchy}
The constants used in the terminal concentration argument are chosen in the following order.
The compact velocity lower bounds \(\eta_0\) and \(\eta_v\) are supplied by
the imported setup results and the local concentration input of the companion
setup paper, respectively.  The local small-data threshold
\(\eps_{\rm sd}\) is then fixed with
\[
   0<\eps_{\rm sd}\le \min\{\eta_0,\eta_v,\eps_{\CKN}\},
\]
where \(\eps_{\CKN}\) denotes the universal CKN regularity threshold whenever
that threshold is invoked.  The retained descendant velocity threshold
\(\eta_*\) is chosen after the pressure gauges and compactness topology are fixed, with
\[
   0<\eta_*\le \frac12\eps_{\rm sd}.
\]
The pressure-inclusive unit CKN activity threshold \(\eps_*\) is then
chosen with
\[
   0<\eps_*\le \frac12\eps_{\rm sd}.
\]
Whenever a compactness or lower-semicontinuity passage loses a harmless fixed
factor, the relevant threshold, \(\eta_*\) for velocity retention or
\(\eps_*\) for pressure-inclusive CKN activity, is replaced once and for all
by a smaller positive value.  No later argument chooses a new concentration
threshold depending on a previously extracted iterated profile extraction.
\end{definition}

\begin{definition}[Refined residual candidate]
\label{def:refined-residual-candidate}
A refined residual candidate is a bounded normalized Seregin ancient limit
obtained from the admissible Seregin extraction theorem in the companion setup
paper, which carries the positive compact local velocity
concentration in \eqref{eq:positive-velocity-main}, and which is not in any alternative
already closed before the refined residual decomposition is entered.  The
definition is local: the only persistent lower bound attached to the candidate by
the setup paper is the fixed compact velocity \(L^3\) lower bound.
\end{definition}

\begin{remark}[No hidden global \(L^3\) assertion]
\label{rem:no-hidden-global-L3}
The terminal arguments below do not claim that a generic residual profile is
globally \(L^3\)-tight or globally bounded in \(L^\infty_\tau L^3_y\).  The
only nonlocal critical-norm conclusion used at the endpoint is the weaker
sequence-\(L^3\) completeness of retained indecomposable concentration
profiles, and that conclusion is derived from local terminal alternatives rather than from an
a priori global estimate on the original residual class.
\end{remark}

\subsection{Notation, threshold, and gauge directories}
\label{subsec:notation-directories}

The following compact directories collect the persistent symbols, thresholds,
and gauges used in the residual closure.  They are not new definitions; each
entry refers to the location where the symbol is first defined.

\paragraph{Velocity, pressure, and profile-space symbols.}
\begin{itemize}[leftmargin=1.6em]
   \item \((V,P)\): centered Seregin ancient profile in \eqref{eq:centered-main};
      first used in \Cref{sec:introduction}.
   \item \((U,Q)\): physical ancient representative when \((V,P)\) is treated
      under physical pullback \eqref{eq:physical-pullback-main}; also used in
      \(\calR_z\)-recentered form throughout
      \Cref{sec:R4-tail-geometry,sec:terminal-concentration}.  The intended
      meaning is fixed locally in each section by the explicit definition of
      the immediately preceding rescaling.
   \item \((W,Q_W)\): blown-down, co-rotating, or tail profile, depending on
      the section; defined locally in
      \Cref{sec:R2-reduction,subsec:hidden-scale-young-variance}.
   \item \(\calE_r(V,P;z_0)\): quotient, scale-invariant parabolic CKN
      concentration, see \Cref{def:local-critical-quantity}.
   \item \(\Vel_r(V;z_0)\): velocity part of the scale-invariant local
      concentration; retained descendants and nonzero limiting profiles are
      certified with this quantity.
   \item \(\calE_1\): pressure-normalized unit CKN concentration, fixed at
      \(r=1\); used for local regularity, CKN smallness, and pressure estimates
      once the velocity lower bound has been identified.
   \item \(\mu_n\): atlas-dependent Radon defect measure in
      \eqref{eq:raw-defect-measure-main}; used for concentration-set and defect
      compactness, not as a scale-invariant quantity.
   \item \(\mathcal A_n(z,r)\): spatial/logarithmic averaged affine-quotient
      oscillation with \(r^{-3}\) normalization, distinct from \(\calE_r\).
   \item \(\mathfrak X_M\): bounded terminal profile class with pressure gauges,
      \Cref{def:terminal-state-space}.
   \item \(\mathfrak X_M^\#\): compact normalized descendant profile class,
      \Cref{def:compact-normalized-state-space}.
   \item \(\Tcrit\): critical-tail compactification of Young/log profiles together
      with pressure-defect, viscous-defect, affine-quotient, virial, and support
      coordinates; abbreviated as \((\bar W,\mathcal Y,\Lambda)\) only when the
      extra coordinates are zero or irrelevant, \Cref{def:critical-diffuse-tail-strata}.
   \item \(\Caff\): affine/parasitic reduction/quotient stratum,
      \Cref{def:affine-constant-lower-stratum}; not a usual nonempty/empty PDE
      class but an exclusion label, see
      \Cref{rem:affine-constant-use,thm:oscillatory-entry-normalization}.
   \item \(\oscsp\): spatial oscillation modulo time-dependent homogeneous
      gauges, \Cref{def:multiscale-affine-oscillation}.
   \item \(\actthree\) and \(\Phi\): constant-mode/non-affine retained
      oscillation, \Cref{def:observer-oscillation,thm:closed-nonaffine-activity}.
\end{itemize}

\paragraph{Threshold ordering.}
\begin{itemize}[leftmargin=1.6em]
   \item \(\eta_0\): compact velocity \(L^3\) lower bound supplied by the base
      Seregin extraction \Cref{hyp:base-seregin-hypotheses}.
   \item \(\eta_v\): compact velocity lower bound from the local concentration
      input of the setup paper, used as initial mass for terminal extraction;
      see \Cref{def:threshold-hierarchy}.
   \item \(\eps_{\rm sd}\): local small-data threshold,
      \(0<\eps_{\rm sd}\le\min\{\eta_0,\eta_v,\eps_{\CKN}\}\),
      \Cref{def:threshold-hierarchy}.
   \item \(\eta_*\): retained descendant velocity threshold,
      \(0<\eta_*\le\eps_{\rm sd}/2\), used in
      \Cref{def:velocity-retention-unit-activity} and in concentration-set extraction.
   \item \(\eps_*\): unit CKN activity threshold,
      \(0<\eps_*\le\eps_{\rm sd}/2\), used only for pressure-inclusive
      CKN activity and smallness tracking.
   \item \(\eps_0\): minimal mesoscopic-scale oscillation threshold; depends
      only on the hidden-scale defect lower bound and on the finite-overlap
      constants of \Cref{lem:finite-overlap-affine-oscillation}; first appears
      in \Cref{def:first-bad-mesoscopic-scale}.
   \item \(\eta_K\): compact oscillation lower bound on the recurrent concentration core,
      first introduced in \Cref{thm:active-path-space-recurrence}.
\end{itemize}
The thresholds are fixed in the order
\(\eta_0,\eta_v\Rightarrow\eps_{\rm sd}\Rightarrow(\eta_*,\eps_*)\Rightarrow\eps_0,
\eta_K\), and no later argument chooses a new concentration threshold depending
on a previously extracted iterated profile extraction.  Phrases of the form ``after lowering
the threshold'' refer either to the once-and-for-all reduction of \(\eta_*\)
or \(\eps_*\) recorded in \Cref{def:threshold-hierarchy}, or to a slight
reduction of the mesoscopic threshold \(\eps_0\) recorded in
\Cref{lem:h2-failure-first-bad-scale}; in either case the reduction does not
disturb earlier choices.

\paragraph{Pressure and velocity gauges.}
\begin{itemize}[leftmargin=1.6em]
   \item \(c(\tau)\) or \(a_{z_0,r}(\tau)\): local pressure time gauge in
      \eqref{eq:pressure-activity-main} and
      \eqref{eq:atlas-dependent-activity-main};
      measurable in \(\tau\); used because pressure is defined modulo functions
      of time.
   \item \(c_n(s)\): homogeneous velocity gauge at the minimal mesoscopic
      scale, \Cref{def:first-bad-spatial-blowdown}; bounded; either compactifies
      or assigns the descendant to \(\Caff\) by
      \Cref{lem:homogeneous-gauge-compactness}; no pointwise derivative
      \(\partial_s c_n\) is used unless that compactness has already been
      obtained.
   \item \(\beta(s)\cdot y+\alpha(s)\): affine parasitic pressure paired with a
      spatially homogeneous velocity mode; quotient notation introduced in
      \Cref{def:affine-constant-lower-stratum} and used only in the gauged
      formulation of the affine/parasitic stratum.
   \item \(L_n(x,s)=b_n(s)\cdot x+a_n(s)\): scaled affine pressure at the
      minimal mesoscopic scale, \Cref{lem:scaled-pressure-compactness-affine};
      affine in \(x\), measurable in \(s\); discarded in the non-affine quotient
      and assigned to \(\Caff\) when it does not vanish.
   \item \(a_{n,K}(t)\): local Seregin pressure time gauge on a compact
      cylinder \(K\), used by the local pressure-gauge compactness theorem
      inherited from the setup paper and recorded in
      \Cref{prop:local-typeI-terminal-compactness}.
\end{itemize}

\begin{longtable}{@{}L{0.28\textwidth}L{0.66\textwidth}@{}}
\caption{Glossary of residual symbols.}
\label{tab:paperIV-glossary}\\
\hline
\textbf{Symbol} & \textbf{Meaning}\\
\hline
\endfirsthead
\hline
\textbf{Symbol} & \textbf{Meaning}\\
\hline
\endhead
\hline
\endfoot
\(\calE_r\) & Scale-invariant parabolic CKN concentration.\\
\(D_r\) & Pressure quotient part of the CKN concentration, modulo time gauges.\\
\(\Vel_r\) & Scale-invariant velocity concentration used for retained
nonzero descendants.\\
\(\mu_n\) & Raw atlas defect measure for concentration-set and compactness
arguments.\\
\(\mathcal A_n(z,r)\) & Spatial/log-tail averaged affine-quotient oscillation.\\
\(\calH_\infty\) & Full covariant spacetime asymptotic hull.\\
\(\Cax\) & Axisymmetric bounded-circulation residual class.\\
\(\Crot\) & Rotational relative-equilibrium residual class.\\
\(\Cstat\) & Stationary-hull residual class with retained local velocity concentration.\\
\(\Caff\) & Affine/parasitic pressure residual quotient.\\
\(\Clogdiff\) & Log-diffuse critical-tail residual alternative.\\
\(\Cyoung\) & Young/defect critical-tail residual alternative, including
non-Dirac Young measures and nonzero retained critical-tail pressure/viscous
defect coordinates.\\
\(\Chomcrit\) & Coherent homogeneous critical-tail alternative.\\
\(\Clogper\) & Coherent log-periodic critical-tail alternative.\\
\(\Capcrit\) & Coherent aperiodic/minimal critical-tail alternative.\\
\(\Cgen\) & Generic terminal residual class.\\
Retained descendant & Descendant with inherited positive unit velocity concentration.\\
Normalized support profile / support profile & Normalized defect
support profile, not necessarily retained at the original threshold.\\
First-generation recentering & Recentered window of the original setup terminal
sequence; compact only under the local-window condition.\\
Ancient-profile recentering & Covariant recentering of an already extracted
bounded ancient profile; compact by global ancient boundedness.\\
Local-window condition & Requirement that every fixed recentered cylinder
undoes into the original local Type~I setup cylinder.\\
Exterior/diffuse terminal branch & Positive residual mass outside active
neighborhoods, below unit velocity thresholds after the reduction step.\\
Pressure/noncompact defect & Failure of local Seregin compactness or pressure
gauge compactness under a selected recentering.\\
Critical-tail defect & Nonzero normalized annular/logarithmic residual profile
visible only after critical-tail blow-down.\\
\hline
\end{longtable}

\section{Exclusion of the axisymmetric bounded-circulation residual class}
\label{sec:R1-exclusion}

In this section we prove the rigidity statement needed for the first refined
residual class.  The argument uses the scale-invariant angular circulation in
centered variables.  The essential point is that the centered Type~I equation
contains the Ornstein--Uhlenbeck drift.  Thus the angular-circulation equation
has an absorbing boundary at the symmetry axis together with a confining drift
at spatial infinity.  This gives a scalar Liouville theorem for bounded ancient
solutions of the circulation equation.  Related axisymmetric Liouville theorems
for Navier--Stokes, in different settings, are proved in
\cite{LZZ2017,LRZ2022}.

Throughout the section we write cylindrical coordinates as
\((\rho,\theta,Z)\), and an axisymmetric vector field as
\[
   V=V_\rho e_\rho+V_\theta e_\theta+V_Z e_Z .
\]
The centered Navier--Stokes equation is
\begin{equation}\label{eq:centered-R1}
   \partial_\tau V-\Delta V+\frac12 y\cdot\nabla V+\frac12 V
   +(V\cdot\nabla)V+\nabla P=0,
   \qquad \nabla\cdot V=0 .
\end{equation}
For an axisymmetric solution define the centered angular circulation
\begin{equation}\label{eq:centered-Gamma-def}
   \Gamma(\rho,Z,\tau):=\rho V_\theta(\rho,Z,\tau).
\end{equation}
This is the same scale-invariant quantity as the physical circulation
\(rU_\theta\).  If the physical representative is
\[
   U(x,t)=(-t)^{-1/2}V(x/\sqrt{-t},-\log(-t)),
\]
then
\[
   rU_\theta(r,z,t)=\rho V_\theta(\rho,Z,\tau),
   \qquad r=\sqrt{-t}\,\rho,
   \quad z=\sqrt{-t}\,Z.
\]

\begin{lemma}[Centered circulation equation]\label{lem:centered-circulation-equation}
Let \((V,P)\) be a smooth bounded axisymmetric solution of
\eqref{eq:centered-R1}.  Then \(\Gamma=\rho V_\theta\) satisfies
\begin{equation}\label{eq:centered-Gamma-equation}
   \partial_\tau\Gamma
   =
   \partial_\rho^2\Gamma-\frac1\rho\partial_\rho\Gamma
      +\partial_Z^2\Gamma
   -\left(V_\rho+\frac\rho2\right)\partial_\rho\Gamma
   -\left(V_Z+\frac Z2\right)\partial_Z\Gamma
\end{equation}
on \((0,\infty)\times\mathbb R\times\mathbb R\).  Moreover
\begin{equation}\label{eq:axis-boundary-Gamma}
   \Gamma(0,Z,\tau)=0
   \qquad\hbox{for all }(Z,\tau)\in\mathbb R\times\mathbb R .
\end{equation}
\end{lemma}

\begin{proof}
The angular component of \eqref{eq:centered-R1} is
\[
   \partial_\tau V_\theta
   -\left(\Delta-\frac1{\rho^2}\right)V_\theta
   +\frac12(\rho\partial_\rho+Z\partial_Z)V_\theta
   +\frac12V_\theta
   +V_\rho\partial_\rho V_\theta+V_Z\partial_ZV_\theta
   +\frac{V_\rho V_\theta}{\rho}=0 .
\]
Multiplying by \(\rho\), using
\[
   \rho\left(\Delta-\frac1{\rho^2}\right)V_\theta
   =\partial_\rho^2\Gamma-\frac1\rho\partial_\rho\Gamma+
      \partial_Z^2\Gamma,
\]
\[
   \rho\left[\frac12(\rho\partial_\rho+Z\partial_Z)V_\theta
      +\frac12V_\theta\right]
   =\frac12(\rho\partial_\rho+Z\partial_Z)\Gamma,
\]
and
\[
   \rho\left(V_\rho\partial_\rho V_\theta+V_Z\partial_ZV_\theta
      +\frac{V_\rho V_\theta}{\rho}\right)
   =V_\rho\partial_\rho\Gamma+V_Z\partial_Z\Gamma,
\]
gives \eqref{eq:centered-Gamma-equation}.  Finally, smoothness of a bounded
axisymmetric vector field at the axis implies \(V_\theta(\rho,Z,\tau)=O(\rho)\)
as \(\rho\downarrow0\), locally uniformly in \((Z,\tau)\).  Hence
\(\Gamma=\rho V_\theta=O(\rho^2)\), which gives
\eqref{eq:axis-boundary-Gamma}.
\end{proof}

We next record the scalar Liouville theorem used below.  It is stated in a
form that isolates precisely the PDE mechanism: bounded drift perturbations of
the centered circulation operator cannot sustain a bounded ancient solution
which vanishes on the symmetry axis.

\begin{lemma}[One-dimensional killed radial absorption]
\label{lem:one-dimensional-killed-radial-absorption}
Fix \(R>1\), \(M<\infty\), and a compact interval
\(K_\rho=[\delta,R-\delta]\Subset(0,R)\).  Let \(\rho_t\) be any weak solution
on \((0,R)\), killed at \(\{0,R\}\), of
\[
   d\rho_t
   =
   -\left(\frac1{\rho_t}+\frac{\rho_t}{2}+\beta_t\right)\,dt
   +\sqrt2\,dW_t,
   \qquad |\beta_t|\le M,
\]
where \(\beta_t\) may be adapted.  If
\(\tau_\partial:=\inf\{t:\rho_t\in\{0,R\}\}\), then
\[
   \lim_{T\to\infty}
   \sup_{\rho_0\in K_\rho}
   \mathbb P_{\rho_0}(\tau_\partial>T)=0 .
\]
More precisely, there are \(T_0=T_0(R,M,\delta)>0\) and
\(\theta=\theta(R,M,\delta)>0\) such that
\[
   \inf_{\rho_0\in K_\rho}
   \mathbb P_{\rho_0}(\tau_\partial\le T_0)\ge\theta .
\]
\end{lemma}

\begin{proof}
Choose
\[
   \delta_0:=\min\{\delta/2,\,[4(M+2)]^{-1}\}.
\]
We first show that every starting point in \(K_\rho\) reaches the left boundary
layer \((0,\delta_0]\) in a fixed time with a probability bounded from below.
While \(\rho_t\ge\delta_0\), the drift satisfies
\[
   -\left(\frac1{\rho_t}+\frac{\rho_t}{2}+\beta_t\right)\le M .
\]
Fix \(T_1>0\) and the deterministic linear path
\[
   \ell(t):=-\frac{R+MT_1+2\delta_0}{T_1}\,t,\qquad 0\le t\le T_1 .
\]
On the Brownian tube event
\[
   E_1:=
   \left\{\sup_{0\le t\le T_1}
      \left|\sqrt2\,W_t-\ell(t)\right|\le\frac{\delta_0}{4}\right\},
\]
the integral form of the SDE gives, for every \(\rho_0\in K_\rho\), that
\(\rho_t\) reaches \((0,\delta_0]\) before time \(T_1\), unless it has already
hit \(0\).  The probability of \(E_1\) is a strictly positive number
\[
   p_1=p_1(R,M,\delta)>0 .
\]
For completeness, this positivity follows by partitioning \([0,T_1]\) into
finitely many subintervals and using the one-dimensional Dirichlet heat kernel
lower bound in the strip of width \(\delta_0/4\) around the line \(\ell\).

It remains to absorb the process from the layer \((0,\delta_0]\).  As long as
the process stays in \((0,2\delta_0]\),
\[
   -\left(\frac1{\rho_t}+\frac{\rho_t}{2}+\beta_t\right)
   \le M-\frac1{2\delta_0}\le -2 .
\]
Let \(T_2:=\delta_0\).  On the event
\[
   E_2:=
   \left\{\sup_{0\le t\le T_2}
      \left(\sqrt2\,W_t-\frac t4\right)\le\frac{\delta_0}{4}\right\},
\]
which has probability \(p_2=p_2(M,\delta)>0\), the same integral estimate gives
\[
   \rho_t\le \rho_0-2t+\frac{\delta_0}{4}+\frac t4
   \qquad\text{until hitting }0 .
\]
Thus every initial point \(\rho_0\le\delta_0\) hits \(0\) before
\(T_2\), on \(E_2\).  Again \(p_2>0\) is an explicit finite-time heat-kernel
barrier estimate for Brownian motion below the affine boundary
\(\delta_0/4+t/4\).

Combining the two blocks and using the strong Markov property gives
\[
   \inf_{\rho_0\in K_\rho}
   \mathbb P_{\rho_0}(\tau_\partial\le T_1+T_2)
   \ge p_1p_2=:\theta>0 .
\]
With \(T_0:=T_1+T_2\), iteration gives
\[
   \sup_{\rho_0\in K_\rho}
   \mathbb P_{\rho_0}(\tau_\partial>kT_0)
   \le (1-\theta)^k ,
\]
and the asserted decay follows.  The constants are independent of the axial
coordinate and of any truncation in that coordinate.
\end{proof}

\begin{lemma}[Axis-absorption Liouville theorem]
\label{lem:axis-absorption-liouville}
Let \(B=(B_\rho,B_Z)\) be a smooth vector field on
\((0,\infty)\times\mathbb R\times\mathbb R\) satisfying
\[
   \|B\|_{L^\infty}<\infty .
\]
Let \(G\in C^{2,1}((0,\infty)\times\mathbb R\times\mathbb R)\cap L^\infty\)
solve
\begin{equation}\label{eq:abstract-Gamma-equation}
   \partial_\tau G
   =
   \partial_\rho^2G-\frac1\rho\partial_\rho G+
      \partial_Z^2G
   -\left(B_\rho+\frac\rho2\right)\partial_\rho G
   -\left(B_Z+\frac Z2\right)\partial_ZG
\end{equation}
for all \(\tau\in\mathbb R\), and assume
\begin{equation}\label{eq:abstract-axis-zero}
   \lim_{\rho\downarrow0}G(\rho,Z,\tau)=0
   \qquad\hbox{locally uniformly in }(Z,\tau).
\end{equation}
Then \(G\equiv0\).
\end{lemma}

\begin{proof}
Set \(M_B=\|B\|_{L^\infty}\).  For \(R>1\) define
\begin{equation}\label{eq:radial-barrier}
   h_R(\rho)
   :=
   \frac{\displaystyle\int_0^{\min\{\rho,R\}} s
      \exp\left(\frac{s^2}{4}-M_Bs\right)\,ds}
   {\displaystyle\int_0^R s
      \exp\left(\frac{s^2}{4}-M_Bs\right)\,ds} .
\end{equation}
Then \(0\le h_R\le1\), \(h_R(0)=0\), \(h_R(R)=1\), and, on
\(0<\rho<R\),
\begin{equation}\label{eq:hR-ode}
   h_R''+\left(M_B-\frac\rho2-\frac1\rho\right)h_R'=0,
   \qquad h_R'>0 .
\end{equation}
Let
\[
   \mathcal L_\tau f
   :=
   \partial_\rho^2f-\frac1\rho\partial_\rho f+\partial_Z^2f
   -\left(B_\rho+\frac\rho2\right)\partial_\rho f
   -\left(B_Z+\frac Z2\right)\partial_Z f .
\]
Since \(B_\rho\ge -M_B\), \eqref{eq:hR-ode} gives
\begin{equation}\label{eq:hR-superharmonic}
   \mathcal L_\tau h_R
   =h_R''-\left(\frac1\rho+\frac\rho2+B_\rho\right)h_R'
   \le
   h_R''+\left(M_B-\frac\rho2-\frac1\rho\right)h_R'
   =0 .
\end{equation}
Thus \((\partial_\tau-\mathcal L_\tau)h_R\ge0\).

We now compare \(G\) to this barrier on the strip
\(\Omega_R=(0,R)\times\mathbb R\).  Let
\(M_G=\|G\|_{L^\infty}\).  We first record the Dirichlet strip decay estimate used
by the comparison argument.  If \(q_{s,R}\) solves
\[
   \partial_\tau q_{s,R}=\mathcal L_\tau q_{s,R},
   \qquad q_{s,R}=0\hbox{ on }\rho=0,R,
   \qquad q_{s,R}(\rho,Z,s)=1,
\]
then, for every compact \(K\Subset\Omega_R\) and every fixed \(t\),
\begin{equation}\label{eq:killed-strip-decay}
   \lim_{s\to-\infty}\sup_{(\rho,Z)\in K}q_{s,R}(\rho,Z,t)=0 .
\end{equation}
To prove \eqref{eq:killed-strip-decay}, fix \(K\Subset\Omega_R\), and choose
\(\delta>0\) such that
\[
   K\subset[\delta,R-\delta]\times\mathbb R .
\]
We use a one-dimensional absorption estimate for the \(\rho\)-component.  Let
\(\rho_\tau\) be any weak solution of
\[
   d\rho_\tau
   =
   -\left(\frac1{\rho_\tau}+\frac{\rho_\tau}{2}
      +\beta_\tau\right)\,d\tau+\sqrt2\,dW_\tau,
   \qquad |\beta_\tau|\le M_B,
\]
killed when it hits \(0\) or \(R\).  The coefficient \(\beta_\tau\) may be any
adapted process; in the application it is
\(B_\rho(\rho_\tau,Z_\tau,\tau)\).  By
\Cref{lem:one-dimensional-killed-radial-absorption}, there are
\(T=T(R,M_B,\delta)>0\) and \(\theta=\theta(R,M_B,\delta)>0\) such that
\[
   \inf_{\rho_0\in[\delta,R-\delta]}
   \mathbb P_{\rho_0}\{\rho_\tau\hbox{ hits }0\hbox{ or }R
      \hbox{ before time }T\}\ge\theta .
\]

The parabolic maximum principle and the Feynman--Kac representation for the
killed equation give
\[
   0\le q_{s,R}(\rho,Z,t)
   \le
   \mathbb P_{\rho}\{\rho_\tau\hbox{ survives on }[s,t]\}.
\]
Iterating the block estimate yields
\[
   \sup_{(\rho,Z)\in K}q_{s,R}(\rho,Z,t)
   \le (1-\theta)^{\lfloor (t-s)/T\rfloor},
\]
which tends to \(0\) as \(s\to-\infty\).  This proves
\eqref{eq:killed-strip-decay} directly on the infinite strip.

\smallskip\noindent\emph{Truncated Feynman--Kac justification.}
The displayed killed-kernel and survival-probability statements are first
established on the bounded box
\[
   \Omega_{\varepsilon,R,L}
   :=(\varepsilon,R)\times(-L,L)
   \qquad (0<\varepsilon<R<\infty,\ L>0)
\]
with Dirichlet boundary conditions on the entire boundary of
\(\Omega_{\varepsilon,R,L}\), where the standard killed parabolic theory
provides a Feynman--Kac representation of the strip kernel
\(q^{(\varepsilon,R,L)}_{s}(\rho,Z,t)\) by the killed diffusion
\((\rho_\tau,Z_\tau)\) with drift
\((-(\rho^{-1}+\rho/2+\beta_\tau),-(Z/2+\gamma_\tau))\) and additive Brownian
noise.  The block survival estimate of
\Cref{lem:one-dimensional-killed-radial-absorption} applies because the
one-dimensional radial absorption in \(\rho\) does not depend on the
\(Z\)-component or on \(\varepsilon\) or \(L\) once
\(\delta\le\rho\le R-\delta\); hence
\[
   \sup_{(\rho,Z)\in K}q^{(\varepsilon,R,L)}_{s}(\rho,Z,t)
   \le (1-\theta)^{\lfloor (t-s)/T\rfloor}
\]
holds uniformly in \(\varepsilon\) and \(L\) for any compact
\(K\Subset (\delta,R-\delta)\times[-L_0,L_0]\) and all
\(\varepsilon<\delta\), \(L>L_0\).  The truncated comparison
\((W_R^\pm)_+\le q^{(\varepsilon,R,L)}_s\) on
\(\Omega_{\varepsilon,R,L}\) then follows from the parabolic maximum
principle applied to the truncated problem (the \(\rho=\varepsilon\)
boundary contributes nonpositively because both \(G/M_G\) and \(h_R\) are
bounded by \(1\), so the truncated boundary corrector
\(\rho_{\rm tr}(\tau):=2\sup\{|G/M_G|+h_R\}\cdot
   \mathbf 1_{\rho=\varepsilon\,\mathrm{or}\,|Z|=L}\)
is killed by the same Dirichlet kernel and tends to zero on every fixed
compact \(K\) as \(L\to\infty\), \(\varepsilon\downarrow0\)).  Sending
\(L\to\infty\) recovers the half-strip
\((\varepsilon,R)\times\mathbb R\) by monotone convergence of the killed
semigroup; sending \(\varepsilon\downarrow0\) recovers
\((0,R)\times\mathbb R\) using \(h_R(0)=0\) and the axis condition
\eqref{eq:abstract-axis-zero} on \(G\); finally sending \(R\to\infty\)
gives the global statement \eqref{eq:G-bounded-by-hR} and the conclusion
\(G\equiv0\) as in the next paragraph.

Fix \(R>1\), \(s<t\), and set
\[
   W^+_{R}:=\frac{G}{M_G}-h_R,
   \qquad
   W^-_{R}:=-\frac{G}{M_G}-h_R,
\]
with the convention that the conclusion is immediate if \(M_G=0\).
By \eqref{eq:abstract-Gamma-equation} and \eqref{eq:hR-superharmonic},
\[
   (\partial_\tau-\mathcal L_\tau)W_R^\pm\le0
   \qquad\hbox{in }\Omega_R\times(s,t).
\]
On \(\rho=0\), \eqref{eq:abstract-axis-zero} and \(h_R(0)=0\) give
\(W_R^\pm\le0\).  On \(\rho=R\), \(h_R(R)=1\) and \(|G|\le M_G\) give again
\(W_R^\pm\le0\).  The parabolic comparison principle with the Dirichlet strip
kernel therefore yields, for \(\tau=t\),
\[
   (W_R^\pm)_+(\rho,Z,t)
   \le q_{s,R}(\rho,Z,t)
   \qquad ((\rho,Z)\in\Omega_R).
\]
Letting \(s\to-\infty\) and using \eqref{eq:killed-strip-decay} gives
\(W_R^\pm\le0\) on compact subsets of \(\Omega_R\).  Since the compact subset
is arbitrary,
\begin{equation}\label{eq:G-bounded-by-hR}
   |G(\rho,Z,t)|\le M_G h_R(\rho)
   \qquad (0<\rho<R, Z\in\mathbb R,
   \ t\in\mathbb R).
\end{equation}
Finally, for each fixed \(\rho<\infty\), the denominator in
\eqref{eq:radial-barrier} tends to infinity super-exponentially as
\(R\to\infty\), while the numerator remains fixed.  Hence
\(h_R(\rho)\to0\).  Letting \(R\to\infty\) in
\eqref{eq:G-bounded-by-hR} gives \(G(\rho,Z,t)=0\) for every
\(\rho>0\), \(Z\in\mathbb R\), and \(t\in\mathbb R\).  The boundary value at
\(\rho=0\) is already zero by assumption.  Therefore \(G\equiv0\).
\end{proof}

\begin{theorem}[Exclusion of the axisymmetric bounded-circulation residual class]
\label{thm:R1-exclusion}
Let \(\Cax(\mathcal S;I,J)\) be the axisymmetric bounded-circulation residual
class in the refined residual decomposition.  Then
\begin{equation}\label{eq:R1-empty}
   \Cax(\mathcal S;I,J)=\varnothing .
\end{equation}
\end{theorem}

\begin{proof}
Assume, for contradiction, that
\(V\in\Cax(\mathcal S;I,J)\).  By definition of \(\Cax\), the
associated physical ancient representative is axisymmetric and has bounded
angular circulation.  Equivalently, in centered variables,
\[
   \Gamma(\rho,Z,\tau)=\rho V_\theta(\rho,Z,\tau)
\]
is bounded on \((0,\infty)\times\mathbb R\times\mathbb R\).  Since \(V\) is a
Seregin ancient limit, it is smooth and bounded in centered variables.  Hence
\(B=(V_\rho,V_Z)\) is bounded.

By \Cref{lem:centered-circulation-equation}, \(\Gamma\) satisfies
\eqref{eq:centered-Gamma-equation} and vanishes at the axis.  Applying
\Cref{lem:axis-absorption-liouville} with \(G=\Gamma\) and
\(B=(V_\rho,V_Z)\) gives
\[
   \Gamma\equiv0 .
\]
Therefore \(V_\theta\equiv0\) for \(\rho>0\), and by smoothness also on the
axis.  Thus the associated ancient solution is axisymmetric without swirl.
But the axisymmetric no-swirl class is one of the structured alternatives
removed before the residual class is defined.  This contradicts
\(V\in\mathcal R(\mathcal S;I,J)\), and therefore no such \(V\) exists.
\end{proof}

\begin{remark}[Why this does not assert the full bounded-swirl Liouville theorem]
The proof uses the centered Type~I equation, specifically the
Ornstein--Uhlenbeck drift terms \(-\rho\partial_\rho/2\) and
\(-Z\partial_Z/2\) in the scalar equation for \(\Gamma\).  It is therefore a
Type~I Seregin limit argument.  It does not prove that every bounded ancient
axisymmetric solution of the unrenormalized Navier--Stokes equations with
bounded physical circulation is swirl-free.  The latter lacks the confining
centered drift and remains a different shell-transport problem.
\end{remark}


\section{The rotational relative-equilibrium residual class}
\label{sec:R2-reduction}

We next isolate the second refined residual class.  This class consists of
centered Type~I ancient limits whose time dependence is a rigid rotation of a
fixed spatial profile.  Passing to the co-rotating frame converts the centered
Navier--Stokes equation into a stationary Leray profile equation with a
Coriolis transport term.  A whole-space Coriolis--Leray Liouville theorem for
all bounded profiles is not used here.  Instead, the Coriolis term is localized as an annular
flux and is then treated by the local concentration analysis proved in
\Cref{sec:terminal-concentration,sec:separated-profiles-atomic}.

Throughout this section the centered equation is
\begin{equation}\label{eq:centered-R2-safe}
   \partial_\tau V-\Delta V+\frac12 y\cdot\nabla V+\frac12 V
   +(V\cdot\nabla)V+\nabla P=0,
   \qquad \nabla\cdot V=0 .
\end{equation}
For \(\Omega\in\mathfrak{so}(3)\), set
\begin{equation}\label{eq:rotation-group-R2-safe}
   Q(\tau):=\exp(\tau\Omega)\in SO(3).
\end{equation}

\begin{definition}[Rotational relative-equilibrium profile]
\label{def:R2-relative-profile-safe}
A bounded smooth ancient solution \((V,P)\) of \eqref{eq:centered-R2-safe} is a
rotational relative equilibrium if there exist
\(\Omega\in\mathfrak{so}(3)\), a bounded smooth divergence-free vector field
\(W\), and a pressure profile \(\Pi\), such that
\begin{equation}\label{eq:R2-relative-form-safe}
   V(y,\tau)=Q(\tau)W(Q(\tau)^Ty),
   \qquad
   P(y,\tau)=\Pi(Q(\tau)^Ty)
\end{equation}
up to the usual time-dependent pressure gauge.
The class \(\Crot\) consists of such residual Seregin profiles, after the
axisymmetric bounded-circulation class and the previously closed alternatives in
the setup paper have been removed.
\end{definition}

\begin{lemma}[Co-rotating stationary equation]
\label{lem:co-rotating-equation-safe}
Let \((V,P)\) be a rotational relative equilibrium.  Then the co-rotating
profile \((W,\Pi)\) satisfies
\begin{equation}\label{eq:R2-coriolis-profile-safe}
   -\Delta W+\frac12 y\cdot\nabla W+\frac12 W
   +(W\cdot\nabla)W
   +\Omega W-(\Omega y)\cdot\nabla W
   +\nabla\Pi=0,
   \qquad \nabla\cdot W=0 .
\end{equation}
Equivalently,
\begin{equation}\label{eq:R2-coriolis-profile-project-safe}
   \mathbb P\Big[
   -\Delta W+\frac12 y\cdot\nabla W+\frac12 W
   +(W\cdot\nabla)W
   +\Omega W-(\Omega y)\cdot\nabla W
   \Big]=0,
\end{equation}
where \(\mathbb P\) is the whole-space Leray projector.
\end{lemma}

\begin{proof}
Let \(z=Q(\tau)^Ty\).  Since \(Q'(\tau)=\Omega Q(\tau)\) and
\(Q(\tau)^T\Omega Q(\tau)=\Omega\), one has
\[
   \partial_\tau z=-\Omega z.
\]
Therefore
\[
   \partial_\tau V(y,\tau)
   =Q(\tau)\Big(\Omega W(z)-(\Omega z)\cdot\nabla W(z)\Big).
\]
The Laplacian, divergence, convection term, and pressure gradient are invariant
under the orthogonal change of variables.  Substitution into
\eqref{eq:centered-R2-safe}, followed by multiplication by \(Q(\tau)^T\), gives
\eqref{eq:R2-coriolis-profile-safe}.  Applying the Leray projector gives
\eqref{eq:R2-coriolis-profile-project-safe}.
\end{proof}

\subsection{The Bernoulli identity and the local Coriolis flux}

Write
\begin{equation}\label{eq:R2-F-def}
   F(y):=\frac12y-\Omega y,
   \qquad
   A(y):=W(y)+F(y),
   \qquad
   \omega:=\nabla\times W .
\end{equation}
Since \(\Omega\) is skew-symmetric, there is a unique vector
\(\varpi\in\mathbb R^3\) such that
\begin{equation}\label{eq:R2-Omega-vector}
   \Omega y=\varpi\times y .
\end{equation}
Then
\begin{equation}\label{eq:R2-curlF}
   \nabla\times F=-2\varpi .
\end{equation}

\begin{lemma}[Coriolis Bernoulli identity]
\label{lem:R2-coriolis-bernoulli}
Let \((W,\Pi)\) be a smooth solution of
\eqref{eq:R2-coriolis-profile-safe}.  Define
\begin{equation}\label{eq:R2-Bernoulli-def}
   B:=\frac12|W|^2+\Pi+F\cdot W .
\end{equation}
Then
\begin{equation}\label{eq:R2-gradient-B}
   \nabla B+\nabla\times\omega-A\times\omega=0,
\end{equation}
and consequently
\begin{equation}\label{eq:R2-Bernoulli-PDE}
   -\Delta B+A\cdot\nabla B
   =
   -|\omega|^2-\omega\cdot(\nabla\times F)
   =
   -|\omega|^2+2\varpi\cdot\omega .
\end{equation}
\end{lemma}

\begin{proof}
The vector identity
\[
   (W\cdot\nabla)W=\nabla\left(\frac12|W|^2\right)-W\times\omega
\]
and the identity
\[
   (F\cdot\nabla)W+(\nabla F)^TW
   =\nabla(F\cdot W)-F\times\omega
\]
apply because
\((\nabla F)^TW=\frac12W+\Omega W\).  Thus
\eqref{eq:R2-coriolis-profile-safe} can be rewritten as
\[
   -\Delta W+
   \nabla\left(\frac12|W|^2+\Pi+F\cdot W\right)
   -(W+F)\times\omega=0 .
\]
Since \(\nabla\cdot W=0\), one has
\(-\Delta W=\nabla\times\omega\), which gives
\eqref{eq:R2-gradient-B}.

Taking divergence in \eqref{eq:R2-gradient-B} gives
\[
   \Delta B=\nabla\cdot(A\times\omega),
\]
because the divergence of a curl vanishes.  Using
\[
   \nabla\cdot(A\times\omega)
   =\omega\cdot(\nabla\times A)-A\cdot(\nabla\times\omega),
   \qquad
   \nabla\times A=\omega+\nabla\times F,
\]
we obtain
\[
   \Delta B
   =|\omega|^2+\omega\cdot(\nabla\times F)
      -A\cdot(\nabla\times\omega).
\]
On the other hand, by \eqref{eq:R2-gradient-B},
\[
   A\cdot\nabla B
   =A\cdot(A\times\omega)-A\cdot(\nabla\times\omega)
   =-A\cdot(\nabla\times\omega).
\]
Combining the last two identities gives \eqref{eq:R2-Bernoulli-PDE}.  Finally,
\eqref{eq:R2-curlF} follows from \eqref{eq:R2-Omega-vector}.
\end{proof}

\begin{lemma}[Localized Coriolis flux identity]
\label{lem:R2-coriolis-flux-localization}
Let \((W,\Pi)\) be as in \Cref{lem:R2-coriolis-bernoulli}, and let
\(\chi\in C_c^\infty(\mathbb R^3)\).  Then
\begin{equation}\label{eq:R2-local-coriolis-source}
   \int_{\mathbb R^3}\chi
   \left(
      -\Delta B+A\cdot\nabla B+|\omega|^2
   \right)\,dy
   =
   -2\int_{\mathbb R^3}(W\times\varpi)\cdot\nabla\chi\,dy .
\end{equation}
In particular the Coriolis contribution is supported only where the cutoff
varies.
\end{lemma}

\begin{proof}
By \eqref{eq:R2-Bernoulli-PDE},
\[
   -\Delta B+A\cdot\nabla B+|\omega|^2=2\varpi\cdot\omega .
\]
Because \(\varpi\) is constant and \(\omega=\nabla\times W\),
\[
   \varpi\cdot\omega
   =
   \nabla\cdot(W\times\varpi).
\]
Multiplying by \(\chi\), integrating over \(\mathbb R^3\), and integrating by
parts gives
\[
   2\int\chi\,\varpi\cdot\omega
   =
   2\int\chi\,\nabla\cdot(W\times\varpi)
   =
   -2\int(W\times\varpi)\cdot\nabla\chi ,
\]
which proves \eqref{eq:R2-local-coriolis-source}.
\end{proof}

\begin{lemma}[Non-small Coriolis flux creates retained local velocity concentration]
\label{lem:R2-annular-flux-active-concentration}
Let \(\chi\in C_c^\infty(\mathbb R^3)\), let
\(A_\chi:=\operatorname{supp}\nabla\chi\), and put
\[
   C_\chi:=\int_{A_\chi}|\nabla\chi|^{3/2}\,dy .
\]
Fix \(T\in(0,1]\), and let
\[
   \mathcal T_{\chi,T}
   :=
   \{(y,\tau): -T<\tau<0,\ Q(\tau)^Ty\in A_\chi\}
\]
be the swept annular tube in the original centered variables.  Cover
\(\mathcal T_{\chi,T}\) by finitely many unit terminal cylinders
\[
   \mathcal T_{\chi,T}\subset \bigcup_{m=1}^{N_{\chi,T}}Q_1(z_m).
\]
The number \(N_{\chi,T}\) depends only on the compact annulus
\(A_\chi\), the time length \(T\), and \(\Omega\), and the cover is chosen with
bounded overlap inside the normalized terminal domain under consideration.
Let \(\eps_{\rm sd}>0\) be the raw unit-window threshold used in
\Cref{def:terminal-measures}.  If \(\varpi\ne0\) and
\begin{equation}\label{eq:R2-flux-active-threshold}
   \left|
      2\int_{\mathbb R^3}(W\times\varpi)\cdot\nabla\chi\,dy
   \right|
   \ge
   2|\varpi|\,C_\chi^{2/3}
      \left(\frac{N_{\chi,T}\eps_{\rm sd}}{T}\right)^{1/3},
\end{equation}
then at least one cylinder \(Q_1(z_m)\) is raw unit-window active at threshold
\(\eps_{\rm sd}\) for the original centered solution \(V\).
\end{lemma}

\begin{proof}
Write
\[
   \delta_\chi:=
   \left|
      2\int_{\mathbb R^3}(W\times\varpi)\cdot\nabla\chi\,dy
   \right| .
\]
H\"older's inequality gives
\[
\begin{aligned}
   \delta_\chi
   &\le
   2|\varpi|\int_{A_\chi}|W|\,|\nabla\chi|\,dy  \\
   &\le
   2|\varpi|
   \left(\int_{A_\chi}|W|^3\,dy\right)^{1/3}
   C_\chi^{2/3}.
\end{aligned}
\]
Thus \eqref{eq:R2-flux-active-threshold} implies
\[
   \int_{A_\chi}|W|^3\,dy\ge \frac{N_{\chi,T}\eps_{\rm sd}}{T}.
\]
Since \(V(y,\tau)=Q(\tau)W(Q(\tau)^Ty)\) and \(Q(\tau)\) is an isometry,
\[
\begin{aligned}
   \iint_{\mathcal T_{\chi,T}} |V(y,\tau)|^3\,dy\,d\tau
   &=
   \int_{-T}^{0}\int_{Q(\tau)A_\chi}
      |W(Q(\tau)^Ty)|^3\,dy\,d\tau \\
   &=
   T\int_{A_\chi}|W(z)|^3\,dz
   \ge N_{\chi,T}\eps_{\rm sd}.
\end{aligned}
\]
Because the tube is covered by \(N_{\chi,T}\) unit terminal cylinders with
bounded overlap, one of
them satisfies
\[
   \iint_{Q_1(z_m)} |V|^3\,dy\,d\tau\ge \eps_{\rm sd}.
\]
The pressure contribution to the local unit defect measure is nonnegative after the
allowed pressure gauge.  Hence
\[
   \mu(Q_1(z_m))\ge \eps_{\rm sd}.
\]
Thus the corresponding recentering sequence is raw unit-window active at
threshold \(\eps_{\rm sd}\) in the sense of \Cref{def:terminal-measures}.
The retained nonzero descendant later obtained from this observer is certified
by the displayed unit velocity lower bound, while the atlas pressure term is
kept only in the raw defect measure.
\end{proof}

\begin{lemma}[Local Coriolis-flux dichotomy]
\label{lem:R2-flux-reduction-dichotomy}
Fix \(\chi\in C_c^\infty(\mathbb R^3)\) and \(T\in(0,1]\).  Along every
terminal extraction of a rotational relative equilibrium, the following
priority alternatives are applied:
\begin{enumerate}[label=(\roman*)]
   \item \(\varpi\ne0\), the flux across
   \(A_\chi=\operatorname{supp}\nabla\chi\) satisfies the threshold
   \eqref{eq:R2-flux-active-threshold}, and the swept annular tube contains a
   unit-scale velocity concentration;
   \item either \(\varpi=0\), or the threshold
   \eqref{eq:R2-flux-active-threshold} fails; then the flux is a bounded
   compactly supported functional of \(W\) on \(A_\chi\), passes to terminal
   limits by local \(L^3\) convergence, and creates no terminal alternative other than
   the single-profile, finite separated-profile, exterior-escape,
   diffuse-exterior-concentration, local-vanishing, or
   local-compactness-failure alternatives of
      \Cref{thm:terminal-exhaustion-main}.
\end{enumerate}
\end{lemma}

\begin{proof}
Alternative \emph{(i)} is precisely
\Cref{lem:R2-annular-flux-active-concentration}.  If \(\varpi=0\), then the flux
vanishes identically.  If \(\varpi\ne0\) and
\eqref{eq:R2-flux-active-threshold} fails, the flux is bounded by the right hand
side of that threshold.  In both cases the functional
\[
   W\longmapsto -2\int_{\mathbb R^3}(W\times\varpi)\cdot\nabla\chi\,dy
\]
depends only on \(W\) on the compact annular set \(A_\chi\).  Under the local
Seregin compactness supplied by the admissible setup hypotheses, velocities converge
strongly in
\(L^3_{\rm loc}\); since \(\nabla\chi\in L^{3/2}\) and \(\varpi\) is fixed, the
functional is continuous along every terminal subsequence:
\[
   \int (W_n\times\varpi)\cdot\nabla\chi
   \longrightarrow
   \int (W_\infty\times\varpi)\cdot\nabla\chi .
\]
Thus a subthreshold flux is retained only as a local term in the limiting
localized Bernoulli identity \eqref{eq:R2-local-coriolis-source}.  It is not a
new location of local velocity concentration.  Every actual location of compact velocity
\(L^3\) concentration is still one of the terminal concentration-set alternatives listed in
\Cref{thm:terminal-exhaustion-main}.
\end{proof}

\begin{remark}[Why this is not a global Liouville argument]
\label{rem:R2-no-global-liouville}
For \(\Omega=0\), the Coriolis flux in
\Cref{lem:R2-coriolis-flux-localization} vanishes.  For \(\Omega\ne0\), the
same lemma does not give a sign-definite maximum principle for \(B\).  It gives
only the local dichotomy in \Cref{lem:R2-flux-reduction-dichotomy}: either a
chosen annulus already contains a recentered local concentration, or the
annular flux is a subthreshold compact functional which passes to the local
terminal limit.  No tail decay, whole-space \(L^p\) bound, Gaussian energy
estimate, or global pressure normalization is used.
\end{remark}

\subsection{Reduction of the rotational relative-equilibrium class}

\begin{proposition}[Rotational residual reduction]
\label{prop:R2-reduction-to-terminal}
Every rotational relative-equilibrium residual either belongs to a previously
closed lower class or produces a terminal concentration profile realized from
the original blow-up sequence
satisfying the hypotheses of the terminal stratification assembly theorem.
\end{proposition}

\begin{proof}
Let \(V\in\Crot(\mathcal S;I,J)\).  By
\Cref{def:R2-relative-profile-safe} and
\Cref{lem:co-rotating-equation-safe}, there are
\(\Omega\in\mathfrak{so}(3)\) and a bounded smooth divergence-free profile
\((W,\Pi)\) satisfying the Coriolis--Leray equation
\eqref{eq:R2-coriolis-profile-safe}.  The centered representative \(V\) is a
normalized Seregin ancient limit, so by
\Cref{hyp:base-seregin-hypotheses} it carries a fixed positive compact
velocity \(L^3\) concentration.

If \(\Omega=0\), then \(V\) is stationary in the centered frame.  It is reduced
to the stationary setup class if it lies in the already closed stationary
alternative, and otherwise to the stationary-hull residual branch handled by
\Cref{prop:R3-stationary-hull-reduction}.  We therefore assume
\(\Omega\ne0\).

Choose a countable family of smooth cutoffs which is cofinal among compact
annuli in the terminal variables.  For each cutoff and each rational
\(T\in(0,1]\), apply the local dichotomy
\Cref{lem:R2-flux-reduction-dichotomy}.  A superthreshold Coriolis flux produces
a unit-scale velocity concentration and is therefore already part of the local
concentration decomposition.  A subthreshold Coriolis flux is a compactly
supported functional
of \(W\), continuous under the local convergence supplied by the admissible setup
compactness hypotheses, and hence remains inside the localized limiting equation; it carries no
independent local velocity
mass and creates no additional terminal alternative.

We now apply the paired terminal concentration-set/remainder decomposition to the inherited compact
density.  Loss of local compactness, local vanishing, exterior escape, and
diffuse exterior concentration cannot be the sole locations of this density by
\Cref{thm:terminal-inactive-exclusion}.  By the preceding paragraph, the
Coriolis flux has already led either to a local concentration or to a
compact local term in the selected terminal equation.  Thus, after local
concentration extraction, the retained density is located either in a finite separated family
of retained velocity concentration profiles or in a single retained velocity concentration
terminal profile, together with any residual diffuse, noncompact, or
critical-tail remainder selected by the terminal extraction procedure.
This gives a terminal input realized from the original blow-up sequence for the terminal stratification
theorem; no R2 closure is used in this reduction step.
\end{proof}

\begin{proposition}[Already-closed integrable subcase]
\label{prop:R2-integrable-subcase}
Let \(V\) be a rotational relative equilibrium with co-rotating profile \(W\).
If
\begin{equation}\label{eq:R2-L3-subcase}
   W\in L^3(\mathbb R^3),
\end{equation}
then \(V\equiv0\).  More generally, the same conclusion holds under any
condition which places the physical representative in
\(L^\infty_tL^3_x\) on its ancient time interval.
\end{proposition}

\begin{proof}
Rotations and the Navier--Stokes critical scaling preserve the \(L^3\) norm.
Thus \(W\in L^3\) implies
\[
   \sup_{\tau\in\mathbb R}\|V(\cdot,\tau)\|_{L^3(\mathbb R^3)}<\infty .
\]
Equivalently, the physical representative lies in the critical class
\(L^\infty_tL^3_x\).  The endpoint \(L^3\) class is one of the already closed
alternatives used in the setup paper, by the Escauriaza--Seregin--\v Sver\'ak endpoint
regularity theorem \cite{ESS2003} and the stationary
Ne\v cas--R\r u\v zi\v cka--\v Sver\'ak Liouville theorem \cite{NRS1996}.
Therefore
\(V\equiv0\).  A zero profile has zero local velocity concentration, so no such profile
can belong to the residual
nonvanishing class.
\end{proof}


\section{The degenerate stationary-hull residual class}
\label{sec:R3-degenerate-stationary-hull}

The third refined residual class consists of Type~I ancient limits whose
time-translation hull contains a stationary
profile with retained local velocity concentration lying on a non-isolated family of
stationary centered profiles.  The role of the stationary-family condition is to separate genuinely degenerate
stationary profile families from isolated stationary profiles and from the
symmetry-generated classes already removed in \(\Cax\) and \(\Crot\).

The retained-velocity concentration condition is essential.  The mere presence of a
stationary-family element in the hull of a residual ancient solution is insufficient.  A
stationary element may occur as a limit without retained local velocity concentration and
therefore need not retain the local velocity concentration inherited from the
original singular point.  Such a hull element is not placed in a class
that is to be closed by a
stationary-profile exclusion.  Therefore the class below is defined using
\emph{retained-velocity occurrence}: the stationary-family profile must
itself occur as a Seregin limit carrying a positive local velocity concentration.

Throughout this section the centered equation is
\begin{equation}\label{eq:centered-R3}
   \partial_\tau V-\Delta V+\frac12 y\cdot\nabla V+\frac12 V
   +(V\cdot\nabla)V+\nabla P=0,
   \qquad \nabla\cdot V=0 .
\end{equation}
Its projected form is
\begin{equation}\label{eq:projected-centered-R3}
   \partial_\tau V=\mathcal F(V),
   \qquad
   \mathcal F(W)
   :=
   \mathbb P\left[
      \Delta W-\frac12(W+y\cdot\nabla W)-(W\cdot\nabla)W
   \right],
\end{equation}
where \(\mathbb P\) is the whole-space Leray projector.

\subsection{Stationary families and retained-velocity hull limits}

\begin{definition}[Admissible stationary phase space]
\label{def:R3-stationary-phase-space}
An admissible stationary phase space is a Banach space \(\mathfrak B\) of
solenoidal vector fields on \(\mathbb R^3\) such that:
\begin{enumerate}[label=(\roman*)]
   \item \(\mathfrak B\hookrightarrow C^{2,\alpha}_{\loc}(\mathbb R^3)
      \cap L^\infty(\mathbb R^3)\) continuously for some
      \(\alpha\in(0,1)\);
   \item the map \(\mathcal F\) in \eqref{eq:projected-centered-R3} is
      well-defined on an open subset \(\mathcal U\subset\mathfrak B\) and maps
      \(\mathcal U\) into a Banach space \(\mathfrak Y\);
   \item \(\mathcal F:\mathcal U\to\mathfrak Y\) is \(C^1\).
\end{enumerate}
\end{definition}

\begin{definition}[Nontrivial stationary family]
\label{def:R3-stationary-family}
Let \(W\in\mathcal U\subset\mathfrak B\) satisfy \(\mathcal F(W)=0\).  We say
that \(W\) lies on an admissible nontrivial stationary family if there exist
\(\varepsilon>0\) and a \(C^1\) curve
\[
   a\mapsto W_a\in\mathcal U,
   \qquad |a|<\varepsilon,
\]
such that
\begin{equation}\label{eq:R3-stationary-family}
   W_0=W,
   \qquad
   \mathcal F(W_a)=0 \quad (|a|<\varepsilon),
\end{equation}
and such that the family is not locally generated only by the action of the
rigid rotation group.  Equivalently, after possibly reparametrizing the curve,
one may assume that
\begin{equation}\label{eq:R3-nontrivial-tangent}
   \dot W_0\notin
   \{\Omega W-(\Omega y)\cdot\nabla W:\Omega\in\mathfrak{so}(3)\}
\end{equation}
whenever the tangent \(\dot W_0\) exists and is nonzero.
\end{definition}

\begin{definition}[Retained stationary-hull velocity concentration]
\label{def:R3-active-profile}
Let \((W,\Pi)\) be a stationary centered profile.  We say that \((W,\Pi)\) is
retained-stationary concentrating if it has retained velocity concentration in the normalized
unit sense: there is a recentering/rescaling realized from the original blow-up
sequence
\(\mathscr T\) such that
\begin{equation}\label{eq:R3-active-velocity}
   \Vel_1\bigl(\mathscr T W;0\bigr)
   \ge \eta_* .
\end{equation}
Equivalently, before the final unit normalization one may require
\(\Vel_r(W;z_0)\ge\eta_*\) for some parabolic cylinder
\(Q_r(z_0)\).  Unnormalized large-cylinder mass is not a retained velocity concentration
criterion.
It has stationary-hull CKN activity if the same realized
normalization also satisfies
\[
   \calE_1(\mathscr T W,\mathscr T\Pi;0)\ge\eps_* .
\]
Only the velocity condition \eqref{eq:R3-active-velocity} certifies that the
stationary-hull profile is nonzero.  The accompanying CKN activity condition is
used only for CKN smallness and pressure estimates.
\end{definition}

\begin{definition}[Degenerate stationary-hull residual class with retained velocity concentration]
\label{def:R3-active-class}
Fix an admissible stationary phase space \(\mathfrak B\).  The degenerate
stationary-hull residual class with retained local velocity concentration
\[
   \Cstat(\mathcal S;I,J;\mathfrak B)
\]
consists of those \(V\in\mathcal R(\mathcal S;I,J)\), after the ordered
removal of \(\Cax\) and \(\Crot\), for which there exists a
sequence \(s_k\in\mathbb R\) and a stationary profile \((W,\Pi)\) such that
\begin{equation}\label{eq:R3-hull-convergence}
   (V(\cdot,\cdot+s_k),P(\cdot,\cdot+s_k)-c_k(\cdot))
   \longrightarrow (W,\Pi)
\end{equation}
locally smoothly for the velocity and locally in \(L^{3/2}\) for the pressure
modulo admissible gauges, and such that:
\begin{enumerate}[label=(\roman*)]
   \item \(W\in\mathfrak B\) and \(\mathcal F(W)=0\);
   \item \(W\) lies on an admissible nontrivial stationary family in the sense
      of \Cref{def:R3-stationary-family};
   \item \((W,\Pi)\) is retained-stationary concentrating in the sense of
      \Cref{def:R3-active-profile};
   \item \((W,\Pi)\) is not in any previously closed stationary, structured,
      axisymmetric bounded-circulation, or rotational relative-equilibrium
      class.
\end{enumerate}
In the ordered residual decomposition one keeps this class of stationary-hull
limits with retained velocity concentration as the degenerate stationary-hull
alternative.  The hull cases
without retained local velocity concentration
are assigned to the generic terminal residual class, where they are
handled by the terminal concentration-set decomposition rather than by a
stationary-profile exclusion.
\end{definition}

\begin{remark}[Why retained velocity concentration is necessary]
\label{rem:R3-active-necessary}
The condition \eqref{eq:R3-active-velocity} prevents the following logical error.  A
nonstationary ancient solution may have a stationary profile in its hull with
zero local velocity concentration on the relevant compact cylinders.  Excluding that stationary
profile would not exclude the original ancient solution.  The
stationary-profile exclusion used below applies only to stationary profiles
that themselves occur as nontrivial Seregin limits with positive local
concentration.
\end{remark}

\subsection{Hull limits are again Seregin limits}

The next lemma is the basic compatibility between time-translation hulls in
centered variables and rescaling sequences in physical variables.  It is the
reason that a stationary profile in the hull which retains local concentration
may be treated as a blow-up profile of the original solution, not merely as a
limit inside the already extracted ancient solution.

\begin{lemma}[Scale reselection for centered hull limits]
\label{lem:R3-hull-is-seregin}
Let \((V,P)\) be a centered Type~I Seregin limit obtained from a finite-energy
suitable weak solution \((u,p)\) at a singular point \(z_*=(x_*,t_*)\).  Suppose
that, for some sequence \(s_k\in\mathbb R\),
\[
   (V(\cdot,\cdot+s_k),P(\cdot,\cdot+s_k)-c_k(\cdot))
   \to (W,\Pi)
\]
locally in the Seregin topology.  Then, after passing to a diagonal subsequence
and reselecting the physical blow-up scales, \((W,\Pi)\) is itself a centered
Type~I Seregin limit of \((u,p)\) at \(z_*\).  If the convergence to \((W,\Pi)\)
retains a positive local velocity lower bound, then \((W,\Pi)\) is a nontrivial
Seregin limit with the same type of inherited nonvanishing condition.
\end{lemma}

\begin{proof}
Let \(\lambda_n\downarrow0\) be a sequence of scales producing \((V,P)\).  In
centered variables, a time translation by \(s\) is a change of the
physical blow-up scale by a factor depending only on \(s\).  With the convention
used in \eqref{eq:centered-R3}, write this adjusted scale as
\begin{equation}\label{eq:R3-adjusted-scale}
   \lambda_{n,s}:=e^{-s/2}\lambda_n .
\end{equation}
Indeed, writing \(t=-e^{-\tau}\), a shift
\(\tau\mapsto\tau+s\) corresponds to multiplying the physical length scale by
\(e^{-s/2}\), so that
\[
   V_{\lambda e^{-s/2}}(y,\tau)=V_\lambda(y,\tau+s)
\]
with the corresponding pressure gauge.  Thus the rescaled fields built from
\((u,p)\) at scale \(\lambda_{n,s}\) are the same as the fields at scale
\(\lambda_n\), translated by \(s\) in centered time.

Fix an exhaustion of compact cylinders \(K_m\subset\mathbb R^3\times\mathbb R\).
For each \(k\), first require
\[
   \lambda_{n(k)}
   \le k^{-1}e^{s_k/2}
\]
when \(s_k<0\), and otherwise require simply \(\lambda_{n(k)}\le k^{-1}\).
Then \(\lambda_{n(k),s_k}=e^{-s_k/2}\lambda_{n(k)}\to0\).  Increasing
\(n(k)\) if needed, arrange that the rescaled fields at scale
\(\lambda_{n(k)}\), shifted by \(s_k\), are within \(1/k\) of
\((V(\cdot,\cdot+s_k),P(\cdot,\cdot+s_k)-c_k)\) on \(K_k\) in the local
Seregin topology.  This is possible because the original sequence converges
locally to \((V,P)\).

By the assumed hull convergence, the shifted limit
\((V(\cdot,\cdot+s_k),P(\cdot,\cdot+s_k)-c_k)\) converges to \((W,\Pi)\) on each
fixed compact cylinder.  The diagonal sequence of physical rescalings with
scales \(\lambda_{n(k),s_k}\) therefore converges locally to \((W,\Pi)\).  Hence
\((W,\Pi)\) is again a centered Type~I Seregin limit at \(z_*\).  The assertion
about nonvanishing follows from strong local \(L^3\) convergence of the
velocity and the retained velocity lower bound; pressure compactness is used
only to keep the limiting pressure gauge admissible.
\end{proof}

\subsection{Elementary consequences of stationary-family degeneracy}

The stationary-family condition has two immediate PDE consequences.  First, a differentiable
nontrivial family produces a nontrivial element of the kernel of the linearized
stationary Leray operator.  Second, if an ancient trajectory is
constrained to a stationary family, then it is stationary.

\begin{lemma}[Stationary family implies linearized degeneracy]
\label{lem:R3-family-kernel}
Let \(a\mapsto W_a\in\mathfrak B\) be a \(C^1\) stationary family satisfying
\eqref{eq:R3-stationary-family}.  Then
\begin{equation}\label{eq:R3-kernel}
   D\mathcal F(W_0)\dot W_0=0 .
\end{equation}
In particular, if \(\dot W_0\ne0\) modulo infinitesimal rotations, the
linearized stationary Leray operator has a genuine non-symmetry kernel element.
\end{lemma}

\begin{proof}
Differentiate \(\mathcal F(W_a)=0\) at \(a=0\).  Since \(\mathcal F\) is \(C^1\)
on \(\mathfrak B\), this gives \eqref{eq:R3-kernel}.  The final statement is
the nontriviality condition in \Cref{def:R3-stationary-family}.
\end{proof}

\begin{lemma}[Trajectories constrained to a stationary family are stationary]
\label{lem:R3-family-constrained-stationary}
Let \(a\mapsto W_a\in\mathfrak B\) be a \(C^1\) stationary family with
\(\dot W_a\ne0\) on an interval.  Suppose that a smooth solution of
\eqref{eq:projected-centered-R3} has the form
\begin{equation}\label{eq:R3-family-trajectory}
   V(\tau)=W_{a(\tau)}
\end{equation}
for a \(C^1\) function \(a(\tau)\).  Then \(a\) is constant, and hence
\(V\) is stationary.
\end{lemma}

\begin{proof}
Since \(\mathcal F(W_a)=0\) for every \(a\), equation
\eqref{eq:projected-centered-R3} gives
\[
   a'(\tau)\dot W_{a(\tau)}=\partial_\tau V(\tau)=\mathcal F(W_{a(\tau)})=0 .
\]
Because \(\dot W_{a(\tau)}\ne0\), we obtain \(a'(\tau)=0\).  Thus \(a\) is
constant on the interval.
\end{proof}

\begin{corollary}[Exact stationary-family subcase]
\label{cor:R3-exact-family-subcase}
If an element of the residual class is itself an exact trajectory on an
admissible stationary family, then it is a stationary centered profile.  Hence
it is excluded whenever it satisfies the stationary \(L^3\) Liouville
criterion, or more generally whenever the stationary profile belongs to one of
the previously closed stationary or structured classes.
\end{corollary}

\begin{proof}
Let \(V(\tau)=W_{a(\tau)}\) be such an exact stationary-family trajectory.  By
\Cref{lem:R3-family-constrained-stationary}, \(a'(\tau)=0\) on every interval on
which the stationary-family parametrization is nondegenerate.  Hence \(a\) is constant and
\(V\) is a single stationary centered profile.  If that stationary profile
satisfies the stationary \(L^3\) Liouville criterion or lies in one of the
previously closed structured or ordered lower residual classes, the corresponding
prior exclusion removes it from the residual class.
\end{proof}

\subsection{Stationary-hull profiles which occur as Seregin limits}

The preceding lemma removes the possible internal dynamics along a stationary
family: exact stationary-family motion is stationary by
\Cref{lem:R3-family-constrained-stationary}.  The remaining question is
attainability by blow-up: whether a degenerate stationary-family profile can
occur as a Seregin limit of a finite-energy suitable weak solution while
retaining local velocity concentration.  The following statement gives the precise
Seregin-limit occurrence exclusion
required to close the degenerate stationary-hull class with retained local
velocity concentration.  Its proof is deferred until the terminal stratification assembly theorem
has been established.

\begin{proposition}[R3 degenerate stationary-family reduction target]
\label{thm:R3-no-attainable-degenerate-family}
Fix an admissible stationary phase space \(\mathfrak B\).  Let \((W,\Pi)\) be a
stationary centered profile satisfying:
\begin{enumerate}[label=(\roman*)]
   \item \(W\in\mathfrak B\), \(\mathcal F(W)=0\), and \(W\) lies on an
      admissible nontrivial stationary family;
   \item \((W,\Pi)\) is a centered Type~I Seregin limit of a finite-energy
      suitable weak solution at a singular point;
   \item \((W,\Pi)\) is retained-stationary concentrating in the sense of
      \Cref{def:R3-active-profile};
   \item \((W,\Pi)\) is not in the small, stationary \(L^3\), tight, closed
      structured/decay, axisymmetric bounded-circulation, or rotational
      relative-equilibrium classes already removed earlier in the stratification.
\end{enumerate}
Then no such \((W,\Pi)\) exists after the terminal stratification assembly is
inserted.
\end{proposition}

\begin{proof}[Proof deferred]
The proof is given after the terminal stratification assembly theorem in
\Cref{cor:R3-closure}.  The present statement records the exact R3 occurrence
input that must be proved; no R3 exclusion is used before the terminal
assembly theorem is available.
\end{proof}

\begin{remark}[Interpretation of the Seregin limit occurrence theorem]
\label{rem:R3-input-interpretation}
\Cref{thm:R3-no-attainable-degenerate-family} excludes occurrence
as a Type~I Seregin limit.  It does not assert that degenerate stationary profiles cannot
exist as elliptic solutions of the stationary centered equation.  It asserts
only that such profiles cannot occur as Type~I blow-up limits with retained
local velocity concentration after all previously closed classes have been removed.  This
is the formulation required for the residual problem: the elliptic family may
exist, but it cannot be reached by the local blow-up procedure coming from a
finite-energy suitable weak solution.
\end{remark}

\subsection{Reduction of the degenerate stationary-hull concentration class}

\begin{proposition}[R3 stationary-hull reduction]
\label{prop:R3-stationary-hull-reduction}
Every retained stationary-hull residual profile either belongs to one of the
stationary, tight, or structured classes already excluded by the setup paper, or it produces a stationary terminal
concentration profile realized from the original blow-up sequence and satisfying the hypotheses of the terminal
stratification assembly theorem.
\end{proposition}

\begin{proof}
Let \(V\in\Cstat(\mathcal S;I,J;\mathfrak B)\).  By
\Cref{def:R3-active-class}, there is a stationary profile \((W,\Pi)\) in the
centered time-translation hull of \((V,P)\), and \((W,\Pi)\) is
retained-stationary concentrating.  The profile \(W\) belongs to \(\mathfrak B\), satisfies
\(\mathcal F(W)=0\), and lies on an admissible nontrivial stationary family.
Moreover, by the ordered definition of \(\Cstat\), it is not
in any of the previously closed lower classes.

By \Cref{lem:R3-hull-is-seregin}, the hull profile \((W,\Pi)\) is itself a
centered Type~I Seregin limit of the original finite-energy suitable weak
solution at the same singular point, after reselecting scales.  Since local
velocity concentration is retained, \((W,\Pi)\) satisfies every condition in
\Cref{thm:R3-no-attainable-degenerate-family}, except that the actual
contradiction is deferred to the terminal stratification assembly theorem.  Thus the
R3 branch has been reduced to a stationary terminal concentration profile
realized from the original blow-up sequence; no R3 closure is used in this
reduction step.
\end{proof}

\begin{proposition}[Already-closed integrable stationary-family subcase]
\label{prop:R3-L3-subcase}
Let \((W,\Pi)\) be a stationary centered profile lying on an admissible
stationary family.  If
\begin{equation}\label{eq:R3-L3}
   W\in L^3(\mathbb R^3),
\end{equation}
then \(W\equiv0\).  Hence no retained-velocity hull profile satisfying
\eqref{eq:R3-L3} can occur in \(\Cstat\).
\end{proposition}

\begin{proof}
The stationary physical representative associated with \(W\) is a backward
self-similar Navier--Stokes solution.  The assumption \(W\in L^3(\mathbb R^3)\)
places it in the stationary \(L^3\) class already excluded by the classical
Ne\v cas--R\r u\v zi\v cka--\v Sver\'ak Liouville theorem \cite{NRS1996},
equivalently by the endpoint ancient Liouville theorem already recorded in
the setup paper and cited below as \Cref{thm:AB-main}.  Hence \(W\equiv0\).  A
zero profile has zero local velocity concentration, so it cannot be
retained-stationary concentrating.
\end{proof}

\begin{remark}[Closure after terminal stratification assembly]
\label{rem:R3-final-obligation}
The preceding argument reduces the stationary-hull class with retained local
velocity concentration from a time-translation hull question to the question whether a
stationary centered profile in a nontrivial stationary family can occur as a
Seregin limit.
\Cref{thm:R3-no-attainable-degenerate-family} is supplied only after the
terminal stratification assembly theorem has been proved.
\end{remark}

\begin{remark}[Residual decomposition statement]
\label{rem:R3-decomposition-replacement}
For the residual decomposition, the third line is reduced as follows and closed
only after the terminal stratification assembly theorem:
\[
\begin{gathered}
   \text{stationary-hull residual with retained velocity concentration}\\
   \xrightarrow{\;\text{stationary-family Seregin limit reduction}\;}
   \text{terminal stationary concentration profile}\\
   \xrightarrow{\;\text{terminal closure}\;}
   \varnothing .
\end{gathered}
\]
The precise class is the
\emph{degenerate stationary-hull residual class with retained local
velocity concentration}.
\end{remark}


% -----------------------------------------------------------------------------
\section{Reduction of degenerate stationary-hull profiles with retained velocity concentration}
\label{sec:R3-terminal-stratification-closure}

We record the terminal concentration input needed for the degenerate stationary-hull
class.  The actual closure is deferred until after the terminal stratification
theorem.  The argument is not specific to the existence of a
non-isolated stationary family.  Once a stationary-family element occurs as
a Type~I Seregin limit carrying retained local velocity concentration, it is a
stationary terminal concentration profile.  Such profiles are later excluded by the
same terminal stratification assembly and separated-profile exclusion used
for the final residual class.  No global coercive
estimate for the original solution is used.

Throughout this section, a centered Type~I profile satisfies
\begin{equation}\label{eq:R3S-centered}
   \partial_\tau U-\Delta U+\frac12 y\cdot\nabla U+\frac12 U
   +(U\cdot\nabla)U+\nabla \Pi=0,
   \qquad \nabla\cdot U=0
\end{equation}
on \(\mathbb R^3\times\mathbb R\), with the usual local suitability and
pressure-gauge conventions.

\subsection{Deferred terminal stratification corollary for R3}

We record the following specialization of the terminal stratification assembly
theorem.  It is stated here exactly in the form needed for the stationary-hull
closure with retained local velocity concentration as a \emph{reduction
target}; the substantive proof is the terminal stratification proof later in
the paper.  Following the recommendation that only one actual R3 closure
theorem appear after terminal assembly, this statement is recorded as a
proposition rather than a theorem.

\begin{proposition}[R3 reduction target]
\label{thm:R3S-terminal-theorem}
Let \((U,\Pi)\) be a bounded centered Type~I Seregin profile satisfying local
suitability and carrying a retained local velocity concentration, i.e. for some
admissible recentering/rescaling \(\mathscr T_{z_0,r}\),
\begin{equation}\label{eq:R3S-retained-activity}
   \Vel_1\bigl(\mathscr T_{z_0,r}U;0\bigr)
   \ge \eta_* .
\end{equation}
Assume that \((U,\Pi)\) is outside the previously closed small-amplitude,
stationary \(L^3\), tight, closed structured/decay, axisymmetric
bounded-circulation, and rotational relative-equilibrium classes.  Then the
following hold.
\begin{enumerate}[label=(\roman*)]
   \item Every terminal extraction of \((U,\Pi)\) has the paired decomposition
   of \Cref{thm:terminal-exhaustion-main}: a retained velocity concentration set
   together with a residual remainder which is locally vanishing, diffuse, or
   noncompact after neighborhoods of the concentration set are removed.
   \item Local vanishing, exterior escape, diffuse exterior concentration, and
   loss of local compactness cannot be the sole locations of the retained
   compact velocity concentration
   \eqref{eq:R3S-retained-activity}.
   \item Mixed compact--diffuse residual profiles are impossible in the refined
   residual class, while compact--noncompact configurations are routed by the
   terminal recentering protocol to the diffuse, pressure/noncompact, or
   critical-tail branches.  Infinite chains of retained concentrating
   descendants are impossible as well.
   \item Finite separated families of retained velocity concentration profiles are
   impossible.
   \item Any remaining single retained velocity concentration profile is terminally
   indecomposable: it has no retained concentrating tail limit, no diffuse
   exterior concentration, and no loss of local compactness along escaping
   recenterings.
   \item Every retained terminally indecomposable concentration profile has a
   backward sequence of finite critical norm: there exist
   \(\tau_k\to-\infty\) such that
   \begin{equation}\label{eq:R3S-atomic-L3-input}
      \sup_k \|U(\cdot,\tau_k)\|_{L^3(\mathbb R^3)}<\infty .
   \end{equation}
   \item Retained terminal concentration profiles are either affine/parasitic
   lower-stratum profiles or have parasitic-free finite-shift mild physical
   pullbacks.  Therefore the Albritton--Barker endpoint Liouville theorem
   applies to any normalized residual profile satisfying
   \eqref{eq:R3S-atomic-L3-input}.
\end{enumerate}
\end{proposition}

\begin{proof}[Proof deferred]
This is the R3 specialization of
\Cref{thm:terminal-stratification}.  It is stated here only as a
reduction target for the stationary-hull section; its proof is the terminal
stratification proof in
\Cref{sec:terminal-concentration,sec:separated-profiles-atomic}.
\end{proof}

\begin{remark}[Local nature of the theorem]
\label{rem:R3S-local-nature}
The theorem above is a local velocity concentration statement.  It uses the local
suitable-weak-solution compactness and pressure normalization developed from
the Caffarelli--Kohn--Nirenberg theory \cite{CKN1982} together with the
imported setup results: local positive concentration
and Type~I Seregin compactness.  It then uses paired terminal
profile decomposition, sole-source exclusions, critical-tail compactification
with realized critical-tail rigidity, recurrence for concentration chains, and the
finite separated-profile theorem.  It does not assert or require a global a
priori bound for the original physical solution.
\end{remark}

For reference, we record the stationary specialization of the atomic
sequence-\(L^3\) lemma used after the terminal stratification assembly theorem.

\begin{lemma}[Indecomposable local profiles have backward sequence-\(L^3\) control]
\label{lem:R3S-atomic-L3}
Let \((U,\Pi)\) be a bounded smooth ancient solution of
\eqref{eq:R3S-centered} which belongs to the normalized parasitic-free compact
state class \(\mathfrak X_M^\#\) imported from the setup construction.  Suppose
that \((U,\Pi)\) has retained local velocity concentration and is terminally
indecomposable: it has no retained
concentrating tail limit, no diffuse exterior concentration in the parabolic
CKN sense, and no escaping recentering sequence along which local compactness,
suitability, or pressure control fails.  Then there
exists a sequence \(\tau_k\to-\infty\) such that
\begin{equation}\label{eq:R3S-atomic-L3}
   \sup_k \|U(\cdot,\tau_k)\|_{L^3(\mathbb R^3)}<\infty .
\end{equation}
\end{lemma}

\begin{proof}
This is the stationary-profile specialization of
\Cref{lem:atomic-sequence-L3}.  The hypotheses listed here are the
terminal indecomposability hypotheses used there: retained unit velocity concentration, no
retained concentrating tail limit, no diffuse exterior concentration, and no
escaping recentering with local compactness or pressure-gauge failure.  Applying
\Cref{lem:atomic-sequence-L3} gives \eqref{eq:R3S-atomic-L3}.  No separate
exterior activity level is introduced in this stationary specialization; the
only exterior activity selection is the velocity quantity \(m_n^{\rm vel}\) in
\Cref{lem:atomic-sequence-L3}, and pressure-inclusive CKN activity is not used
to certify nontriviality of a tail profile.
\end{proof}

\subsection{Deferred stationary concentration statement}

\begin{remark}[Deferred R3 stationary consequence]
\label{thm:R3S-no-active-stationary-carrier}
Let \((W,\Pi)\) be a stationary solution of the centered equation,
\begin{equation}\label{eq:R3S-stationary-centered}
   -\Delta W+\frac12 y\cdot\nabla W+\frac12 W
   +(W\cdot\nabla)W+\nabla\Pi=0,
   \qquad \nabla\cdot W=0 .
\end{equation}
Assume that:
\begin{enumerate}[label=(\roman*)]
   \item \((W,\Pi)\) is obtained as a centered Type~I Seregin limit of a
   finite-energy suitable weak solution at a singular point;
   \item \((W,\Pi)\) has retained local velocity concentration in the sense of
   \eqref{eq:R3S-retained-activity};
   \item \((W,\Pi)\) is outside all previously closed lower classes listed in
   \Cref{thm:R3S-terminal-theorem}.
\end{enumerate}
Then no such \((W,\Pi)\) exists.  This statement is the stationary case of
the R3 closure after \Cref{thm:terminal-stratification}; it follows from
\Cref{cor:R3-closure} and is recorded here only to name the exact stationary
concentration statement used by the stationary-hull reduction.
\end{remark}

\subsection{Deferred degenerate stationary-hull closure statement}

\begin{remark}[Deferred R3 degenerate-family consequence]
\label{cor:R3S-no-active-degenerate-family}
Fix an admissible stationary phase space \(\mathfrak B\).  There is no stationary
centered profile \((W,\Pi)\) satisfying all of the following conditions:
\begin{enumerate}[label=(\roman*)]
   \item \(W\in\mathfrak B\), \(W\) solves the stationary centered equation
   \eqref{eq:R3S-stationary-centered}, and \(W\) lies on an admissible
   nontrivial stationary family;
   \item \((W,\Pi)\) is a centered Type~I Seregin limit of a finite-energy
   suitable weak solution at a singular point;
   \item \((W,\Pi)\) has retained local velocity concentration;
   \item \((W,\Pi)\) is outside the previously closed lower classes.
\end{enumerate}
This statement is recorded as a deferred R3 consequence; it follows from the
post-terminal R3 closure \Cref{cor:R3-closure}.  The stationary-family
condition in \textup{(i)} is additional structure; once the terminal
stratification assembly theorem is available, the stationary concentration
argument applies to these profiles in particular.
\end{remark}

\begin{proposition}[Stationary-hull reduction to terminal stratification]
\label{prop:R3S-R3-terminal-reduction}
Let \(\Cstat(\mathcal S;I,J;\mathfrak B)\) be the degenerate
stationary-hull residual class with retained local velocity concentration.  Every
element of this class produces a stationary terminal concentration profile
realized from the original blow-up sequence and satisfying the hypotheses of
\Cref{thm:terminal-stratification}.
\end{proposition}

\begin{proof}
Let \(V\in\Cstat(\mathcal S;I,J;\mathfrak B)\).  By the
definition of the degenerate stationary-hull residual class with retained local
velocity concentration, the centered translation hull of \((V,P)\) contains a stationary
profile \((W,\Pi)\) with retained local velocity concentration
which lies in an admissible nontrivial stationary family and is outside the
previously closed lower classes.  The scale-reselection lemma for centered hull
limits shows that \((W,\Pi)\) is itself a centered Type~I Seregin limit of the
original finite-energy suitable weak solution at the same singular point.

Thus \((W,\Pi)\) satisfies the hypotheses of the terminal stratification
assembly theorem in the stationary case.  The actual emptiness conclusion is recorded
after that terminal stratification assembly theorem as \Cref{cor:R3-closure}.
\end{proof}

\begin{remark}[How this enters the residual decomposition]
\label{rem:R3S-decomposition}
The third residual line is reduced as follows and closed only after the terminal
stratification theorem:
\[
\begin{gathered}
   \text{degenerate stationary-hull residual with retained velocity concentration}\\
   \Longrightarrow
   \text{stationary profile with retained velocity concentration}
   \Longrightarrow
   \varnothing .
\end{gathered}
\]
The second implication is proved by the terminal stratification assembly:
decomposable profiles produce either an excluded separated concentration
family, diffuse exterior concentration, or loss of local compactness;
indecomposable profiles have backward sequence-\(L^3\) control and are ruled
out by the endpoint ancient Liouville theorem.
\end{remark}


\section{Tail geometry of the generic terminal residual}
\label{sec:R4-tail-geometry}

We now pass to the generic terminal residual class.  The preliminary point is that
the centered equation is not invariant under ordinary fixed spatial translations.
The correct far-field recenterings are self-similar covariant translations.

\begin{definition}[Covariant translations]
\label{def:covariant-translations}
For \(a\in\R^3\), define
\[
   (T_aV)(y,\tau):=V(y+e^{\tau/2}a,\tau),
   \qquad
   (T_aP)(y,\tau):=P(y+e^{\tau/2}a,\tau).
\]
\end{definition}

\begin{lemma}[Covariant translation invariance]
\label{lem:covariant-translation-invariance}
If \((V,P)\) solves \eqref{eq:centered-main}, then
\((T_aV,T_aP)\) also solves \eqref{eq:centered-main}.
\end{lemma}

\begin{proof}
Set \(z=y+e^{\tau/2}a\).  Then
\[
   \partial_\tau(T_aV)(y,\tau)
   =
   (\partial_\tau V)(z,\tau)
   +\frac12e^{\tau/2}a\cdot\nabla V(z,\tau).
\]
Moreover,
\[
   \frac12 y\cdot\nabla(T_aV)(y,\tau)
   =
   \frac12 z\cdot\nabla V(z,\tau)
   -\frac12e^{\tau/2}a\cdot\nabla V(z,\tau).
\]
The two additional transport terms cancel in
\(\partial_\tau T_aV+\frac12y\cdot\nabla T_aV\).  The Laplacian, the
zeroth-order term, the convection term, the pressure gradient, and the
divergence constraint are translation covariant.  Substitution into
\eqref{eq:centered-main} proves the claim.
\end{proof}

\begin{remark}[Fixed translates create a spurious drift]
\label{rem:raw-translates}
If one sets \(V_n(y,\tau)=V(y+x_n,\tau)\) with fixed \(x_n\), then the centered
equation acquires the additional transport term
\[
   -\frac12 x_n\cdot\nabla V_n .
\]
Thus fixed translates do not define a compact invariant far-field hull for the
renormalized dynamics.
\end{remark}

\begin{definition}[Spacetime covariant action]
\label{def:spacetime-covariant-action}
For \(s\in\R\), define the centered time translation
\[
   (\Theta_sV)(y,\tau):=V(y,\tau+s),
   \qquad
   (\Theta_sP)(y,\tau):=P(y,\tau+s).
\]
For a tail center \(x\in\R^3\) observed at centered time \(\sigma\), define the
covariant tail recentering
\[
   \mathscr T_{(x,\sigma)}V(y,s)
   :=V(y+e^{s/2}x,\sigma+s),
   \qquad
   \mathscr T_{(x,\sigma)}P(y,s)
   :=P(y+e^{s/2}x,\sigma+s).
\]
Equivalently,
\[
   \mathscr T_{(x,\sigma)}
   =
   \Theta_\sigma T_{e^{-\sigma/2}x}.
\]
Thus \(\mathscr T_{(x,\sigma)}\) is a composition of exact symmetries of the
centered equation.  Ordinary parabolic recenterings are used below only for local
terminal extraction before one has passed to a centered ancient profile; all
escaping tail limits of centered profiles use \(\mathscr T_{(x,\sigma)}\).
\end{definition}

\begin{lemma}[Semidirect covariance of the centered action]
\label{lem:semidirect-covariance}
The centered time translations and spatial covariant translations satisfy
\[
   \Theta_sT_a\Theta_{-s}=T_{e^{s/2}a},
   \qquad s\in\R,\ a\in\R^3.
\]
Equivalently, if \(G=\mathbb R^3\rtimes\mathbb R\) acts by
\[
   \mathscr G_{(a,s)}:=\Theta_sT_a,
\]
then the group law in this convention is
\begin{equation}\label{eq:semidirect-group-law}
   (a,s)\cdot(b,t)
   =
   (e^{-t/2}a+b,\ s+t),
   \qquad
   \mathscr G_{(a,s)}\mathscr G_{(b,t)}
   =
   \mathscr G_{(a,s)\cdot(b,t)} .
\end{equation}
In particular the subgroup \(\{T_a:a\in\R^3\}\) is normal in the full
spacetime covariant action, and the subspace of observables invariant under all
spatial covariant translations is invariant under every \(\Theta_s\).
\end{lemma}

\begin{proof}
For a profile \(V\),
\[
   (\Theta_sT_a\Theta_{-s}V)(y,\tau)
   =
   V(y+e^{(\tau+s)/2}a,\tau)
   =
   (T_{e^{s/2}a}V)(y,\tau).
\]
The pressure transforms identically.  The asserted invariance of the
spatially invariant observable subspace is the same identity written on the
induced \(L^2\)-representation.
The group law \eqref{eq:semidirect-group-law} follows from
\[
   \Theta_sT_a\Theta_tT_b
   =
   \Theta_{s+t}T_{e^{-t/2}a+b}.
\]
\end{proof}

\begin{theorem}[Covariant observer calculus]
\label{thm:covariant-observer-calculus}
Let \(Z:=\mathbb R^3\times\mathbb R\) denote the centered observer space, with
points \(z=(\xi,\sigma)\).  For \(a\in\mathbb R^3\) and
\(\sigma\in\mathbb R\),
\[
   (T_aU)(y,\tau)=U(y+e^{\tau/2}a,\tau),
   \qquad
   (\Theta_\sigma U)(y,\tau)=U(y,\tau+\sigma).
\]
Then
\[
   \Theta_\sigma T_a=T_{e^{\sigma/2}a}\Theta_\sigma,
   \qquad
   \Theta_\sigma T_a\Theta_{-\sigma}=T_{e^{\sigma/2}a}.
\]
For \(z=(\xi,\sigma)\), define the covariant observer
\[
   \calR_zU
   :=
   \Theta_\sigma T_{e^{-\sigma/2}\xi}U .
\]
Equivalently,
\[
   (\calR_{(\xi,\sigma)}U)(y,s)
   =
   U(y+e^{s/2}\xi,\sigma+s).
\]
The covariant tube pulled back from \(Q_r(0,0)\) is
\[
   \mathcal Q_r^{\rm cov}(\xi,\sigma)
   :=
   \left\{
      (Y,\tau):
      |\tau-\sigma|<r^2,\ 
      |Y-e^{(\tau-\sigma)/2}\xi|<r
   \right\}.
\]
The map
\[
   (y,s)\mapsto(Y,\tau)
   =
   (y+e^{s/2}\xi,\sigma+s)
\]
is a bijection from \(Q_r(0,0)\) onto
\(\mathcal Q_r^{\rm cov}(\xi,\sigma)\) with spacetime Jacobian one.  The same
formulas apply to the pressure, modulo the admissible time-dependent gauges.
\end{theorem}

\begin{proof}
The semidirect identities are
\Cref{lem:semidirect-covariance}; written without conjugation,
\[
   (\Theta_\sigma T_aU)(y,\tau)
   =
   U(y+e^{(\tau+\sigma)/2}a,\tau+\sigma)
   =
   (T_{e^{\sigma/2}a}\Theta_\sigma U)(y,\tau).
\]
Taking \(a=e^{-\sigma/2}\xi\) gives
\[
   (\Theta_\sigma T_{e^{-\sigma/2}\xi}U)(y,s)
   =
   U(y+e^{(s+\sigma)/2}e^{-\sigma/2}\xi,\sigma+s)
   =
   U(y+e^{s/2}\xi,\sigma+s).
\]
If \(Y=y+e^{s/2}\xi\) and \(\tau=\sigma+s\), then
\[
   y=Y-e^{(\tau-\sigma)/2}\xi,\qquad s=\tau-\sigma,
\]
so \(|s|<r^2\) and \(|y|<r\) are equivalent to the defining inequalities for
\(\mathcal Q_r^{\rm cov}(\xi,\sigma)\).  The transformation is a
time-dependent spatial translation, hence has Jacobian one.
\end{proof}

\begin{definition}[Terminal profile class and local topology]
\label{def:terminal-state-space}
For \(M<\infty\), let \(\mathfrak X_M\) be the class of pairs
\((U,\Pi)\), modulo addition of arbitrary functions of time to \(\Pi\), such
that:
\begin{enumerate}[label=\textup{(\roman*)}]
   \item \(U\) is a smooth bounded ancient solution of
   \eqref{eq:centered-main} on \(\mathbb R^3\times\mathbb R\), with
   \(\|U\|_{L^\infty}\le M\);
   \item \((U,\Pi)\) is locally suitable on every compact terminal cylinder;
   \item on every compact cylinder the pressure has an admissible
   \(L^{3/2}\)-representative after subtracting a time-dependent gauge.
\end{enumerate}
The topology on \(\mathfrak X_M\) is
\[
   U_n\to U \quad\hbox{in }C^\infty_{\loc},
   \qquad
   \Pi_n-c_n(\tau)\rightharpoonup \Pi
      \quad\hbox{in }L^{3/2}_{\loc}
\]
for suitable time-dependent gauges \(c_n\).  All profile limits, tail
descendants, concentrating descendant hulls, and recurrent-core limits below are taken in
this topology.
\end{definition}

\begin{definition}[Observer oscillation and spatial oscillation]
\label{def:observer-oscillation}
Let \(Q=B\times I\) be a bounded covariant terminal cylinder.  For a velocity
field \(U\), define the constant-in-time oscillation functional
\[
   \actthree(U;Q)
   :=
   \inf_{c\in\mathbb R^3}\iint_Q |U-c|^3\,dy\,d\tau .
\]
Define also the spatial oscillation modulo time-dependent homogeneous gauges
\[
   \oscsp(U;Q)
   :=
   \inf_{c\in L^3(I;\mathbb R^3)}
      \iint_Q |U(y,\tau)-c(\tau)|^3\,dy\,d\tau .
\]
Thus \(\actthree\) detects nonconstant time-dependent homogeneous modes, while
\(\oscsp\) is the quotient seminorm used to remove spatially homogeneous
parasitic modes.
A covariant observer \(z\in Z\) is \(\eta\)-oscillatory for \(U\) if
\[
   \actthree\bigl(U;\mathcal Q_1^{\rm cov}(z)\bigr)\ge\eta .
\]
Fix once and for all a proper covariant metric \(d_Z\) on \(Z\), and choose a
separated covariant unit-net \(\{z_k\}_{k\in\mathbb N}\subset Z\) such that
\[
   Z\subset \bigcup_{k\in\mathbb N} B_Z(z_k,1/10),
   \qquad
   \#\{k:z_k\in B_Z(z,2)\}\le N_0
\]
for every \(z\in Z\).  The associated covariant unit cylinders therefore have
bounded overlap with constant depending only on \(N_0\).  The oscillation
measure on observer-center sets \(E\subset Z\) is the atomic measure on this
fixed net,
\[
   \mu_U^{\rm osc}(E)
   :=
   \sum_{z_k\in E}
      \actthree\bigl(U;\mathcal Q_1^{\rm cov}(z_k)\bigr) .
\]
For a sequence \(U_n\) we write \(\mu_n^{\rm osc}\) for this measure.  The
fixed packing constant \(N_0\) is used below in the moving-window estimate for
escaping observer regions.  The concentration-set extraction uses \(\actthree\),
because after parasitic-free normalization it detects genuine non-affine
oscillation.  Mesoscopic cylinders use \(\oscsp\), equivalently the functional
\(\mathcal A_n(z,r)\) of
\Cref{def:multiscale-affine-oscillation}, because the minimal-scale quotient
must ignore time-dependent homogeneous gauges.  Thus the concentration-set
construction and the mesoscopic affine quotient are kept distinct.
\end{definition}

\begin{lemma}[Local compactness and continuity of the profile transformations]
\label{lem:state-space-local-compactness}
Let \((U_n,\Pi_n)\subset\mathfrak X_M\) be a sequence with the local suitability
and pressure-gauge bounds supplied by the setup paper on every fixed compact
cylinder.
Then every subsequence has a further subsequence converging in the topology of
\(\mathfrak X_M\).  Moreover the actions \(T_a\), \(\Theta_s\), the covariant
tail recenterings \(\mathscr T_{(x,\sigma)}\) are continuous on
\(\mathfrak X_M\).  The affine maps satisfy the bound-tracking continuity
\[
   S_c:\mathfrak X_M\longrightarrow \mathfrak X_{M+|c|},
\]
again with the pressure interpreted modulo the corresponding transformed time
gauges.
\end{lemma}

\begin{proof}
On each compact cylinder \(K\), the setup local suitability, dissipation,
strong local \(L^3\)-velocity compactness, and pressure-gauge compactness give
a Seregin subsequential limit in the local topology of \(\mathfrak X_M\).
Diagonalizing over the countable compact exhaustion gives the asserted
subsequence.  With the setup pressure-atlas bounds and local suitability fixed,
interior parabolic estimates for \eqref{eq:centered-main} upgrade the velocity
convergence to \(C^\infty_{\loc}\); the pressure representatives converge
weakly in \(L^{3/2}_{\loc}\) after the local gauges.  The formulas defining \(T_a\),
\(\Theta_s\), \(\mathscr T_{(x,\sigma)}\), and \(S_c\) are smooth changes of
variables and affine pressure changes on each fixed compact cylinder.  The
first three preserve the \(L^\infty\) bound \(M\).  The affine map \(S_c\)
subtracts the constant vector \(c\), so its velocity bound is at most
\(M+|c|\).  All four maps preserve the local topology and transform admissible
gauges into admissible gauges.  No global spatial norm or global pressure
representative is used.
\end{proof}

\begin{definition}[Compact normalized descendant profile class]
\label{def:compact-normalized-state-space}
Let \(\mathfrak X_M^\#\subset\mathfrak X_M\) denote the sequential closure, in
the local topology of \(\mathfrak X_M\), of all normalized Seregin profiles from
the setup paper and all profiles obtained from them by the limiting operations used
in this paper: centered time translations, covariant translations, covariant
tail recenterings, terminal profile limits, concentrating descendant limits, recurrent-core
limits, and the canonical affine normalization whenever the affine parameter is
bounded.  The closure is always taken with the local pressure gauges supplied by
the normalized Seregin compactness theorem.
\end{definition}

\begin{lemma}[Compactness and continuity on the normalized profile class]
\label{lem:compact-normalized-state-space}
The space \(\mathfrak X_M^\#\) is compact in the local topology inherited from
\(\mathfrak X_M\).  The maps \(T_a\), \(\Theta_s\), and the covariant
recenterings used to form descendants act continuously on
\(\mathfrak X_M^\#\).  The affine map satisfies the bound-tracking continuity
\[
   S_c:\mathfrak X_M^\#\longrightarrow \mathfrak X_{M+|c|}^\# .
\]
\end{lemma}

\begin{proof}
By definition, every sequence in \(\mathfrak X_M^\#\) is a diagonal limit of
normalized Seregin descendants with the same local \(L^\infty\), local
suitability, pressure-gauge, and interior smooth bounds on each fixed compact
cylinder.  Applying \Cref{lem:state-space-local-compactness} on an exhaustion of
compact cylinders gives a further subsequential limit in \(\mathfrak X_M\).  The
closure defining \(\mathfrak X_M^\#\) contains that limit, so
\(\mathfrak X_M^\#\) is sequentially compact; the local topology is metrizable on
bounded sets by the standard countable exhaustion, hence compact.

Continuity of the limiting operations is the continuity already proved in
\Cref{lem:state-space-local-compactness}, restricted to the closed normalized
profile class.  The affine map changes only the velocity bound and the explicit
affine pressure gauge, so it maps \(\mathfrak X_M^\#\) continuously into
\(\mathfrak X_{M+|c|}^\#\).
\end{proof}

\begin{definition}[Affine homogeneous symmetry]
\label{def:affine-symmetry}
For \(c\in\R^3\), define
\[
   (S_cV)(y,\tau):=V(y-2c,\tau)-c,
   \qquad
   (S_cP)(y,\tau):=P(y-2c,\tau)+\frac12 c\cdot y .
\]
\end{definition}

\begin{lemma}[Affine renormalized Galilean symmetry]
\label{lem:affine-symmetry}
If \((V,P)\) solves \eqref{eq:centered-main}, then
\((S_cV,S_cP)\) also solves \eqref{eq:centered-main}.
\end{lemma}

\begin{proof}
Set \(U(y,\tau)=V(y-2c,\tau)-c\),
\(\Pi(y,\tau)=P(y-2c,\tau)+\frac12c\cdot y\), and \(z=y-2c\).  Then
\[
   \partial_\tau U=(\partial_\tau V)(z,\tau),
   \qquad \Delta U=(\Delta V)(z,\tau),
   \qquad \nabla U=(\nabla V)(z,\tau).
\]
The nonlinear term is
\[
   (U\cdot\nabla)U=(V-c)\cdot\nabla V
   =(V\cdot\nabla)V-c\cdot\nabla V,
\]
where all quantities on the right are evaluated at \((z,\tau)\).  The drift
term satisfies
\[
   \frac12 y\cdot\nabla U
   =\frac12(z+2c)\cdot\nabla V
   =\frac12z\cdot\nabla V+c\cdot\nabla V.
\]
The two \(c\cdot\nabla V\) terms cancel.  Finally,
\[
   \frac12U+\nabla\Pi
   =
   \frac12(V-c)+\nabla P+\frac12c
   =
   \frac12V+\nabla P.
\]
Substitution into \eqref{eq:centered-main} at the point \((z,\tau)\) gives the
equation for \((U,\Pi)\).  The divergence condition is unchanged by translation
and subtracting a constant vector.
\end{proof}

\begin{definition}[Pressure gauges and affine quotient]
\label{def:pressure-gauge-affine-quotient}
On every ordinary local cylinder \(Q\), pressure representatives are identified
modulo time functions:
\[
   P_Q\sim P_Q+c_Q(\tau).
\]
At the affine/parasitic quotient level, and only there, one also identifies
pressure representatives modulo the affine spatial functions paired with
spatially homogeneous parasitic velocities:
\[
   P_Q\sim P_Q+b_Q(\tau)\cdot y+c_Q(\tau).
\]
The affine term is allowed exactly when it corresponds to a removable
homogeneous mode in the velocity equation, as in
\Cref{def:affine-constant-lower-stratum}.  Generic terminal residual branches
are always considered after this quotient has either been fixed by the
canonical parasitic-free representative or assigned to \(\Caff\).
\end{definition}

\begin{lemma}[Pressure atlas compatibility]
\label{lem:paperIV-pressure-atlas-compatibility}
Let \(Q,Q'\) be overlapping local cylinders in a normalized Seregin extraction.
Before affine quotienting, admissible pressure representatives satisfy
\[
   P_Q-P_{Q'}=c_{Q,Q'}(\tau)
\]
on \(Q\cap Q'\).  At the affine/parasitic quotient level they satisfy
\[
   P_Q-P_{Q'}=b_{Q,Q'}(\tau)\cdot y+c_{Q,Q'}(\tau)
\]
with the affine term paired with the corresponding removable homogeneous
velocity mode.  The velocity equation is unchanged after the time-gauge
identification, and after the affine identification it is interpreted in the
quotient equation with the parasitic correction removed or assigned to
\(\Caff\).
\end{lemma}

\begin{proof}
The time-gauge statement is the setup pressure atlas compatibility imported in
\Cref{thm:imported-setup-results}.  The affine statement follows by substituting
a spatially homogeneous velocity \(b(\tau)\) into the centered equation:
the only pressure gradient it can create is
\(-(\partial_\tau b+\frac12b)\), hence an affine pressure in \(y\), modulo a
time function.  This is the quotient convention in
\Cref{def:pressure-gauge-affine-quotient}.  All terminal limiting operations
below choose their pressure gauges before taking compactness limits, and if the
affine term is nontrivial the profile is assigned to \(\Caff\).
\end{proof}

\begin{definition}[Affine/parasitic lower stratum and residual normalization]
\label{def:affine-constant-lower-stratum}
The affine/parasitic lower stratum \(\Caff\) consists of the spatially
homogeneous modes which are invisible to the spatial quotient seminorm
\(\oscsp\).  It has
two equivalent realizations.
\begin{enumerate}[label=\textup{(\roman*)}]
   \item A genuine centered profile belongs to \(\Caff\) if, after one affine
   homogeneous symmetry \(S_c\), its velocity is spatially constant.  In the
   parasitic-free representative this is the zero profile.
   \item A quotient limit belongs to \(\Caff\) if it is represented by a
   locally bounded spatially homogeneous velocity \(b(\tau)\).  In the weak
   centered equation this mode is paired with the affine pressure functional
   \[
      \bigl(-\partial_\tau b(\tau)-\tfrac12 b(\tau)\bigr)\cdot y
   \]
   in \(\mathcal D'_\tau\mathcal D'_y\).  If the mode is realized as a smooth
   suitable profile, this affine pressure is an ordinary local pressure; if it
   arises only as a measurable homogeneous gauge \(c_n(\tau)\), it is used only
   in the quotient weak formulation and never as an independent
   \(L^{3/2}_{\loc}\) pressure.
\end{enumerate}
In the ordered residual decomposition this stratum is part of the previously
closed lower classes, not part of \(\Cgen\).  Equivalently, every generic
residual profile is considered after the canonical affine normalization which
removes its spatially homogeneous mild mode.  Retained non-affine oscillation in
the residual class is therefore always measured after this normalization.
\end{definition}

\begin{remark}[How affine constants are used]
\label{rem:affine-constant-use}
When a compact recurrent tail core is shown below to consist only of spatially
homogeneous parasitic modes, the contradiction is not that an unnormalized
nonzero constant has zero local \(L^3\)-mass.  The contradiction is that these
modes have already been assigned to the closed lower stratum in
\Cref{def:affine-constant-lower-stratum}; after the canonical affine
normalization they are the zero profile and cannot represent a generic residual
profile carrying retained velocity concentration or non-affine oscillation.
\end{remark}

\begin{lemma}[Homogeneous gauges generate only the affine/parasitic stratum]
\label{lem:homogeneous-gauges-are-affine-stratum}
Let \(c_n:I\to\mathbb R^3\) be bounded in \(L^\infty(I)\), and let
\[
   C_n(y,\tau):=c_n(\tau).
\]
Every local weak limit, time translate, or normalized noncompactness limit of
\((C_n)\) has vanishing spatial quotient oscillation on every compact cylinder
and belongs to \(\Caff\).  Conversely, after the canonical
parasitic-free normalization, no nonzero spatially homogeneous mode remains in
the residual profile class.
\end{lemma}

\begin{proof}
For every compact cylinder \(Q\),
\[
   \inf_{c(\tau)}
   \iint_Q |C_n(y,\tau)-c(\tau)|^3\,dy\,d\tau=0,
\]
with the choice \(c(\tau)=c_n(\tau)\).  The same identity is stable under local
weak limits and time recenterings.  Hence every such limit has zero spatial
quotient oscillation on every compact cylinder and is, by
\Cref{def:affine-constant-lower-stratum}, an affine/parasitic quotient mode.
If a homogeneous mode is realized as a smooth centered solution, substituting
\(U=b(\tau)\) into \eqref{eq:centered-main} gives
\[
   \partial_\tau b+\frac12b+\nabla P=0,
\]
so the pressure is affine in \(y\) up to a time gauge, exactly as in the
definition of \(\Caff\).  If \(b\) is only a bounded measurable gauge, the same
identity is understood distributionally and only in the quotient formulation.
The residual representative is defined after removal of this quotient
direction; therefore a nonzero retained effect of \(b\) is assigned to the
closed lower stratum, while the parasitic-free residual profile has no such mode.
\end{proof}

\begin{theorem}[Closed non-affine oscillation functional]
\label{thm:closed-nonaffine-activity}
Let \(Q\Subset\mathbb R^3\times\mathbb R\) be a bounded connected parabolic
cylinder.  For \(U\in\mathfrak X_M\), define
\[
   d_Q(U)
   :=
   \actthree(U;Q)^{1/3}.
\]
Then
\[
   |d_Q(U)-d_Q(V)|
   \le
   \|U-V\|_{L^3(Q)} .
\]
In particular \(d_Q\) is continuous under the local topology of
\(\mathfrak X_M\).  Moreover,
\[
   d_Q(U)=0
   \quad\Longleftrightarrow\quad
   U \hbox{ is spatially and temporally constant on } Q .
\]
Fix once and for all a nested exhaustion
\[
   Q_1\Subset Q_2\Subset\cdots,\qquad
   \bigcup_{j=1}^{\infty}Q_j=\mathbb R^3\times\mathbb R,
\]
by bounded connected parabolic cylinders.  Define
\[
   \Phi(U)
   :=
   \sum_{j=1}^{\infty}2^{-j}\min\{1,d_{Q_j}(U)\}.
\]
Then \(\Phi\) is continuous on \(\mathfrak X_M\), and
\[
   \Phi(U)=0
   \quad\Longleftrightarrow\quad
   U \hbox{ is globally constant}.
\]
Consequently, for every \(\eta>0\), the non-affine oscillation set
\[
   \mathfrak A_\eta
   :=
   \{(U,\Pi)\in\mathfrak X_M:\Phi(U)\ge\eta\}
\]
is closed in the local profile topology.
\end{theorem}

\begin{proof}
For every \(c\in\mathbb R^3\), the triangle inequality in \(L^3(Q)\) gives
\[
   \|U-c\|_{L^3(Q)}
   \le
   \|U-V\|_{L^3(Q)}+\|V-c\|_{L^3(Q)}.
\]
Taking the infimum over \(c\) gives
\[
   d_Q(U)\le \|U-V\|_{L^3(Q)}+d_Q(V).
\]
Interchanging \(U\) and \(V\) proves the Lipschitz estimate.  Local smooth
convergence implies local \(L^3\)-convergence, so \(d_Q\) is continuous in the
topology of \(\mathfrak X_M\).

If \(d_Q(U)=0\), choose \(c_k\in\mathbb R^3\) such that
\(\|U-c_k\|_{L^3(Q)}\to0\).  The sequence \(c_k\) is bounded; otherwise
\[
   \|U-c_k\|_{L^3(Q)}
   \ge |c_k|\,|Q|^{1/3}-\|U\|_{L^3(Q)}
   \to\infty ,
\]
which is impossible.  Passing to a subsequence, \(c_k\to c\), and hence
\(U=c\) almost everywhere on \(Q\).  Since \(U\) is smooth, \(U\equiv c\)
pointwise on \(Q\).  The converse is immediate.

Each summand in the definition of \(\Phi\) is continuous and bounded by
\(2^{-j}\).  The Weierstrass test gives uniform convergence of the defining
series, hence continuity of \(\Phi\).  If \(\Phi(U)=0\), then
\(d_{Q_j}(U)=0\) for every \(j\), so \(U\) is constant on every \(Q_j\).  The
cylinders are nested and overlap, so the constants agree and \(U\) is globally
constant.  The converse is again immediate.  Closedness of
\(\mathfrak A_\eta\) follows from continuity of \(\Phi\).
\end{proof}

\begin{corollary}[Compact normalized non-affine oscillation sets]
\label{cor:compact-normalized-activity-set}
For every \(\eta>0\),
\[
   \mathfrak A_\eta^\#
   :=
   \{(U,\Pi)\in\mathfrak X_M^\#:\Phi(U)\ge\eta\}
\]
is compact in the normalized profile topology.
\end{corollary}

\begin{proof}
By \Cref{lem:compact-normalized-state-space}, \(\mathfrak X_M^\#\) is compact.
By \Cref{thm:closed-nonaffine-activity}, \(\Phi\) is continuous in the local
topology.  Therefore \(\mathfrak A_\eta^\#\) is a closed subset of a compact
space.
\end{proof}

\begin{theorem}[Affine normalization dichotomy]
\label{thm:affine-normalization-dichotomy}
Let \(S_c\) be the affine homogeneous symmetry of
\Cref{def:affine-symmetry}, and suppose the ordered lower strata include all
constant centered profiles as in \Cref{def:affine-constant-lower-stratum}.
For every \((U,\Pi)\in\mathfrak X_M\), exactly one of the following alternatives
holds:
\begin{enumerate}[label=\textup{(\roman*)}]
   \item \(\oscsp(U;Q)=0\) on every compact cylinder \(Q\).  Then \(U\) is a
   spatially homogeneous quotient mode, and \((U,\Pi)\) belongs to the
   affine/parasitic lower stratum \(\Caff\).
   \item For the parasitic-free representative, after an admissible affine
   symmetry \(S_c\), one has \(\Phi(S_cU)>0\).  Equivalently, for some compact
   cylinder \(Q\), \(\actthree(S_cU;Q)>0\).
\end{enumerate}
In the retained normalized residual space, alternative \textup{(i)} has already
been removed by the canonical affine/parasitic normalization, so only
alternative \textup{(ii)} can remain.
Moreover, if \(K\subset\mathfrak X_M\) is compact and
\[
   K\cap\{(U,\Pi):\Phi(U)=0\}=\varnothing ,
\]
then there exists \(\eta_K>0\) such that
\[
   \Phi(U)\ge\eta_K
   \qquad\hbox{for every }(U,\Pi)\in K .
\]
\end{theorem}

\begin{proof}
If \(\oscsp(U;Q)=0\) on every compact cylinder, then on each time slice \(U\) is
spatially constant in the distributional sense.  The local constants are
compatible on overlapping cylinders, so \(U(y,\tau)=b(\tau)\) locally in time
for a measurable bounded function \(b\).  By
\Cref{def:affine-constant-lower-stratum}, this is precisely the
affine/parasitic quotient stratum; if it is realized by a smooth profile, the
corresponding affine pressure is the one described there, and otherwise it is
used only in the quotient weak formulation.

If the spatial oscillation does not vanish identically, \(U\) is not a
spatially homogeneous quotient mode.  After choosing the canonical
parasitic-free representative and applying any admissible affine symmetry
\(S_c\), the resulting velocity is not globally constant.  Hence, by
\Cref{thm:closed-nonaffine-activity}, \(\Phi(S_cU)>0\).  Conversely,
if \(\Phi(S_cU)=0\), then \(S_cU\) is globally constant, and the original
profile is in \(\Caff\).  This proves the dichotomy and the retained normalized
statement.  For the compact statement, \(\Phi\) is continuous and
strictly positive on \(K\), so it attains a positive minimum
\[
   \eta_K:=\min_{(U,\Pi)\in K}\Phi(U)>0 .
\]
\end{proof}

\begin{definition}[Escaping covariant hull]
\label{def:escaping-covariant-hull}
Let \(V\) be a bounded smooth ancient solution of \eqref{eq:centered-main}.
Define
\[
   \calH_\infty(V)
   :=
   \left\{
      W:\ \exists (a_n,s_n)\in G,\ |(a_n,s_n)|_{\rm cov}\to\infty,
      \mathscr G_{(a_n,s_n)}V\to W
      \text{ in } C^\infty_{\loc}(\R^3\times\R)
   \right\}.
\]
Here \(G=\mathbb R^3\rtimes\mathbb R\) is the covariant spacetime group in
\Cref{lem:semidirect-covariance}, and any proper left-invariant metric may be
used for \(|(a,s)|_{\rm cov}\).  Thus the recurrent tail hull is invariant
under the full covariant action, not merely under ordinary spatial translations.
\end{definition}

\begin{theorem}[Compactness of the normalized escaping covariant hull]
\label{thm:escaping-covariant-hull-compactness}
Let \((V,P)\in\mathfrak X_M^\#\) be represented after canonical
affine/parasitic normalization.  Assume that every escaping covariant sequence
for \((V,P)\) has a compact subsequence in the normalized local topology with
the setup pressure atlas fixed; if this assumption fails, the sequence is
assigned instead to the local-compactness or pressure-gauge-loss terminal
alternative.  In the compactness branch, \(\calH_\infty(V)\) is nonempty and
compact in \(C^\infty_{\loc}(\R^3\times\R)\).  Every
\(W\in\calH_\infty(V)\) is a bounded smooth ancient solution of
\eqref{eq:centered-main} with
\(\|W\|_{L^\infty}\le M\).
\end{theorem}

\begin{proof}
For any escaping covariant sequence \((a_n,s_n)\), the transformed solutions
\(\mathscr G_{(a_n,s_n)}V\) solve the same equation and have the same
\(L^\infty\) bound by
\Cref{lem:covariant-translation-invariance,lem:semidirect-covariance}.  The
\(L^\infty\) bound alone is not used to obtain a time modulus or derivative
bound.  By the compactness-branch assumption, the transformed profiles have a
subsequence converging in the normalized pressure-atlas topology, hence locally
smoothly in velocity.  Passing to the limit in the equation gives another
bounded ancient solution.  Nonemptiness follows by applying this compactness
branch to any escaping sequence in \(G\).

It remains to record closedness of the escaping hull.  Suppose
\[
   W_j\in\calH_\infty(V),\qquad W_j\to W
   \quad\hbox{in }C^\infty_{\loc}.
\]
For each \(j\), choose an escaping sequence \(g_{j,k}\in G\) with
\(\mathscr G_{g_{j,k}}V\to W_j\).  Choose \(k(j)\) so large that
\[
   |g_{j,k(j)}|_{\rm cov}\ge j,
   \qquad
   d_j\!\left(\mathscr G_{g_{j,k(j)}}V,W_j\right)\le j^{-1},
\]
where \(d_j\) is a metric for \(C^j\)-convergence on
\(\overline{B_j}\times[-j,j]\).  Then
\[
   \mathscr G_{g_{j,k(j)}}V\to W
   \quad\hbox{in }C^\infty_{\loc},
   \qquad
   |g_{j,k(j)}|_{\rm cov}\to\infty .
\]
Hence \(W\in\calH_\infty(V)\).  This proves closedness.

To see sequential precompactness of the whole hull, let
\(W_j\in\calH_\infty(V)\) be arbitrary.  For each \(j\), choose an escaping
sequence \(g_{j,k}\) with \(\mathscr G_{g_{j,k}}V\to W_j\), and choose
\(k(j)\) so that
\[
   |g_{j,k(j)}|_{\rm cov}\ge j,\qquad
   d_j\!\left(\mathscr G_{g_{j,k(j)}}V,W_j\right)\le j^{-1}.
\]
The sequence \(\mathscr G_{g_{j,k(j)}}V\) is an escaping transformed sequence
in the compactness branch, so it has a locally smooth convergent subsequence.
Along that subsequence, \(W_j\) converges to the same limit.  The closedness
argument places the limit in \(\calH_\infty(V)\).
Thus every sequence in the hull has a convergent subsequence with limit still
in the hull; metrizability of the local smooth topology gives compactness.
\end{proof}

\begin{lemma}[Full covariant invariance of the escaping covariant hull]
\label{lem:escaping-covariant-hull-invariance}
For every \(g\in G\),
\[
   \mathscr G_g\calH_\infty(V)=\calH_\infty(V).
\]
In particular, if \(W\in\calH_\infty(V)\), then both \(T_aW\) and
\(\Theta_sW\) belong to \(\calH_\infty(V)\) for all \(a\in\mathbb R^3\) and
\(s\in\mathbb R\).
\end{lemma}

\begin{proof}
Let \(W=\lim_n\mathscr G_{g_n}V\), with \(|g_n|_{\rm cov}\to\infty\).  Then
\[
   \mathscr G_gW
   =
   \lim_n\mathscr G_g\mathscr G_{g_n}V
   =
   \lim_n\mathscr G_{gg_n}V .
\]
Since the metric on \(G\) is proper and left-invariant, \(|gg_n|_{\rm cov}\to
\infty\).  Hence \(\mathscr G_gW\in\calH_\infty(V)\).  Applying the same
argument to \(g^{-1}\) gives equality.
\end{proof}

\begin{definition}[Tail alternatives]
\label{def:R4-tail-strata}
Let \(\calC\) denote the set of constant vector fields.  Define
the following classes only in the compact normalized hull branch, after the
local-compactness and pressure-gauge-loss alternatives identified in
\Cref{thm:escaping-covariant-hull-compactness}:
\[
   {\Cgen}_{\mathrm{tail\text{-}const}}
   :=
   \{V\in\Cgen:\calH_\infty(V)\subset\calC\},
\]
and
\[
   {\Cgen}_{\mathrm{tail\text{-}prof}}
   :=
   \{V\in\Cgen:\calH_\infty(V)\hbox{ contains a nonconstant profile}\}.
\]
\end{definition}

\begin{proposition}[Tail dichotomy and tail-zero normalization]
\label{prop:tail-dichotomy}
Every \(V\in\Cgen\) belongs to
\({\Cgen}_{\mathrm{tail\text{-}const}}\cup
{\Cgen}_{\mathrm{tail\text{-}prof}}\).  If
\(\calH_\infty(V)=\{C_\infty\}\) is a singleton constant hull, then
\[
   U:=S_{C_\infty}V
\]
solves \eqref{eq:centered-main} and satisfies
\[
   \calH_\infty(U)=\{0\}.
\]
\end{proposition}

\begin{proof}
The dichotomy is the definition.  Since \(T_a\) and \(S_c\) commute,
\[
   T_a(S_{C_\infty}V)=S_{C_\infty}(T_aV).
\]
Thus every far-field limit of \(S_{C_\infty}V\) is
\[
   S_{C_\infty}(C_\infty)=0.
\]
The equation is preserved by \Cref{lem:affine-symmetry}.
\end{proof}

\begin{definition}[Tail-minimal recurrent core]
\label{def:tail-minimal-core}
A nonempty compact set \(\calM\subset\calH_\infty(V)\) is a tail-minimal
recurrent core if it is invariant under all covariant translations \(T_a\),
invariant under all time translations
\[
   (\Theta_sW)(y,\tau):=W(y,\tau+s),
\]
and contains no proper nonempty compact subset with the same two invariance
properties.
\end{definition}

\begin{lemma}[Pressure gradients in invariant local averages]
\label{lem:recurrent-core-pressure-gauges}
Let \(\calM\) be a compact set of smooth bounded solutions of
\eqref{eq:centered-main}, compact in the locally smooth topology and invariant
under the covariant translations \(T_a\) and the time translations
\(\Theta_s\).  Let \(\mu\) be a Borel probability measure on \(\calM\)
invariant under both actions.  For each \(W\in\calM\), let \(P_W\) be any
smooth local pressure near \((0,0)\) satisfying \eqref{eq:centered-main}.  Put
\[
   u_i(W):=W_i(0,0),
   \qquad
   p_i(W):=\partial_iP_W(0,0),
   \qquad
   b_{ij}(W):=u_i(W)u_j(W).
\]
Let \(\Pi_0\) denote the orthogonal projection in
\(L^2(\calM,\mu)\) onto the closed subspace of functions invariant under all
spatial covariant translations \(T_a\), and set
\[
   u_i^0:=\Pi_0u_i,\qquad p_i^0:=\Pi_0p_i.
\]
Then the pressure contribution in the recurrent-core energy average is
gauge-independent and satisfies
\begin{equation}\label{eq:pressure-average-core}
   \int_{\calM}u_i(W)p_i(W)\,d\mu(W)
   =
   \int_{\calM}u_i^0(W)p_i^0(W)\,d\mu(W).
\end{equation}
Equivalently, the nonzero spatial-frequency part of the pressure gradient
does no work in invariant average:
\begin{equation}\label{eq:pressure-nonzero-modes-core}
   \int_{\calM}u_i(W)\bigl(p_i(W)-p_i^0(W)\bigr)\,d\mu(W)=0 .
\end{equation}
\end{lemma}

\begin{proof}
The additive pressure gauge does not enter, because only \(\nabla P_W\) appears.
The pressure gradient is fixed by the velocity through the equation
\[
   \nabla P_W
   =
   \Delta W-\frac12(W+y\cdot\nabla W)
   -\partial_\tau W-(W\cdot\nabla)W .
\]
Thus \(W\mapsto p_i(W)\) is a continuous bounded local observable on
\(\calM\).  The local pressure reconstruction step in
the imported setup results supplies the local pressure representatives; the present
argument uses only their gradients, so the additive gauges never enter.

Let \(U_a\) be the unitary representation of spatial covariant translations
on \(L^2(\calM,\mu)\),
\[
   (U_aF)(W):=F(T_aW).
\]
Let \(D_1,D_2,D_3\) be its skew-adjoint infinitesimal generators.  For a smooth
local observable \(F\), \(D_kF\) is obtained by differentiating
\(F(T_{he_k}W)\) at \(h=0\).  In particular, at \(\tau=0\),
\[
   D_k u_i(W)=\partial_kW_i(0,0).
\]
Hence
\begin{equation}\label{eq:state-div-free}
   D_i u_i=0
   \qquad\hbox{in }L^2(\calM,\mu),
\end{equation}
because each \(W\) is divergence-free.

Taking the spatial divergence of \eqref{eq:centered-main} gives the local
pressure Poisson equation
\begin{equation}\label{eq:local-pressure-poisson-core}
   -\Delta_yP_W=\partial_i\partial_j(W_iW_j).
\end{equation}
Indeed, the divergence of \(\partial_\tau W\), \(\Delta W\), and \(W\) is
zero; the divergence of \(y\cdot\nabla W\) is
\(\partial_i(y_j\partial_jW_i)=\partial_jW_j+y_j\partial_j\partial_iW_i=0\);
and \((W\cdot\nabla)W_i=\partial_j(W_jW_i)\).  Applying the infinitesimal
translation generators at \((0,0)\) gives
\begin{equation}\label{eq:state-pressure-divergence}
   D_i p_i=-D_iD_jb_{ij}.
\end{equation}
The pressure gradient is curl-free, hence
\begin{equation}\label{eq:state-pressure-curl}
   D_kp_i=D_ip_k .
\end{equation}

Let
\[
   A:=-\sum_{k=1}^3D_k^2.
\]
By the spectral theorem for the commuting unitary group \(U_a\), \(A\) is
self-adjoint and nonnegative, and \(\Pi_0\) is the spectral projection
of \(A\) at the eigenvalue \(0\).  Define the nonzero spatial part
\[
   q_i:=p_i-p_i^0=(I-\Pi_0)p_i .
\]
Equations \eqref{eq:state-pressure-divergence} and
\eqref{eq:state-pressure-curl} imply, on the nonzero spectral subspace,
\[
   q_k
   =
   D_kA^{-1}(I-\Pi_0)D_iD_jb_{ij}
\]
in the standard regularized sense.  Explicitly, with
\[
   q_{k,\varepsilon}
   :=
   D_k(\varepsilon+A)^{-1}(I-\Pi_0)D_iD_jb_{ij},
   \qquad \varepsilon>0,
\]
the joint spectral resolution of \(D_1,D_2,D_3\) gives
\[
   q_{k,\varepsilon}\longrightarrow q_k
   \quad\hbox{in }L^2(\calM,\mu)
   \quad\hbox{as }\varepsilon\downarrow0.
\]
The convergence is precisely the multiplier convergence
\[
   \frac{\xi_k\xi_i\xi_j}{\varepsilon+|\xi|^2}
      \longrightarrow
   \frac{\xi_k\xi_i\xi_j}{|\xi|^2}
   \qquad (\xi\ne0),
\]
together with \((I-\Pi_0)\), which removes the point \(\xi=0\).  The
curl-free relation \eqref{eq:state-pressure-curl} says that the nonzero
spatial spectral part of \(p\) is parallel to \(\xi\), and
\eqref{eq:state-pressure-divergence} fixes its longitudinal coefficient;
there is therefore no additional harmonic component.

For each \(\varepsilon>0\), skew-adjointness of \(D_k\) and
\eqref{eq:state-div-free} give
\begin{align*}
   \int_{\calM}u_k q_{k,\varepsilon}\,d\mu
   &=
   \int_{\calM}
      u_kD_k(\varepsilon+A)^{-1}(I-\Pi_0)D_iD_jb_{ij}\,d\mu \\
   &=
   -\int_{\calM}
      (D_ku_k)(\varepsilon+A)^{-1}(I-\Pi_0)D_iD_jb_{ij}\,d\mu
   =0 .
\end{align*}
Passing to the limit in \(L^2\) yields
\[
   \int_{\calM}u_k q_k\,d\mu=0,
\]
which is \eqref{eq:pressure-nonzero-modes-core}.  Since \(p_i^0\) lies in the
range of \(\Pi_0\) and \(u_i-u_i^0\) is orthogonal to that range,
\[
   \int_{\calM}(u_i-u_i^0)p_i^0\,d\mu=0 .
\]
Combining this with \eqref{eq:pressure-nonzero-modes-core} proves
\eqref{eq:pressure-average-core}.
\end{proof}

\begin{theorem}[Rigidity of recurrent profile tail cores]
\label{thm:recurrent-tail-core-rigidity}
Let \(\calM\) be a compact tail-minimal recurrent core for
\eqref{eq:centered-main} arising from the normalized Seregin setup in
the companion setup paper,
with the affine zero mode interpreted by the affine symmetry
\Cref{lem:affine-symmetry}.  Equivalently, after applying one fixed affine
normalization \(S_c\), the spatially homogeneous part of the core satisfies the
bounded mild centered formulation recorded in
the imported setup results.  Then
every \(W\in\calM\) is spatially and temporally constant.
More precisely, there exists \(m\in\R^3\) such that
\[
   W(y,\tau)\equiv m
   \qquad
   \hbox{for every }W\in\calM.
\]
After applying the affine symmetry \(S_m\), the core is reduced to the zero
solution.
\end{theorem}

\begin{proof}
The group generated by covariant translations and time translations is
amenable on the compact invariant set \(\calM\).  Choose an invariant Borel
probability measure \(\mu\) with full support on a minimal component, and write
\[
   \langle F\rangle:=\int_{\calM}F(W)\,d\mu(W)
\]
for local observables.  Let \(P_W\) be a local pressure associated to \(W\).
As in \Cref{lem:recurrent-core-pressure-gauges}, set
\[
   u_i(W):=W_i(0,0),
   \qquad
   p_i(W):=\partial_iP_W(0,0),
\]
and let \(\Pi_0\) be the orthogonal projection onto observables invariant under
all spatial covariant translations.  Write
\[
   u_i^0:=\Pi_0u_i,
   \qquad
   p_i^0:=\Pi_0p_i .
\]
Invariance gives
\[
   \langle \partial_\tau F(W)(0,0)\rangle=0,
   \qquad
   \langle \partial_j F(W)(0,0)\rangle=0
\]
for every smooth local observable \(F\).  These identities follow first for
difference quotients generated by the group action and then by smoothness.

Evaluate \eqref{eq:centered-main} at \((0,0)\).  The drift term vanishes at
\(y=0\), which gives
\[
   \partial_\tau W(0,0)-\Delta W(0,0)+\frac12 W(0,0)
   +(W(0,0)\cdot\nabla)W(0,0)+\nabla P_W(0,0)=0 .
\]
Project this identity onto the spatially invariant subspace.  The spatial
   derivative terms vanish under \(\Pi_0\): the Laplacian is
   \(D_jD_j u_i\), and the nonlinear term is \(D_j(u_ju_i)\) because
   \(D_ju_j=0\).  Time translations do not literally commute with the
   covariant spatial translations; by \Cref{lem:semidirect-covariance} they
   normalize them:
   \[
      \Theta_sT_a\Theta_{-s}=T_{e^{s/2}a}.
   \]
   Hence the subspace of observables invariant under all \(T_a\) is invariant
   under \(\Theta_s\), and the orthogonal projection \(\Pi_0\) commutes with the
   time-translation representation.  Therefore
   \(\Pi_0\partial_\tau u_i=\partial_\tau u_i^0\).  Thus
\begin{equation}\label{eq:zero-mode-momentum-core}
   \partial_\tau u_i^0+\frac12u_i^0+p_i^0=0
   \qquad\hbox{in }L^2(\calM,\mu).
\end{equation}
Taking the inner product of \eqref{eq:zero-mode-momentum-core} with \(u_i^0\),
summing in \(i\), and using invariance under time translations,
\[
   \left\langle u_i^0\partial_\tau u_i^0\right\rangle
   =
   \frac12\left\langle\partial_\tau |u^0|^2\right\rangle
   =0,
\]
gives
\begin{equation}\label{eq:zero-mode-pressure-work-core}
   \left\langle u_i^0p_i^0\right\rangle
   =
   -\frac12\left\langle|u^0|^2\right\rangle .
\end{equation}

Testing the equation against \(W\), evaluating at \((0,0)\), and averaging, the time
term is
\[
   \left\langle W\cdot\partial_\tau W\right\rangle
   =\frac12\left\langle\partial_\tau |W|^2\right\rangle=0 .
\]
The diffusion term gives
\[
   -W\cdot\Delta W
   =
   |\nabla W|^2-\divergence(W\cdot\nabla W),
\]
and therefore
\[
   -\left\langle W\cdot\Delta W\right\rangle
   =
   \left\langle|\nabla W|^2\right\rangle .
\]
The nonlinear term is a divergence:
\[
   W\cdot(W\cdot\nabla)W
   =
   \frac12W\cdot\nabla |W|^2
   =
   \frac12\divergence(|W|^2W),
\]
so its average is zero.  Decompose
\[
   p_i=p_i^0+(p_i-p_i^0).
\]
The nonzero spatial part of the pressure gradient does no work in invariant
average by \Cref{lem:recurrent-core-pressure-gauges}, and the remaining
spatial zero mode is evaluated by
\eqref{eq:zero-mode-pressure-work-core}.  Hence
\[
   \left\langle u_ip_i\right\rangle
   =
   \left\langle u_i^0p_i^0\right\rangle
   =
   -\frac12\left\langle|u^0|^2\right\rangle .
\]
Thus
\begin{equation}\label{eq:mean-energy-core}
   \left\langle|\nabla W|^2\right\rangle
   +\frac12\left\langle|W|^2\right\rangle
   -\frac12\left\langle|u^0|^2\right\rangle=0.
\end{equation}
Both terms are nonnegative.  The second nonnegative quantity is
\[
   \left\langle|u|^2\right\rangle-\left\langle|u^0|^2\right\rangle
   =
   \left\langle|u-u^0|^2\right\rangle,
\]
because \(\Pi_0\) is an orthogonal projection in \(L^2(\calM,\mu)\).
Therefore
\[
   \left\langle|\nabla W|^2\right\rangle=0,
   \qquad
   \left\langle|u-u^0|^2\right\rangle=0 .
\]
Since \(\mu\) has full support and the observables are continuous in the local
smooth topology, \(\nabla W(0,0)=0\) for every \(W\in\calM\).  Invariance
moves \((0,0)\) to any \((y,\tau)\), so \(\nabla W(y,\tau)=0\) everywhere.
Thus every element of \(\calM\) is spatially constant:
\[
   W(y,\tau)=a_W(\tau).
\]
It remains only to identify the allowed spatially homogeneous zero mode.  This
is not a pressure-gauge assertion.  By the normalization in the setup paper
quoted in the statement, there is one fixed affine normalization \(S_c\) for which
the spatially homogeneous representatives satisfy the bounded mild centered
formulation recorded in the imported setup results.  For a spatially constant
mild centered solution \(B(\tau)\), the nonlinear Duhamel term in the mild
formulation vanishes, and the Ornstein--Uhlenbeck--Stokes semigroup acts on
constants by
\[
   S(\tau-\sigma)B(\sigma)=e^{-(\tau-\sigma)/2}B(\sigma).
\]
Thus
\[
   B(\tau)=e^{-(\tau-\sigma)/2}B(\sigma)
   \qquad(\sigma<\tau).
\]
Letting \(\sigma\to-\infty\) and using boundedness of the ancient trajectory
gives \(B(\tau)=0\) for every \(\tau\).  Undoing the fixed affine
normalization gives a single vector \(m=c\) and
\[
   W(y,\tau)\equiv m
   \qquad(W\in\calM).
\]
The affine symmetry \(S_m\) then maps the core to the zero solution.
\end{proof}

\section{Terminal concentration profiles and separated centers}
\label{sec:terminal-concentration}

The terminal part of the proof is local in parabolic cylinders.  It does not require a
global \(L^3\) bound, a global CKN budget, a finite critical-mass budget, or a
global nonlinear profile decomposition for the original residual profile.
Everything is formulated on fixed compact terminal cylinders after applying the
recentering sequence under consideration.

\begin{definition}[Parabolic recentering sequence]
\label{def:parabolic-recentering-sequence}
Let \((V_n,P_n)\) be a terminal extraction sequence on expanding domains
\(D_n\uparrow \R^3\times\R\), with local suitability and pressure-gauge
control.  A parabolic recentering sequence is a sequence
\[
   \mathfrak c=(z_n)_{n\ge1},\qquad z_n=(y_n,\tau_n)\in D_n .
\]
Two recentering sequences are equivalent if
\[
   \sup_n\left(|y_n-y_n'|+|\tau_n-\tau_n'|^{1/2}\right)<\infty,
\]
and asymptotically separated if the same parabolic distance tends to infinity.
The recentered fields associated to \(\mathfrak c\) are denoted
\[
   V_n^{\mathfrak c}(y,s):=V_n(y+y_n,s+\tau_n),
   \qquad
   P_n^{\mathfrak c}(y,s):=
   P_n(y+y_n,s+\tau_n)-a_n^{\mathfrak c}(s),
\]
where \(a_n^{\mathfrak c}\) is an admissible time-dependent local pressure
gauge.  All convergence assertions for a recentering sequence are made on fixed
compact sets in the variables \((y,s)\).
\end{definition}

\begin{definition}[Combined parabolic-covariant descendant recentering]
\label{def:combined-raw-covariant-recentering}
Let \(g_n=(y_n,\tau_n)\) be an ordinary parabolic recentering sequence for a terminal
extraction \((V_n,P_n)\), and let \(\zeta_n=(\xi_n,\sigma_n)\) be a covariant
tail recentering sequence in the centered limit generated by \(g_n\).  The
combined parabolic-covariant descendant recentering is denoted
\[
   d_n=(g_n;\xi_n,\sigma_n)
\]
and acts on the original terminal sequence by
\[
   V_n^{d}(y,s)
   :=
   V_n(y+y_n+e^{s/2}\xi_n,\ s+\tau_n+\sigma_n),
\]
with the pressure pulled back by the same formula and then adjusted by an
admissible time-dependent gauge.  Its relative distance from \(g_n\) is
\[
   |\xi_n|+|\sigma_n|^{1/2}.
\]
Thus \(d_n\) is separated from \(g_n\) exactly when this relative distance tends
to infinity.  This definition is used only to lift descendants of a centered
profile back to the original terminal extraction; ordinary first-generation
terminal recenterings remain the ordinary parabolic recenterings of
\Cref{def:parabolic-recentering-sequence}.
\end{definition}

\begin{definition}[Raw atlas defect measure and raw unit-window active recenterings]
\label{def:terminal-measures}
After choosing local pressure gauges \(a_n(\tau)\), define
\[
   d\mu_n
   :=
   \left(|V_n|^3+|P_n-a_n(\tau)|^{3/2}\right)\,dy\,d\tau,
   \qquad
   d\nu_n:=|\nabla V_n|^2\,dy\,d\tau .
\]
The measure \(\mu_n\) is the atlas-dependent Radon measure from
\Cref{def:local-critical-quantity}.  It is raw, gauge-dependent, and used on
fixed unit observer cylinders for local compactness and concentration-set
extraction.  It is not a scale-invariant CKN measure.  Scale-invariant
threshold and regularity statements at arbitrary radius are written with
\(\Vel_r\), \(D_r\), or \(\calE_r\).
The raw mass condition below is only a selection criterion for observer
windows.  It never certifies that a limiting profile is nonzero: retained
nonzero profiles are certified by \(\Vel_1\) or by the explicitly stated
non-affine oscillation functional.  If the selected raw mass is pressure-only,
it is treated by \Cref{lem:no-pressure-only-retained-profile}.
The local mass of a recentering sequence \(\mathfrak c=(z_n)\) is
\[
   m(\mathfrak c):=\limsup_{n\to\infty}\mu_n(Q_1(z_n)).
\]
For a combined parabolic-covariant descendant recentering
\(d_n=(g_n;\xi_n,\sigma_n)\), the same notation means the mass of the pulled
back defect measure in the combined coordinates:
\[
   m(\mathfrak d):=
   \limsup_{n\to\infty}\mu_n^{\mathfrak d}(Q_1(0,0)),
\]
where \(\mu_n^{\mathfrak d}\) is obtained from \(\mu_n\) by the coordinate
change in \Cref{def:combined-raw-covariant-recentering} and by the admissible
pressure gauges used in that recentered frame.  Thus first-generation
recentering sequences and lifted covariant descendants use the same local
unit-cylinder criterion, expressed in their own local coordinates.
For a fixed raw unit-window threshold \(\eps_{\rm sd}>0\), the recentering sequence
is raw unit-window active at threshold \(\eps_{\rm sd}\) if
\(m(\mathfrak c)\ge\eps_{\rm sd}\), and locally vanishing at unit scale
otherwise.  It has loss of local compactness if
\[
   \limsup_{n\to\infty}\nu_n(Q_2(z_n))=\infty
\]
or if the local pressure gauges fail to be uniformly controlled in
\(L^{3/2}(Q_2(z_n))\).
\end{definition}

\begin{lemma}[Gauge compatibility and retained-velocity semicontinuity]
\label{lem:gauge-compatibility}
Let \((U_n,\Pi_n)\to(U,\Pi)\) in the topology of
\(\mathfrak X_M\) on a fixed compact cylinder containing \(Q_1(0,0)\).  After
choosing the pressure gauges used in that convergence, the following hold.
\begin{enumerate}[label=\textup{(\roman*)}]
   \item The pressure gauges are compatible with centered time shifts,
   covariant translations, covariant tail recenterings, affine normalization,
   and diagonal extraction of descendant profiles: applying any of these operations only
   changes the gauge by another admissible function of time.
   \item The local CKN quantity is lower semicontinuous in the only direction
   used for pressure-inclusive limits:
   \[
      \calE_1(U,\Pi;0)
      \le
      \liminf_{n\to\infty}\calE_1(U_n,\Pi_n;0).
   \]
   Consequently, if the limit profile has
   \(\calE_1(U,\Pi;0)\ge\eta\), then all sufficiently far approximating
   profiles have CKN mass at least any prescribed \(\eta'<\eta\).
   \item If retained nontriviality is carried by the velocity part, then it is
   closed under the same convergence:
   \[
      \liminf_{n\to\infty}\Vel_1(U_n;0)\ge\eta
      \quad\Longrightarrow\quad
      \Vel_1(U;0)\ge\eta .
   \]
   In applications where the lower bound is first obtained for velocity alone
   and then converted into a CKN lower bound, the descendant threshold is chosen
   below the resulting positive constant as in
   \Cref{def:threshold-hierarchy}.
\end{enumerate}
\end{lemma}

\begin{proof}
Time shifts and covariant translations compose the pressure with a smooth local
change of variables; affine normalization adds the explicit linear pressure
\(\frac12c\cdot y\); and every local pressure is defined only up to a
time-dependent function.  Thus admissible gauges remain admissible under all
operations listed in (i).  For (ii), strong local convergence of \(U_n\) gives
\[
   \iint_{Q_1}|U_n|^3\to\iint_{Q_1}|U|^3 .
   \]
For the pressure part, pass to the quotient space
\[
   L^{3/2}(Q_1)\big/
   L^{3/2}((-1,0);\{\hbox{spatial constants}\}).
\]
The quotient norm is convex and weakly lower semicontinuous; equivalently,
after choosing any admissible atlas gauges realizing the local convergence,
\[
   D_1(\Pi;0)\le \liminf_{n\to\infty}D_1(\Pi_n;0).
\]
Adding this to the strong \(L^3\) convergence of the velocity gives the claimed
lower semicontinuity direction for \(\calE_1\).  The velocity-only assertion
follows directly from strong \(L^3(Q_1)\) convergence.  The statement is
entirely local on \(Q_1\).
\end{proof}

\begin{lemma}[No pressure-only retained active limit]
\label{lem:no-pressure-only-retained-profile}
Let \((U_n,\Pi_n)\to(U,\Pi)\) in the local Seregin topology on \(Q_2(0,0)\),
after passing to compatible pressure gauges whenever such gauges are available.
Assume that, on the cylinders where pressure compactness is invoked, the
pressure representatives are the normalized whole-space Riesz pressures modulo
functions of time and modulo the affine/parasitic quotient of
\Cref{def:pressure-gauge-affine-quotient}.
If \(U=0\) on \(Q_2(0,0)\), then
\[
   D_1(\Pi;0)=0 .
\]
Consequently, suppose
\[
   \liminf_{n\to\infty}\calE_1^{\mathfrak a}(U_n,\Pi_n;0)\ge\eta>0
   \qquad\hbox{but}\qquad
   \Vel_1(U;0)=0 .
\]
Then the velocity part of the approximating concentration vanishes in the
limit, and the remaining lower bound is pressure-only.  At least one of the
following ordered alternatives holds; when more than one applies, the proof
uses the first applicable alternative:
\begin{enumerate}[label=\textup{(\roman*)}]
   \item the pressure gauges or pressure masses are not strongly compact in the
   local \(L^{3/2}\) quotient topology on \(Q_2(0,0)\), so the recentering
   belongs to the pressure-loss part of the noncompact terminal alternative;
   \item the limiting pressure defect is affine in space modulo functions of
   time and is assigned to the affine/parasitic quotient stratum \(\Caff\);
   \item the local Riesz part on a slightly larger cylinder carries the
   pressure lower bound, and hence yields a positive velocity lower bound on
   a finite family of neighboring unit cylinders.
\end{enumerate}
Consequently no pressure-only lower bound is used below to certify a retained
nonzero terminal profile.  Retained descendants are certified by \(\Vel_1\), or
by the non-affine oscillation functional \(\Phi\), and pressure-inclusive
\(\calE_1\) is used only after one of the three alternatives above has been
resolved.
\end{lemma}

\begin{proof}
If \(U=0\) on \(Q_2\), then the limiting centered equation gives, in
distributions on \(Q_2\),
\[
   \nabla\Pi
   =
   -\partial_\tau U+\Delta U-\frac12y\cdot\nabla U-\frac12U
   -(U\cdot\nabla)U
   =0 .
\]
Thus \(\Pi=c(\tau)\) on \(Q_2\), and its quotient pressure concentration on
\(Q_1\) vanishes:
\[
   D_1(\Pi;0)=0 .
\]

If the gauges are not compact, or if the pressure representatives are only
weakly compact and retain a nonzero \(L^{3/2}\) defect in the local quotient,
this is alternative \textup{(i)} by definition of the pressure-loss part of the
noncompact terminal alternative.  Assume therefore that, after subtracting the
chosen time gauges, the pressures are strongly compact in
\(L^{3/2}_{\loc}\) modulo spatial constants.  Since
\(\Vel_1(U;0)=0\) and \(U_n\to U\) strongly in \(L^3(Q_1)\), the velocity
contribution to \(\calE_1^{\mathfrak a}(U_n,\Pi_n;0)\) tends to zero.  Any
remaining positive lower bound is therefore represented by the strong
\(L^{3/2}\) quotient pressure limit modulo gauges.

If \(U=0\) on \(Q_2\), the first paragraph shows that this quotient pressure
limit has zero \(D_1\)-mass, a contradiction.  Hence a genuine pressure lower
bound with compact pressure gauges can only survive through either an
affine/parasitic pressure mode, through a local Riesz part generated by
velocity products on the finite enlargement used by the pressure
representation, or through failure of the chosen normalized whole-space
pressure representative to remain compact in the local quotient.

On \(Q_2\), the pressure satisfies the local decomposition
\[
   \Pi_n-a_n(\tau)
   =
   \mathcal R_i\mathcal R_j(U_{n,i}U_{n,j})+h_n,
\]
where \(h_n\) is harmonic in the spatial variables on the smaller cylinder,
after the standard cutoff used in the setup pressure atlas.  If the velocity
products on the finite enlargement used by this cutoff decomposition carry a
positive \(L^{3/2}\) lower bound, then a finite unit-cylinder cover gives
\textup{(iii)}.  Otherwise the Riesz source tends to zero strongly in
\(L^{3/2}\) on that enlargement, and the Riesz part of the pressure converges
strongly to zero in \(L^{3/2}(Q_1)\).  Any remaining nonzero pressure defect is
therefore harmonic in space on the unit cylinder.  Its affine part is
the parasitic pressure mode allowed in \(\Caff\), giving \textup{(ii)}.  If the
harmonic representative is not affine after quotienting by time functions,
then it is not an admissible compact pressure-only remnant under the normalized
whole-space Riesz representation modulo time and affine gauges assumed in the
lemma.  It is therefore recorded as pressure-gauge or pressure-representation
failure, namely alternative \textup{(i)}.  This proves that a retained terminal
profile is never certified by pressure alone.
\end{proof}

\begin{lemma}[Local pressure decomposition and pressure-only alternatives]
\label{lem:local-pressure-decomposition}
Let \((u_n,p_n)\) be locally suitable on a parabolic cylinder
\(Q_4=B_4\times(-16,0)\) with \(\|u_n\|_{L^\infty(Q_4)}\le M\) and pressure
gauges \(a_n(t)\) compatible with \Cref{def:pressure-gauge-affine-quotient}.
Fix a compactly supported cutoff \(\chi\in C_c^\infty(B_3)\) with
\(\chi\equiv1\) on \(B_2\).  Then on \(B_2\)
\begin{equation}\label{eq:local-pressure-split}
   p_n-a_n(t)=p_n^{\rm loc}+h_n,
   \qquad
   p_n^{\rm loc}:=\mathcal R_i\mathcal R_j(\chi u_{n,i}u_{n,j}),
\end{equation}
where \(h_n\) is harmonic in the spatial variables on \(B_2\) for almost
every \(t\in(-4,0)\) and depends only on the values of
\((u_n,p_n-a_n)\) on \(Q_4\setminus Q_2\) (the boundary data of the cutoff
representation).  The following four properties hold.
\begin{enumerate}[label=\textup{(\roman*)}]
   \item \emph{Local Riesz strong continuity.} If
   \(u_n\to0\) strongly in \(L^3(Q_4)\), then
   \(p_n^{\rm loc}\to0\) strongly in \(L^{3/2}(Q_2)\), with
   \(\|p_n^{\rm loc}\|_{L^{3/2}(Q_2)}\le C\|u_n\|_{L^3(Q_4)}^2\).
   \item \emph{Harmonic moment compactness.} For every multi-index
   \(\alpha\) with \(|\alpha|\le K\) and every Schwartz test function
   \(\psi\) compactly supported in \(B_1\),
   \[
      t\mapsto \int (\partial^\alpha h_n)(x,t)\,\psi(x)\,dx
   \]
   is bounded uniformly in \(n\) by a constant depending only on
   \(\|u_n\|_{L^3(Q_4)}\), \(\|p_n-a_n\|_{L^{3/2}(Q_4)}\), \(K\), and
   \(\psi\).
   \item \emph{Affine quotient record.} The quotient of \(h_n\) by spatial
   affine functions
   \[
      h_n(x,t)\sim h_n(x,t)+(\alpha_n(t)+\beta_n(t)\cdot x)
   \]
   is bounded in any local norm by the same data; the
   quotient \([h_n]\in D_1\) is the affine quotient of
   \Cref{def:pressure-gauge-affine-quotient}.  If the affine part of \(h_n\)
   is the only carrier of a positive pressure lower bound, the corresponding
   recentering belongs to the affine/parasitic stratum \(\Caff\).
   \item \emph{Non-affine harmonic record.} If \([h_n]\) has a nonzero
   limit \([h]\) in the quotient \(D_1\), then either
   \([h]\) is purely affine (case \textup{(ii)} of
   \Cref{lem:no-pressure-only-retained-profile}), or the non-affine harmonic
   remainder is not used as retained activity.  If a nonzero local Riesz part
   is present on the finite enlargement, it gives neighboring velocity
   activity by \textup{(i)}.  If no such Riesz part is present, the
   non-affine harmonic remainder is recorded as failure of the normalized
   whole-space pressure representative, modulo time and affine gauges, to be
   compact in the local quotient; this is the pressure/noncompact alternative.
\end{enumerate}
The conclusion of \textup{(iv)} is the precise mechanism behind the reduction
of the pressure-only case: only the local Riesz component is allowed to produce
neighboring velocity activity.  A genuinely non-affine harmonic quotient is
not converted into a velocity source by itself; outside the affine quotient it
is treated as pressure-representation or pressure-compactness failure.
\end{lemma}

\begin{proof}
The decomposition \eqref{eq:local-pressure-split} is the standard local
pressure cutoff: writing
\[
   -\Delta(p_n-a_n)=\partial_i\partial_j(u_{n,i}u_{n,j})
\]
on \(Q_4\) and applying the Riesz operators
\(\mathcal R_i\mathcal R_j\) to the cutoff right-hand side
\(\partial_i\partial_j(\chi u_{n,i}u_{n,j})\) gives a local representative of
the velocity-product part of the pressure, and the difference
\(h_n=p_n-a_n-p_n^{\rm loc}\) is then a harmonic function in space on
\(B_2\) at almost every time.  This is the same construction used in the
setup paper for the local pressure atlas
(\Cref{thm:imported-setup-results}\,(3)).

\textup{(i)} The boundedness of Riesz transforms on \(L^{3/2}\) gives
\[
   \|p_n^{\rm loc}\|_{L^{3/2}(Q_2)}
   \le C\|\chi u_{n,i}u_{n,j}\|_{L^{3/2}(Q_4)}
   \le C\|u_n\|_{L^3(Q_4)}^2 .
\]
Strong \(L^3\)-vanishing of \(u_n\) on \(Q_4\) thus forces strong
\(L^{3/2}\)-vanishing of \(p_n^{\rm loc}\) on \(Q_2\).

\textup{(ii)} The standard interior estimates for harmonic functions in
\(B_2\) bound every spatial derivative of \(h_n\) on \(B_1\) by an
\(L^{3/2}\)-norm of \(h_n\) on \(B_2\), which in turn is bounded by
\(\|p_n-a_n\|_{L^{3/2}(Q_4)}+\|p_n^{\rm loc}\|_{L^{3/2}(Q_2)}\).  The
displayed bound on time-integrated moments follows.

\textup{(iii)} The affine subspace of harmonic functions on \(B_2\) is
finite-dimensional; the projection onto it is bounded for any local norm.
The quotient \([h_n]\) is therefore bounded.  By the affine pressure quotient
fixed in \Cref{def:pressure-gauge-affine-quotient}, the affine part of \(h_n\)
is the parasitic pressure mode allowed in \(\Caff\); a recentering with
positive lower bound carried only by this mode is the affine/parasitic
recentering and is recorded in \(\Caff\).

\textup{(iv)} Assume the limit \([h]\) is non-affine in the quotient
\(D_1\).  A harmonic function may be non-affine while still having no local
Poisson source in \(B_2\); therefore this fact alone is not used to produce a
velocity lower bound.  If the localized Riesz component
\(p_n^{\rm loc}\) has a nonzero lower bound on the finite enlargement entering
the cutoff, \textup{(i)} and boundedness of the Riesz transforms give a
neighboring unit cylinder with positive velocity-product mass, hence the
velocity alternative of \Cref{lem:no-pressure-only-retained-profile}.  If
\(p_n^{\rm loc}\to0\) strongly and \([h]\) is non-affine, then the surviving
pressure lower bound is neither affine/parasitic nor represented by the
normalized whole-space Riesz pressure modulo the allowed gauges.  It is
therefore assigned to the pressure-representation or pressure-compactness
failure alternative, not to retained velocity concentration.
\end{proof}

\begin{definition}[Local-window condition for first-generation recenterings]
\label{def:first-generation-local-window-condition}
Let \((V_n,P_n)\) be a first-generation terminal sequence obtained by undoing
the setup paper local Type~I rescaling on a fixed physical cylinder
\(\mathcal Q_{\rm loc}\).  A parabolic recentering sequence
\(\mathfrak c=(z_n)\), \(z_n=(y_n,\tau_n)\), satisfies the
\emph{local-window condition} if, for every compact cylinder
\[
   K_{R,S}:=B_R(0)\times[-S,S]\Subset\mathbb R^3\times\mathbb R,
\]
the cylinder obtained by writing \(K_{R,S}\) in the \(\mathfrak c\)-centered
coordinates and then undoing the setup rescaling is contained in
\(\mathcal Q_{\rm loc}\) for all sufficiently large \(n\).  Equivalently, every
fixed compact cylinder in the recentered variables lies in the part of the
first-generation extraction where the setup pressure atlas, local energy
inequality, and endpoint Type~I bounds are available.
\end{definition}

\begin{lemma}[Local-window compactness for first-generation recenterings]
\label{lem:first-generation-local-window-compactness}
Let \((V_n,P_n)\) be a first-generation admissible terminal sequence from the
setup extraction, and let \(\mathfrak c=(z_n)\) satisfy the local-window
condition of \Cref{def:first-generation-local-window-condition}.  If
\[
   \liminf_{n\to\infty}\Vel_1(V_n;z_n)>0,
\]
then the recentered sequence admits a subsequence and admissible pressure gauges
for which
\[
   (V_n^{\mathfrak c},P_n^{\mathfrak c})
   \longrightarrow (U,\Pi)
\]
locally in the Seregin topology.  The limit is locally suitable, pressure
normalized, bounded on compact terminal windows by the imported setup bounds,
and nonzero:
\[
   \Vel_1(U;0)>0
\]
after lowering the velocity threshold once.
\end{lemma}

\begin{proof}
Fix \(R,S<\infty\).  By the local-window condition, the physical image of
\(K_{2R,2S}\) in the \(\mathfrak c\)-frame lies inside the setup local Type~I
cylinder for all large \(n\).  The imported setup results therefore gives the
local energy, dissipation, velocity \(L^3\), and pressure-gauge bounds on this
physical image, uniformly in \(n\), after the pressure gauges supplied by the
setup atlas.  Rescaling back to the \(\mathfrak c\)-coordinates gives the
standard Seregin compactness bounds on \(K_{R,S}\): strong local \(L^3\)
compactness for the velocity, weak local \(L^{3/2}\) compactness for pressure
modulo time functions, and stability of the local energy inequality.

Diagonalizing over \(R,S\in\mathbb N\) gives a locally suitable terminal limit.
The retained lower bound is carried by the velocity part, so strong
\(L^3_{\loc}\) convergence gives a positive \(\Vel_1\) lower bound in the
limit after the fixed threshold loss.  No compactness is asserted outside the
local-window region.
\end{proof}

\begin{lemma}[Covariant compactness or pressure-atlas loss for descendants]
\label{lem:ancient-profile-covariant-compactness}
Let \((U,\Pi)\in\mathfrak X_M^\#\) be represented after canonical
affine/parasitic normalization.  For every covariant recentering sequence
\(g_n\in G\), either the transforms
\[
   \mathscr G_{g_n}(U,\Pi)
\]
have no compact closure in the normalized local topology with the setup
pressure atlas fixed, in which case the sequence is assigned to the
local-compactness or pressure-gauge-loss terminal alternative, or they admit a
locally compact pressure-gauge subsequence in the Seregin topology.  In the
compactness branch, if
\[
   \liminf_{n\to\infty}\Vel_1(\mathscr G_{g_n}U;0)>0,
\]
then the limiting descendant is nonzero.
\end{lemma}

\begin{proof}
The covariant action is an exact symmetry of the centered equation by
\Cref{lem:semidirect-covariance}.  Hence the transformed velocities remain
bounded by \(\|U\|_{L^\infty}\) and solve the same centered equation on all of
\(\mathbb R^3\times\mathbb R\).  This boundedness alone is not used as a
compactness theorem.  If the setup pressure-atlas representatives and local
compactness bounds fail to have a compact subsequence, this is precisely the
noncompact/pressure-gauge terminal alternative.  Otherwise the transformed
profiles lie in a compact subset of \(\mathfrak X_M^\#\), so a diagonal
subsequence converges in the normalized local topology to a bounded ancient
suitable descendant.  Strong local \(L^3\) convergence transfers any velocity
lower bound to the limit.
\end{proof}

\begin{lemma}[Escaping first-generation recenterings are assigned to terminal alternatives]
\label{lem:escaping-first-generation-alternatives}
Let \((V_n,P_n)\) be a first-generation terminal sequence and let
\(\mathfrak c=(z_n)\) carry retained unit velocity concentration:
\[
   \liminf_{n\to\infty}\Vel_1(V_n;z_n)>0 .
\]
If \(\mathfrak c\) fails the local-window condition, then it is not handled
by local compactness.  Instead it is assigned, according to the
compact-window-to-expanding-region protocol of
\Cref{lem:compact-window-expanding-region-protocol}, to one of the terminal
exterior alternatives: an exterior retained velocity profile, a diffuse defect,
a pressure/noncompact defect, or a critical-tail defect.
\end{lemma}

\begin{proof}
Failure of the local-window condition means that some fixed recentered compact
cylinder leaves the physical setup cylinder after undoing the first-generation
scaling.  The imported local Seregin compactness theorem is therefore not
available on that cylinder, so the recentering cannot be treated by
\Cref{lem:first-generation-local-window-compactness}.  The retained unit
velocity concentration gives a positive residual measure in the corresponding selected
expanding terminal region.  Applying
\Cref{lem:compact-window-expanding-region-protocol} to that region selects one
of the protocol's priority alternatives: a selected unit observer with velocity
concentration, a residual measure escaping to the observer boundary, a failure
of local compactness or pressure-gauge compactness, or a nonzero normalized
critical-tail scale.  These are reduction alternatives, not compactness
conclusions.
\end{proof}

\begin{corollary}[Local compactness or alternative selection for retained recenterings]
\label{lem:retained-recentering-alternatives}
Let \((V_n,P_n)\) be a terminal sequence obtained either as a first-generation
setup sequence or as a descendant of an already extracted bounded
ancient profile.  Any recentering sequence carrying retained unit velocity concentration is
of exactly one of the following types.
\begin{enumerate}[label=\textup{(\roman*)}]
   \item It is first-generation and satisfies the local-window condition; then
   \Cref{lem:first-generation-local-window-compactness} gives a nonzero locally
   compact Seregin limit.
   \item It acts on an already extracted normalized ancient profile and its
   pressure-atlas orbit has compact closure; then
   \Cref{lem:ancient-profile-covariant-compactness} gives a nonzero locally
   compact descendant.
   \item It is either first-generation and fails the local-window condition, or
   it is an ancient-profile recentering whose pressure-atlas orbit has no
   compact closure; then
   \Cref{lem:escaping-first-generation-alternatives,lem:ancient-profile-covariant-compactness}
   assigns it to an exterior, diffuse, pressure/noncompact, or critical-tail
   terminal branch.
\end{enumerate}
Thus retained unit velocity concentration under an arbitrary recentering is never used by itself as a
compactness hypothesis.
\end{corollary}

\begin{proof}
First-generation recenterings either satisfy or fail the local-window condition.
The two cases are handled by
\Cref{lem:first-generation-local-window-compactness,lem:escaping-first-generation-alternatives}.
Recenterings of extracted ancient profiles act on a bounded ancient profile defined on the full
centered ancient space.  They are handled by
\Cref{lem:ancient-profile-covariant-compactness}: compact pressure-atlas orbits
give locally compact descendants, while noncompact or pressure-gauge-loss
orbits are assigned to the corresponding terminal alternative.  These cases exhaust the
recenterings used in the terminal extraction arguments.
\end{proof}

\begin{theorem}[Local concentration-set extraction]
\label{thm:maximal-active-family}
Let \(Z=\mathbb R^3\times\mathbb R\) be the centered covariant observer space,
and let \(K_m\Subset Z\), \(m\ge1\), be a compact exhaustion.  Let
\(\mu_n\) be the raw atlas defect measures of a terminal sequence, with local finite
mass on compact observer windows:
\[
   \sup_n\mu_n(K_m^+)<\infty
   \qquad(m\ge1),
\]
where \(K_m^+\) is a compact enlargement such that every covariant unit
cylinder \(\mathcal Q_1^{\rm cov}(z)\), \(z\in K_m\), lies in \(K_m^+\).
The compact-window concentration sets are
\[
   \widetilde A_{n,m}^{\eta}
   :=
   \{z\in K_m:\mu_n(\mathcal Q_1^{\rm cov}(z))\ge\eta\}.
\]
For the Hausdorff--Fell extraction we use the equivalent smoothed version below;
thresholds are adjusted only by fixed cutoff constants.  Fix
\(\chi\in C_c^\infty(Q_1(0,0))\), \(0\le\chi\le1\), with
\(\chi\equiv1\) on \(Q_{1/2}(0,0)\).  For \(z\in Z\), put
\[
   \chi_z(Y,\tau)
   :=
   \chi\bigl(\calR_z^{-1}(Y,\tau)\bigr),
   \qquad
   M_n(z):=\int \chi_z\,d\mu_n .
\]
For \(\eta>0\), define the compact-window concentration set
\[
   A_{n,m}^{\eta}
   :=
   \{z\in K_m:M_n(z)\ge\eta\}.
\]
After passing to a subsequence, for every \(m\) and every rational
\(\eta>0\), the compact sets \(A_{n,m}^{\eta}\) converge in the
Hausdorff--Fell topology to compact sets \(A_m^\eta\subset K_m\).  Moreover,
for every open neighborhood \(O\subset K_m\) of \(A_m^\eta\),
\[
   \limsup_{n\to\infty}
   \sup_{z\in K_m\setminus O}M_n(z)
   \le \eta .
\]
Thus outside every neighborhood of the concentration set \(A_m^\eta\), no
\(\eta\)-concentrating covariant unit tube remains in the compact observer window
\(K_m\).  The concentration set is the closed locus \(A_m^\eta\), not a finite or
countable enumeration of tubes.
\end{theorem}

\begin{proof}
For fixed \(n\), the map \(z\mapsto\chi_z\) is continuous in \(C_c\) on
\(K_m\), and all supports are contained in the compact enlargement \(K_m^+\).
Since \(\mu_n\) is a Radon measure finite on \(K_m^+\), \(M_n\) is continuous
on \(K_m\).  Therefore \(A_{n,m}^{\eta}\) is compact.

The hyperspace of compact subsets of the compact metric space \(K_m\) is
compact in the Hausdorff--Fell topology.  Hence, for each fixed pair
\((m,\eta)\), a subsequence exists along which \(A_{n,m}^{\eta}\) converges.
A diagonal argument over \(m\in\mathbb N\) and
\(\eta\in\mathbb Q_+\) gives simultaneous convergence.

Let \(O\subset K_m\) be an open neighborhood of \(A_m^\eta\).  If the stated
residual bound failed, there would exist \(\delta>0\), a subsequence
\(n_k\), and points \(z_{n_k}\in K_m\setminus O\) such that
\[
   M_{n_k}(z_{n_k})\ge\eta+\delta .
\]
After replacing \(O\) by a slightly smaller open neighborhood of \(A_m^\eta\),
we may assume \(K_m\setminus O\) is compact.  Passing to a further subsequence,
\[
   z_{n_k}\to z_*\in K_m\setminus O .
\]
For all large \(k\),
\[
   z_{n_k}\in A_{n_k,m}^{\eta+\delta/2}.
\]
Hausdorff--Fell convergence gives \(z_*\in A_m^{\eta+\delta/2}\).  Since
\(A_m^{\eta+\delta/2}\subset A_m^\eta\subset O\), this contradicts
\(z_*\in K_m\setminus O\).  The residual estimate follows.  All bounds used
here are local on \(K_m^+\); no global CKN budget or countability of concentrating
profiles is assumed.
\end{proof}

\begin{lemma}[Closed-active-locus reduction]
\label{lem:closed-active-locus-reduction}
Let \(A_m^\eta\subset K_m\) be the compact concentration loci produced by
\Cref{thm:maximal-active-family} for a rational threshold
\(\eta\in(0,\eps_{\CKN}/2)\) and a compact observer window \(K_m\Subset Z\).
Then \(A_m^\eta\) is a compact subset of \(K_m\) for which the following
reduction holds.  For every open neighborhood \(O\Supset A_m^\eta\) in
\(K_m\) there is \(N(O)\in\mathbb N\) such that for all \(n\ge N(O)\) the
following hold simultaneously.
\begin{enumerate}[label=\textup{(\arabic*)}]
   \item \emph{Sole carrier of compact retained concentration.}  Every
   covariant unit cylinder
   \(\mathcal Q_1^{\rm cov}(z)\subset K_m\) with center \(z\in K_m\setminus O\)
   carries less than \(\eta\) defect mass:
   \(M_n(z)<\eta\).
   \item \emph{Locality of the residual.}  The residual measure
   \(\mu_n\lfloor (K_m\setminus O)\) satisfies
   \(\sup_{z\in K_m\setminus O} M_n(z)<\eta\); equivalently, no covariant unit
   cylinder centered in \(K_m\setminus O\) is \(\eta\)-concentrating.
   \item \emph{No additional compact concentration locus.}  Any further
   Hausdorff--Fell limit point of compact \(\eta\)-concentration sets in
   \(K_m\) is contained in \(A_m^\eta\); the compact concentration locus is
   the unique maximal closed locus of compact retained concentration in
   \(K_m\).
\end{enumerate}
Consequently, the only candidate compact carriers of retained terminal
concentration on \(K_m\) at threshold \(\eta\) are the points of \(A_m^\eta\),
and the trichotomy alternatives of
\Cref{def:active-removed-residual-measure} (vanishing, diffuse, or
noncompact residual) describe the entirety of the residual after a
\(\rho\)-neighborhood of \(A_m^\eta\) is removed.
\end{lemma}

\begin{proof}
Compactness of \(A_m^\eta\) is recorded in
\Cref{thm:maximal-active-family}.  Statement \textup{(1)}--\textup{(2)} is
the displayed residual estimate of that theorem applied with neighborhood
\(O\) of \(A_m^\eta\).  Statement \textup{(3)} follows because every other
compact \(\eta\)-concentration limit \(A'\subset K_m\) along the same
extraction is, by Hausdorff--Fell limit-point characterization, contained in
\(A_m^\eta\): otherwise a point \(z'\in A'\setminus A_m^\eta\) would carry a
covariant unit defect mass \(\ge\eta\) and could not be evicted from
\(K_m\setminus O\) for \(O\) any open neighborhood of \(A_m^\eta\),
contradicting \textup{(1)}.  The trichotomy then exhausts the residual on
\(K_m\setminus U_\rho(A_m^\eta)\) by definition of
\Cref{def:active-removed-residual-measure}.
\end{proof}

\begin{lemma}[Closed active loci reduce to single, finite separated, or chain alternatives]
\label{lem:closed-active-locus-single-finite-chain}
Let \(A_m^\eta\subset K_m\) be a nonempty compact active locus from
\Cref{thm:maximal-active-family}.  For each \(z\in A_m^\eta\), choose
centers \(z_n\in A_{n,m}^{\eta'}\), with \(\eta'<\eta\), converging to \(z\)
in the Hausdorff--Fell realization and extract the corresponding locally
compact retained velocity profile whenever the local-window/ancient-profile
compactness hypotheses of \Cref{lem:retained-recentering-alternatives} hold.  Declare two
realized centers preliminarily related if their realizing recentering sequences
have bounded covariant parabolic distance and converge, after the admissible
pressure gauges and after passing to a common subsequence, to the same retained
terminal profile class.  The non-separated local profile equivalence relation
is the transitive closure of this preliminary relation.  Then the compact active
locus contributes to the terminal assembly in exactly one of the following
priority ways.
\begin{enumerate}[label=\textup{(\roman*)}]
   \item All realized centers in \(A_m^\eta\) lie in one equivalence class; the
   compact active locus represents a single retained profile class.
   \item There are at least two separated equivalence classes; choosing one
   realizing sequence from each of finitely many such classes gives a finite
   separated family of retained velocity concentration profiles in the sense of
   \Cref{def:finite-separated-profile-family}.
   \item After every finite separated subfamily has been selected, a further
   retained active descendant remains separated from the selected family.
   Iterating the realized descendant construction and using
   \Cref{lem:diagonal-lifting-descendants,thm:descendant-heredity} produces an
   infinite retained concentration chain.
\end{enumerate}
Thus a continuum active set is not a separate terminal alternative: after
quotienting by non-separated local profile equivalence, it enters the terminal
assembly as a single profile class, a finite separated family, or an infinite
retained chain.
\end{lemma}

\begin{proof}
The relation is an equivalence relation by construction.  If all realized
centers are equivalent, the active locus has one non-separated local profile
class, proving \textup{(i)}.

If not, choose two inequivalent realized center sequences.  Their covariant
distance cannot stay bounded along a subsequence that realizes the same local
profile class, because then they would be related by the preliminary relation.
Thus, after passing to the realizing subsequences defining the inequivalent
classes, they are separated in the sense of
\Cref{def:finite-separated-profile-family}.  The same argument selects any
finite number of separated inequivalent classes, giving \textup{(ii)} whenever
the process stops after finitely many selections.

If the process does not stop, construct inductively a sequence of retained
active descendants, each separated from the finite family already chosen.  The
diagonal lifting lemma realizes each descendant as a recentering of the
original terminal extraction, and descendant heredity keeps the profile in the
same normalized retained residual class unless it has already fallen into a
closed lower stratum.  Removing the closed lower strata, the inductive sequence
is an infinite retained concentration chain.  This is \textup{(iii)}.
\end{proof}

\begin{definition}[Residual measure after concentration-set removal]
\label{def:active-removed-residual-measure}
Let \(U_\rho(A_{n,m}^{\eta})\) be the \(\rho\)-neighborhood of the concentration set in
the compact observer window \(K_m\).  The residual measure after deleting
positive-threshold concentration neighborhoods is
\[
   \mu_{n,m,\eta,\rho}^{\rm rem}
   :=
   \mu_n\lfloor
   \left(K_m\setminus U_\rho(A_{n,m}^{\eta})\right).
\]
The terminal trichotomy uses this restricted measure:
vanishing means
\(\mu_{n,m,\eta,\rho}^{\rm rem}(K_m)\to0\) along the selected extraction;
diffuse defect means positive residual mass remains while no unit cylinder
exceeds the concentration threshold after deletion; noncompact defect means that an
escaping residual recentering loses local compactness or pressure-gauge
control under its own normalization.
\end{definition}

\begin{definition}[Canonical countable family of expanding terminal regions]
\label{def:canonical-expanding-family}
The \emph{canonical countable family} \(\mathcal E\subset 2^Z\) is the closure
under finite intersection of the family of regions of one of the following
elementary types:
\begin{enumerate}[label=\textup{(\roman*)},leftmargin=2.2em]
   \item exterior balls \(\{|y|>N\}\) with \(N\in\mathbb N\);
   \item dyadic spatial annuli \(\{2^k<|y|<2^{k+1}\}\) with
      \(k\in\mathbb Z\);
   \item time slabs \(\{|\tau-q|<\ell\}\) with \(q\in\mathbb Q\) and
      \(\ell\in\mathbb Q_{>0}\);
   \item logarithmic annuli \(\{2^k<|y-y_0|<2^{k+1}\}\) with
      \(y_0\in\mathbb Q^3\) and \(k\in\mathbb Z\),
\end{enumerate}
together with their images under finitely many covariant observer shifts
\(T_a\) with \(a\in\mathbb Q^3\rtimes\mathbb Q\) and parabolic rescalings
\(\Theta_s\) with \(s\in\mathbb Q\).  Equivalently, \(\mathcal E\) is the set
of all finite intersections of the rationally parametrized images of
\textup{(i)}--\textup{(iv)}.  This family is countable.

Every concrete exterior or tail window used in the residual proof is of one
of the following two forms: \textup{(a)} an element of \(\mathcal E\) up to
covariant action; \textup{(b)} a multiplicative or additive scale-shift of an
element of \(\mathcal E\) by an index-dependent parameter \(R_n\to\infty\) or
\(R_n\downarrow0\).  In case \textup{(a)} the inclusion in \(\mathcal E\) is
direct; in case \textup{(b)} the relevant scale-shifted region is one of the
canonical dyadic annuli or logarithmic annuli of \(\mathcal E\) after
substitution \(R_n\mapsto2^{k(n)}\) for an integer index \(k(n)\) chosen
within a bounded distance from \(\log_2R_n\), at the cost of an absolute
multiplicative loss in the captured mass.

For convenience we write \(E_n^{(j)}=g_n^{(j)}\cdot E^{(j)}\), where
\(E^{(j)}\in\mathcal E\) and \(g_n^{(j)}\) is the index-dependent covariant
shift or scale change selecting the \(j\)-th window at time \(n\).
\end{definition}

\begin{theorem}[Compact-window to expanding-region terminal protocol]
\label{lem:compact-window-expanding-region-protocol}
Let \(K_m\uparrow Z\) be the compact observer exhaustion used in
\Cref{thm:maximal-active-family}, and let \(\mu_n\) be the raw terminal defect
measures.  Let \(\mathcal E\) be the canonical countable family of expanding
terminal regions of \Cref{def:canonical-expanding-family}, indexed by
\(j\in\mathbb N\), and write \(E_n^{(j)}\subset Z\) for its members along the
extraction.  After one diagonal subsequence, for every \(j\), every rational
\(\eta>0\), and every rational \(\rho>0\), and for every CKN threshold
\(\eta\le\eps_{\CKN}/2\), exactly one of the following \emph{priority
alternatives} occurs in the order listed for the concentration-deleted
measures
\[
   \mu_n\lfloor\bigl(E_n^{(j)}\setminus U_\rho(A_{n,m}^{\eta})\bigr).
\]
\begin{enumerate}[label=\textup{(\alph*)}]
   \item \emph{Expanding-region vanishing.}  For every \(m\),
   \[
      \mu_n\bigl(E_n^{(j)}\cap K_m\setminus U_\rho(A_{n,m}^{\eta})\bigr)
      \longrightarrow0,
   \]
   and the remaining mass outside \(K_m\) tends to zero as \(m\to\infty\) along
   the selected protocol.
   \item \emph{Compact subthreshold regular residual.}
   A positive amount of residual mass remains in some fixed compact
   observer window \(K_m\) after deleting \(U_\rho(A_{n,m}^{\eta})\), every
   unit cylinder inside the residual region has CKN activity \(\calE_1<\eta\),
   pressure gauges are compact, and no escape outside \(K_m\) takes place.
   \item \emph{Boundary diffuse defect.}  A positive amount of residual mass
   leaves every fixed \(K_m\); normalized restrictions converge on
   \(\overline Z\) to a nonzero probability measure supported on
   \(\overline Z\setminus Z\).
   \item \emph{Noncompact or pressure-gauge defect.}  There exists a recentering
   sequence \(z_n\in E_n^{(j)}\) carrying positive velocity or atlas CKN mass for
   which the local Seregin compactness or pressure-gauge compactness required in
   \Cref{def:terminal-measures} fails.
   \item \emph{Critical-tail defect.}  No unit observer carries a retained
   velocity threshold and no compactness failure occurs, but there are scales
   \(R_n\to\infty\) or \(R_n\downarrow0\), according to the observer convention,
   for which the normalized residual measure on the associated annular/logarithmic
   windows has nonzero limiting mass.  This is the critical-tail compactification
   input of \Cref{def:critical-diffuse-tail-strata}.
\end{enumerate}
The priority order is: \textup{(a)} preempts the others; failing
\textup{(a)}, the alternatives \textup{(b)}--\textup{(e)} are intrinsically
distinct after deleting unit cylinders above the CKN threshold, with
\textup{(b)} being the only
compact-window alternative in which no unit cylinder has CKN activity above
\(\eta\).  By \Cref{lem:compact-subthreshold-regular-residual}, alternative
\textup{(b)} carries no retained terminal velocity concentration and is
removed from the residual decomposition before any later use; the terminal
stratification theorem
\Cref{thm:terminal-stratification} therefore lists only the alternatives
\textup{(c)}--\textup{(e)} together with the four compact retained
alternatives proved later.
\end{theorem}

\begin{proof}
The extraction in \Cref{thm:maximal-active-family} is already diagonal over the
countable pairs \((m,\eta)\).  Refine it over the countable selected regions
\(E_n^{(j)}\) and rational deletion radii \(\rho\); since \(\mathcal E\) is
itself countable by \Cref{def:canonical-expanding-family}, the resulting
diagonal extraction is a countable refinement.  If alternative
\textup{(a)} fails, then for some \(j,\eta,\rho\) either a positive amount of
mass remains in some fixed compact window after deleting
\(U_\rho(A_{n,m}^{\eta})\), or positive mass escapes every fixed compact
window.

In the first case, if a unit observer inside the residual set has velocity
activity above the velocity threshold, it is an active velocity recentering and
its observer cylinder is raw-measure active by
\Cref{lem:oscillation-active-implies-raw-active}; but every raw-measure active
observer at threshold \(\eta\) has been deleted by construction, so this case
does not occur after deletion.  If the local compactness or pressure gauges
fail along the selected observer sequence, this is \textup{(d)}.  If
compactness holds but the lower bound is pressure-only,
\Cref{lem:no-pressure-only-retained-profile} places it either in \textup{(d)},
to the affine/parasitic quotient, or to a neighboring velocity observer.  The
remaining sub-alternative is \textup{(b)}: every unit cylinder in
\(K_m\setminus U_\rho(A_{n,m}^{\eta})\) has CKN activity \(\calE_1<\eta\),
pressure gauges are compact, and the residual mass is contained in
\(K_m\) for some fixed \(m\) (no escape).

In the second case, normalize the residual restrictions outside \(K_m\) and use
weak compactness of probability measures on \(\overline Z\).  The limit has no
mass on any compact subset of \(Z\), hence is supported on the boundary and gives
\textup{(c)}.  Finally suppose that the residual mass is invisible at every
fixed unit observer but is not lost in the compactification.  Cover the selected
tail region by dyadic annular/logarithmic windows and normalize the residual
restriction in that region to a probability measure.  Apply
\Cref{lem:annular-tight-vs-log-escape-dichotomy}.  If a fixed-width logarithmic
window captures a positive amount of the normalized mass after a suitable
logarithmic recentering, then that window gives a fixed multiplicative annulus
at macroscopic scales \(R_n\), hence the macroscopic critical-tail input of
\Cref{def:critical-diffuse-tail-strata}.  If no such window exists, the
log-radius mass is diffuse along the logarithmic axis; its weak-\(*\) limit on
the logarithmic compactification is recorded as the log-diffuse critical-tail
coordinate \(\Lambda\ne0\).  Both subcases are alternative \textup{(e)}.  This
selection is local in the chosen windows and uses no global mass budget.
\end{proof}

\begin{lemma}[Compact subthreshold regular residual cannot carry retained velocity concentration]
\label{lem:compact-subthreshold-regular-residual}
Let the residual measure
\(\mu_{n,m,\eta,\rho}^{\rm rem}\) of
\Cref{def:active-removed-residual-measure} fall into alternative \textup{(b)}
of \Cref{lem:compact-window-expanding-region-protocol} at the threshold
\(\eta\le\eps_{\CKN}/2\), where \(\eps_{\CKN}\) is the CKN smallness constant
of \Cref{thm:imported-setup-results}\,(1).  Then the corresponding portion of
the residual measure is supported on a relatively compact set of
\(\eps_{\CKN}\)-regular points of the underlying suitable solutions, and the
following two conclusions hold.
\begin{enumerate}[label=\textup{(\roman*)}]
   \item \emph{Local regularity.}  The compact-window residual mass is
   carried by velocity fields which are uniformly bounded and locally
   smooth on a slight enlargement of the selected residual region; in
   particular it cannot accumulate retained terminal velocity concentration
   at any compact observer.
   \item \emph{Disjointness from singular cores.}  After alternative
   \textup{(b)} has been recorded and removed, no surviving residual
   measure with retained velocity concentration intersects the same compact
   region.  Equivalently, alternative \textup{(b)} cannot coexist with a
   retained singular concentration core inside the same compact observer
   window after the CKN threshold has been chosen below \(\eps_{\CKN}\).
\end{enumerate}
\end{lemma}

\begin{proof}
Fix the compact window \(K_m\) and the deletion radius \(\rho\) provided by
the alternative.  Let \(R_m=\dist(K_m,\partial K_{m+1})>0\).  Cover
\(K_m\) by a finite collection of unit covariant cylinders
\(\mathcal Q_1^{\rm cov}(z_i)\), \(i=1,\dots,N_m\), arranged so that the
slightly enlarged cylinders \(\mathcal Q_{1+\delta}^{\rm cov}(z_i)\) with
\(\delta\in(0,R_m/4)\) still lie in \(K_{m+1}\setminus
U_{\rho/2}(A_{n,m}^{\eta})\) for all sufficiently large \(n\), using that the
CKN-threshold locus has been deleted at radius \(\rho\) and the cylinders are
chosen at distance \(\rho/2\) from it.

By the compact subthreshold hypothesis, every such cylinder satisfies
\[
   \calE_1(U_n,\Pi_n;z_i)
   <\eta\le \eps_{\CKN}/2 .
\]
The pressure gauge attached to \((U_n,\Pi_n)\) on \(K_{m+1}\) is compact in
the relevant atlas
(\Cref{def:terminal-measures}).  By the imported \(\eps\)-regularity criterion
of Caffarelli--Kohn--Nirenberg, recorded in
\Cref{thm:imported-setup-results}\,(1), each cylinder
\(\mathcal Q_1^{\rm cov}(z_i)\) is \(\eps_{\CKN}\)-regular for the
underlying suitable solution \((U_n,\Pi_n)\) along the diagonal extraction:
the velocity is bounded by an absolute constant on
\(\mathcal Q_{1-\delta}^{\rm cov}(z_i)\) and is locally smooth there.  Taking the
finite union over \(i\) and over the at most countably many windows in
alternative \textup{(b)} gives a uniformly bounded smooth representative of
the velocity on a slight enlargement of the residual support, proving
\textup{(i)}.

For \textup{(ii)}, a retained singular concentration core inside the same
compact window has, by definition of the CKN threshold below
\(\eps_{\CKN}\), at least one unit cylinder with CKN activity exceeding
\(\eta\).  All such cylinders are removed by passing to the
concentration-deleted measure at radius \(\rho\).  The compact subthreshold
alternative is recorded only on the complement of these cylinders.  Hence
the compact subthreshold residual and any retained singular core are
disjointly supported on the compact window after deletion.  In particular
the compact subthreshold residual carries no retained terminal velocity
concentration, and is removed from the residual decomposition before the
terminal stratification theorem
\Cref{thm:terminal-stratification} is applied.
\end{proof}

\begin{lemma}[Annular tight/log-escape dichotomy]
\label{lem:annular-tight-vs-log-escape-dichotomy}
Let \(\lambda_n\) be normalized residual probability measures on dyadic
annuli or logarithmic windows in escaping tail regions, with
\(\lambda_n(\mathbb R^3\setminus\{0\})=1\) and supports tending to infinity
or to zero in the radial sense.  Write
\[
   \Lambda_n:=(\log|\cdot|)_*\lambda_n .
\]
Then, after passing to a subsequence, the following priority alternatives
exhaust the possibilities.
\begin{enumerate}[label=\textup{(\arabic*)},leftmargin=2.2em]
   \item \emph{Macroscopic annular tightness.}  There exist macroscopic
   scales \(R_n\) (with \(R_n\to\infty\) or \(R_n\downarrow0\) according to
   the observer convention) and constants \(0<a<b<\infty\) and \(\delta>0\)
   such that
   \[
      \lambda_n\bigl(\{aR_n<|x|<bR_n\}\bigr)\ge\delta
      \qquad\hbox{for all sufficiently large }n.
   \]
   \item \emph{Logarithmic escape.}  The pushforwards
   \(\Lambda_n\) have no positive mass concentration in any fixed-width
   logarithmic window after arbitrary recentering:
   \[
      \lim_{n\to\infty}
      \sup_{\ell\in\mathbb R}
      \Lambda_n([\ell-L,\ell+L])=0
      \qquad\hbox{for every }L<\infty .
   \]
   In this case every weak-\(*\) limit of \(\Lambda_n\) on the logarithmic
   compactification
   \(\overline{\mathbb R}_{\log}=[-\infty,+\infty]\) is supported at the
   logarithmic boundary and gives the log-diffuse coordinate
   \(\Lambda\ne0\).
\end{enumerate}
\end{lemma}

\begin{proof}
For \(L>0\), define the logarithmic concentration function
\[
   Q_n(L):=\sup_{\ell\in\mathbb R}\Lambda_n([\ell-L,\ell+L]).
\]
If for some \(L<\infty\),
\[
   \limsup_{n\to\infty}Q_n(L)>0,
\]
then, after passing to a subsequence, there are logarithmic centers
\(\ell_n\) and \(\delta>0\) such that
\[
   \Lambda_n([\ell_n-L,\ell_n+L])\ge\delta .
\]
With \(R_n=e^{\ell_n}\), this is exactly
\[
   \lambda_n\bigl(\{e^{-L}R_n\le |x|\le e^LR_n\}\bigr)\ge\delta
\]
for all sufficiently large \(n\).  Since the original measures live on
escaping tail regions, the selected \(R_n\) tends to infinity or to zero
according to the observer convention after passing to a further subsequence.
Thus alternative \textup{(1)} holds with \(a=e^{-L}\) and \(b=e^L\).

If no such \(L\) exists, then \(Q_n(L)\to0\) for every rational \(L>0\), hence
for every finite \(L\) by monotonicity.  This is the stated logarithmic
diffuseness condition.  Compactness of probability measures on
\(\overline{\mathbb R}_{\log}\) gives a weak-\(*\) limit point.  The vanishing
of \(Q_n(L)\) for every finite \(L\) rules out mass at any finite logarithmic
point, because a neighborhood of such a point contains some fixed interval
\([\ell-L,\ell+L]\).  Since each \(\Lambda_n\) has total mass one, the limiting
measure is nonzero and is supported on \(\{-\infty,+\infty\}\).  This is the
log-diffuse alternative.
\end{proof}

\begin{remark}[How the dichotomy enters the protocol]
\label{rem:annular-dichotomy-protocol}
The dichotomy
\Cref{lem:annular-tight-vs-log-escape-dichotomy} is invoked exactly at the
last step of the proof of
\Cref{lem:compact-window-expanding-region-protocol}: the normalized residual
measure on the selected dyadic annular window is either macroscopically
annular tight, in which case it gives the macroscopic critical-tail input of
\Cref{def:critical-diffuse-tail-strata}, or it has nonzero logarithmic
escape, in which case it yields the log-defect measure \(\Lambda\) used in
the critical-tail compactification of
\Cref{def:critical-diffuse-tail-strata}.  In both cases alternative
\textup{(e)} of the protocol is realized.
\end{remark}

\begin{lemma}[Oscillatory observers are raw-measure active]
\label{lem:oscillation-active-implies-raw-active}
Let \(\mu_n\) be the raw atlas defect measure used in
\Cref{thm:maximal-active-family}, so that \(\mu_n\) dominates the velocity
measure \(|U_n|^3\,dy\,d\tau\).  If a covariant unit observer \(z\) satisfies
\[
   \actthree(U_n;\mathcal Q_1^{\rm cov}(z))\ge\eta ,
\]
then the same observer is raw-measure active at threshold \(\eta\):
\[
   \mu_n(\mathcal Q_1^{\rm cov}(z))\ge\eta .
\]
Consequently, after the raw atlas concentration sets of
\Cref{thm:maximal-active-family} are removed at every positive threshold, no
positive constant-in-time oscillatory unit observer remains outside those sets.
\end{lemma}

\begin{proof}
By definition of the constant-in-time oscillation,
\[
   \actthree(U_n;\mathcal Q_1^{\rm cov}(z))
   =
   \inf_{c\in\mathbb R^3}
      \iint_{\mathcal Q_1^{\rm cov}(z)}|U_n-c|^3
   \le
   \iint_{\mathcal Q_1^{\rm cov}(z)}|U_n|^3 .
\]
The raw atlas defect measure contains the velocity contribution, hence
\[
   \iint_{\mathcal Q_1^{\rm cov}(z)}|U_n|^3
   \le
   \mu_n(\mathcal Q_1^{\rm cov}(z)).
\]
Thus oscillation implies raw-measure activity with the same threshold.  The
last statement follows by applying the local concentration-set extraction theorem to
all rational positive thresholds and then lowering thresholds by harmless fixed
cover and cutoff constants when the smoothed bump convention is used.
\end{proof}

\begin{lemma}[Spatial quotient defect measures are locally CKN-controlled]
\label{lem:spatial-quotient-defect-cKN-control}
On the same bounded-overlap covariant unit-cylinder cover, define the spatial
quotient oscillation measure
\[
   \nu_U^{\rm sp}(E)
   :=
   \sum_{z_k\in E}
      \oscsp\bigl(U;\mathcal Q_1^{\rm cov}(z_k)\bigr).
\]
Then \(\nu_U^{\rm sp}\) is locally finite on compact observer windows and is
dominated, up to the fixed cover multiplicity, by the local velocity part of
the defect measure.  In particular, every compactness argument in
\Cref{thm:diffuse-defect-compactness} applies without change to
\(\nu_U^{\rm sp}\) after concentration-set removal.
\end{lemma}

\begin{proof}
For each unit covariant cylinder,
\[
   \oscsp\bigl(U;\mathcal Q_1^{\rm cov}(z_k)\bigr)
   \le
   \iint_{\mathcal Q_1^{\rm cov}(z_k)}|U|^3 .
\]
Summing over a bounded-overlap cover gives local finiteness and domination by
the local velocity contribution, hence by the raw atlas defect measure.  The proof of
\Cref{thm:diffuse-defect-compactness} uses only positivity, local finiteness,
bounded-overlap domination, and weak compactness of normalized restrictions on
the observer compactification.  These properties hold for
\(\nu_U^{\rm sp}\) exactly as they hold for the constant-in-time oscillation
measure \(\mu_U^{\rm osc}\).  In particular,
\[
   \nu_U^{\rm sp}(\{z_k\})
   =
   \oscsp\bigl(U;\mathcal Q_1^{\rm cov}(z_k)\bigr)
   \le
   \actthree\bigl(U;\mathcal Q_1^{\rm cov}(z_k)\bigr).
\]
Hence, on escaping observer windows where all positive-threshold unit
\(\actthree\)-observers have been removed, the atomic masses of
\(\nu_U^{\rm sp}\) tend to zero, so the finite atomic window-selection lemma
applies to \(\nu_U^{\rm sp}\) exactly as it does to \(\mu_U^{\rm osc}\).
\end{proof}

\begin{lemma}[Separation of recentered cylinders]
\label{lem:separated-recentered-cylinders}
Let \(\mathfrak c_i=(z_{i,n})\) and \(\mathfrak c_j=(z_{j,n})\) be
asymptotically separated recentering sequences.  For every fixed \(R<\infty\), the
cylinder \(Q_R(z_{j,n})\), written in the coordinates of \(\mathfrak c_i\),
leaves every compact subset of the recentered variables.  Consequently
no positive compact velocity \(L^3\) concentration assigned to \(\mathfrak c_j\) is visible on any
fixed compact cylinder of the \(\mathfrak c_i\)-frame.
\end{lemma}

\begin{proof}
In the \(\mathfrak c_i\)-coordinates the center of \(Q_R(z_{j,n})\) is
\[
   \bigl(y_{j,n}-y_{i,n},\,\tau_{j,n}-\tau_{i,n}\bigr).
\]
The parabolic distance of this point from the origin tends to infinity by
asymptotic separation.  Hence, for every fixed compact terminal cylinder
\(K\), one has
\[
   K\cap\bigl(Q_R(z_{j,n})-z_{i,n}\bigr)=\varnothing
\]
for all sufficiently large \(n\).  This is a purely local statement about the
relative positions of the two recentering sequences.  It does not estimate the solution on
all of space and does not use a global norm.
\end{proof}

\begin{theorem}[Paired terminal concentration-profile decomposition]
\label{thm:terminal-exhaustion-main}
Every terminal concentration sequence admits a paired decomposition
\[
   \hbox{retained velocity concentration set}
   \quad+\quad
   \hbox{residual remainder}.
\]
More precisely, after passing to a subsequence, the compact observer-window
concentration sets \(A_m^\eta\) of \Cref{thm:maximal-active-family} exist
simultaneously for every \(m\) and every rational \(\eta>0\).  Every observer
sequence that remains in these sets at a fixed positive velocity threshold and has no
loss of local compactness generates a bounded retained terminal profile after
local pressure gauges.  Outside every neighborhood of the concentration set in each
fixed compact observer window, the residual raw unit-window mass is below the chosen
threshold in the precise concentration-deleted sense of
\(\mu_{n,m,\eta,\rho}^{\rm rem}\) from
\Cref{def:active-removed-residual-measure}.  Applying the compact-window to
expanding-region protocol
\Cref{lem:compact-window-expanding-region-protocol}, the compact subthreshold
regular residual is removed first by
\Cref{lem:compact-subthreshold-regular-residual}.  The surviving residual
remainder is then of one of the following priority types:
\[
   \mathcal R_{\rm van},\qquad
   \mathcal R_{\rm diff},\qquad
   \mathcal R_{\rm noncomp},\qquad
   \mathcal R_{\rm crit}.
\]
Here \(\mathcal R_{\rm van}\) means the concentration-deleted residual measures vanish
on compact observer windows and then on the selected expanding regions;
\(\mathcal R_{\rm diff}\) means positive concentration-deleted residual mass
persists at the boundary of the observer compactification while every fixed
unit cylinder is below threshold; and
\(\mathcal R_{\rm noncomp}\) means that some escaping residual recentering loses
the local Seregin compactness or pressure-gauge control; and
\(\mathcal R_{\rm crit}\) means that residual mass is first visible only after
annular/logarithmic critical-tail normalization.  Thus mixed configurations such
as one compact concentration profile plus diffuse exterior remainder are part of the
decomposition; they are not silently identified with the pure diffuse stratum.
\end{theorem}

\begin{proof}
Apply the local concentration-set extraction theorem,
\Cref{thm:maximal-active-family}, on the compact observer exhaustion
\(\{K_m\}\).  It produces closed concentration sets \(A_m^\eta\) and the residual
estimate
\[
   \limsup_{n\to\infty}
   \sup_{z\in K_m\setminus O}M_n(z)\le\eta
\]
outside every neighborhood \(O\) of \(A_m^\eta\).  This is the precise
replacement for a maximal finite or countable tube family: concentrating centers are
recorded by closed loci, and the residual is defined by restricting away from
their neighborhoods.

If an observer sequence stays in a positive-threshold concentration set and has no
loss of the local compactness specified in \Cref{def:terminal-measures}, then
\Cref{lem:retained-recentering-alternatives} gives a bounded retained terminal profile:
the diagonal compactness is supplied there either by the first-generation
local-window compactness lemma or by covariant compactness for bounded ancient profiles.  If such an observer
sequence loses local dissipation or pressure-gauge control, the corresponding
remainder is \(\mathcal R_{\rm noncomp}\).

It remains to classify what is outside the concentration sets.  Apply the
compact-window-to-expanding-region protocol
\Cref{lem:compact-window-expanding-region-protocol} to the selected exterior and
tail regions used in the terminal extraction.  Its vanishing alternative is
\(\mathcal R_{\rm van}\).  Its boundary-measure alternative is
\(\mathcal R_{\rm diff}\).  Its compactness or pressure-gauge failure
alternative is \(\mathcal R_{\rm noncomp}\), while its critical-scale
alternative is recorded for the later critical-tail branch.  These residual
alternatives are mutually exclusive once the concentration sets and selected protocol
are fixed, and they exhaust the residual possibilities.  No step asserts that a
diffuse or exterior remainder cannot coexist with compact concentration
profiles.
\end{proof}

\begin{corollary}[Persistent local velocity concentration forces retained velocity concentration]
\label{cor:persistent-forces-active}
If a terminal sequence satisfies
\[
   \liminf_{n\to\infty}
   \iint_{Q_1(0,0)} |V_n|^3\,dy\,d\tau\ge\eta_0>0
\]
with \(\eta_0\ge\eps_{\rm sd}\), then the terminal decomposition contains a
local concentration profile unless the distinguished observer sequence has
loss of local compactness.  In the generic residual after local vanishing and
local compactness failure are excluded, the persistent compact velocity concentration is
therefore carried by the retained velocity concentration set of
\Cref{thm:terminal-exhaustion-main}; the locus may contain one, finitely many,
infinitely many, or continuum many concentrating observer classes, and the residual
remainder may still be diffuse until the mixed compact--diffuse stratum is
excluded.
\end{corollary}

\begin{proof}
The raw atlas defect measure \(\mu_n\) dominates the velocity contribution.  Hence the
origin recentering sequence has \(\mu\)-mass at least \(\eta_0\), and is raw
unit-window active at threshold \(\eps_{\rm sd}\) unless it has loss of local compactness.  If it
has no loss of local compactness, concentration-set extraction places its observer
class in a positive-threshold concentration set \(A_m^\eta\) of
\Cref{thm:terminal-exhaustion-main}.
\end{proof}

\begin{lemma}[Local vanishing and exterior alternatives do not carry compact concentration]
\label{lem:radiative-stratum-exclusion}
Let a terminal extraction inherit the positive compact velocity \(L^3\) concentration
supplied by the imported setup results.  Then the retained
compact velocity concentration
cannot be carried solely by local vanishing, exterior escape, or diffuse
exterior concentration.
\end{lemma}

\begin{proof}
The assertion is local.  By the local no-escape and concentration input in
the imported setup results, failure of local vanishing at a singular terminal
point gives a positive local velocity concentration sequence.  The setup paper
carries this through the Type~I Seregin extraction as the compact velocity
lower bound \eqref{eq:positive-velocity-main}.  Cover the fixed compact
cylinder
\(K\) in \eqref{eq:positive-velocity-main} by finitely many unit terminal
cylinders.  At least one cylinder in the cover carries a positive velocity
lower bound, say \(\eta_K>0\).  Since \(\mu_n\) dominates
\(|V_n|^3\,dy\,d\tau\), choosing \(\eps_{\rm sd}\le\eta_K\) makes that unit
cylinder raw unit-window active at threshold \(\eps_{\rm sd}\).  Thus the distinguished compact
concentration cannot be locally vanishing.

Local vanishing after removal of raw unit-window active recentering sequences means
that the residual profile has raw unit-window mass below \(\eps_{\rm sd}\) on each fixed
compact terminal cylinder; it therefore cannot contain the fixed compact
velocity lower bound.  An exterior escaping concentration is visible only after
its own recentering; relative to the distinguished compact terminal cylinder it
escapes every bounded window.  By
\Cref{lem:separated-recentered-cylinders}, it cannot be the sole location of a
lower bound that remains in the distinguished compact window.  Finally, diffuse
exterior concentration is defined by persistence through expanding regions
while all fixed unit cylinders are locally vanishing.  If such a sequence
contained the compact lower-bound threshold from the imported setup results, the
positive concentration input in that setting would select a unit terminal
cylinder with mass
at least the velocity threshold, contradicting local vanishing.  Thus none of
these alternatives can be the unique location of the retained compact velocity
concentration.
\end{proof}

\begin{theorem}[Local exclusion of noncompact and nonconcentrating terminal alternatives]
\label{thm:terminal-inactive-exclusion}
For terminal sequences generated by the Seregin setup theorem in
the companion setup paper, loss of local compactness, local vanishing,
exterior escape, and
diffuse exterior concentration are not possible sole locations of the retained
compact velocity concentration.
\end{theorem}

\begin{proof}
Retained recentering sequences with loss of local compactness are first routed
by \Cref{lem:retained-recentering-alternatives} to one of the terminal
branches.  If the routed branch is locally compact, it is handled by the
retained-profile protocol.  If it is diffuse, pressure/noncompact, or
critical-tail, it is excluded only after applying the corresponding branch
closure theorem; thus the routing lemma is not itself used as a contradiction.
Local vanishing, exterior escape, and diffuse exterior concentration are
excluded by \Cref{lem:radiative-stratum-exclusion}.  The proof uses only the
positive concentration input and the local compactness hypotheses in the
imported setup results, and fixed compact terminal cylinders.
\end{proof}

\begin{definition}[Observer compactification]
\label{def:observer-compactification}
We fix a metrizable compactification \(\overline Z\) of the observer space
\(Z=\mathbb R^3\times\mathbb R\) with the following properties.
\begin{enumerate}[label=\textup{(\roman*)}]
   \item Compact observer balls of finite covariant radius embed continuously.
   \item The covariant observer maps generated by \(T_a\) and \(\Theta_s\)
   extend continuously to \(\overline Z\).
   \item Unit-cylinder concentration evaluation maps are upper semicontinuous after
   the fixed smoothing convention of \Cref{thm:maximal-active-family}.
   \item The normalized restrictions of locally finite defect measures are
   tight as probability measures on \(\overline Z\).
\end{enumerate}
Such a compactification is obtained by embedding \(Z\) into the product of
compact intervals generated by a countable separating family of bounded
continuous observer functions, including the smoothed unit-cylinder evaluation
functions and the coordinate cutoffs for a compact exhaustion, and then taking
the closure.
\end{definition}

\begin{theorem}[Diffuse-defect compactness]
\label{thm:diffuse-defect-compactness}
Let \(\mu_n^{\rm osc}\) be the residual constant-in-time oscillation measures
of \Cref{def:observer-oscillation}, after removing neighborhoods of all
positive-threshold concentration sets from \Cref{thm:maximal-active-family}.  By
\Cref{lem:oscillation-active-implies-raw-active}, this removal also removes all
positive-threshold oscillatory unit observers on compact windows.
Assume there are observer regions \(E_n\subset Z\), escaping every compact
subset of \(Z\), and a number \(m_0>0\) such that
\[
   0<m_0\le \mu_n^{\rm osc}(E_n)<\infty ,
\]
while no unit observer in \(E_n\) carries a positive oscillation threshold:
\[
   \sup_{z\in E_n}
      \actthree(U_n;\mathcal Q_1^{\rm cov}(z))
   \longrightarrow0 .
\]
Let \(\overline Z\) be the compact metrizable observer compactification of
\Cref{def:observer-compactification}, and define probability
measures
\[
   \lambda_n
   :=
   \frac{\mu_n^{\rm osc}\lfloor E_n}{\mu_n^{\rm osc}(E_n)}
   \qquad\hbox{on }\overline Z .
\]
After passing to a subsequence,
\[
   \lambda_n\rightharpoonup\lambda
\]
weakly as probability measures on \(\overline Z\).  The limit satisfies
\[
   \lambda(\overline Z)=1,\qquad
   \lambda(B)=0\quad\hbox{for every }B\Subset Z .
\]
Thus the limit is a nonzero diffuse-defect profile supported at the boundary of
the observer compactification, not a local concentration profile.
\end{theorem}

\begin{proof}
The theorem is stated directly in the finite-mass normalization regime needed
to define \(\lambda_n\).  If one starts only from a positive lower bound on an
escaping oscillation region, \Cref{lem:finite-diffuse-window-selection}
provides an escaping finite-mass net-cell subset to which the present argument
applies.

The measures \(\lambda_n\) are probability measures on the compact metric space
\(\overline Z\), so weak compactness gives a convergent subsequence.  Testing
against the constant function \(1\) preserves total mass and gives
\(\lambda(\overline Z)=1\).

Let \(B\Subset Z\).  Since \(E_n\) escapes every compact subset of \(Z\), one
has \(B\cap E_n=\varnothing\) for all sufficiently large \(n\).  Hence
\(\lambda_n(B)=0\) eventually, and the portmanteau theorem gives
\(\lambda(B)=0\).  The additional unit-scale diffuseness assumption gives the
stronger moving-window estimate
\[
   \sup_{z\in Z}
      \lambda_n\bigl(B_Z(z,1)\bigr)
   \le
   \frac{N_0}{m_0}
   \sup_{\zeta\in E_n\cap\{z_k\}}
      \actthree(U_n;\mathcal Q_1^{\rm cov}(\zeta))
   \longrightarrow0 .
\]
Indeed, \(\lambda_n\) is atomic on the fixed observer net of
\Cref{def:observer-oscillation}, and every unit ball for \(d_Z\) contains at
most \(N_0\) net points.  Therefore
\[
\begin{aligned}
   \lambda_n\bigl(B_Z(z,1)\bigr)
   &=
   \frac{1}{\mu_n^{\rm osc}(E_n)}
   \sum_{z_k\in E_n\cap B_Z(z,1)}
      \actthree\bigl(U_n;\mathcal Q_1^{\rm cov}(z_k)\bigr)
   \\
   &\le
   \frac{N_0}{m_0}
   \sup_{\zeta\in E_n\cap\{z_k\}}
      \actthree\bigl(U_n;\mathcal Q_1^{\rm cov}(\zeta)\bigr).
\end{aligned}
\]
Equivalently, after replacing \(E_n\) by its net-cell representative, every
unit observer ball contains at most \(N_0\) contributing atoms of
\(\mu_n^{\rm osc}\).
Therefore no positive unit-scale local oscillation profile is hidden inside the
escaping windows.  The proof is local on the chosen observer windows and uses no
global critical-mass or energy budget.
\end{proof}

\begin{lemma}[Finite diffuse-window selection]
\label{lem:finite-diffuse-window-selection}
Let \(\mu_n^{\rm osc}\) be the observer oscillation measures from
\Cref{def:observer-oscillation}.  Suppose \(E_n\subset Z\) escapes every compact
subset of \(Z\) and
\[
   \liminf_{n\to\infty}\mu_n^{\rm osc}(E_n)>0,
   \qquad
   \sup_{z\in E_n}\actthree(U_n;\mathcal Q_1^{\rm cov}(z))\to0 .
\]
Then for every \(m_0>0\) below the surviving mass there is, after passing to a
subsequence, an escaping net-cell subset \(E_n'\subset E_n\cap\{z_k\}_{k\in\mathbb N}\)
such that
\[
   m_0\le \mu_n^{\rm osc}(E_n')\le 2m_0
\]
for all large \(n\).
\end{lemma}

\begin{proof}
Since \(m_0\) is below the surviving mass, after passing to a subsequence we
may assume
\[
   \mu_n^{\rm osc}(E_n)\ge m_0
\]
for every \(n\).  Since the unit-scale oscillation threshold on \(E_n\) tends
to zero, after passing to a further subsequence we may also assume
\[
   \sup_{z\in E_n}\actthree(U_n;\mathcal Q_1^{\rm cov}(z))\le m_0
\]
for every \(n\).  Replace \(E_n\) by its net-cell representative
\(E_n\cap\{z_k\}_{k\in\mathbb N}\), which does not change
\(\mu_n^{\rm osc}(E_n)\) because the measure is atomic on the fixed observer net.

Enumerate the net points in \(E_n\) arbitrarily and add them one by one until
the accumulated \(\mu_n^{\rm osc}\)-mass first reaches or exceeds \(m_0\).  The
partial sum at the preceding step is \(<m_0\), and the last added atom has mass
at most \(m_0\), so the resulting finite subset \(E_n'\subset E_n\) satisfies
\[
   m_0\le \mu_n^{\rm osc}(E_n')\le 2m_0 .
\]
Because \(E_n\) escapes every compact subset of \(Z\), every subset
\(E_n'\subset E_n\) also escapes every compact subset.  Thus the selected
finite windows retain positive normalized oscillation mass while avoiding any
division by an infinite total mass.
\end{proof}

\begin{lemma}[Finite window selection for atomic observer defects]
\label{lem:finite-atomic-observer-window-selection}
Let \(\gamma_n\) be a nonnegative atomic measure on the fixed observer net
\(\{z_k\}_{k\in\mathbb N}\).  Suppose \(E_n\subset Z\) escapes every compact
subset of \(Z\),
\[
   \liminf_{n\to\infty}\gamma_n(E_n)>0,
   \qquad
   \sup_{z_k\in E_n}\gamma_n(\{z_k\})\to0 .
\]
Then, for every \(m_0>0\) below the surviving mass, after passing to a
subsequence there is an escaping finite net-cell subset
\(E_n'\subset E_n\cap\{z_k\}_{k\in\mathbb N}\) such that
\[
   m_0\le \gamma_n(E_n')\le 2m_0
\]
for all large \(n\).
\end{lemma}

\begin{proof}
Since \(m_0\) lies below the surviving mass, after passing to a subsequence we
may assume \(\gamma_n(E_n)\ge m_0\) for every \(n\).  Since the atomic masses
on \(E_n\) tend to zero, after passing to a further subsequence we may also
assume
\[
   \sup_{z_k\in E_n}\gamma_n(\{z_k\})\le m_0
\]
for every \(n\).  Replacing \(E_n\) by \(E_n\cap\{z_k\}_{k\in\mathbb N}\) does
not change \(\gamma_n(E_n)\) because \(\gamma_n\) is supported on the fixed
observer net.

Enumerate the atoms of \(E_n\) arbitrarily and add them one by one until the
accumulated \(\gamma_n\)-mass first reaches or exceeds \(m_0\).  The partial
sum at the previous step is \(<m_0\), and the last added atom has mass at most
\(m_0\), so the resulting finite subset \(E_n'\subset E_n\) satisfies
\[
   m_0\le \gamma_n(E_n')\le 2m_0 .
\]
Because \(E_n\) escapes every compact subset of \(Z\), every subset
\(E_n'\subset E_n\) also escapes every compact subset.
\end{proof}

\begin{lemma}[Support transfer for the diffuse-defect compactification]
\label{lem:diffuse-defect-support-transfer}
Let \(X_{\rm def}\) be the compact metrizable diffuse-defect space of
\Cref{thm:diffuse-defect-compactification-construction}, and let
\[
   \lambda_n=\frac{\mu_n^{\rm osc}\lfloor E_n}{\mu_n^{\rm osc}(E_n)}
\]
be the normalized realized diffuse-defect measures from
\Cref{thm:diffuse-defect-compactness}, viewed in \(X_{\rm def}\).  If
\(\lambda_n\rightharpoonup\lambda\) and \(A\subset X_{\rm def}\) is a closed
union of already excluded diffuse alternatives, then support of \(\lambda\) in
\(A\) implies that no diffuse defect mass survives outside \(A\).  Conversely,
surviving normalized diffuse mass outside \(A\) produces a sequence-realized
support profile in \(X_{\rm def}\setminus A\).
\end{lemma}

\begin{proof}
This is \Cref{lem:compact-defect-support-transfer} with
\(X=X_{\rm def}\), \(\mu_n=\mu_n^{\rm osc}\), and the escaping observer
regions \(E_n\).  The construction theorem makes \(X_{\rm def}\) compact
metrizable and keeps the limit inside the closed sequence-realized subset, so
the support point obtained outside \(A\) is a realized diffuse-defect profile.
\end{proof}

\begin{definition}[Defect observables]
\label{def:diffuse-defect-observables}
Let \(M\ge 1\) be the velocity bound of the underlying admissible class.
Fix the following six countable families of bounded continuous observables,
each defined on diffuse-defect candidates produced by
\Cref{thm:diffuse-defect-compactness}.
\begin{enumerate}[label=\textup{(\arabic*)},leftmargin=2.2em]
   \item \emph{Observer test functions.} A countable dense family
   \(\{\varphi_j\}_{j\in\mathbb N}\subset C(\overline Z)\), separating
   probability measures on the compact metrizable observer compactification
   \(\overline Z\) of \Cref{def:observer-compactification}.
   \item \emph{Smoothed unit-cylinder activity observables.} For each rational
   \(z\in Z\) and each \(\eta\in\mathbb Q_{>0}\), the smooth bounded functional
   \(\Psi_{z,\eta}\) recording the convolved CKN activity \(\calE_1\) of a
   defect representative on the unit covariant cylinder
   \(\mathcal Q_1^{\rm cov}(z)\), thresholded by a smooth cutoff at
   \(\eta\).  This family is countable.
   \item \emph{Compact-exhaustion cutoff observables.} For each
   \(m\in\mathbb N\), a smooth cutoff \(\zeta_m\) supported in \(K_{m+1}\)
   and equal to \(1\) on \(K_m\), recording the mass that the defect measure
   places on the compact exhaustion of \Cref{thm:maximal-active-family}.
   \item \emph{Local velocity observables.} For each compact cylinder
   \(K_m\) and each finite list \(F_1,\dots,F_N\) of polynomials of total
   degree \(\le m\) in the variables \(W,\nabla W,\dots,\nabla^m W\), the
   localized observable \(\int F_j(W,\nabla W,\dots)\zeta_m\,dy\,d\tau\).
   \item \emph{Pressure-gradient observables modulo gauges.} For each
   compactly supported divergence-free test field
   \(\Psi\in C_c^\infty(K_m,\mathbb R^3)\) with rational coefficients in a
   countable dense subspace, the bounded observable
   \[
      \mathcal P_\Psi(W,P)=\langle \nabla P,\Psi\rangle ,
   \]
   computed for representatives modulo the time-only and affine pressure
   gauges of \Cref{def:pressure-gauge-affine-quotient}.
   \item \emph{Log-annular observables.} For each rational \(\rho_0\) and
   \(A>0\), the cutoffs
   \(\chi_{\rho_0,A}\) of \Cref{def:log-annular-compactification} applied
   to the log-radial pushforward of any defect representative.
\end{enumerate}
The combined family is countable; let \(\{\mathcal O_k\}_{k\in\mathbb N}\)
enumerate it after fixing a uniform bound \(\|\mathcal O_k\|_\infty\le C_k\).
Fix also a countable dense subgroup \(G_{\mathbb Q}\subset G=\mathbb R^3\rtimes
\mathbb R\), and enlarge the observable family by all rational pullbacks
\(\mathcal O_k\circ q\), \(q\in G_{\mathbb Q}\).  The enlarged family is still
countable and is again denoted by \(\{\mathcal O_k\}_{k\in\mathbb N}\).
\end{definition}

\begin{theorem}[Construction of the diffuse-defect compactification]
\label{thm:diffuse-defect-compactification-construction}
There exists a compact metrizable space \(X_{\rm def}\), a continuous action
\(\sigma\colon G\times X_{\rm def}\to X_{\rm def}\) of
\(G=\mathbb R^3\rtimes\mathbb R\), and a closed \(G\)-invariant subset
\[
   X_{\rm def}^{\rm anc}\subset X_{\rm def} ,
\]
such that every normalized diffuse-defect sequence generated by the original
terminal sequence has a subsequential limit in \(X_{\rm def}^{\rm anc}\).
\end{theorem}

\begin{proof}
Let \(\{\mathcal O_k\}_{k\in\mathbb N}\) be the observable family of
\Cref{def:diffuse-defect-observables} with uniform bound
\(\|\mathcal O_k\|_\infty\le C_k\).  Choose an exhaustion
\[
   C_1\Subset C_2\Subset\cdots\Subset G
\]
by compact sets with \(e\in C_1\) and such that for every compact
\(H\Subset G\) and every \(m\) there exists an index \(m(H,m)\) satisfying
\(C_mH\subset C_{m(H,m)}\).

For each diffuse-defect candidate \(\omega\) and each pair \((k,m)\), define
the orbit map
\[
   \iota_{k,m}(\omega)(g):=\mathcal O_k(g\cdot\omega),
   \qquad g\in C_m .
\]
By \Cref{lem:state-space-local-compactness,thm:covariant-observer-calculus},
the family \(\{\iota_{k,m}(\omega):\omega\}\) is uniformly bounded by \(C_k\)
and equicontinuous on \(C_m\).  Arzel\`a--Ascoli therefore gives a compact
metrizable closure
\[
   K_{k,m}:=\overline{\{\iota_{k,m}(\omega):\omega\}}^{\,C(C_m)}
   \subset C(C_m,[-C_k,C_k]).
\]
Set
\[
   \mathcal X_0:=\prod_{k,m\in\mathbb N}K_{k,m},
   \qquad
   \iota(\omega):=\bigl(\iota_{k,m}(\omega)\bigr)_{k,m},
\]
and define
\[
   X_{\rm def}:=\overline{\iota(\text{all defect candidates})}
   \subset \mathcal X_0 .
\]
Then \(\mathcal X_0\) is compact metrizable as a countable product of compact
metrizable spaces, so \(X_{\rm def}\) is compact metrizable as a closed subset.

For every defect candidate the orbit maps satisfy the compatibility relation
\[
   \iota_{k,m'}(\omega)|_{C_m}=\iota_{k,m}(\omega)
   \qquad (m'\ge m).
\]
Since the restriction maps \(C(C_{m'})\to C(C_m)\) are continuous, every point
\(x=(f_{k,m})\in X_{\rm def}\) satisfies the same compatibility.  Now fix
\(h\in G\).  Given \(m\), choose any index \(m'\) with \(C_mh\subset C_{m'}\)
and define
\[
   (\sigma_hx)_{k,m}(g):=f_{k,m'}(gh),
   \qquad g\in C_m .
\]
Compatibility makes this independent of the choice of \(m'\).  For every
defect candidate \(\omega\),
\[
   (\sigma_h\iota(\omega))_{k,m}(g)
   =\mathcal O_k(gh\cdot\omega)
   =\iota_{k,m}(h\cdot\omega)(g),
\]
so \(\sigma_h\iota(\omega)=\iota(h\cdot\omega)\).  Hence \(\sigma_h\) maps the
dense embedded set \(\iota(\text{all defect candidates})\) into itself and
therefore preserves the closure \(X_{\rm def}\).

To check joint continuity, fix a compact \(H\Subset G\) and a coordinate
\((k,m)\).  Choose \(m_H\) with \(C_mH\subset C_{m_H}\).  Then the map
\[
   (h,f)\longmapsto \bigl[g\mapsto f(gh)\bigr]
\]
is continuous from \(H\times C(C_{m_H})\) to \(C(C_m)\) in the sup norm,
because right translation on the compact set \(C_m\) depends continuously on
\(h\in H\).  Therefore each coordinate of \((h,x)\mapsto\sigma_hx\) is
continuous on \(H\times X_{\rm def}\); since the product topology is generated by
finitely many coordinates at a time, \((h,x)\mapsto\sigma_hx\) is jointly
continuous on \(G\times X_{\rm def}\).  The group law holds on the dense
embedded set \(\iota(\text{all defect candidates})\) and hence on all of
\(X_{\rm def}\) by continuity.

The sequence-realized class:
\[
   X_{\rm def}^{\rm anc}
   :=
   \overline{\iota(\text{sequence-realized defect profiles})}
   \subset X_{\rm def} ,
\]
where ``sequence-realized'' means realized from the original blow-up
sequence in the sense of
\Cref{def:sequence-realized-residual-object}.  This gives a closed subset of
the compact product, and the diagonal argument below verifies that this closure
is still represented by limits of realized states rather than by an abstract
compactification artifact.

It remains to check \(G\)-invariance of \(X_{\rm def}^{\rm anc}\) and the
diagonal-closure property.  \(G\)-invariance follows because the covariant
observer maps belong to the allowed realized operations.  For the
diagonal-closure property, suppose \(\omega_j\in X_{\rm def}^{\rm anc}\) and
\(\omega_j\to\omega\) in \(X_{\rm def}\).  By definition of the closure each
\(\omega_j\) is a limit point of a sequence \(\iota(\omega_j^{(n)})\) with
\(\omega_j^{(n)}\) realized from the original blow-up sequence.  Choose
indices \(n=n(j)\) so large that \(d(\iota(\omega_j),\iota(\omega_j^{(n(j))}))
<2^{-j}\).  Then \(\iota(\omega_j^{(n(j))})\to\iota(\omega)\), and the
\(\omega_j^{(n(j))}\) are realized.  Hence \(\omega\) is in the closure of
the realized states, i.e.\ \(\omega\in X_{\rm def}^{\rm anc}\).

Finally, every normalized diffuse-defect sequence generated by the original
terminal sequence is a sequence \(\omega_n\) of realized states.  Their
images \(\iota(\omega_n)\) lie in the compact metrizable space
\(X_{\rm def}\) and have a convergent subsequence; the limit lies in
\(X_{\rm def}^{\rm anc}\) by the closure argument above.
\end{proof}

\begin{definition}[Diffuse-defect profile class]
\label{def:diffuse-defect-compactification}
The compact metrizable space \(X_{\rm def}=\mathfrak D_M\) and its closed
\(G\)-invariant sequence-realized subset \(X_{\rm def}^{\rm anc}=
\mathfrak D_M^{\rm anc}\) are those constructed in
\Cref{thm:diffuse-defect-compactification-construction}.  The
diffuse-defect mass functional is normalized to be identically one on
\(\mathfrak D_M\).  The covariant observer group
\(G=\mathbb R^3\rtimes\mathbb R\) acts on \(\mathfrak D_M\) continuously
through the action \(\sigma\) of
\Cref{thm:diffuse-defect-compactification-construction}.
\end{definition}

\begin{proposition}[Amenable averaging protocol for realized diffuse defects]
\label{prop:amenable-averaging-realized-diffuse-defects}
Let \(G=\mathbb R^3\rtimes\mathbb R\) act continuously on the compact metrizable
diffuse-defect compactification \(X_{\rm def}=\mathfrak D_M\).  Let \(\nu_0\) be a Borel
probability measure supported on the closed sequence-realized class
\(\mathfrak D_M^{\rm anc}\).  Then there exists a \(G\)-invariant Borel
probability measure \(\nu\) on \(X_{\rm def}\), obtained as a weak-* limit of
Følner averages
\[
   \nu_\alpha
   :=
   \frac1{|F_\alpha|}
   \int_{F_\alpha} g_\#\nu_0\,dg,
\]
where \((F_\alpha)\) is a left Følner net in \(G\).  Moreover
\(\nu(\mathfrak D_M^{\rm anc})=1\), so the invariant diffuse defect remains
sequence-realized.
\end{proposition}

\begin{proof}
The semidirect group \(G=\mathbb R^3\rtimes\mathbb R\) is solvable and hence
amenable; choose a left Følner net \((F_\alpha)\) for a left Haar measure.  The
action map \(g\mapsto g_\#\nu_0\) is weak-* continuous because the action on
\(\mathfrak D_M\) is continuous.  Therefore each \(\nu_\alpha\) is a Borel
probability measure.  Compactness of \(X_{\rm def}\) implies weak-* compactness
of \(\mathcal P(X_{\rm def})\), so a subnet converges to a probability measure
\(\nu\).

Let \(h\in G\) and let \(F\in C(X_{\rm def})\).  Then
\[
\begin{aligned}
   \int F\,d(h_\#\nu_\alpha)-\int F\,d\nu_\alpha
   &=
   \frac1{|F_\alpha|}
   \int_{F_\alpha}\left[
      \int F(hg\omega)\,d\nu_0(\omega)
      -
      \int F(g\omega)\,d\nu_0(\omega)
   \right]dg .
\end{aligned}
\]
The absolute value is bounded by
\[
   2\|F\|_{L^\infty}
   \frac{|hF_\alpha\triangle F_\alpha|}{|F_\alpha|},
\]
which tends to zero by the Følner property.  Passing to the weak-* limit gives
\(h_\#\nu=\nu\).  Since \(\mathfrak D_M^{\rm anc}\) is closed and \(G\)-invariant,
each averaged measure \(\nu_\alpha\) is supported on
\(\mathfrak D_M^{\rm anc}\).  Weak-* limits preserve support on closed sets, so
\(\nu(\mathfrak D_M^{\rm anc})=1\).
\end{proof}

\begin{theorem}[Recurrent diffuse-defect measure]
\label{thm:recurrent-diffuse-defect-measure}
The covariant observer action of \(G\) is continuous on \(\mathfrak D_M\).
Every nonzero diffuse-defect profile
\(\lambda\in\mathfrak D_M\) generates a \(G\)-invariant Borel probability
measure \(\mathbb P\) on \(\mathfrak D_M\).  The normalized diffuse-defect mass
remains one under this averaging.
\end{theorem}

\begin{proof}
The continuity of the covariant observer action is the conclusion of the
construction theorem
\Cref{thm:diffuse-defect-compactification-construction}.  In that construction,
the action is first defined on realized defect profiles by the covariant
observer maps and then extended continuously to the compact product closure of
the countable observable family.
Apply \Cref{prop:amenable-averaging-realized-diffuse-defects} to
\(\nu_0=\delta_\lambda\).  The resulting invariant probability measure
\(\mathbb P\) is supported on the closed sequence-realized diffuse-defect class.
Since the diffuse-defect mass is identically one on \(\mathfrak D_M\), its
average against \(\mathbb P\) is also one.
\end{proof}

\begin{theorem}[Diffuse-defect trichotomy]
\label{thm:diffuse-defect-trichotomy}
Let \((U_n,\Pi_n)\subset\mathfrak X_M\) be a normalized terminal residual
sequence with \(\|U_n\|_{L^\infty}\le M\).  Suppose that, after removing the
concentration sets at every positive affine-normalized threshold, there are escaping
observer windows \(E_n\) such that
\[
   \mu_n^{\rm osc}(E_n)\ge m_0>0 .
\]
Then, after passing to a subsequence, the residual diffuse defect falls into
one of the following ordered alternatives:
\begin{enumerate}[label=\textup{D\arabic*.}]
   \item \emph{Regenerated concentration.}  There exist covariant observers \(z_n\)
   and a threshold \(\eta>0\) such that
   \[
      \actthree(U_n;\mathcal Q_1^{\rm cov}(z_n))\ge\eta .
   \]
   This produces a retained concentrating descendant.
   \item \emph{Affine/parasitic collapse.}  Every velocity tangent associated
   with the diffuse-defect hull is locally constant after the admissible affine
   normalization; the defect belongs to the affine/parasitic lower stratum of
   \Cref{def:affine-constant-lower-stratum}.
   \item \emph{Critical diffuse tail.}  There is a nonzero recurrent
   diffuse-defect profile supported on \(\overline Z\setminus Z\), carrying
   normalized diffuse oscillation mass one, annihilating every fixed unit
   concentration observable, and not contained in the affine/parasitic lower stratum.
\end{enumerate}
Equivalently, residual diffuse defects are contained in the union of the regenerated
concentration alternative, the affine/parasitic stratum, and the critical diffuse-tail
class \(\Ccritdiff\).
\end{theorem}

\begin{proof}
If D1 occurs, the asserted trichotomy holds.  Assume D1 does not occur.  Then, for
every fixed \(\eta>0\), all sufficiently large \(n\) satisfy
\[
   \sup_{z\in E_n}
      \actthree(U_n;\mathcal Q_1^{\rm cov}(z))<\eta .
\]
Since D1 does not occur, the atom size of \(\mu_n^{\rm osc}\) on \(E_n\)
tends to zero.  By \Cref{lem:finite-diffuse-window-selection}, after passing to
a subsequence we may replace \(E_n\) by an escaping net-cell subset
\(E_n'\subset E_n\) such that
\[
   m_0\le \mu_n^{\rm osc}(E_n')\le 2m_0 .
\]
Relabel \(E_n'\) as \(E_n\).  The hypotheses of
\Cref{thm:diffuse-defect-compactness} now apply to
\(\mu_n^{\rm osc}\lfloor E_n\), and hence produce a diffuse-defect profile
\(\lambda\) with total mass one and zero mass on each compact subset of \(Z\).
By \Cref{thm:recurrent-diffuse-defect-measure}, the orbit closure of
\(\lambda\) carries a \(G\)-invariant probability measure \(\mathbb P\).

If every local velocity tangent in the support of \(\mathbb P\), after the
canonical affine normalization, is locally constant, then
\Cref{def:affine-constant-lower-stratum} places the defect in the affine
lower stratum.  This is D2.  If not, the invariant diffuse boundary profile has
positive normalized spatial quotient oscillation mass, no retained unit-scale
velocity concentration observable, and non-affine velocity content.  This is precisely the critical diffuse-tail
alternative D3.  The alternatives are ordered by first extracting all regenerated
concentration and then all affine/parasitic collapse; after those removals, the
remaining nonzero diffuse profile is \(\Ccritdiff\).
\end{proof}

\begin{definition}[Critical-tail compactification]
\label{def:critical-diffuse-tail-strata}
The class \(\Ccritdiff\) consists of recurrent diffuse-defect profiles produced
by \Cref{thm:diffuse-defect-trichotomy} which are supported on the boundary of
the covariant observer compactification, carry normalized spatial quotient oscillation
mass one, annihilate every fixed unit concentration observable, and are not
affine/parasitic.

The critical-tail compactification \(\Tcrit\) is the compact profile class of
tuples
\[
   \mathcal T=
   (\bar W,\mathcal Y,\Lambda;
    \widehat{\mathscr P},\mathscr D,
    \mathfrak q_{\rm aff},\mathfrak v_{\rm vir},\mathfrak s)
\]
generated by such critical diffuse profiles after macroscopic blow-down.  Here
\(\bar W\) is the affine-normalized barycentric macroscopic velocity,
\(\mathcal Y_{x,s}\) is the Young measure generated by the affine-normalized
fields \(W_n=V_n-c_n(s)\) and recording unresolved bounded macroscopic
oscillation, with velocity support in \(\overline{B_{M_{\rm tail}}}\) where
\(M_{\rm tail}=2M\), and
\(\Lambda\) is the log-scale diffuse measure recording mass escaping every fixed
multiplicative annulus.  The coordinates
\(\widehat{\mathscr P}\), \(\mathscr D\), \(\mathfrak q_{\rm aff}\),
\(\mathfrak v_{\rm vir}\), and \(\mathfrak s\) record, respectively, the actual
pressure-defect pairings, viscous-defect pairings, affine quotient data, optional
virial-tight coordinates, and sequence-realized support data.  When these extra
coordinates are zero or irrelevant, we suppress them and write simply
\((\bar W,\mathcal Y,\Lambda)\).  In the coarse log-diffuse case
\(\Lambda\ne0\), annular Young, pressure, viscous, and virial coordinates are
inactive until one restricts to normalized log-window support profiles, where
the annular compactification is applied to the selected window.  The exact PDE
and pressure/viscous Young balance are
written instead for the ungauged blow-downs \(V_n\) and their ungauged Young
measure, denoted \(\mathcal Y^V\); outside \(\Caff\), the homogeneous gauge
bridge of \Cref{lem:critical-tail-homogeneous-gauge-bridge} transfers
Dirac/non-Dirac and hidden-scale alternatives between \(\mathcal Y^V\) and the
affine-normalized Young measure \(\mathcal Y\).  The coherent subspace
\(\Tcoh\) is the special case
\[
\begin{aligned}
   &\mathcal Y_{x,s}=\delta_{\bar W(x,s)},\qquad
   \Lambda=0,\qquad
   \widehat{\mathscr P}=0,\qquad
   \mathscr D=0,\\
   &\text{and no retained affine or pressure/noncompact defect.}
\end{aligned}
\]
The purely log-diffuse subspace is denoted \(\Tlogdiff\), and the Young/defect
tail subspace is denoted \(\Tyoung\).
A sequence
\[
   \mathcal T_n=(\bar W_n,\mathcal Y_n,\Lambda_n;
   \widehat{\mathscr P}_n,\mathscr D_n,
   \mathfrak q_{{\rm aff},n},\mathfrak v_{{\rm vir},n},\mathfrak s_n)
\]
with actual annular pressure representatives \(\widehat Q_n^{\rm act}\), defined
modulo time-only gauges, and affine-quotient representatives \(Q_n^\#\),
converges to \(\mathcal T\), with limiting annular pressure representative \(Q\),
in \(\Tcrit\) if:
\begin{enumerate}[label=\textup{(\roman*)}]
   \item \(\bar W_n\rightharpoonup^\ast\bar W\) in
   \(L^\infty_{\loc}\) on macroscopic annular cylinders;
   \item the affine-normalized Young measures generated by \(W_n\) converge
   against every bounded continuous velocity observable with compact support in
   the macroscopic variables: for every
   \(\psi\in C_c^\infty\) supported in a macroscopic annular cylinder and every
   \(F\in C_b(\mathbb R^3)\),
   \[
      \iint \psi(x,s)
      \int F(\xi)\,d\mathcal Y_{n,x,s}(\xi)\,dx\,ds
      \longrightarrow
      \iint \psi(x,s)
      \int F(\xi)\,d\mathcal Y_{x,s}(\xi)\,dx\,ds ;
   \]
   \item the log-scale defect measures \(\Lambda_n\) converge vaguely on the
   logarithmic annular compactification;
   \item the actual pressure gradients
   \(\nabla\widehat Q_n^{\rm act}\) converge annularly distributionally after
   the time-dependent gauges fixed in
   \Cref{def:critical-pressure-gradient-convergence}, and for every pair
   \((F,\psi)\) in the fixed countable \(C^2\)-test family of
   \Cref{def:critical-pressure-gradient-convergence} the scalar defect
   coordinates \(\widehat{\mathscr P}_{n,F}[\psi]\) and
   \(\mathscr D_{n,F}[\psi]\) converge;
   \item on the scale-critical virial-tight recurrent subhull of
   \Cref{def:scale-critical-virial-tight-subhull}, the local virial,
   dissipation, and pressure-flux coordinates on the fixed rational family of
   log-windows are recorded as bounded compactification data rather than
   asserted to be continuous for the weak velocity-pressure topology alone;
   \item the affine quotient data, including the quotient pressures
   \(Q_n^\#\) and any retained affine-force coordinates, converge in the
   finite-dimensional local quotient topology; and
   \item when the profiles are sequence-realized, the limiting profile remains in
   the closed sequence-realized support class generated by the original terminal
   sequence; equivalently, the associated support data converge in the
   diffuse-defect state space \(X_{\rm def}^{\rm anc}\).
\end{enumerate}

For coherent elements, there is an ungauged representative
\((V,Q^{\rm act})\) satisfying on punctured annular cylinders
\(\{0<r_0<|x|<r_1<\infty\}\times I\)
\[
   \partial_s V+\frac12 x\cdot\nabla V+\frac12V+\nabla Q^{\rm act}=0,
   \qquad
   \nabla\cdot V=0
\]
in distributions.  The affine-normalized representative is
\(W=V-c(s)\).  The pair \((W,Q)\in\Ccohcrit\) is understood modulo the
affine/parasitic pressure direction: equivalently, \(W\) satisfies the same
transport equation only in the quotient by spatially homogeneous velocity modes
and affine pressure forces.  If one chooses to write a literal equation for
\(W\), the compensating affine pressure
\((c'(s)+\tfrac12c(s))\cdot x\) is allowed as a distribution in time.  The
coherent tail has nonzero spatial quotient oscillation and recurrence under the
induced logarithmic dilation flow.
\end{definition}

\begin{theorem}[Compactness of the critical-tail compactification]
\label{thm:critical-tail-compactification-compactness}
For every sequence-realized family of critical-tail profiles, either an affine
pressure descendant or a pressure/noncompact descendant is generated and routed
to the corresponding lower stratum, or else, after those alternatives have been
routed away, the pressure and viscous defect coordinates are uniformly bounded
on the fixed countable test family and the family has a convergent subsequence
in \(\Tcrit\).  If all profiles in the retained bounded sequence are
sequence-realized, then the limit is sequence-realized.
\end{theorem}

\begin{proof}
First apply the affine-pressure and pressure/noncompact routing alternatives.  A
nonzero retained affine-force coordinate is assigned to \(\Caff\).  In the
complementary retained branch the affine-pressure coefficient dichotomy gives
\(b_n\to0\) on compact time intervals, while
\Cref{lem:scaled-pressure-compactness-affine} gives
\(\|q_n^{\rm na}\|_{L^p(K)}\le C_{p,K,M}R_n^{-1}\) on every fixed annular
cylinder.  The retained alternative in
\Cref{lem:local-energy-blow-down} gives
\[
   R_n^{-2}\iint |\nabla V_n|^2\phi\le C_{M,K,\phi}.
\]
Consequently, for every pair in the fixed countable test family, the viscous
coordinates are bounded by this local-energy bound and the bounded Hessian of
\(F\).  For the pressure coordinates, write
\(\widehat Q_n^{\rm act}=q_n^{\rm na}+b_n(s)\cdot x+a_n(s)\).  The time-only
gauge has zero spatial gradient; the affine part is bounded, and in fact
vanishes in the retained branch, because \(b_n\to0\); and the non-affine part is
bounded after integrating by parts, using the finite-\(p\) pressure estimate and
the same local-energy bound.  More explicitly, on a fixed annular cylinder \(K\),
choose \(p=2\) in \Cref{lem:scaled-pressure-compactness-affine}; then
\[
   \left|\iint q_n^{\rm na}
      \operatorname{div}_x(\psi\nabla_\xi F(V_n))\right|
   \le
   C\|q_n^{\rm na}\|_{L^2(K)}
   \bigl(1+\|\nabla V_n\|_{L^2(K)}\bigr)
   \le C_{F,\psi,K,M},
\]
because \(\|q_n^{\rm na}\|_{L^2(K)}=O(R_n^{-1})\) and
\(\|\nabla V_n\|_{L^2(K)}=O(R_n)\) in the retained local-energy branch.  Thus
all pressure and viscous defect coordinates on the fixed countable
critical-tail test family are uniformly bounded.  We therefore work in this
retained bounded-coordinate branch.

The barycentric velocities are bounded in \(L^\infty_{\loc}\), so
Banach--Alaoglu gives local weak-\(*\) subsequential convergence.  The
fundamental compactness theorem for Young measures gives convergence of the
oscillation measures against bounded continuous velocity observables.  The
logarithmic annular compactification is metrizable on the countable family of
windows used in the extraction, hence Prokhorov compactness gives a vaguely
convergent subsequence of the log-scale measures.  The pressure components are
recorded only through distributional gradients after the time and affine gauges;
local pressure compactness modulo these gauges gives distributional
subsequential limits.  The pressure-defect coordinates
\(\widehat{\mathscr P}_{n,F}[\psi]\) and viscous-defect coordinates
\(\mathscr D_{n,F}[\psi]\) are part of the defining topology from
\Cref{def:critical-pressure-gradient-convergence,def:critical-diffuse-tail-strata};
on the fixed countable family they are uniformly bounded in the retained branch
by the finite-\(n\) Young identity and the retained local-energy alternative in
\Cref{lem:local-energy-blow-down}, so a diagonal extraction gives convergent
subsequences for all of them simultaneously.  For realized recurrent
profiles lying in the scale-critical virial-tight subhull of
\Cref{def:scale-critical-virial-tight-subhull}, the virial coordinates are
likewise uniformly bounded on the countable rational family by
\Cref{lem:realized-log-window-virial-compactness}, hence admit simultaneous
limits as coordinate data of the recurrent virial hull.  The affine quotient data
are finite-dimensional on each fixed time window, and the quotient topology was
chosen so that bounded sequences have subsequential limits.  Diagonalizing over
the countable annular exhaustion gives convergence in all components listed in
\Cref{def:critical-diffuse-tail-strata}.  For sequence-realized sequences,
choose for each approximating profile a realizing subsequence from the original
terminal sequence and diagonalize over the annular exhaustion and the countable
observables in the topology above.  The diagonal realizing sequence converges
to the same limit, so the limiting critical-tail profile remains
sequence-realized.
\end{proof}

\begin{lemma}[Support transfer for the critical-tail compactification]
\label{lem:critical-tail-support-transfer}
Let \(\Tcrit\) be the compact metrizable critical-tail compactification and let
\[
   \lambda_n=\frac{\mu_n\lfloor E_n}{\mu_n(E_n)}
\]
be normalized realized critical-tail defect measures, where \(\mu_n(E_n)>0\).
If \(\lambda_n\rightharpoonup\lambda\in\mathcal P(\Tcrit)\) and
\(A\subset\Tcrit\) is closed, then support of \(\lambda\) in \(A\) removes all
surviving critical-tail defect mass outside \(A\).  If positive normalized
critical-tail mass survives outside \(A\), then a sequence-realized support
profile in \(\Tcrit\setminus A\) exists.
\end{lemma}

\begin{proof}
This is \Cref{lem:compact-defect-support-transfer} with \(X=\Tcrit\).  The
compactness theorem above ensures that subsequential limits of realized
critical-tail profiles remain sequence-realized, so the support point outside
the closed set is not an abstract compactification artifact.
\end{proof}

\begin{definition}[Log-annular compactification]
\label{def:log-annular-compactification}
Set \(\rho=\log|x|\) and write \(\overline{\mathbb R}_{\log}=[-\infty,+\infty]\)
for the two-point compactification of the log-radius axis.  Fix the countable
family of log-window cutoffs
\[
   \chi_{\rho_0,A}(\rho):=\chi\!\Bigl(\tfrac{\rho-\rho_0}{A}\Bigr),
   \qquad \rho_0\in\mathbb Q,\ A\in\mathbb Q_{>0},
\]
where \(\chi\) is the fixed compactly supported even cutoff of
\Cref{subsec:log-annular-virial-critical-defects}.  The log-annular space
\(\overline{\mathbb R}_{\log}\) is endowed with the metrizable vague topology
generated by these cutoffs.  A sequence \(\Lambda_n\) of finite Borel measures
on \(\overline{\mathbb R}_{\log}\) converges to \(\Lambda\) vaguely if
\[
   \int \chi_{\rho_0,A}(\rho)\,d\Lambda_n(\rho)
   \longrightarrow
   \int \chi_{\rho_0,A}(\rho)\,d\Lambda(\rho)
\]
for every cutoff in the countable family above and additionally for the test
functions
\[
   \rho\mapsto \chi(\rho/A),\quad
   \rho\mapsto \chi((\rho\mp A)/1)
\]
witnessing the masses at \(\pm\infty\).
\end{definition}

\begin{definition}[Young-measure convergence at the macroscopic critical scale]
\label{def:young-measure-convergence}
Let \(\overline{B_M}\subset\mathbb R^3\) be the closed velocity ball of radius
\(M\), and set
\[
   M_{\rm tail}:=2M .
\]
The ungauged exact PDE fields \(V_n\) take values in \(\overline{B_M}\), while
the affine-normalized fields \(W_n=V_n-c_n(s)\), with \(|c_n(s)|\le M\), take
values in \(\overline{B_{M_{\rm tail}}}\).  A sequence of affine-normalized Young
measures \(\mathcal Y_n\) on the macroscopic spacetime annulus
\(\{a<|x|<b\}\times I\) with values in \(\overline{B_{M_{\rm tail}}}\)
converges to \(\mathcal Y\) if for every \(\psi\in C_c^\infty(\{a<|x|<b\}\times
I)\) and every \(F\in C(\overline{B_{M_{\rm tail}}})\),
\[
   \iint \psi(x,s)\!\int F(\xi)\,d\mathcal Y_{n,x,s}(\xi)\,dx\,ds
   \longrightarrow
   \iint \psi(x,s)\!\int F(\xi)\,d\mathcal Y_{x,s}(\xi)\,dx\,ds .
\]
This is the fundamental theorem of Young measures applied to the bounded
sequence of macroscopic representatives \(W_n\), with associated Young
measures \(\mathcal Y_{n,x,s}\) recording the unresolved bounded oscillations
at \((x,s)\).  The countable critical-tail pressure and viscous test family in
\Cref{def:critical-pressure-gradient-convergence} remains on
\(\overline{B_M}\), because those identities are tested against the ungauged
exact PDE variables \(V_n\).  Equivalently, \(\mathcal Y_n\to\mathcal Y\) in the
\(\sigma(L^\infty,L^1)\)-topology induced by tensor test functions.
\end{definition}

\begin{definition}[Distributional pressure-gradient convergence at critical scale]
\label{def:critical-pressure-gradient-convergence}
Fix once and for all a countable dense family \(\mathfrak F_{\rm crit}\) of test
pairs \((F,\psi)\), where \(F\in C^2(\overline{B_M})\) and
\(\psi\in C_c^\infty(\{a<|x|<b\}\times I)\) has rational coefficients on the
chosen countable annular exhaustion.  Let \(\mathfrak F_{\rm crit}^{\rm conv}\)
denote the subfamily of pairs in this exhaustion for which
\(\psi\ge0\) and \(\nabla_\xi^2F\ge0\) on \(\overline{B_M}\), enlarged if necessary
so that, for every rational vector \(q\in\mathbb Q^3\), it contains the convex
quadratic observable
\[
   F_q(\xi)=\frac12(q\cdot\xi)^2,
   \qquad
   \nabla_\xi^2F_q=q\otimes q .
\]
The countable family \(\{q\otimes q:q\in\mathbb Q^3\}\) determines every positive
semidefinite matrix-valued Radon defect measure by density and polarization.
Thus vanishing of all positive convex coordinates implies vanishing of the full
viscous defect matrix, and hence all signed viscous coordinates vanish.

Let \(\widehat Q_n^{\rm act}\) be actual macroscopic pressure representatives,
defined modulo time-only gauges on the same annular cylinder, and let
\[
   Q_n^\#(x,s):=\widehat Q_n^{\rm act}(x,s)-b_n(s)\cdot x
\]
be their affine-quotient representatives.  We say that
\(\nabla\widehat Q_n^{\rm act}\) converges \emph{annularly distributionally}
to \(\nabla Q\), written
\(\nabla\widehat Q_n^{\rm act}\to\nabla Q\) in \(\mathcal D'_{\rm ann}\), if
for every compactly supported divergence-free test vector field
\(\Psi\in C_c^\infty(\{a<|x|<b\}\times I,\mathbb R^3)\)
\[
   \langle\nabla\widehat Q_n^{\rm act},\Psi\rangle
   =-\langle\widehat Q_n^{\rm act},\nabla\cdot\Psi\rangle
   =0
\]
and for every general compactly supported \(\Phi\in C_c^\infty(\{a<|x|<b\}\times
I,\mathbb R^3)\) one has
\[
   \langle\nabla\widehat Q_n^{\rm act},\Phi\rangle
   \to\langle\nabla Q,\Phi\rangle .
\]
In addition, for every \((F,\psi)\in\mathfrak F_{\rm crit}\) the scalar defect
coordinates
\[
   \widehat{\mathscr P}_{n,F}[\psi]
   :=
   \iint \psi\,\nabla_\xi F(V_n)\cdot\nabla\widehat Q_n^{\rm act}
\]
and
\[
   \mathscr D_{n,F}[\psi]
   :=
   R_n^{-2}\iint
   \psi\,\nabla V_n^T\nabla_\xi^2F(V_n)\nabla V_n
\]
are required to converge along the selected compactification subsequence.  If
one passes to the affine-quotient representative
\(Q_n^\#=\widehat Q_n^{\rm act}-b_n(s)\cdot x\), then
\[
   \widehat{\mathscr P}_{n,F}[\psi]
   =
   \mathscr P^\#_{n,F}[\psi]+\mathscr B_{n,F}[\psi],
\]
where
\[
   \mathscr P^\#_{n,F}[\psi]
   :=
   \iint \psi\,\nabla_\xi F(V_n)\cdot\nabla Q_n^\#
\]
and
\[
   \mathscr B_{n,F}[\psi]
   :=
   \iint \psi\,\nabla_\xi F(V_n)\cdot b_n(s).
\]
A nonzero retained affine-force coordinate \(\mathscr B_{n,F}[\psi]\) is
assigned to \(\Caff\).  Outside \(\Caff\), the critical-tail Young identity is
written using the actual pressure coordinates
\(\widehat{\mathscr P}_{n,F}[\psi]\).
\end{definition}

\begin{lemma}[Critical-tail Young balance with pressure and viscous defects]
\label{lem:critical-tail-young-log-transport}
Let \(V_n\in L^\infty\) be uniformly bounded by \(M\) and satisfy the exact
critical-tail blow-down equation
\begin{equation}\label{eq:critical-blow-down-equation}
   \partial_s V_n
   -R_n^{-2}\Delta V_n
   +\frac12 x\cdot\nabla V_n+\frac12 V_n
   +R_n^{-1}V_n\cdot\nabla V_n+\nabla \widehat Q_n^{\rm act}=0,
   \qquad \nabla\cdot V_n=0
\end{equation}
on a fixed macroscopic annular cylinder \(\{a<|x|<b\}\times I\), with
\(R_n\to\infty\), and pressure gradients satisfying the time-gauge convention of
\Cref{def:critical-pressure-gradient-convergence}.  Let
\((\bar W,\mathcal Y,\Lambda)\) be the limiting critical-tail profile in
\(\Tcrit\), where \(\mathcal Y\) is the Young measure generated by the
affine-normalized fields \(W_n=V_n-c_n(s)\).  Let
\(\mathcal Y^V_{x,s}\) denote the ungauged Young measure generated by \(V_n\).
After passing to the subsequence defining \(\Tcrit\), the actual pressure and
viscous pairings
\[
   \widehat{\mathscr P}_{n,F}[\psi]
   :=
   \iint \psi(x,s)\,
      \nabla_\xi F(V_n(x,s))\cdot
      \nabla\widehat Q_n^{\rm act}(x,s)\,dx\,ds
\]
\[
   \mathscr D_{n,F}[\psi]
   :=
   R_n^{-2}\iint
      \psi(x,s)\,
      \nabla_xV_n(x,s)^T
      \nabla_\xi^2F(V_n(x,s))
      \nabla_xV_n(x,s)\,dx\,ds
\]
\[
   \mathscr E_{n,F}[\psi]
   :=
   -R_n^{-2}\iint F(V_n)\,\Delta\psi\,dx\,ds
   -R_n^{-1}\iint F(V_n)V_n\cdot\nabla\psi\,dx\,ds .
\]
converge for every \(\psi\in C_c^\infty\) and every
\(F\in C^2(\overline{B_M})\) in the countable dense test class used to define
the critical-tail topology.  Denote the limits by
\(\widehat{\mathscr P}_F[\psi]\) and \(\mathscr D_F[\psi]\).  Define also the
explicit error \(\mathscr E_{n,F}[\psi]\) above.  For each finite \(n\), the exact identity
\begin{equation}\label{eq:young-transport-identity-finite-n}
\begin{aligned}
   &\iint
   \Bigl(\partial_s\psi+\tfrac12 x\cdot\nabla_x\psi+\tfrac32\psi\Bigr)
   F(V_n)\,dx\,ds\\
   &\qquad =
   \frac12\iint \psi\,\nabla_\xi F(V_n)\cdot V_n\,dx\,ds
   +\widehat{\mathscr P}_{n,F}[\psi]
   +\mathscr D_{n,F}[\psi]+\mathscr E_{n,F}[\psi].
\end{aligned}
\end{equation}
holds.  Moreover \(\mathscr E_{n,F}[\psi]\to0\) as \(n\to\infty\).  If one
rewrites the pressure with the affine-quotient representative
\(Q_n^\#=\widehat Q_n^{\rm act}-b_n(s)\cdot x\), then
\(\widehat{\mathscr P}_{n,F}[\psi]=\mathscr P^\#_{n,F}[\psi]+
\mathscr B_{n,F}[\psi]\); a nonzero retained affine-force coordinate
\(\mathscr B_{n,F}[\psi]\) is routed to \(\Caff\).  Consequently the Young
measure \(\mathcal Y^V\) satisfies the exact weak
balance
\begin{equation}\label{eq:young-transport-identity}
\begin{aligned}
   &\iint
   \Bigl(\partial_s\psi+\tfrac12 x\cdot\nabla_x\psi+\tfrac32\psi\Bigr)
   \int F(\xi)\,d\mathcal Y^V_{x,s}(\xi)\,dx\,ds\\
   &\qquad =
   \frac12\iint \psi
   \int \xi\cdot\nabla_\xi F(\xi)\,d\mathcal Y^V_{x,s}(\xi)\,dx\,ds
   +\widehat{\mathscr P}_F[\psi]+\mathscr D_F[\psi].
\end{aligned}
\end{equation}
with the usual approximation from \(C^2\) to bounded continuous observables on
\(\overline{B_M}\) whenever the velocity derivatives and the two defect
pairings are represented by uniform limits.  In particular, if \(F\) is
homogeneous of degree \(p\) and
\(\widehat{\mathscr P}_F[\psi]=\mathscr D_F[\psi]=0\), then
\[
   \iint
   \Bigl(\partial_s\psi+\tfrac12x\cdot\nabla_x\psi
      +\tfrac{3-p}{2}\psi\Bigr)
   \int F(\xi)\,d\mathcal Y^V_{x,s}(\xi)\,dx\,ds=0 .
\]
For the affine-normalized Young measure \(\mathcal Y\), the corresponding
observable is \(F(\xi+c(s))\), not \(F(\xi)\), unless the homogeneous gauge has
already been removed in the parasitic-free representative.
The log-scale component \(\Lambda\), which records the covariant positions of
the selected annular windows rather than a pressure-renormalized velocity
observable, is invariant under the induced log-translation
\(\rho\mapsto\rho+\tfrac12 s\) on \(\overline{\mathbb R}_{\log}\).  If one
also records the affine-normalized Young measure \(\mathcal Y\) generated by
\(W_n\), then outside \(\Caff\) the bridge lemma
\Cref{lem:critical-tail-homogeneous-gauge-bridge} identifies \(\mathcal Y^V\)
and \(\mathcal Y\) by deterministic spatially homogeneous translation, so their
Dirac/non-Dirac and hidden-scale alternatives agree.
\end{lemma}

\begin{proof}
Test \eqref{eq:critical-blow-down-equation} against \(\psi(x,s)F'(V_n(x,s))\)
for \(F\in C^2(\overline{B_M})\) and \(\psi\in C_c^\infty\).  Each term is
integrated by parts inside the chosen macroscopic annular cylinder where
\(\psi\) has compact support.  The viscous term contributes
\[
   -R_n^{-2}\iint \psi\,\nabla_\xi F(V_n)\cdot\Delta V_n\,dx\,ds .
\]
Integrating by parts gives
\[
\begin{aligned}
&-R_n^{-2}\iint \psi\,\nabla_\xi F(V_n)\cdot\Delta V_n\,dx\,ds \\
&\quad =
R_n^{-2}\iint
\psi\,\nabla_xV_n^T\nabla_\xi^2F(V_n)\nabla_xV_n\,dx\,ds
+
R_n^{-2}\iint
\nabla_xF(V_n)\cdot\nabla\psi\,dx\,ds .
\end{aligned}
\]
We are in the retained non-affine branch of
\Cref{lem:local-energy-blow-down}; otherwise the sequence has already been
routed to the affine or pressure/noncompact alternative and no \(\Tcrit\)
Young-balance extraction is being performed.  Since \(F'\) and \(F''\) are
bounded on \(\overline{B_M}\), \(V_n\) is bounded by \(M\), and the local energy
bound from
\Cref{lem:local-energy-blow-down} gives \(R_n^{-2}\iint
|\nabla V_n|^2\psi\,dx\,ds\le C\), the first term is exactly
\(\mathscr D_{n,F}[\psi]\) and is uniformly bounded on the fixed countable test
family.  The second term equals
\(-R_n^{-2}\iint F(V_n)\Delta\psi\,dx\,ds\), hence is \(O(R_n^{-2})\) and
tends to \(0\).

The convective term, after the standard chain-rule manipulation, equals
\[
   R_n^{-1}\!\iint F'(V_n)(V_n\cdot\nabla V_n)\psi\,dx\,ds
   =R_n^{-1}\!\iint (V_n\cdot\nabla)F(V_n)\,\psi\,dx\,ds
   =-R_n^{-1}\!\iint F(V_n)V_n\cdot\nabla\psi\,dx\,ds ,
\]
using \(\nabla\cdot V_n=0\).  Since \(F\) and \(V_n\) are bounded, this term
is \(O(R_n^{-1})\) and tends to \(0\).

The drift term \(\tfrac12 x\cdot\nabla V_n+\tfrac12 V_n\) tested against
\(F'(V_n)\psi\) gives, after the chain rule and integration by parts in
\(x\),
\[
   \iint\Bigl(\tfrac12 x\cdot\nabla F(V_n)+\tfrac12 F'(V_n)V_n\Bigr)\psi\,dx\,ds
   =-\iint F(V_n)\Bigl(\tfrac12\nabla\cdot(x\psi)\Bigr)\,dx\,ds
   +\iint \tfrac12 F'(V_n)V_n\,\psi\,dx\,ds ,
\]
where the boundary term vanishes by compact support of \(\psi\).  Using
\(\tfrac12\nabla\cdot(x\psi)=\tfrac32\psi+\tfrac12 x\cdot\nabla\psi\) gives
\[
   -\iint F(V_n)
   \Bigl(\tfrac32\psi+\tfrac12 x\cdot\nabla\psi\Bigr)\,dx\,ds
   +\iint \tfrac12 F'(V_n)V_n\,\psi\,dx\,ds .
\]
Combining with the remaining \(\partial_s\) term, which gives
\(-\iint F(V_n)\partial_s\psi\,dx\,ds\) after integration by parts in time,
and keeping the pressure term explicitly, we obtain the exact finite-\(n\)
identity
\[
\begin{aligned}
   &\iint
   F(V_n)
   \Bigl(\partial_s\psi+\tfrac12 x\cdot\nabla\psi+\tfrac32\psi\Bigr)\,dx\,ds
   -\iint \tfrac12 F'(V_n)V_n\,\psi\,dx\,ds \\
   &\qquad =
   \widehat{\mathscr P}_{n,F}[\psi]
   +\mathscr D_{n,F}[\psi]+\mathscr E_{n,F}[\psi],
\end{aligned}
\]
which is exactly \eqref{eq:young-transport-identity-finite-n}.  The remainder
\(\mathscr E_{n,F}[\psi]\) tends to zero because the two terms estimated above
are \(O(R_n^{-2})\) and \(O(R_n^{-1})\), respectively.  Since the two terms on
the left-hand side are uniformly bounded by \(M\|\psi\|_{C^1}\) and
\(\mathscr D_{n,F}[\psi]\) is uniformly bounded as above, the finite-\(n\)
identity also gives a uniform bound for the actual pressure pairings
\(\widehat{\mathscr P}_{n,F}[\psi]\) on the fixed countable family.  These scalar
coordinates are therefore available for diagonal extraction in the definition of
\(\Tcrit\).

Choosing \(F(\xi)=\xi\cdot e_k\) recovers the linear barycenter equation;
choosing more general \(F\) and passing to the ungauged Young-measure limit using
\Cref{def:young-measure-convergence} gives the exact weak identity
\[
\begin{aligned}
   &\iint
   \Bigl(\partial_s\psi+\tfrac12 x\cdot\nabla\psi+\tfrac32\psi\Bigr)
   \int F(\xi)\,d\mathcal Y^V_{x,s}(\xi)\,dx\,ds\\
   &\qquad =
   \iint \tfrac12 \psi
   \int F'(\xi)\xi\,d\mathcal Y^V_{x,s}(\xi)\,dx\,ds
   +\widehat{\mathscr P}_F[\psi]+\mathscr D_F[\psi].
\end{aligned}
\]
This is \eqref{eq:young-transport-identity}.  The homogeneous-observable
special case follows from Euler's identity
\(\xi\cdot\nabla_\xi F=pF\) whenever the pressure and viscous defect pairings
vanish for that observable.  If either pairing does not vanish, the nonzero
defect is retained as pressure-gradient or Young/viscous defect data in
\(\Tcrit\) and is assigned to the pressure/noncompact, affine quotient, or
Young/defect critical-tail alternatives rather than silently discarded.

The same calculation applied to the radial pushforward
\(\rho\mapsto\log|x|\) of \(V_n\), or directly to the dyadic annular
\(\Lambda_n\), gives the invariance of \(\Lambda\) under the induced
log-translation \(\rho\mapsto\rho+\tfrac12 s\).
\end{proof}

\subsection{Log-annular virial identities for recurrent critical defects}
\label{subsec:log-annular-virial-critical-defects}

Throughout this subsection let
\(\chi\in C_c^\infty((-2,2))\) satisfy \(0\le\chi\le1\) and
\(\chi\equiv1\) on \((-1,1)\).  For \(\rho\in\mathbb R\) and \(A\ge1\), set
\[
   \phi_{\rho,A}(y)
   :=
   \chi\!\left(\frac{\log |y|-\rho}{A}\right),
   \qquad R:=e^\rho .
\]
The cutoff is used only on annuli avoiding the origin.  When differentiating
with respect to \(\rho\), \(A\) is kept fixed.

\begin{lemma}[Log-annular oscillation virial identity]
\label{lem:log-annular-oscillation-virial}
Let \((U,\Pi)\in\mathfrak X_M\) be a smooth bounded solution of
\eqref{eq:centered-main}.  Fix \(A\ge1\), \(\rho\in\mathbb R\), and put
\(\phi=\phi_{\rho,A}\).  Define the weighted affine gauge
\[
   c_{\rho,A}(\tau)
   :=
   \frac{\int \phi(y)U(y,\tau)\,dy}{\int \phi(y)\,dy},
   \qquad
   V_{\rho,A}(y,\tau):=U(y,\tau)-c_{\rho,A}(\tau),
\]
so that
\[
   \int \phi V_{\rho,A}\,dy=0 .
\]
Set
\[
   E_{\rho,A}(\tau):=\int \phi |V_{\rho,A}|^2\,dy,
   \qquad
   \mathcal V_{\rho,A}(\tau):=R^{-1}E_{\rho,A}(\tau),
\]
\[
   D_{\rho,A}(\tau):=\int \phi |\nabla U|^2\,dy,
   \qquad
   \mathcal D_{\rho,A}(\tau):=R^{-1}D_{\rho,A}(\tau).
\]
Then \(\mathcal V_{\rho,A}\) is locally absolutely continuous in \(\tau\), and
for a.e. \(\tau\), after subtracting an arbitrary pressure gauge \(\pi(\tau)\),
one has
\begin{equation}
\label{eq:log-annular-virial-identity}
\begin{aligned}
   \partial_\tau \mathcal V_{\rho,A}
   +2\mathcal D_{\rho,A}
   +\frac12\partial_\rho \mathcal V_{\rho,A}
   &=
   R^{-1}\int (\Delta\phi)|V_{\rho,A}|^2\,dy                         \\
   &\quad
   +R^{-1}\int |V_{\rho,A}|^2U\cdot\nabla\phi\,dy                    \\
   &\quad
   +2R^{-1}\int (\Pi-\pi(\tau))V_{\rho,A}\cdot\nabla\phi\,dy .
\end{aligned}
\end{equation}
Moreover
\begin{equation}
\label{eq:log-cutoff-derivative-bounds}
   |\nabla\phi_{\rho,A}(y)|\le \frac{C_\chi}{A|y|},
   \qquad
   |\Delta\phi_{\rho,A}(y)|\le
   \frac{C_\chi}{A|y|^2}+\frac{C_\chi}{A^2|y|^2},
\end{equation}
and
\[
   y\cdot\nabla\phi_{\rho,A}(y)
   =
   \frac1A\chi'\!\left(\frac{\log |y|-\rho}{A}\right),
   \qquad
   \partial_\rho\phi_{\rho,A}=-y\cdot\nabla\phi_{\rho,A}.
\]
\end{lemma}

\begin{proof}
Since \(c_{\rho,A}(\tau)\) is the weighted least-squares affine gauge,
\[
   \int\phi V_{\rho,A}\,dy=0 .
\]
Hence the derivative of the gauge drops out:
\[
   \frac{d}{d\tau}\int\phi |U-c_{\rho,A}|^2
   =
   2\int\phi V_{\rho,A}\cdot\partial_\tau U .
\]
Substitute the centered equation
\[
   \partial_\tau U
   =
   \Delta U-(U\cdot\nabla)U-\frac12y\cdot\nabla U-\frac12U-\nabla\Pi .
\]
The diffusion term gives
\[
   2\int\phi V\cdot\Delta U
   =
   -2\int\phi|\nabla U|^2+\int(\Delta\phi)|V|^2 .
\]
The convective term gives, using \(U\cdot\nabla U=U\cdot\nabla V\) and
\(\nabla\cdot U=0\),
\[
   -2\int\phi V\cdot(U\cdot\nabla U)
   =
   -\int\phi U\cdot\nabla |V|^2
   =
   \int |V|^2U\cdot\nabla\phi .
\]
The centered drift and linear terms give
\[
\begin{aligned}
   -\int\phi V\cdot(y\cdot\nabla U)-\int\phi V\cdot U
   &=
   -\frac12\int\phi y\cdot\nabla |V|^2
   -\int\phi |V|^2                                      \\
   &=
   \frac12\int y\cdot\nabla\phi |V|^2
   +\frac12\int\phi |V|^2 .
\end{aligned}
\]
The terms involving \(c_{\rho,A}\) vanish because
\(\int\phi V=0\).  Finally, for any scalar gauge \(\pi(\tau)\),
\[
   -2\int\phi V\cdot\nabla\Pi
   =
   2\int(\Pi-\pi(\tau))V\cdot\nabla\phi,
\]
because \(\nabla\cdot V=0\) and
\[
   \int \pi(\tau)V\cdot\nabla\phi\,dy
   =
   \pi(\tau)\int\nabla\cdot(\phi V)\,dy=0 .
\]
Combining these identities yields
\[
\begin{aligned}
   \partial_\tau E_{\rho,A}+2D_{\rho,A}
   &=
   \int(\Delta\phi)|V|^2
   +\int |V|^2U\cdot\nabla\phi
   +2\int(\Pi-\pi)V\cdot\nabla\phi                         \\
   &\quad
   +\frac12\int\phi |V|^2
   +\frac12\int y\cdot\nabla\phi |V|^2 .
\end{aligned}
\]
Since
\[
   \partial_\rho\phi_{\rho,A}=-y\cdot\nabla\phi_{\rho,A}
\]
and the weighted least-squares condition also removes the
\(\partial_\rho c_{\rho,A}\) term,
\[
   \partial_\rho E_{\rho,A}
   =
   \int \partial_\rho\phi_{\rho,A}|V|^2
   =
   -\int y\cdot\nabla\phi_{\rho,A}|V|^2 .
\]
Therefore
\[
   \frac12 R^{-1}\left(
      \int\phi |V|^2+
      \int y\cdot\nabla\phi |V|^2
   \right)
   =
   -\frac12\partial_\rho(R^{-1}E_{\rho,A})
   =
   -\frac12\partial_\rho\mathcal V_{\rho,A}.
\]
Thus, after multiplying the energy identity by \(R^{-1}\), the centered drift
and zeroth-order contribution is
\(-\frac12\partial_\rho\mathcal V_{\rho,A}\).  Moving it to the left gives
\eqref{eq:log-annular-virial-identity}.  The cutoff bounds follow by direct
differentiation in polar coordinates.
\end{proof}

\begin{remark}[Virial normalization]
\label{rem:virial-normalization-warning}
The normalization \(R^{-1}\int\phi |U-c|^2\) is the correct scale for an actual
spatial tail of size \(|y|^{-1}\).  It is not automatically a compact observable
for an arbitrary bounded residual profile.  In the critical-tail reduction
below, one uses it either on selected scale-critical defect windows or as an
observable on the compact recurrent defect hull after the required
scale-critical virial tightness has been verified.  Without that compactness,
the virial identity remains valid for smooth annular cutoffs but cannot by
itself close the critical-tail alternatives.
\end{remark}

\begin{definition}[Scale-critical virial-tight recurrent subhull]
\label{def:scale-critical-virial-tight-subhull}
A realized recurrent critical-tail defect lies in the scale-critical virial-tight
recurrent subhull if, along a realizing recurrent sequence and for every
rational log-window \((\rho,A)\),
\[
   \sup_n
   \Bigl(
      |\mathcal V_{\rho,A}^{(n)}|
      +|\mathcal D_{\rho,A}^{(n)}|
      +|\mathcal F_{\rho,A}^{(n)}|
   \Bigr)<\infty ,
\]
where \(\mathcal F_{\rho,A}^{(n)}\) denotes the sum of the three right-hand-side
flux contributions in \eqref{eq:log-annular-virial-identity}.  On this subhull
the virial, dissipation, and flux quantities are adjoined to the realized
recurrent compactification as bounded scalar coordinates on the fixed rational
family of log-windows.
\end{definition}

\begin{lemma}[Realized log-window virial compactness]
\label{lem:realized-log-window-virial-compactness}
Let \(\mathcal D\in\Tcrit\) be a recurrent critical-tail defect generated by a
normalized terminal sequence realized from the original blow-up sequence after concentration-set removal and
affine normalization, and assume that \(\mathcal D\) lies in the scale-critical
virial-tight recurrent subhull of
\Cref{def:scale-critical-virial-tight-subhull}.  Then, along the selected
defect subsequence and after the admissible local pressure gauges of
\Cref{prop:local-typeI-terminal-compactness}, the annular virial observables
\(\mathcal V_{\rho,A}\), the local dissipation \(\mathcal D_{\rho,A}\), and the
right-hand side flux contributions
\[
   R^{-1}\int(\Delta\phi)|V|^2 ,\qquad
   R^{-1}\int|V|^2 U\cdot\nabla\phi ,\qquad
   2R^{-1}\int(\Pi-\pi)V\cdot\nabla\phi
\]
appearing in \eqref{eq:log-annular-virial-identity} are \emph{not} asserted to
be continuous for the weak velocity-pressure topology alone.  Instead, on the
fixed countable rational family of log-windows \((\rho,A)\), these virial,
dissipation, and pressure-flux quantities are included as scalar coordinates of
the realized recurrent critical-tail compactification.  Along every realized
sequence they are uniformly bounded, hence admit simultaneous subsequential
limits by diagonal extraction.  The resulting coordinate functions are
continuous by definition on the corresponding compact virial hull.  In the
strong-compactness subcase, these coordinate limits agree with the ordinary
annular integrals of the limiting coherent profile.
\end{lemma}

\begin{proof}
The required boundedness is precisely the defining property of the
scale-critical virial-tight subhull.  Therefore, on the fixed countable
rational family of log-windows, the virial, dissipation, and flux observables
admit simultaneous subsequential limits by diagonal extraction.  These
quantities are not claimed to be continuous under weak \(L^\infty\)-velocity
and distributional-pressure convergence alone: quadratic energy is only weakly
lower semicontinuous, dissipation is not weakly continuous, and the pressure
flux requires more than distributional convergence against fixed test fields.
Instead, on the virial-tight subhull one adjoins these bounded scalar
coordinates to the recurrent compactification.  By construction the resulting
coordinate maps are continuous on the corresponding compact virial hull.  If at
any stage the homogeneous gauge sequence loses compactness,
\Cref{lem:homogeneous-gauge-compactness} assigns the descendant to \(\Caff\).
In the complementary case the observability hypothesis of
\Cref{lem:recurrent-virial-equality-reduction} is verified by these recorded
coordinates, and the reduction provided by the minimal mesoscopic-scale theorem
\Cref{thm:first-bad-mesoscopic-reduction} closes the diagram.  When the annular
blow-downs converge strongly, the recorded coordinate limits coincide with the
usual virial integrals of the limiting coherent profile.
\end{proof}

\begin{lemma}[Recurrent virial equality case and reduction]
\label{lem:recurrent-virial-equality-reduction}
Let \(\mathcal D\in\Tcrit\) be a recurrent critical-tail defect generated by a
bounded normalized terminal sequence after concentration-set removal and affine
normalization.  Assume that \(\mathcal D\) lies in the scale-critical
virial-tight recurrent subhull of
\Cref{def:scale-critical-virial-tight-subhull}.  Thus the log-annular virial,
dissipation, and flux coordinates in
\Cref{lem:log-annular-oscillation-virial} are bounded continuous observables on
the chosen recurrent hull.  If the defect is not in this virial-tight subhull,
the failure of virial tightness is assigned, by definition of the ordered
critical-tail compactification, to the log-diffuse, pressure/noncompact, or
hidden-scale variance branch before the recurrent virial equality is invoked.
Let \(\mathbb P\) be an
invariant probability measure on that hull for the joint time-translation and
log-translation flow.  Assume also that the averaged pressure gauges are chosen
by a local annular pressure decomposition, so that the three right-hand side
terms in \eqref{eq:log-annular-virial-identity} have limits as defect
observables.  Then the following conclusions hold.

\begin{enumerate}[label=\textup{(\roman*)}]
   \item \emph{Zero-flux equality case.}  If
   \begin{equation}
   \label{eq:zero-virial-flux-assumption}
      \lim_{A\to\infty}
      \int
      \left|
         R^{-1}\int (\Delta\phi)|V|^2
         +R^{-1}\int |V|^2U\cdot\nabla\phi
         +2R^{-1}\int(\Pi-\pi)V\cdot\nabla\phi
      \right|\,d\mathbb P=0,
   \end{equation}
   then every local tangent in the support of \(\mathbb P\) has zero spatial
   gradient.  Consequently
   \[
      \mathcal D\in\Caff .
   \]
   \item \emph{Tight Dirac flux reduction.}  Suppose
   \eqref{eq:zero-virial-flux-assumption} fails, and suppose the nonzero
   limiting flux is neither concentrated on a positive-threshold unit concentrating
   observer nor lost to logarithmic infinity.  Suppose moreover that the
   annular blow-downs selected by this nonzero flux generate a Dirac Young
   measure.  Then the selected annular blow-downs converge strongly in
   \(L^3_{\loc}\) on macroscopic annular cylinders to a nonzero coherent
   critical tail.  Hence
   \[
      \mathcal D\in\Ccohcrit .
   \]
\end{enumerate}
The minimal mesoscopic-scale reduction theorem proved below supplies the tight-Dirac
alternative for realized recurrent defects.  Combined with that theorem, the
preceding alternatives give the inclusion
\[
   \mathcal D\in\Cactive\cup\Caff\cup\Ccohcrit .
\]
\end{lemma}

\begin{proof}
Because we work on the virial-tight recurrent hull, the observables
\(\mathcal V_{\rho,A}\), \(\mathcal D_{\rho,A}\), and \(\mathcal F_{\rho,A}\) are
bounded and continuous.  Hence averaging against the invariant measure
\(\mathbb P\) is legitimate.  We do not invoke infinitesimal generator-domain
theory for \(\mathcal V_{\rho,A}\).  Instead, fix a nonnegative smooth cutoff
\(\eta\in C_c^\infty(\mathbb R^2)\) with \(\eta\equiv1\) on \([-1,1]^2\) and
support in \([-2,2]^2\), and set
\[
   \eta_T(\tau,\zeta):=\eta(\tau/T,\zeta/T),
   \qquad
   N_T:=\iint_{\mathbb R^2}\eta_T(\tau,\zeta)\,d\tau\,d\zeta .
\]
Translate the finite-level identity \eqref{eq:log-annular-virial-identity} by
the joint time/log action, multiply by \(\eta_T\), integrate over
\((\tau,\zeta)\in\mathbb R^2\), and then integrate over the recurrent hull
against \(\mathbb P\).  Because the virial coordinates are recorded as bounded
continuous observables on the virial-tight hull, this averaged identity passes
to the hull level.  Moving the derivative terms onto \(\eta_T\) gives
\[
\begin{aligned}
   &\frac1{N_T}\iint
   \bigl(\partial_\tau\mathcal V_{\rho+\zeta,A}
   +\tfrac12\partial_\rho\mathcal V_{\rho+\zeta,A}\bigr)
   \,\eta_T(\tau,\zeta)\,d\tau\,d\zeta \\
   &\qquad =
   -\frac1{N_T}\iint
   \mathcal V_{\rho+\zeta,A}
   \Bigl(\partial_\tau\eta_T(\tau,\zeta)
      +\tfrac12\partial_\zeta\eta_T(\tau,\zeta)\Bigr)
   \,d\tau\,d\zeta .
\end{aligned}
\]
After integrating also over \(\mathbb P\), invariance of \(\mathbb P\) under the
joint translation action makes the average of each translated observable equal
to its \(\mathbb P\)-mean.  Since
\(\sup|\mathcal V_{\rho,A}|<\infty\) on the virial-tight hull and
\(\|\nabla\eta_T\|_{L^1}/N_T\to0\), the averaged derivative contribution tends
to zero as \(T\to\infty\).  Therefore the averaged identity yields
\[
   2\int \mathcal D_{\rho,A}\,d\mathbb P
   =
   \int \mathcal F_{\rho,A}\,d\mathbb P,
\]
where \(\mathcal F_{\rho,A}\) denotes the sum of the three right-hand side flux
terms in \eqref{eq:log-annular-virial-identity}.  Under
\eqref{eq:zero-virial-flux-assumption} and by nonnegativity of the dissipation,
\[
   \lim_{A\to\infty}\int \mathcal D_{\rho,A}\,d\mathbb P=0 .
\]
Choose a countable exhausting family of ordinary unit covariant cylinders.  For
each cylinder, choose \(A\) large enough that the corresponding cutoff is
identically one on a fixed enlargement.  The preceding limit then implies that,
for \(\mathbb P\)-a.e. profile, the local dissipation
\(\int_Q|\nabla U|^2\) vanishes on every cylinder in the exhausting family.
Hence \(\nabla U=0\) distributionally on every compact cylinder.  Thus every
local tangent is spatially constant.  By the affine normalization dichotomy,
this is precisely the affine/parasitic lower stratum \(\Caff\).

We now prove the tight Dirac flux reduction.  Since the averaged flux does not
vanish and is not concentrated on a unit concentrating observer, there are logarithmic
annuli with nonzero normalized flux but with no positive-threshold unit
constant-in-time oscillatory subcylinder.  Since the flux is not lost to logarithmic
infinity, these annuli may be selected with a tight macroscopic scale
\(R_n\to\infty\).  Passing to the corresponding macroscopic variables gives
bounded blown-down fields
\[
   \widetilde U_n:=\calR_{(y_n,\sigma_n)}U_n,\qquad
   V_n(x,s)=\widetilde U_n(R_nx,s),\qquad
   W_n(x,s)=\widetilde U_n(R_nx,s)-c_n(s)
\]
on every fixed annular cylinder, with pressure and affine gauges normalized by
the exact blow-down definitions of \Cref{def:macroscopic-critical-tail-blowdown}.
Terminal compactness gives a
weak-\(*\) \(L^\infty_{\loc}\) limit and an associated Young measure.  By the
Dirac hypothesis on the selected flux, the Young measure is
\(\delta_W\) almost everywhere.  Uniform boundedness then upgrades convergence
to strong \(L^3_{\loc}\) convergence on annular cylinders for the
affine-normalized representatives, and for the ungauged representatives up to
the deterministic homogeneous translation supplied by the gauge bridge.  Passing
to the limit in the exact ungauged equation
\eqref{eq:critical-tail-exact-blowdown-equation}, the viscous and nonlinear
terms vanish by the factors \(R_n^{-2}\) and \(R_n^{-1}\), and the normalized
pressure gradients converge in annular distributions.  The ungauged limit
therefore satisfies
\[
   \partial_s V+\frac12x\cdot\nabla V+\frac12V+\nabla Q^{\rm act}=0,
   \qquad \nabla\cdot V=0,
\]
on macroscopic annular cylinders.  The field \(W=V-c(s)\) is the corresponding
affine-quotient representative; any term \((c'(s)+\tfrac12c(s))\cdot x\) is
understood as an affine pressure distribution and is not used as an ordinary
pressure without additional time regularity.  The selected flux is nonzero, so
the quotient representative has nonzero spatial quotient oscillation.  Thus
\((W,Q)\in\Ccohcrit\) in the quotient sense of
\Cref{def:critical-diffuse-tail-strata}.
\end{proof}

\subsection{Logarithmic escape as a hidden-scale obstruction}
\label{subsec:logarithmic-escape-hidden-scale}

The logarithmic escape alternative is not ruled out below by a separate global
tail estimate.  The point of the next reduction is narrower and more useful:
once tight annular descendants have no hidden-scale variance, a nonzero
logarithmic escape component cannot survive.  Thus the log-diffuse alternative and
the Young alternative are both reduced to the same minimal mesoscopic-scale
theorem.

\begin{definition}[Log-diffuse recurrent critical defect]
\label{def:log-diffuse-recurrent-critical-defect}
A realized recurrent critical-tail defect is log-diffuse if, after projecting
the normalized diffuse oscillation measures to the logarithmic radial variable
\(\rho=\log |y|\), their log marginals \(\lambda_n\) satisfy
\[
   \sup_{\rho_0\in\mathbb R}
   \lambda_n([\rho_0-A,\rho_0+A])\longrightarrow0
   \qquad\hbox{for every fixed }A<\infty .
\]
Equivalently, the critical-tail compactification has a nonzero logarithmic
escape component \(\Lambda\).  This alternative is denoted by \(\Clogdiff\).
\end{definition}

\begin{definition}[Log-window normalized support profile]
\label{def:log-window-descendant}
Let \(\lambda_n\) be the normalized logarithmic defect marginals of a
log-diffuse critical defect.  A log-window normalized support profile, also
called a log-window support profile, is obtained by choosing
numbers \(A\ge1\), centers \(\rho_n\), and intervals
\[
   I_n=[\rho_n-A,\rho_n+A]
\]
with \(\lambda_n(I_n)>0\), restricting the defect measures to \(I_n\), and
renormalizing:
\[
   \widetilde\lambda_n
   :=
   \frac{\lambda_n\lfloor I_n}{\lambda_n(I_n)} .
\]
After applying the logarithmic recentering \(\rho\mapsto\rho-\rho_n\), the
renormalized measures are tight in the fixed interval \([-A,A]\).  Any
critical-tail profile obtained as a subsequential limit of these normalized
restrictions is called a log-window normalized support profile of the original
log-diffuse defect.  It is a retained descendant only if a separate unit
velocity lower bound is proved.
\end{definition}

\begin{lemma}[Nonzero logarithmic escape has tight log-window support profiles]
\label{lem:log-window-descendant-selection}
Let \(\mathcal D\in\Clogdiff\) be a nonzero realized recurrent critical-tail
defect after concentration-set removal and affine normalization.  Then, for some
fixed \(A\ge1\), \(\mathcal D\) has a nonzero log-window normalized support
profile.  This profile is a tight macroscopic annular critical-tail realization and
retains normalized spatial quotient oscillation mass.
\end{lemma}

\begin{proof}
The logarithmic marginals \(\lambda_n\) are probability measures.  Fix any
\(A\ge1\).  Cover \(\mathbb R_\rho\) by the intervals
\[
   I_{k,A}:=[2Ak-A,2Ak+A],\qquad k\in\mathbb Z .
\]
For each \(n\), at least one interval has positive \(\lambda_n\)-mass.  Choose
one such interval and denote it by \(I_n=[\rho_n-A,\rho_n+A]\).  Restricting to
\(I_n\), renormalizing, and translating \(\rho_n\) to the origin gives
probability measures supported in the fixed compact interval \([-A,A]\).
Compactness of the observer boundary profile class and the local Seregin
compactness of \(\mathfrak X_M\) give a subsequential limit.  By construction,
the limiting profile has total normalized defect mass one, hence is nonzero.

The logarithmic recentering fixes the selected multiplicative annulus, so the
resulting critical-tail realization is tight in logarithmic scale.  The
original diffuse defect was measured after affine normalization by the
oscillation defect measure; the normalized restriction preserves total
oscillation mass equal to one.  Hence the normalized support profile carries nonzero
spatial quotient oscillation.  No positive lower bound for the unnormalized
mass \(\lambda_n(I_n)\) is required, because the profile is a normalized
boundary profile.
\end{proof}

\begin{lemma}[Support transfer for log-window defect compactification]
\label{lem:log-window-support-transfer}
Let \(X_{\log}\) be the compact metrizable log-window defect space obtained by
restricting the critical-tail compactification \(\Tcrit\) to normalized
logarithmic windows, and let
\[
   \lambda_n=\frac{\Lambda_n\lfloor I_n}{\Lambda_n(I_n)}
   \qquad(\Lambda_n(I_n)>0)
\]
be normalized realized log-window restrictions.  If
\(\lambda_n\rightharpoonup\lambda\in\mathcal P(X_{\log})\) and
\(A\subset X_{\log}\) is closed, then support of \(\lambda\) in \(A\) removes
all surviving log-window defect mass outside \(A\); any positive surviving
mass outside \(A\) gives a normalized log-window support profile outside \(A\).
\end{lemma}

\begin{proof}
Apply \Cref{lem:compact-defect-support-transfer} in the compact metrizable
space \(X_{\log}\).  The normalization by \(\Lambda_n(I_n)\) is exactly the
realized support-profile normalization; the size of \(\Lambda_n(I_n)\) is
irrelevant after passing to probability measures.
\end{proof}

\begin{lemma}[Renormalized log-window heredity]
\label{lem:renormalized-log-window-heredity}
Let \(\mathcal D\in\Clogdiff\) be a log-diffuse defect realized from the original blow-up sequence, and
let \(\mathcal F\subset\Tcrit\) be a closed union of alternatives which is
closed under retained descendants and normalized support profiles.  Suppose that every nonzero
log-window normalized support profile of \(\mathcal D\) belongs to
\(\mathcal F\).  Then the
original log-diffuse defect itself has no surviving mass outside
\(\mathcal F\).  In particular, if \(\mathcal F\) has already been excluded
by the lower-stratum closure theorems, then \(\mathcal D\) is impossible.
\end{lemma}

\begin{proof}
Let \(\Lambda\) be the compactified logarithmic defect measure associated with
\(\mathcal D\).  Suppose that \(\Lambda(\Tcrit\setminus\mathcal F)>0\).  Since
\(\mathcal F\) is closed, choose a compact set
\[
   K\subset \Tcrit\setminus\mathcal F
\]
with \(\Lambda(K)>0\).  Pick a support point
\(x_*\in K\cap\supp\Lambda\).  By the definition of support in the compact
log-window defect space and by
\Cref{lem:log-window-support-transfer}, there are logarithmic windows \(I_n\)
with \(\Lambda_n(I_n)>0\) whose normalized restrictions converge to a support
profile represented by \(x_*\).  This limit profile lies in
\(\Tcrit\setminus\mathcal F\), contradicting the hypothesis that every nonzero
log-window normalized support profile of \(\mathcal D\) belongs to
\(\mathcal F\).  Hence \(\Lambda\) is supported on \(\mathcal F\).  Since
\(\mathcal F\) is closed under retained descendants and normalized support
profiles, the original log-diffuse defect has no surviving mass outside
\(\mathcal F\).
\end{proof}

\begin{theorem}[Logarithmic escape reduces to hidden-scale compactness]
\label{thm:log-diffuse-from-hidden-scale}
Assume \Cref{thm:no-hidden-scale-variance-realized-defects}.  Let
\(\mathcal D\in\Clogdiff\) be a nonzero realized log-diffuse critical-tail
defect generated after concentration-set removal and affine normalization.
Then every nonzero log-window normalized support profile of \(\mathcal D\)
belongs to
\[
   \Cactive\cup\Caff\cup\Ccohcrit .
\]
Consequently, after the concentration, affine/parasitic, and coherent critical-tail
alternatives have been removed, no log-diffuse alternative remains.
\end{theorem}

\begin{proof}
By
\Cref{lem:log-window-descendant-selection}, nonzero tight log-window
normalized support profiles exist.  Let \(\widetilde{\mathcal D}\) be any such
profile.  If this support profile regenerates a
positive-threshold unit concentrating observer, then it belongs to \(\Cactive\).  If
its spatial quotient oscillation vanishes, then
\Cref{thm:affine-normalization-dichotomy} places it in \(\Caff\).  It remains
to consider the case in which it is neither concentrating nor affine.

Because \(\widetilde{\mathcal D}\) was obtained by restriction to a fixed
logarithmic window, its logarithmic escape component is zero.  The only
remaining possible obstructions in the critical-tail compactification are a
non-Dirac annular Young defect or a nonzero pressure/viscous defect coordinate.
By
\Cref{thm:no-hidden-scale-variance-realized-defects}, every tight annular
blow-down generated by the descendant has no hidden-scale variance.  Hence
\Cref{lem:hidden-variance-young-obstruction} makes the Young measure Dirac.
The pressure/viscous defect-coordinate alternatives vanish in the tight Dirac
branch by \Cref{lem:critical-tail-defect-coordinate-closure}, after the already
closed affine, pressure/noncompact, concentration, and hidden-scale branches have
been removed.  In the remaining full \(\Tcoh\) case, the coherent extraction
\Cref{cor:coherent-critical-tail-extraction} then gives a nonzero coherent
critical tail; thus \(\widetilde{\mathcal D}\in\Ccohcrit\).

Since \(\widetilde{\mathcal D}\) was arbitrary, every nonzero log-window
normalized support profile of \(\mathcal D\) lies in the closed excluded set
\[
   \mathcal F:=\Cactive\cup\Caff\cup\Ccohcrit .
\]
The set \(\mathcal F\) is closed in the log-window compactification by the
closedness of the positive-threshold activity observables, the closed
affine-normalized oscillation functional, and the coherent-tail extraction
topology used in \Cref{cor:coherent-critical-tail-extraction}.
By \Cref{lem:renormalized-log-window-heredity}, if every log-window normalized
support profile of the original log-diffuse defect lies in \(\mathcal F\), then the
original defect has no surviving mass outside \(\mathcal F\).  Once the concentration,
affine/parasitic, and coherent critical-tail alternatives are removed, this
makes the original log-diffuse defect impossible.
\end{proof}

\begin{corollary}[Log-diffuse alternative excluded by minimal mesoscopic scale reduction]
\label{cor:log-diffuse-branch-exclusion}
Inside the retained residual decomposition,
\[
   \Clogdiff=\varnothing .
\]
Equivalently, after affine/parasitic lower-stratum removal, logarithmic hidden
escape is excluded by the same hidden-scale compactness theorem that removes
the Young alternative.
\end{corollary}

\begin{proof}
By \Cref{thm:first-bad-mesoscopic-reduction},
\Cref{thm:no-hidden-scale-variance-realized-defects} applies to all tight
annular descendants.  Hence \Cref{thm:log-diffuse-from-hidden-scale} reduces
every nonzero log-diffuse defect to a descendant in
\(\Cactive\cup\Caff\cup\Ccohcrit\).  The concentration alternative has been removed by
terminal concentration-set extraction, the affine stratum by
the routing theorem \Cref{thm:oscillatory-entry-normalization}, and the coherent critical-tail
alternative by \Cref{cor:coherent-critical-tail-branch-closure}.  Hence no
log-diffuse alternative remains.
\end{proof}

\subsection{Hidden-scale Young variance}
\label{subsec:hidden-scale-young-variance}

\begin{definition}[Hidden-scale variance]
\label{def:hidden-scale-variance}
Let \(W_n\) be a macroscopic critical-tail blow-down on a compact annular
cylinder
\[
   K\Subset\{0<|x|<\infty\}\times I .
\]
For \(0<\delta<1\), set
\[
   \mathfrak H_K(W_n,\delta)
   :=
   \sup_{\substack{|h|<\delta\\ |\theta|<\delta}}
   \|W_n(\cdot+h,\cdot+\theta)-W_n\|_{L^3(K)} ,
\]
where the translated functions are evaluated on a fixed slightly larger
annular cylinder containing all \((x+h,s+\theta)\) with
\((x,s)\in K\) and \(|h|+|\theta|<\delta\).  The sequence has no hidden-scale
variance on \(K\) if
\[
   \lim_{\delta\downarrow0}\limsup_{n\to\infty}
   \mathfrak H_K(W_n,\delta)=0 .
\]
Because the blow-down is spatial/logarithmic rather than parabolic, the time
shift \(\theta\) is measured in the centered variable itself; only the spatial
displacement can later be magnified into a mesoscopic spatial radius by
multiplication with \(R_n\).
It has no hidden-scale variance if this holds on every compact macroscopic
annular cylinder.
\end{definition}

\begin{lemma}[Hidden-scale variance is the Young obstruction]
\label{lem:hidden-variance-young-obstruction}
Let \(W_n\) be uniformly bounded in \(L^\infty_{\rm loc}\) on macroscopic
annular cylinders and suppose
\[
   W_n\rightharpoonup^\ast \bar W
   \qquad\hbox{in }L^\infty_{\rm loc}.
\]
On a compact annular cylinder \(K\), the following are equivalent:
\begin{enumerate}[label=\textup{(\roman*)}]
   \item \(W_n\to\bar W\) strongly in \(L^3(K)\);
   \item the generated Young measure is Dirac on \(K\):
   \[
      \mathcal Y_{x,s}=\delta_{\bar W(x,s)}
      \qquad\hbox{for a.e. }(x,s)\in K;
   \]
   \item \(W_n\) has no hidden-scale variance on \(K\).
\end{enumerate}
Consequently the non-Dirac Young-measure subcase of a realized Young/defect
critical tail is precisely a failure of the hidden-scale variance condition on
some compact annular cylinder; the pressure/viscous subcases are closed
separately by \Cref{lem:critical-tail-defect-coordinate-closure}.
\end{lemma}

\begin{proof}
The uniform \(L^\infty(K)\) bound gives a uniform \(L^3(K)\) bound.  The
Fr\'echet--Kolmogorov compactness criterion on \(L^3(K)\), applied after
extending the functions to a fixed slightly larger cylinder, says that this
bounded family is relatively compact in \(L^3(K)\) if and only if the
translation modulus in \Cref{def:hidden-scale-variance} tends to zero.  Thus
\textup{(iii)} is equivalent to relative compactness in \(L^3(K)\).  Any strong
subsequential limit must agree with the weak-* limit \(\bar W\), so relative
compactness is equivalent to strong convergence to \(\bar W\).  This proves
\textup{(i)}\(\Leftrightarrow\)\textup{(iii)}.

Strong \(L^3(K)\) convergence implies convergence in measure, hence the
associated Young measure is \(\delta_{\bar W}\).  Conversely, if the generated
Young measure is Dirac and \(\{W_n\}\) is uniformly bounded in \(L^\infty(K)\),
then \(W_n\to\bar W\) in measure on \(K\).  Uniform boundedness gives uniform
integrability of \(|W_n-\bar W|^3\), so Vitali convergence yields strong
\(L^3(K)\) convergence.  Hence
\textup{(i)}\(\Leftrightarrow\)\textup{(ii)}.
\end{proof}

\begin{definition}[Multiscale affine oscillation]
\label{def:multiscale-affine-oscillation}
Let \((U_n,\Pi_n)\) be a terminal sequence.  For a covariant observer
\(z=(\xi,\sigma)\) and \(r\ge1\), define the critical-tail cylinder
\[
   \mathcal Q_{r,\mathrm{ct}}^{\rm cov}(z)
   :=
   \calR_z\bigl(B_r(0)\times(-1,1)\bigr).
\]
Define
\[
   \mathcal A_n(z,r)
   :=
   r^{-3}
   \oscsp\bigl(\calR_zU_n;B_r(0)\times(-1,1)\bigr).
\]
The exponent \(3\) is the spatial dimension; the centered time interval is kept
at fixed size because critical-tail blow-downs are spatial/logarithmic
blow-downs, not parabolic rescalings.  Allowing the minimizing constant to
depend on time removes the harmless spatially homogeneous affine mode.  Unit
concentration-set removal is formulated with \(\actthree\); after the
affine/parasitic quotient has been applied, the mesoscopic residual analysis is
formulated with \(\oscsp\), and \(\sup_z\mathcal A_n(z,1)\) is below the retained
spatial-oscillation threshold on the residual exterior window, after the fixed
cover constants are absorbed into the threshold.
\end{definition}

\begin{definition}[Minimal mesoscopic oscillation scale]
\label{def:first-bad-mesoscopic-scale}
Let \(R_n\to\infty\) be the macroscopic scale of a tight critical-tail
blow-down.  A minimal mesoscopic oscillation scale is a sequence of covariant centers
\(z_n\), radii \(r_n\), and a number \(\varepsilon_0>0\) such that
\[
   1\ll r_n\ll R_n,
   \qquad
   \mathcal A_n(z_n,r_n)\ge\varepsilon_0,
\]
and, after replacing \(r_n\) by a dyadic comparable scale if necessary,
\[
   \sup_{z}\sup_{1\le r\le r_n/2}\mathcal A_n(z,r)
   <\varepsilon_0/2
\]
on the selected residual exterior window.  Thus all smaller retained
subscales are affine-flat, while the scale \(r_n\) is the first scale at which
non-affine oscillation is visible.
\end{definition}

\begin{lemma}[Finite-overlap control of affine oscillation]
\label{lem:finite-overlap-affine-oscillation}
Let \(I\Subset\mathbb R\), let \(f\in L^\infty_{\loc}(\mathbb R^3\times I)\)
with \(\|f\|_{L^\infty}\le M\), and let
\[
   B_1(a_0),B_1(a_1),\dots,B_1(a_N)
\]
be a finite chain of unit balls such that
\[
   |B_1(a_j)\cap B_1(a_{j+1})|\ge\kappa>0
   \qquad (j=0,\dots,N-1).
\]
Assume that for each \(j\) there exists
\(c_j\in L^\infty(I;\mathbb R^3)\), \(|c_j|\le M\), such that
\[
   \int_I\int_{B_1(a_j)}|f(y,s)-c_j(s)|^3\,dy\,ds\le \delta .
\]
Then, for \(0\le j,k\le N\),
\begin{equation}\label{eq:finite-overlap-constants}
   \int_I |c_j(s)-c_k(s)|^3\,ds
   \le C(N,\kappa)\delta ,
\end{equation}
and
\begin{equation}\label{eq:finite-overlap-transfer}
   \int_I\int_{B_1(a_j)}|f(y,s)-c_k(s)|^3\,dy\,ds
   \le C(N,\kappa)\delta .
\end{equation}
Consequently, if all unit cylinders in a bounded chain have affine-normalized
oscillation \(o(1)\), then all their local affine gauges agree up to \(o(1)\)
in \(L^3(I)\).
\end{lemma}

\begin{proof}
For neighboring indices put \(E_j=B_1(a_j)\cap B_1(a_{j+1})\).  Since
\(|E_j|\ge\kappa\), the pointwise inequality
\[
   |c_j(s)-c_{j+1}(s)|^3
   \le
   4\kappa^{-1}\int_{E_j}
   \bigl(|f(y,s)-c_j(s)|^3+|f(y,s)-c_{j+1}(s)|^3\bigr)\,dy
\]
holds for a.e. \(s\).  Integrating in \(s\) gives
\[
   \int_I|c_j-c_{j+1}|^3\,ds\le C_\kappa\delta .
\]
Summing along the chain and using
\[
   \left|\sum_{\ell=j}^{k-1}(c_\ell-c_{\ell+1})\right|^3
   \le |j-k|^2\sum_{\ell=j}^{k-1}|c_\ell-c_{\ell+1}|^3
\]
proves \eqref{eq:finite-overlap-constants}.  Finally,
\[
   |f-c_k|^3\le4|f-c_j|^3+4|c_j-c_k|^3
\]
on \(B_1(a_j)\times I\), and
\eqref{eq:finite-overlap-transfer} follows from
\eqref{eq:finite-overlap-constants}.
\end{proof}

\begin{lemma}[Spacetime finite-overlap translation control]
\label{lem:spacetime-finite-overlap-translation}
Let \(K\Subset\mathbb R^3\times I\), let
\((h_n,\theta_n)\to0\), and let \(f_n\) be uniformly bounded in \(L^\infty\) on
a fixed neighborhood of
\[
   K\cup\bigl(K+(h_n,\theta_n)\bigr).
\]
Assume that every unit covariant cylinder in a finite spacetime chain covering
the segment joining \((x,s)\in K\) to \((x+h_n,s+\theta_n)\) has spatial
quotient oscillation \(o(1)\).  Then there are
spatially homogeneous functions \(a_n,b_n\in L^\infty(I;\mathbb R^3)\) such
that
\[
   \|f_n-a_n\|_{L^3(K)}
   +\|f_n(\cdot+h_n,\cdot+\theta_n)-b_n\|_{L^3(K)}
   =o(1),
\]
and
\[
   \|f_n(\cdot+h_n,\cdot+\theta_n)-f_n-(b_n-a_n)\|_{L^3(K)}=o(1).
\]
Consequently, if the translated difference has a positive \(L^3(K)\) lower
bound while no unit concentrating observer is present, then that lower bound is
carried by the spatially homogeneous time-dependent mode \(b_n-a_n\).  In the
ordered residual decomposition this descendant belongs to the affine/parasitic stratum
\(\Caff\).
\end{lemma}

\begin{proof}
Cover \(K\) and \(K+(h_n,\theta_n)\) by finitely many unit spacetime cylinders
contained in the fixed neighborhood.  Refine the cover so that any two
successive cylinders in the chain have spatial overlap bounded below and time
overlap bounded below by constants depending only on \(K\).  On each cylinder
choose a near-minimizing spatially homogeneous gauge.  Applying
\Cref{lem:finite-overlap-affine-oscillation} on each time-overlap interval and
summing along the finite chain shows that all gauges on the chain agree in
\(L^3\) on their common time intervals up to \(o(1)\).  Gluing the finitely many
local gauges by this overlap relation gives global-in-space functions
\(a_n(s)\) on the part of the chain covering \(K\) and \(b_n(s)\) on the part
covering \(K+(h_n,\theta_n)\).  The first displayed estimate follows from the
local small oscillation and finite overlap.  Subtracting the two approximations
gives the second displayed estimate.

If the translated difference remains bounded below, then
\(\|b_n-a_n\|_{L^3(\operatorname{proj}_sK)}\) stays bounded below.  The
corresponding descendant has zero \(\oscsp\) on every compact cylinder but a
nonzero homogeneous time-dependent component.
That is the affine/parasitic lower stratum in the ordered profile
decomposition.
\end{proof}

\begin{lemma}[Translation defect produces concentration, affine collapse, or a mesoscopic radius]
\label{lem:translation-defect-to-bad-radius}
Let \(K\Subset\{0<|x|<\infty\}\times I\), let
\[
   \widetilde U_n:=\calR_{(y_n,\sigma_n)}U_n,\qquad
   W_n(x,s)=\widetilde U_n(R_nx,s)-c_n(s),
   \qquad R_n\to\infty,
\]
be a bounded macroscopic blow-down, and suppose that for some
\((h_n,\theta_n)\to0\) and \(\eta_0>0\),
\[
   \|W_n(\cdot+h_n,\cdot+\theta_n)-W_n\|_{L^3(K)}\ge\eta_0 .
\]
Since \((h_n,\theta_n)\to0\), the time displacement remains inside a uniformly
finite chain of centered unit time windows.  Only the spatial displacement is
magnified by the spatial/logarithmic blow-down.  Set
\[
   \ell_n:=R_n|h_n|.
\]
Then, after passing to a subsequence, exactly one of the following alternatives
is selected by the ordered priority rule.
\begin{enumerate}[label=\textup{(\roman*)}]
   \item A positive-threshold unit concentrating observer is regenerated.
   \item The defect is carried only by a spatially homogeneous time-dependent
   affine mode; equivalently the descendant is assigned to \(\Caff\).
   \item \(1\ll\ell_n\ll R_n\), and there are covariant centers \(z_n\), radii
   \(r_n\sim\ell_n\), and \(\varepsilon_0=\varepsilon_0(\eta_0,K)>0\) such that
   \[
      \mathcal A_n(z_n,r_n)\ge\varepsilon_0 .
   \]
\end{enumerate}
\end{lemma}

\begin{proof}
Since \((h_n,\theta_n)\to0\), the time displacement remains inside a uniformly
finite chain of centered unit time windows, while \(\ell_n/R_n=|h_n|\to0\).
First assume \(\ell_n\) is bounded.  The two points compared in the original
centered variables are then connected, on the compact set \(K\), by a
uniformly finite spacetime chain of covariant unit cylinders.  If a unit
cylinder in that chain has a positive constant-in-time oscillation lower bound,
alternative (i) holds.  Otherwise the spatial quotient oscillation on every
cylinder in the chain tends to zero.
By \Cref{lem:spacetime-finite-overlap-translation}, the displayed translation
defect can persist only through a spatially homogeneous time-dependent mode.
This is alternative (ii).  Thus, if neither (i) nor (ii) occurs,
\(\ell_n\to\infty\).

In that remaining case \(1\ll\ell_n\ll R_n\).  Pulling back a finite
macroscopic cover of \(K\cup(K+(h_n,\theta_n))\) gives centered cylinders with
spatial radius comparable to \(\ell_n\); the time shift \(\theta_n\) contributes
only a uniformly finite overlap in the centered time variable.  If all such
cylinders had spatial quotient oscillation below a sufficiently small
\(\varepsilon_0=\varepsilon_0(\eta_0,K)\), then
\Cref{lem:spacetime-finite-overlap-translation}, after rescaling the finite
overlap constants, would again force the translated fields to differ by only an
affine/parasitic time mode.  This is alternative (ii), contrary to the present
case.  Therefore one such cylinder has
\(\mathcal A_n(z_n,r_n)\ge\varepsilon_0\) with \(r_n\sim\ell_n\), proving
(iii).
\end{proof}

\begin{lemma}[Failure of hidden-scale compactness produces a minimal mesoscopic scale]
\label{lem:h2-failure-first-bad-scale}
Let \(\widetilde U_n=\calR_{(y_n,\sigma_n)}U_n\) and
\[
   W_n(x,s)=\widetilde U_n(R_nx,s)-c_n(s)
\]
be a tight macroscopic
annular blow-down after concentration-set removal.  If \(W_n\) fails the
hidden-scale variance condition on some compact annular cylinder \(K\), then
either a positive-threshold unit concentrating observer is regenerated, the affine
stratum \(\Caff\) is reached by a descendant, or there is a minimal
mesoscopic oscillation scale in the sense of
\Cref{def:first-bad-mesoscopic-scale}.
\end{lemma}

\begin{proof}
Failure of hidden-scale variance gives a compact annular cylinder \(K\), a
number \(\eta_0>0\), and translations \((h_n,\theta_n)\to0\) such that, after
passing to a subsequence,
\[
   \|W_n(\cdot+h_n,\cdot+\theta_n)-W_n\|_{L^3(K)}\ge \eta_0 .
\]
By \Cref{lem:translation-defect-to-bad-radius}, the translation defect either
regenerates a unit concentrating observer, reaches \(\Caff\), or gives centers
\(z_n\), radii \(\rho_n\) with \(1\ll\rho_n\ll R_n\), and
\(\varepsilon_0>0\) such that
\[
   \mathcal A_n(z_n,\rho_n)\ge\varepsilon_0 .
\]
In particular, pure time-translation defects are already routed to the first
two alternatives and do not create a fictitious mesoscopic spatial radius.
In the last case, let \(r_n\) be the smallest dyadic radius in
\([1,C\rho_n]\) for which
\[
   \sup_z \mathcal A_n(z,r_n)\ge\varepsilon_0
\]
on the selected residual exterior window.  Unit concentration-set removal excludes
\(r_n=O(1)\), so \(r_n\to\infty\).  Since \(r_n\lesssim \rho_n\) and
\(\rho_n/R_n\to0\), one has \(r_n/R_n\to0\).  By dyadic minimality, all
smaller scales up to \(r_n/2\) have oscillation below \(\varepsilon_0/2\), after
slightly lowering \(\varepsilon_0\).  This is a minimal mesoscopic oscillation scale.
\end{proof}

\begin{definition}[Minimal-scale spatial blow-down]
\label{def:first-bad-spatial-blowdown}
Let \(z_n\) and \(r_n\) be a minimal mesoscopic oscillation scale.  Write
\[
   U_n^{z_n}:=\calR_{z_n}U_n,\qquad
   \Pi_n^{z_n}:=\calR_{z_n}\Pi_n .
\]
Choose measurable velocity gauges \(c_n:(-T,T)\to\mathbb R^3\), with
\(|c_n|\le M\), which are minimizing, or within \(o(1)\) of minimizing, the
affine oscillation on each compact time interval.  The minimal-scale spatial
blow-down is
\[
   V_n(x,s):=U_n^{z_n}(r_nx,s),
   \qquad
   W_n(x,s):=V_n(x,s)-c_n(s).
\]
Thus \(V_n\) is the ungauged blow-down and \(W_n\) is the affine-normalized
velocity used to measure non-affine oscillation.  The associated scaled
pressure is
\[
   Q_n^{\rm act}(x,s):=r_n^{-1}\Pi_n^{z_n}(r_nx,s),
\]
defined modulo time-only gauges.  For the affine-quotient topology one writes
\[
   Q_n^\#(x,s):=Q_n^{\rm act}(x,s)-L_n(x,s),
   \qquad L_n(x,s)=b_n(s)\cdot x+a_n(s),
\]
but this affine subtraction is used only in the quotient topology, not inside
the local energy inequality.  No pointwise time derivative of \(c_n\) is
required.
This is not the Navier--Stokes parabolic scaling.  It is the
spatial/logarithmic tail blow-down adapted to the centered residual equation.
The actual pressure normalization is chosen so that
\(\nabla_xQ_n^{\rm act}\) is the rescaled pressure-gradient term appearing in
the centered blow-down equation.
\end{definition}

\begin{remark}[First-bad pressure and gauge convention]
\label{rem:first-bad-pressure-gauge-rule}
Throughout the minimal mesoscopic-scale argument, all distributional equations,
pressure Poisson equations, and local energy inequalities are written for the
ungauged field \(V_n\) and its genuine scaled pressure \(Q_n^{\rm act}\).  The field
\(W_n=V_n-c_n(s)\) is used only for quotient quantities: spatial oscillation,
Young measures, and variance after the homogeneous gauge sequence has either
compactified in time or has been assigned to the affine/parasitic stratum
\(\Caff\).  Thus an ``affine pressure paired with \(c_n\)'' is quotient
notation.  It is not asserted to be an \(L^{3/2}_{\loc}\) pressure, and no
pointwise derivative \(\partial_s c_n\) is used unless compactness of \(c_n\)
has already been obtained.
\end{remark}

\begin{lemma}[Minimal-scale blow-down equation]
\label{lem:first-bad-blowdown-equation}
Let \(V_n,W_n,Q_n^{\rm act},c_n\) be the minimal-scale blow-downs of
\Cref{def:first-bad-spatial-blowdown}.  On every compact cylinder
\(K\Subset\mathbb R^3\times\mathbb R\), \(V_n\) and \(W_n\) are
divergence-free, \(\|V_n\|_{L^\infty(K)}\le M\),
\(\|W_n\|_{L^\infty(K)}\le2M\), and the ungauged field satisfies
\begin{equation}\label{eq:first-bad-scaled-equation}
   \partial_sV_n
   -r_n^{-2}\Delta V_n
   +\frac12x\cdot\nabla V_n+\frac12V_n
   +r_n^{-1}(V_n\cdot\nabla)V_n
   +\nabla Q_n^{\rm act}
   =0
\end{equation}
in distributions.  The affine-normalized field \(W_n=V_n-c_n(s)\) is used for
oscillation and Young-measure compactness.  Its equation is understood only in
the quotient sense made precise in
\Cref{lem:first-bad-quotient-formulation}; no pointwise derivative of \(c_n\)
and no \(L^{3/2}_{\loc}\) affine pressure for \(W_n\) is assumed.
Moreover
\begin{equation}\label{eq:first-bad-pressure-poisson}
   -\Delta Q_n^{\rm act}
   =
   r_n^{-1}\partial_i\partial_j
   \bigl(V_{n,i}V_{n,j}\bigr)
\end{equation}
in distributions, after the same gauges.
\end{lemma}

\begin{proof}
The covariant recentering \(\calR_{z_n}\) is an exact symmetry of the centered
equation, so \((U_n^{z_n},\Pi_n^{z_n})\) solves
\eqref{eq:centered-main}.  Since centered
time is not rescaled,
\[
   \partial_sV_n=(\partial_\tau U_n^{z_n})(r_nx,s),\qquad
   \Delta_yU_n^{z_n}=r_n^{-2}\Delta_xV_n,
\]
\[
   y\cdot\nabla_yU_n^{z_n}=x\cdot\nabla_xV_n,\qquad
   (U_n^{z_n}\cdot\nabla_y)U_n^{z_n}
   =
   r_n^{-1}(V_n\cdot\nabla_x)V_n .
\]
If \(Q_n^{\rm act}=r_n^{-1}\Pi_n^{z_n}(r_nx,s)\), then
\[
   \nabla_xQ_n^{\rm act}=(\nabla_y\Pi_n^{z_n})(r_nx,s).
\]
Substituting these identities gives the scaled equation for \(V_n\).  We do not
treat \(W_n=V_n-c_n(s)\) as carrying an ordinary pressure unless the gauges have
additional time regularity.  Instead, all PDE identities below are first written
for the ungauged field \(V_n\), and the passage to \(W_n\) is made only after
quotienting out spatially homogeneous modes.  Taking the divergence
of the ungauged scaled equation gives
\eqref{eq:first-bad-pressure-poisson}; the linear drift has zero divergence
contribution because \(V_n\) is divergence-free.  The \(L^\infty\) bounds follow
from \(\|U_n\|_{L^\infty}\le M\) and \(|c_n|\le M\).
\end{proof}

\begin{lemma}[Local energy inequality under spatial/logarithmic blow-down]
\label{lem:local-energy-blow-down}
Let \((U_n,\Pi_n)\) be locally suitable centered profiles realized from the
original blow-up sequence in the sense of
\Cref{def:sequence-realized-residual-object}, let \(z_n\) be covariant
centers, and let \(r_n\to\infty\).  Define
\[
   V_n(x,s)=U_n^{z_n}(r_nx,s),
   \qquad
   Q_n^{\rm act}(x,s)=r_n^{-1}\Pi_n^{z_n}(r_nx,s),
\]
Then \((V_n,Q_n^{\rm act})\) satisfy the exact local energy inequality under
spatial/logarithmic blow-down on every relatively compact spacetime domain
\(K\Subset\mathbb R^3\times\mathbb R\): for every nonnegative
\(\phi\in C_c^\infty(K)\),
\begin{equation}\label{eq:blow-down-local-energy}
\begin{aligned}
   r_n^{-2}\!\iint |\nabla V_n|^2\phi\,dx\,ds
   &\le
   \iint \frac{|V_n|^2}{2}
   \Bigl(
      \partial_s\phi
      +r_n^{-2}\Delta\phi
      +\frac12 x\cdot\nabla\phi
      +\frac12\phi
   \Bigr)\,dx\,ds\\
   &\quad
   +r_n^{-1}\!\iint
   \frac{|V_n|^2}{2}V_n\cdot\nabla\phi\,dx\,ds
   +\iint Q_n^{\rm act} V_n\cdot\nabla\phi\,dx\,ds .
\end{aligned}
\end{equation}
Moreover, exactly one of the following alternatives holds:
\begin{enumerate}[label=\textup{(\roman*)},leftmargin=2.2em]
\item a retained affine-pressure descendant or pressure/noncompact defect is
generated and the sequence is routed to the corresponding lower stratum;
\item after time-only pressure gauges, the non-affine branch satisfies
\begin{equation}\label{eq:blow-down-local-energy-bound}
   r_n^{-2}\!\iint |\nabla V_n|^2\phi\,dx\,ds
   \le
   C\!\left(M,\|\phi\|_{C^2},|K|\right)
\end{equation}
with a finite constant uniform in \(n\), and the pressure flux has the stated
distributional limit in the local energy inequality.
\end{enumerate}
All later uses of \eqref{eq:blow-down-local-energy-bound} are in alternative
\textup{(ii)}, after the affine and pressure/noncompact alternatives have been
routed away.
\end{lemma}

\begin{proof}
The setup paper local energy inequality (recalled inside the imported results
\Cref{thm:imported-setup-results}\,(3)) states that for every locally suitable
centered profile \((U,\Pi)\) and every nonnegative
\(\Phi\in C_c^\infty(\mathbb R^3\times\mathbb R)\),
\begin{equation}\label{eq:centered-local-energy-original}
   \iint |\nabla_y U|^2\Phi\,dy\,d\tau
   \le
   \iint \frac{|U|^2}{2}\bigl(\partial_\tau\Phi+\Delta_y\Phi
      +\tfrac12 y\cdot\nabla_y\Phi+\tfrac12\Phi\bigr)\,dy\,d\tau
   +\iint
   \Bigl(\tfrac{|U|^2}{2}+\Pi\Bigr)U\cdot\nabla_y\Phi\,dy\,d\tau .
\end{equation}
This identity is preserved by the exact covariant recentering
\(\calR_{z_n}\), so it applies to \(U_n^{z_n}\) and \(\Pi_n^{z_n}\) with the
same constants.

Apply \eqref{eq:centered-local-energy-original} to \(U_n^{z_n},\Pi_n^{z_n}\)
with the test function
\[
   \Phi(y,\tau):=r_n^{-3}\,\phi(y/r_n,\tau)
\]
where \(\phi\in C_c^\infty(K)\) is fixed.  Change variables
\(y=r_nx\), keeping \(\tau=s\) (centered time is not rescaled).  Then
\[
   \iint|\nabla_yU_n^{z_n}|^2\Phi\,dy\,d\tau
   =\iint |\nabla_yU_n^{z_n}(r_nx,s)|^2 r_n^{-3}\phi(x,s)\,r_n^3\,dx\,ds
   =r_n^{-2}\iint|\nabla_xV_n|^2\phi\,dx\,ds ,
\]
because \(|\nabla_yU_n^{z_n}|^2(r_nx,s)=r_n^{-2}|\nabla_xV_n|^2(x,s)\).
Similarly,
\[
   \iint\frac{|U_n^{z_n}|^2}{2}\partial_\tau\Phi\,dy\,d\tau
   =\iint\frac{|V_n|^2}{2}\partial_s\phi\,dx\,ds ,
\]
\[
   \iint\frac{|U_n^{z_n}|^2}{2}\Delta_y\Phi\,dy\,d\tau
   =r_n^{-2}\iint\frac{|V_n|^2}{2}\Delta_x\phi\,dx\,ds ,
\]
\[
   \iint\frac{|U_n^{z_n}|^2}{2}\tfrac12 y\cdot\nabla_y\Phi\,dy\,d\tau
   =\iint\frac{|V_n|^2}{2}\tfrac12 x\cdot\nabla_x\phi\,dx\,ds ,
\]
\[
   \iint\frac{|U_n^{z_n}|^2}{2}\tfrac12\Phi\,dy\,d\tau
   =\iint\frac{|V_n|^2}{2}\tfrac12\phi\,dx\,ds .
\]
For the convective and pressure-flux terms,
\[
   \iint
   \Bigl(\tfrac{|U_n^{z_n}|^2}{2}+\Pi_n^{z_n}\Bigr)
   U_n^{z_n}\cdot\nabla_y\Phi\,dy\,d\tau
   =
   r_n^{-1}\!\iint
   \tfrac{|V_n|^2}{2}V_n\cdot\nabla_x\phi\,dx\,ds
   +\iint Q_n^{\rm act}V_n\cdot\nabla_x\phi\,dx\,ds ,
\]
using \(Q_n^{\rm act}=r_n^{-1}\Pi_n^{z_n}(r_nx,s)\),
\(\nabla_y\Phi=r_n^{-4}\nabla_x\phi(\cdot/r_n,\cdot)\), and
\(dy=r_n^3dx\).  Substituting
all changes into \eqref{eq:centered-local-energy-original} gives
\eqref{eq:blow-down-local-energy}.

The exact inequality \eqref{eq:blow-down-local-energy} holds before routing.  To
obtain the uniform bound, use the affine pressure coefficient dichotomy for
spatial/logarithmic blow-downs,
\Cref{lem:affine-pressure-coefficient-dichotomy-general}.  If a retained affine
force or pressure/noncompact component occurs, we are in alternative
\textup{(i)}.  Otherwise, in alternative \textup{(ii)},
\(Q_n^{\rm act}=q_n^{\rm na}+a_n(s)+o(1)\) with the affine spatial coefficient
vanishing and the non-affine part converging strongly to zero after time-only
gauges.  The time-only gauge contributes nothing to the pressure flux because
\(\nabla\cdot V_n=0\).  The remaining right-hand side is then bounded by the
\(L^\infty\) bound \(|V_n|\le M\), by \(r_n^{-1}\le1\) and \(r_n^{-2}\le1\) for
all large \(n\), and by H\"older's inequality on the compact support of \(\phi\),
with a constant depending only on \(M\), \(\|\phi\|_{C^2}\), and \(|K|\).  Since
\(r_n\to\infty\), the viscous and convective terms carry the vanishing factors
\(r_n^{-2}\) and \(r_n^{-1}\).  After passing to the time-only pressure-gauge
subsequence and to the bounded-velocity subsequence, the remaining drift and
pressure-flux terms pass to their distributional limits.  This applies both to
minimal mesoscopic blow-downs and to macroscopic critical-tail blow-downs, with
\(r_n=R_n\) in the latter case.
\end{proof}

\begin{lemma}[Minimal-scale quotient formulation]
\label{lem:first-bad-quotient-formulation}
Let \(V_n,W_n,c_n,Q_n^{\rm act},Q_n^\#\) be as in
\Cref{lem:first-bad-blowdown-equation}.  The minimal-scale analysis uses the
following quotient formulation.
\begin{enumerate}[label=\textup{(\roman*)}]
   \item The distributional equation and local energy inequality are always
   applied to the ungauged field \(V_n\) and its actual pressure
   \(Q_n^{\rm act}\), modulo time-only gauges.
   \item The affine-normalized field \(W_n=V_n-c_n(s)\) is used only for
   oscillation functionals and Young measures, which are invariant under
   spatially homogeneous shifts after the gauge sequence is either compact or
   assigned to \(\Caff\).
   \item The affine-quotient pressure \(Q_n^\#=Q_n^{\rm act}-L_n\), with
   \(L_n(x,s)=b_n(s)\cdot x+a_n(s)\), is used only in the quotient topology
   after the affine component has either been routed to \(\Caff\) or discarded
   from the non-affine quotient description.
   \item If \(c_n\) is not compact in \(L^3_{\loc}\) in time, the homogeneous
   gauge sequence itself generates an affine/parasitic descendant.  If
   \(c_n\to c\) strongly in \(L^3_{\loc}\), all Young-measure variances and
   strong compactness statements pass from \(V_n\) to \(W_n\) by deterministic
   translation of the generated Young measures.
\end{enumerate}
Thus the symbol ``affine pressure paired with \(c_n\)'' means a quotient
functional, not an asserted \(L^{3/2}_{\loc}\) pressure inside the local energy
inequality unless the time regularity of \(c_n\) has been separately proved.
\end{lemma}

\begin{proof}
Statement (i) is the construction in
\Cref{lem:first-bad-blowdown-equation}.  The local energy inequality is stable
under the minimal-scale spatial blow-down for \(V_n\), because \(V_n\) is obtained
from an actual centered solution by an exact covariant recentering and a spatial
scaling; no subtraction of \(c_n\) is involved.

For (ii) and (iii), observe that subtracting a spatially homogeneous function
does not change any \(\oscsp\)-type seminorm in which the same class of
homogeneous gauges is minimized over.  If the sequence \(c_n\) is not
compact in \(L^3_{\loc}\), \Cref{lem:homogeneous-gauge-compactness} places the
resulting homogeneous descendant in \(\Caff\).  If \(c_n\to c\) strongly, then
\Cref{lem:variance-homogeneous-shift} identifies the Young measure of \(W_n\)
with the deterministic translate of the Young measure of \(V_n\), and the
variance densities agree.  Hence all compactness and variance conclusions used
for \(W_n\) follow from identities proved for \(V_n\), without requiring
\(\partial_s c_n\) to be a function.
\end{proof}

\begin{lemma}[Covariance of the whole-space Riesz pressure]
\label{lem:whole-space-riesz-pressure-covariance}
At the finite-energy Leray level one may choose the standard whole-space
pressure representative
\[
   p(\cdot,t)=\mathcal R_i\mathcal R_j(u_i u_j)(\cdot,t)
\]
modulo a function of time.  This representative is preserved, modulo the
corresponding time gauge, by physical Type~I scaling, the centered
self-similar change of variables, centered time translation, spatial
translation, and the covariant translations used in this paper.  Local pressure
representatives are restrictions of this whole-space representative, followed
only by the admissible addition of functions of time.
\end{lemma}

\begin{proof}
Let \(K_{ij}\) be the Calder\'on--Zygmund pressure kernel.  It is homogeneous of
degree \(-3\):
\[
   K_{ij}(\lambda x)=\lambda^{-3}K_{ij}(x),\qquad \lambda>0 .
\]
Therefore, if
\[
   u^\lambda(x,t)=\lambda u(x_0+\lambda x,t_0+\lambda^2t),
   \qquad
   p^\lambda(x,t)=\lambda^2p(x_0+\lambda x,t_0+\lambda^2t),
\]
then
\[
   p^\lambda(\cdot,t)
   =
   \mathcal R_i\mathcal R_j
   (u^\lambda_i u^\lambda_j)(\cdot,t)
\]
modulo the scaled time gauge.  Spatial translations commute with convolution by
\(K_{ij}\), and time translations only relabel the time variable.  The centered
change of variables is a time-dependent parabolic scaling with the
pressure factor \(P(y,\tau)=(-t)p(x,t)\); hence the same homogeneity gives the
centered Riesz representative at each fixed centered time.  The covariant
translations used below are compositions of these exact spatial/time changes
and the admissible pressure gauges recorded in the centered equation.  Since
all additional gauges are functions of time, their gradients vanish and their
restrictions to compact cylinders are precisely the local pressure gauges used
in the Seregin topology.
\end{proof}

\begin{lemma}[Harmonic remainders in normalized pressure limits with blow-down coefficient]
\label{lem:local-harmonic-remainders-affine}
Let \((U_n,\Pi_n)\) be a sequence of normalized profiles obtained from the
whole-space representative in
\Cref{lem:whole-space-riesz-pressure-covariance}, before applying any
homogeneous velocity quotient.  Suppose that, before the affine quotient, each
pressure representative has the form
\[
   \Pi_n(\cdot,s)
   =
   \kappa_n\mathcal R_i\mathcal R_j(U_{n,i}U_{n,j})(\cdot,s)+\alpha_n(s),
\]
where \(\kappa_n\to\kappa\in\{0,1\}\), with \(\kappa=1\) for natural-scale
profiles and \(\kappa=0\) for spatial/logarithmic blow-downs.  If
\(U_n\to U\) locally smoothly and, after time-dependent gauges,
\(\Pi_n\to\Pi\) locally in \(L^{3/2}\) or weakly in \(L^{3/2}\), then
\[
   \Pi(\cdot,s)-\kappa\mathcal R_i\mathcal R_j(U_iU_j)(\cdot,s)
\]
is a time-dependent constant.  After the affine/parasitic quotient is applied,
the only additional admissible harmonic pressure direction is a linear function
\(\beta(s)\cdot y\); this linear part is the pressure paired with the
spatially homogeneous parasitic velocity mode and is absent in the
parasitic-free representative.
\end{lemma}

\begin{proof}
Before the affine quotient, set
\[
   P_n^\ast(\cdot,s):=\kappa_n\mathcal R_i\mathcal R_j(U_{n,i}U_{n,j})(\cdot,s).
\]
Since the velocities are uniformly bounded in \(L^\infty\), these
representatives have uniformly bounded BMO seminorms on \(\mathbb R^3\):
\[
   \|P_n^\ast(\cdot,s)\|_{\mathrm{BMO}}
   \le C M^2 .
\]
The BMO bound is stable under the exact scalings and translations in
\Cref{lem:whole-space-riesz-pressure-covariance}, and \(|\kappa_n|\le1\).  Hence,
after passing to a subsequence at each fixed time interval and using a diagonal
extraction, the representatives \(P_n^\ast\) converge weakly-* in local
BMO-duality to a whole-space BMO pressure \(P^\ast\), modulo time gauges.  The
local \(L^{3/2}\)-pressure limit \(\Pi\) agrees with \(P^\ast\) on each compact
cylinder, up to the same time-dependent gauge, because both are local limits of
the same normalized representatives.  Passing to the limit in the pressure
Poisson equation gives
\[
   -\Delta P^\ast(\cdot,s)
   =
   \kappa\,\partial_i\partial_j(U_iU_j)(\cdot,s)
\]
in distributions.  Thus the difference
\[
   h(\cdot,s):=P^\ast(\cdot,s)-\kappa\mathcal R_i\mathcal R_j(U_iU_j)(\cdot,s)
\]
is harmonic in space and belongs to \(\mathrm{BMO}(\mathbb R^3)\) modulo an
additive constant.  An entire harmonic BMO function is constant: indeed, for
each ball \(B_R(x)\), the interior estimate for harmonic functions and the
John--Nirenberg control of mean oscillation give
\[
   |\nabla h(x,s)|
   \le
   C R^{-1}\|h(\cdot,s)\|_{\mathrm{BMO}(\mathbb R^3)}
\]
with \(C\) independent of \(R\); letting \(R\to\infty\) yields
\(\nabla h(\cdot,s)=0\).  Since \(\Pi\) and \(P^\ast\) differ only by the local
time gauge, before affine quotienting the difference between \(\Pi\) and
\(\kappa\mathcal R_i\mathcal R_j(U_iU_j)\) is spatially constant.

Equivalently, if a non-affine harmonic component were present on some compact
ball and time interval, elliptic regularity would give a point where a
nonzero spatial derivative of order at least one remains.  Testing against
compactly supported functions with zero spatial mean, and with zero first
moments when isolating the affine quotient, would detect this component.  Such
a component cannot arise from the uniformly BMO whole-space Riesz
representatives plus time-only gauges described in
\Cref{lem:whole-space-riesz-pressure-covariance}.

The affine/parasitic normalization is the only subsequent operation that
changes the velocity by a spatially homogeneous function.  For a smooth
homogeneous mode \(b(s)\), substitution into the centered equation gives the
affine pressure gradient
\[
   \nabla_y\bigl[(-\partial_s b-\tfrac12 b)\cdot y\bigr].
\]
For merely bounded measurable homogeneous gauges this expression is used only
as the quotient functional specified in
\Cref{def:affine-constant-lower-stratum,rem:first-bad-pressure-gauge-rule,lem:first-bad-quotient-formulation}.
Thus the normalized profile class contains no non-affine harmonic pressure
remainders: only time gauges and, inside \(\Caff\), the explicit affine
parasitic pressure direction are allowed.
\end{proof}

\begin{proposition}[Normalized pressure representation with blow-down coefficient]
\label{prop:normalized-pressure-representation}
Every normalized pressure representative used in the residual profile class has
the form
\begin{equation}\label{eq:normalized-pressure-representation}
   \Pi(\cdot,s)
   =
   \kappa\,\mathcal R_i\mathcal R_j(U_iU_j)(\cdot,s)
   +\beta(s)\cdot y+\alpha(s),
\end{equation}
where \(\alpha\) is a time gauge, \(\beta(s)\cdot y\) is the affine pressure
paired with the spatially homogeneous parasitic velocity mode, and the scalar
\(\kappa\) is determined by the operation which produced the profile.

For natural-scale Seregin profiles, terminal profiles, exact covariant
descendants, and genuine parabolic concentrating descendants, \(\kappa=1\).
For finite spatial/logarithmic blow-down representatives, the non-affine Riesz
part carries coefficient \(r_n^{-1}\) or \(R_n^{-1}\).  For their limiting
blow-down profiles, \(\kappa=0\).  Outside the affine/parasitic
stratum \(\Caff\), the representative is taken with \(\beta=0\).  In
particular, for spatial/logarithmic blow-downs the canonical non-affine Riesz
part carries coefficient \(r_n^{-1}\) or \(R_n^{-1}\) and therefore vanishes in
the limit.  The actual pressure representative used in the local energy
inequality may still contain a spatially affine harmonic component.  Such an
affine component is not subtracted inside the local energy inequality; it is
either shown to vanish by the affine-coefficient dichotomy or routed to
\(\Caff\).
\end{proposition}

\begin{proof}
At natural scale the coefficient is \(1\) by the whole-space Riesz pressure
representative and its covariance, recorded in
\Cref{lem:whole-space-riesz-pressure-covariance}.  The exclusion of non-affine
harmonic pressure remainders in normalized limits, with the limiting
coefficient tracked explicitly, is
\Cref{lem:local-harmonic-remainders-affine}.  Under the spatial/logarithmic
blow-down \(x\mapsto r_nx\), the actual pressure is normalized as
\[
   Q_n^{\rm act}(x,s)=r_n^{-1}\Pi_n(r_nx,s),
\]
so the non-affine Riesz part is
\[
   r_n^{-1}\mathcal R_i\mathcal R_j(V_{n,i}V_{n,j})(x,s).
\]
The same formula with \(R_n^{-1}\) holds for macroscopic critical-tail
blow-downs.  The BMO bound for the Riesz representative gives local
\(L^{3/2}\)-compactness modulo constants, and the prefactor tending to zero
forces the non-affine Riesz part to vanish on compact blow-down windows.  The
only possible surviving harmonic direction is affine, and that direction is
precisely the affine/parasitic quotient.  This yields
\eqref{eq:normalized-pressure-representation} with the stated values of
\(\kappa\).
\end{proof}

\begin{lemma}[Scaled pressure compactness modulo affine gauges]
\label{lem:scaled-pressure-compactness-affine}
Let \((U_n,\Pi_n)\) be bounded centered solutions with
\(\|U_n\|_{L^\infty}\le M\), let \(z_n\) be covariant centers, let
\(r_n\to\infty\), and set
\[
   \widetilde U_n(y,s):=(\calR_{z_n}U_n)(y,s),
   \qquad
   \widetilde\Pi_n(y,s):=(\calR_{z_n}\Pi_n)(y,s).
\]
For
\[
   \widetilde Q_n(x,s):=r_n^{-1}\widetilde\Pi_n(r_nx,s)
\]
and every compact cylinder
\(K=B_R\times I\Subset\mathbb R^3\times\mathbb R\), there are affine functions
\[
   L_n(x,s)=b_n(s)\cdot x+a_n(s)
\]
such that the non-affine pressure part
\[
   q_n^{\rm na}:=\widetilde Q_n-L_n
\]
satisfies, for every finite \(1\le p<\infty\),
\begin{equation}\label{eq:scaled-pressure-nonaffine-Lp}
   \|q_n^{\rm na}\|_{L^p(K)}
   \le C_{p,K,M}r_n^{-1} .
\end{equation}
In particular,
\begin{equation}\label{eq:scaled-pressure-compactness-affine}
   \widetilde Q_n-L_n\to0
   \qquad\hbox{strongly in }L^p(K)
   \quad\text{for every finite }p.
\end{equation}
The same statement holds for the actual minimal-scale pressures
\(Q_n^{\rm act}\); after subtracting the affine gauges, it becomes the quotient
statement for \(Q_n^\#\) in the sense of
\Cref{lem:first-bad-quotient-formulation}.  The same finite-\(p\) non-affine
estimate also holds for macroscopic critical-tail actual pressures
\(\widehat Q_n^{\rm act}\), with \(r_n\) replaced by the critical-tail scale
\(R_n\), on every fixed annular cylinder.  If the affine parts
\(L_n\) carry a nonzero retained descendant, then that descendant lies in the
affine/parasitic stratum \(\Caff\).  Otherwise the affine parts may be
discarded in the affine-normalized quotient, and the non-affine pressure
converges strongly to zero.
\end{lemma}

\begin{proof}
For each \(n\), the covariant recentering is an exact symmetry, so
\((\widetilde U_n,\widetilde\Pi_n)\) is again a bounded centered solution with
\(\|\widetilde U_n\|_{L^\infty}\le M\).  For a.e. \(s\), the pressure
Poisson equation is
\[
   -\Delta\widetilde\Pi_n(\cdot,s)
   =\partial_i\partial_j(\widetilde U_{n,i}\widetilde U_{n,j})(\cdot,s)
\]
in \(\mathbb R^3\).  Let
\[
   \Pi_n^{\rm loc}(\cdot,s)
   :=\mathcal R_i\mathcal R_j(\widetilde U_{n,i}\widetilde U_{n,j})(\cdot,s)
\]
be the whole-space Calder\'on--Zygmund representative.  Since
\(\|\widetilde U_n\|_{L^\infty}\le M\),
\[
   \|\Pi_n^{\rm loc}(\cdot,s)\|_{\mathrm{BMO}(\mathbb R^3)}
   \le C M^2
\]
uniformly in \(s\).  Choose the measurable constant representative
\(a_{n,\rho}(s)=(\Pi_n^{\rm loc})_{B_\rho}(s)\), the spatial average on
\(B_\rho\).  Hence, by the John--Nirenberg inequality, for every ball
\(B_\rho\) and every finite \(1\le p<\infty\), the same constants
\(a_{n,\rho}(s)\) satisfy
\[
   \|\Pi_n^{\rm loc}(\cdot,s)-a_{n,\rho}(s)\|_{L^p(B_\rho)}
   \le C_{p} M^2 |B_\rho|^{1/p}.
\]
After the scaling \(x\mapsto r_nx\), this gives
\[
   \|r_n^{-1}\Pi_n^{\rm loc}(r_n\cdot,s)-a_n(s)\|_{L^p(B_R)}
   \le C_{p,R,M}r_n^{-1}\to0 .
\]
where \(a_n(s)=r_n^{-1}a_{n,r_nR}(s)\).  The bound is uniform for \(s\in I\),
hence the same measurable averages \(a_n(s)\) obey, for every finite \(p\),
\[
   \|r_n^{-1}\Pi_n^{\rm loc}(r_n\cdot,s)-a_n(s)\|_{L^p(B_R\times I)}
   \le C_{p,R,M,I}r_n^{-1}\to0 .
\]

By \Cref{prop:normalized-pressure-representation}, after the spatial blow-down
the non-affine Riesz part carries the explicit coefficient \(r_n^{-1}\), while
the only possible surviving harmonic component is affine.  Thus for each \(n\)
there exist measurable functions \(\beta_n:I\to\mathbb R^3\) and
\(\alpha_n:I\to\mathbb R\) such that
\[
   \widetilde\Pi_n(\cdot,s)
   =\Pi_n^{\rm loc}(\cdot,s)+\beta_n(s)\cdot y+\alpha_n(s)
\]
on \(\mathbb R^3\).  After scaling,
\[
   r_n^{-1}\bigl(\beta_n(s)\cdot r_nx+\alpha_n(s)\bigr)
   =
   \beta_n(s)\cdot x+r_n^{-1}\alpha_n(s),
\]
which is an affine function of \(x\).  Taking
\[
   L_n(x,s):=\beta_n(s)\cdot x+a_n(s)+r_n^{-1}\alpha_n(s),
\]
the remaining pressure is the scaled non-affine Riesz representative already estimated
above.  This proves \eqref{eq:scaled-pressure-nonaffine-Lp} and
\eqref{eq:scaled-pressure-compactness-affine}.  The minimal-scale pressure is used
only up to the affine quotient described in
\Cref{lem:first-bad-quotient-formulation}; in that quotient the same estimate
applies.  The macroscopic critical-tail statement is identical, with the spatial
blow-down parameter denoted \(R_n\), because the normalized pressure
representation of \Cref{prop:normalized-pressure-representation} gives the same
explicit non-affine coefficient \(R_n^{-1}\) on fixed annular windows.
\end{proof}

\begin{lemma}[Affine pressure coefficient dichotomy for spatial/logarithmic blow-downs]
\label{lem:affine-pressure-coefficient-dichotomy}
\label{lem:affine-pressure-coefficient-dichotomy-general}
Assume that on every compact cylinder \(K\Subset\mathbb R^3\times\mathbb R\)
the actual pressure at a spatial/logarithmic blow-down scale \(r_n\to\infty\)
has the affine expansion
\[
   Q_n^{\rm act}(x,s)
   =
   b_n(s)\cdot x+a_n(s)+o_{L^{3/2}(K)}(1)
\]
as supplied by scaled pressure compactness.  Then exactly one of the following
alternatives occurs:
\begin{enumerate}[label=\textup{(\roman*)}]
   \item for some compact time interval \(I\Subset\mathbb R\),
   \[
      \limsup_{n\to\infty}\|b_n\|_{L^{3/2}(I)}>0,
   \]
   and the associated affine pressure/velocity quotient descendant is assigned
   to \(\Caff\);
   \item for every compact time interval \(I\Subset\mathbb R\),
   \[
      b_n\to0
      \qquad\text{in }L^{3/2}(I;\mathbb R^3),
   \]
   and hence, after time-only gauges,
   \[
      Q_n^{\rm act}-a_n(s)\to0
      \qquad\text{strongly in }L^{3/2}_{\loc}.
   \]
\end{enumerate}
\end{lemma}

\begin{proof}
Scaled pressure compactness gives that, on each compact cylinder, the actual
pressure differs from an affine function \(b_n(s)\cdot x+a_n(s)\) by a term
converging strongly to zero in \(L^{3/2}\).  At minimal mesoscopic scales this is
\Cref{lem:scaled-pressure-compactness-affine}; at macroscopic critical-tail
scales it is the affine part of the actual/quotient pressure split in
\Cref{def:critical-pressure-gradient-convergence}.  If \(b_n\) does not converge to
zero on some compact time interval \(I\), then a nontrivial affine pressure
force remains on that interval.  After the same affine normalization used in
the spatial/logarithmic quotient, this produces a nonzero spatially homogeneous
pressure-force descendant; by the affine/parasitic routing convention this
descendant belongs to \(\Caff\).  This is alternative \textup{(i)}.

If alternative \textup{(i)} does not occur, then for every compact interval
\(I\) one must have \(b_n\to0\) in \(L^{3/2}(I)\).  On a compact cylinder
\(K=B_R\times I\),
\[
   \|b_n(s)\cdot x\|_{L^{3/2}(K)}
   \le C_R\|b_n\|_{L^{3/2}(I)}\to0,
\]
so the affine spatial part vanishes in \(L^{3/2}(K)\).  The remaining term is
the time gauge \(a_n(s)\), proving alternative \textup{(ii)}.
\end{proof}

\begin{lemma}[Pressure--velocity product convergence]
\label{lem:pressure-velocity-product-convergence}
Let \(K\Subset\mathbb R^3\times\mathbb R\).  Suppose
\[
   Q_n\to Q\quad\hbox{strongly in }L^{3/2}(K),
   \qquad
   W_n\rightharpoonup^\ast W\quad\hbox{in }L^\infty(K),
\]
with \(\sup_n\|W_n\|_{L^\infty(K)}<\infty\).  Then
\[
   Q_nW_n\to QW
   \qquad\hbox{in }\mathcal D'(K).
\]
\end{lemma}

\begin{proof}
Let \(\varphi\in C_c^\infty(K;\mathbb R^3)\).  Then
\[
   \iint_K (Q_nW_n-QW)\cdot\varphi
   =
   \iint_K (Q_n-Q)W_n\cdot\varphi
   +\iint_K Q(W_n-W)\cdot\varphi .
\]
The first term tends to zero by the strong \(L^{3/2}\) convergence of \(Q_n\)
and the uniform \(L^\infty\) bound on \(W_n\).  Since \(Q\varphi\in L^1(K)\),
the second term tends to zero by weak-* convergence in \(L^\infty(K)\).
\end{proof}

\begin{lemma}[Pressure compactness at a minimal mesoscopic scale]
\label{lem:first-bad-pressure-compactness}
Let \(Q_n^\#\) be the affine-quotient pressures associated to the actual scaled
pressures \(Q_n^{\rm act}\) from \Cref{lem:first-bad-blowdown-equation}.  Assume in
addition the \emph{normalized whole-space Riesz-pressure hypothesis}: each
\(Q_n^{\rm act}\) is the scaled whole-space Riesz-pressure representative, with
blow-down coefficient \(r_n^{-1}\), of the minimal-scale rescaled solution,
normalized modulo time-only functions in the sense of the local pressure atlas
imported in \Cref{thm:imported-setup-results}\,(3); equivalently, each
\(Q_n^{\rm act}\) is the canonical scaled representative produced by
\Cref{lem:local-pressure-decomposition} before affine subtraction, and
\(Q_n^\#\) is obtained from \(Q_n^{\rm act}\) only after passing to the affine
quotient topology.
For every compact cylinder \(K\Subset\mathbb R^3\times\mathbb R\), exactly
one of the following holds:
\begin{enumerate}[label=\textup{(\roman*)}]
   \item the affine pressure/velocity gauges produce a retained descendant in
   \(\Caff\);
   \item after discarding those affine gauges in the affine-normalized quotient,
   there exists a pressure limit \(Q\) such that
\[
   Q_n^\#\to Q
   \qquad\hbox{strongly in }L^{3/2}(K).
\]
   In fact, in the non-affine quotient one may take \(Q=0\).  Consequently, if
   \(W_n\rightharpoonup^\ast \bar W\) in \(L^\infty_{\loc}\), then
\[
   Q_n^\#W_n\rightharpoonup Q\bar W
   \qquad\hbox{in distributions on }K .
\]
\end{enumerate}
\end{lemma}

\begin{proof}
This is \Cref{lem:scaled-pressure-compactness-affine} applied first to the
actual minimal-scale pressures \(Q_n^{\rm act}\), and only afterwards to the
affine-quotient representatives \(Q_n^\#\).  If the affine gauges retain nonzero descendant concentration,
we are in alternative \textup{(i)}.  Otherwise the affine gauges are precisely
the quotient directions removed by affine normalization, and the remaining
pressure converges strongly to \(0\) in \(L^{3/2}(K)\).  Strong convergence of
the pressure and weak-* convergence of the uniformly bounded velocity give the
product convergence in distributions by
\Cref{lem:pressure-velocity-product-convergence}.
\end{proof}

\begin{lemma}[Actual pressure in the minimal-scale energy inequality]
\label{lem:actual-pressure-minimal-scale-energy}
In the minimal-scale blow-down argument, the local energy inequality is always
applied to the ungauged field \(V_n\) and the actual scaled pressure
\(Q_n^{\rm act}\), modulo time-only gauges.  Exactly one of the following
holds:
\begin{enumerate}[label=\textup{(\roman*)}]
\item the affine pressure/velocity component produces a retained
affine-parasitic descendant, hence the profile is assigned to \(\Caff\);
\item after time-only gauge normalization, the actual non-affine pressure
satisfies
\[
   Q_n^{\rm act}-a_n(s)\to0
   \qquad\text{strongly in }L^{3/2}_{\rm loc},
\]
and the pressure flux in the local energy inequality satisfies
\[
   \iint Q_n^{\rm act}V_n\cdot\nabla\phi\,dx\,ds
   \longrightarrow0
   \qquad\text{for every }\phi\in C_c^\infty.
\]
Equivalently,
\[
   \divergence(Q_n^{\rm act}V_n)\to0
   \qquad\text{in distributions}
\]
in the non-affine branch.
\end{enumerate}
No spatially affine pressure subtraction is used inside the local energy
inequality unless its contribution has first been routed to \(\Caff\).
\end{lemma}

\begin{proof}
The actual scaled pressure \(Q_n^{\rm act}=r_n^{-1}\Pi_n^{z_n}(r_n\cdot,s)\)
is the pressure entering \eqref{eq:first-bad-scaled-equation} and the local
energy inequality.  Its non-affine part is the scaled whole-space Riesz
pressure with coefficient \(r_n^{-1}\), so the BMO/John--Nirenberg estimate in
\Cref{lem:scaled-pressure-compactness-affine} and the affine-coefficient
dichotomy of \Cref{lem:affine-pressure-coefficient-dichotomy} imply exactly the
      required two-way alternative.  If an affine pressure/velocity mode survives
after the quotient normalization, then by definition that mode is routed to
\(\Caff\), which is alternative \textup{(i)}.

In the complementary case there are measurable time gauges \(a_n(s)\) such
that
\[
   Q_n^{\rm act}-a_n(s)\to0
   \qquad\hbox{strongly in }L^{3/2}_{\rm loc}.
\]
Because \(V_n\) is uniformly bounded in \(L^\infty_{\loc}\),
\((Q_n^{\rm act}-a_n(s))V_n\to0\) in distributions by
\Cref{lem:pressure-velocity-product-convergence}.  This product convergence is
used only after pairing with \(\nabla\phi\) in the local energy pressure flux.  The
time-only gauge itself does not contribute to that flux, since for every compactly
supported scalar \(\phi\),
\[
   \iint a_n(s)V_n\cdot\nabla\phi
   =-\iint a_n(s)\phi\,\nabla\cdot V_n=0.
\]
Thus the pressure flux is determined by the actual non-affine pressure alone,
and \(\divergence(Q_n^{\rm act}V_n)\to0\) distributionally in the non-affine
quotient.  This is alternative
\textup{(ii)}.
\end{proof}

\begin{lemma}[Homogeneous gauges: compactness modulo the affine stratum]
\label{lem:homogeneous-gauge-compactness}
Let \(c_n:I\to\mathbb R^3\) be uniformly bounded in \(L^\infty(I)\) on every
compact interval \(I\Subset\mathbb R\).  Then, after passing to a subsequence,
exactly one of the following alternatives is used in the ordered minimal-scale
argument:
\begin{enumerate}[label=\textup{(\roman*)}]
   \item \(c_n\to c\) strongly in \(L^3_{\loc}(\mathbb R;\mathbb R^3)\);
   \item the spatially homogeneous sequence
   \[
      C_n(y,s):=c_n(s)
   \]
   has a noncompact homogeneous quotient defect.  This defect has zero spatial
   quotient oscillation on every compact cylinder and is assigned to the
   affine/parasitic stratum \(\Caff\).
\end{enumerate}
\end{lemma}

\begin{proof}
If \(c_n\) has a strongly convergent subsequence in \(L^3_{\loc}\), we pass to
that subsequence and are in alternative \textup{(i)}.  If no such subsequence
exists, then the family is not relatively compact in \(L^3(I)\) for some
compact interval \(I\Subset\mathbb R\).  Since the \(L^\infty(I)\) bound gives
uniform boundedness and tightness on this fixed interval, the only possible
failure in the Fr\'echet--Kolmogorov criterion is time-translation
equicontinuity.  Hence there are translations \(\theta_n\to0\) and
\(\eta>0\) such that
\[
   \|c_n(\cdot+\theta_n)-c_n\|_{L^3(I)}\ge\eta .
\]
The associated spatially homogeneous profiles \(C_n(y,s)=c_n(s)\) have zero
spatial quotient oscillation on every compact cylinder.  The displayed lower
bound records noncompactness only in the homogeneous time-dependent quotient;
it does not assert the existence of a strongly convergent nonzero homogeneous
descendant.  By the affine/parasitic routing convention and
\Cref{lem:homogeneous-gauges-are-affine-stratum}, this quotient defect belongs
to \(\Caff\).  This is alternative \textup{(ii)}.
\end{proof}

\begin{lemma}[Variance is invariant under homogeneous velocity gauges]
\label{lem:variance-homogeneous-shift}
Let \(V_n\) be uniformly bounded in \(L^\infty_{\loc}\), let
\(c_n\to c\) strongly in \(L^3_{\loc}(\mathbb R)\), and set
\[
   W_n(x,s):=V_n(x,s)-c_n(s).
\]
If \(V_n\) generates a Young measure \(\nu_{x,s}\) with barycenter
\(\bar V\), then \(W_n\) generates the translated Young measure
\[
   \widetilde\nu_{x,s}(A)=\nu_{x,s}(A+c(s))
\]
with barycenter \(\bar W=\bar V-c(s)\).  The variance densities agree:
\[
   \langle|\xi-\bar V(x,s)|^2\rangle_{\nu_{x,s}}
   =
   \langle|\xi-\bar W(x,s)|^2\rangle_{\widetilde\nu_{x,s}} .
\]
Consequently any distributional variance inequality proved for \(V_n\) passes
unchanged to the affine-normalized fields \(W_n\).
\end{lemma}

\begin{proof}
After passing to a subsequence, \(c_n(s)\to c(s)\) a.e. and in measure on every
compact time interval.  For every continuous compactly supported function
\(F(\xi)\) and every compactly supported test function \(\varphi(x,s)\),
\[
   \iint \varphi F(W_n)\,dx\,ds
   =
   \iint \varphi F(V_n-c_n(s))\,dx\,ds
   \to
   \iint\varphi\int F(\xi-c(s))\,d\nu_{x,s}(\xi)\,dx\,ds .
\]
This is precisely the translated Young measure.  Translation by a deterministic
vector preserves centered second moments, giving the variance identity.  The
distributional inequality involves only this variance density and therefore is
unchanged.
\end{proof}

\begin{lemma}[Support transfer for the Young variance defect]
\label{lem:young-variance-support-transfer}
Let \(K\Subset\mathbb R^3\times\mathbb R\) be a compact macroscopic cylinder
and let \(M_{\rm tail}=2M\).  Set
\(X_Y(K)=K\times\mathcal P(\overline{B_{M_{\rm tail}}})\), where
\(\overline{B_{M_{\rm tail}}}\) is the compact velocity ball for the
affine-normalized fields \(W_n=V_n-c_n(s)\).  Let
\[
   \lambda_n=\frac{\mu_n^Y\lfloor E_n}{\mu_n^Y(E_n)}
   \qquad(\mu_n^Y(E_n)>0)
\]
be normalized realized Young-measure variance defects on \(X_Y(K)\).  If
\(\lambda_n\rightharpoonup\lambda\) and \(A\subset X_Y(K)\) is closed, then
support of \(\lambda\) in \(A\) removes all surviving Young variance defect
outside \(A\); positive surviving variance outside \(A\) produces a normalized
Young support profile outside \(A\).
\end{lemma}

\begin{proof}
The space \(X_Y(K)\) is compact metrizable because \(K\) is compact and
\(\mathcal P(\overline{B_{M_{\rm tail}}})\) is compact metrizable in the weak
topology.
The conclusion is \Cref{lem:compact-defect-support-transfer} applied to this
space.  This is the only measure-theoretic transfer used when the hidden-scale
Young defect is removed after the variance inequality.
\end{proof}

\begin{lemma}[Renormalization for the coherent barycenter equation]
\label{lem:coherent-barycenter-renormalization}
Let \(Z\in L^\infty_{\loc}(\mathbb R^3\times\mathbb R)\),
\(Q\in L^{3/2}_{\loc}(\mathbb R^3\times\mathbb R)\), and
\(\nabla\cdot Z=0\).  Suppose
\[
   \partial_s Z+\frac12 x\cdot\nabla Z+\frac12 Z+\nabla Q=0
\]
in distributions.  Then
\[
   \partial_s\frac{|Z|^2}{2}
   +\frac12x\cdot\nabla\frac{|Z|^2}{2}
   +\frac12|Z|^2
   +\divergence(QZ)=0
\]
in distributions.
\end{lemma}

\begin{proof}
Let \(\varrho_\varepsilon\) be a standard space--time mollifier and write
\(Z_\varepsilon=\varrho_\varepsilon*Z\), \(Q_\varepsilon=\varrho_\varepsilon*Q\)
on a compact cylinder \(K\), after extending across a slightly larger compact
cylinder.  Since \(x/2\) is smooth and Lipschitz on \(K\),
\[
   \varrho_\varepsilon*((x/2)\cdot\nabla Z)
   -(x/2)\cdot\nabla Z_\varepsilon\to0
   \qquad\hbox{in }L^p(K)
\]
for every finite \(p\).  The mollified equation is therefore
\[
   \partial_s Z_\varepsilon+\frac12x\cdot\nabla Z_\varepsilon
   +\frac12Z_\varepsilon+\nabla Q_\varepsilon=R_\varepsilon,
   \qquad R_\varepsilon\to0\quad\hbox{in }L^p_{\loc}.
\]
Multiplying by \(Z_\varepsilon\) gives
\[
   \partial_s\frac{|Z_\varepsilon|^2}{2}
   +\frac12x\cdot\nabla\frac{|Z_\varepsilon|^2}{2}
   +\frac12|Z_\varepsilon|^2
   +Z_\varepsilon\cdot\nabla Q_\varepsilon
   =R_\varepsilon\cdot Z_\varepsilon .
\]
Because \(\nabla\cdot Z=0\), also \(\nabla\cdot Z_\varepsilon=0\), and hence
\[
   Z_\varepsilon\cdot\nabla Q_\varepsilon
   =\divergence(Q_\varepsilon Z_\varepsilon).
\]
Now \(Z_\varepsilon\to Z\) in \(L^p_{\loc}\) for every finite \(p\), while
\(Z\in L^\infty_{\loc}\), and \(Q_\varepsilon\to Q\) in
\(L^{3/2}_{\loc}\).  Thus \(Q_\varepsilon Z_\varepsilon\to QZ\) in
distributions, and \(R_\varepsilon\cdot Z_\varepsilon\to0\) in distributions.
Letting \(\varepsilon\downarrow0\) gives the claimed renormalized identity.
\end{proof}

\begin{lemma}[Barycenter equation and variance inequality]
\label{lem:first-bad-young-variance}
Let \(V_n,W_n,Q_n^{\rm act},Q_n^\#,c_n\) be a minimal-scale blow-down, where
\(Q_n^{\rm act}\) is the actual scaled pressure entering
\eqref{eq:first-bad-scaled-equation} and \(Q_n^\#\) is its affine-quotient
representative.  Then either the sequence has
a retained affine/parasitic descendant in \(\Caff\), or, after passing to a
subsequence, \(c_n\to c\) strongly in \(L^3_{\loc}\),
\[
   V_n\rightharpoonup^\ast \bar V,\qquad
   W_n\rightharpoonup^\ast \bar W
   \qquad\hbox{in }L^\infty_{\loc},
\]
and \(V_n,W_n\) generate Young measures \(\nu_{x,s}\) and
\(\widetilde\nu_{x,s}\), respectively.  Set
\[
   E(x,s):=\langle |\xi|^2\rangle_{\widetilde\nu_{x,s}},
   \qquad
   \mathcal V(x,s):=E(x,s)-|\bar W(x,s)|^2 .
\]
Then \(\bar V\) satisfies the coherent transport equation with pressure \(Q\).
The affine-normalized barycenter \(\bar W=\bar V-c(s)\) satisfies the same
equation only as a quotient modulo spatially homogeneous modes; after passing
to the parasitic-free representative outside \(\Caff\), this quotient equation
is represented as
\begin{equation}\label{eq:first-bad-barycenter-equation}
   \partial_s\bar W+\frac12x\cdot\nabla\bar W+\frac12\bar W+\nabla Q=0,
   \qquad \nabla\cdot\bar W=0,
\end{equation}
and the nonnegative variance density \(\mathcal V\in L^\infty_{\loc}\)
satisfies
\begin{equation}\label{eq:first-bad-variance-ineq}
   \partial_s\mathcal V+\frac12x\cdot\nabla\mathcal V+\mathcal V\le0
\end{equation}
in distributions.
\end{lemma}

\begin{proof}
By \Cref{lem:homogeneous-gauge-compactness}, either the homogeneous gauge
defect belongs to \(\Caff\), or \(c_n\to c\) strongly in \(L^3_{\loc}\) after
passing to a subsequence.  We work in the second case.
The weak-* compactness of \(V_n,W_n\) and generation of Young measures are
standard consequences of the uniform \(L^\infty_{\loc}\) bound.  Testing
\eqref{eq:first-bad-scaled-equation} against a compactly supported
divergence-free vector field, the viscous term is bounded by
\[
   Cr_n^{-2}\|V_n\|_{L^\infty},
\]
and the nonlinear term by
\[
   Cr_n^{-1}\|V_n\|_{L^\infty}^2.
\]
Both tend to zero.  The quotient-pressure compactness in
\Cref{lem:actual-pressure-minimal-scale-energy} and the coefficient form of
\Cref{prop:normalized-pressure-representation} give either a retained affine
descendant in \(\Caff\), or time gauges \(a_n(s)\) such that
\(Q_n^{\rm act}-a_n(s)\to0\) strongly in \(L^{3/2}_{\loc}\).  Since time gauges
have zero spatial gradient, passing to the limit in
\eqref{eq:first-bad-scaled-equation} first yields the barycenter equation for
\(\bar V\) with the actual pressure representative \(Q=0\) in the non-affine
quotient.  Only after this step do we pass to the affine-quotient description
recorded by \(Q_n^\#\).  Thus \(\bar V\) satisfies the coherent transport
equation.  Since
\(\bar W=\bar V-c(s)\), passing from \(\bar V\) to \(\bar W\) changes only the
spatially homogeneous quotient component.  We do not differentiate \(c\) as a
function.  Instead, \Cref{lem:first-bad-quotient-formulation} says that this is
an affine/parasitic quotient direction.  If it has retained contribution, the profile
is in \(\Caff\); otherwise the canonical parasitic-free representative removes
it, and the representative satisfies
\eqref{eq:first-bad-barycenter-equation}.

We next pass to the local energy inequality for the ungauged fields \(V_n\).
The appropriate form is the scaled suitable inequality of
\Cref{lem:local-energy-blow-down}, applied to nonnegative compactly supported
test functions.  This avoids using a formal pointwise multiplication of the
scaled equation by \(V_n\).  There is no affine pressure work term in this
inequality; the homogeneous gauge has not been subtracted.  The local energy
inequality is applied to \(V_n\) with the actual pressure
\(Q_n^{\rm act}\), not with the affine-quotient representative.  By
\Cref{lem:actual-pressure-minimal-scale-energy}, either the affine pressure
component has already routed the sequence to \(\Caff\), or
\(Q_n^{\rm act}-a_n(s)\to0\) strongly in \(L^{3/2}_{\rm loc}\), and the
time-only gauge contributes zero to \(\divergence(Q_n^{\rm act}V_n)\).  In the
latter case the pressure flux converges distributionally to
\(\divergence(Q\bar V)\), with \(Q=0\) in the non-affine quotient.  In
\eqref{eq:blow-down-local-energy}, the viscous diffusion and nonlinear-flux
terms vanish in the limit because of the factors \(r_n^{-2}\) and
\(r_n^{-1}\), while the dissipation term is nonnegative and may be dropped.
The strong actual-pressure convergence modulo time-only gauges gives
\[
   \divergence(Q_n^{\rm act}V_n)\to\divergence(Q\bar V).
\]
Therefore
\[
   \partial_s\frac{E_V}{2}
   +\frac12x\cdot\nabla\frac{E_V}{2}
   +\frac12E_V
   +\divergence(Q\bar V)
   \le0
\]
in distributions, where
\[
   E_V(x,s):=\langle|\xi|^2\rangle_{\nu_{x,s}}
\]
is the second moment of the Young measure generated by \(V_n\).  On the other
hand, applying \Cref{lem:coherent-barycenter-renormalization} to the
barycenter equation for \(\bar V\) gives
\[
   \partial_s\frac{|\bar V|^2}{2}
   +\frac12x\cdot\nabla\frac{|\bar V|^2}{2}
   +\frac12|\bar V|^2
   +\divergence(Q\bar V)=0 .
\]
Subtracting gives the variance inequality for the variance of \(V_n\).  By
\Cref{lem:variance-homogeneous-shift}, the same variance density is
\(\mathcal V\), the variance of \(W_n\).  This proves
\eqref{eq:first-bad-variance-ineq}.
\end{proof}

\begin{lemma}[Bounded ancient variance Liouville theorem]
\label{lem:bounded-variance-liouville}
Let \(F\in L^\infty_{\loc}(\mathbb R^3\times\mathbb R)\) satisfy
\[
   0\le F\le C<\infty
\]
and
\[
   \partial_sF+\frac12x\cdot\nabla F+F\le0
\]
in distributions on \(\mathbb R^3\times\mathbb R\).  Then \(F\equiv0\).
\end{lemma}

\begin{proof}
\emph{Step 1: renormalization.}  Set \(G(x,s):=e^{s}F(x,s)\), so that
\(0\le G\le Ce^{s}\).  Since \(G=e^sF\), the product rule gives, in
distributions,
\[
   \partial_sG+\frac12x\cdot\nabla G
   =
   e^s\left(\partial_sF+\frac12x\cdot\nabla F+F\right)\le0.
\]
Equivalently, for every smooth nonnegative compactly supported test function
\(\psi\), the distributional hypothesis applied to \(e^s\psi\ge0\) yields
\begin{equation}\label{eq:variance-renormalized-transport}
   \partial_sG+\frac12 x\cdot\nabla G\le 0
   \qquad\hbox{in distributions on }\mathbb R^3\times\mathbb R .
\end{equation}

\emph{Step 2: backward adjoint semigroup.}  The forward characteristic flow of
\(\widetilde L:=\partial_s+\frac12 x\cdot\nabla\) starting from \((x_0,s_0)\) is
\(x(\sigma)=e^{(\sigma-s_0)/2}x_0\).  Define, for \(T\ge 0\), the backward
adjoint operator on test functions
\[
   (\mathcal S_T\psi)(x,s)
   :=
   e^{-3T/2}\psi\!\bigl(e^{-T/2}x,\,s-T\bigr) .
\]
The prefactor \(e^{-3T/2}\) is the Jacobian of \(x\mapsto e^{-T/2}x\)
on \(\mathbb R^3\), so \(\int\mathcal S_T\psi\,dx\,ds=\int\psi\,dx\,ds\) and
the family \(\{\mathcal S_T\}_{T\ge 0}\) is the adjoint of the forward
characteristic flow on \(L^1(\mathbb R^3\times\mathbb R)\).  In particular, by
the substitution \(x=e^{T/2}x'\), \(s=s'+T\) (Jacobian \(e^{3T/2}\)),
\begin{equation}\label{eq:variance-pairing-substitution}
   \langle G,\mathcal S_T\psi\rangle
   =\iint G(e^{T/2}x',s'+T)\,\psi(x',s')\,dx'\,ds' .
\end{equation}

\emph{Step 3: monotonicity in \(T\).}  Fix \(\psi\ge 0\) smooth and compactly
supported.  A direct computation using the chain rule shows that
\[
   \frac{d}{dT}\mathcal S_T\psi
   =
   \widetilde L^\ast(\mathcal S_T\psi)
\]
where \(\widetilde L^\ast=-\partial_s-\frac12 x\cdot\nabla-\frac32\).
Therefore
\[
   \frac{d}{dT}\langle G,\mathcal S_T\psi\rangle
   =\langle G,\widetilde L^\ast(\mathcal S_T\psi)\rangle
   =\langle \widetilde L G,\mathcal S_T\psi\rangle .
\]
Since \(\mathcal S_T\psi\ge 0\) and \(\widetilde L G\le 0\) distributionally by
\eqref{eq:variance-renormalized-transport},
\begin{equation}\label{eq:variance-monotone}
   \frac{d}{dT}\langle G,\mathcal S_T\psi\rangle\le 0 .
\end{equation}

\emph{Step 4: test-function vanishing and conclusion.}  Let
\(\varphi\in C_c^\infty(\mathbb R^3\times\mathbb R)\) be nonnegative and, for
\(T\ge0\), define
\[
   \psi_T(x,s)
   :=
   e^{3T/2}\varphi(e^{T/2}x,s+T).
\]
Then \(\psi_T\ge0\), \(\psi_T\in C_c^\infty\), and a direct computation gives
\(\mathcal S_T\psi_T=\varphi\).  Applying \eqref{eq:variance-monotone} with
\(\psi=\psi_T\) and evaluating at time \(T\) yields
\[
   \langle G,\varphi\rangle
   =
   \langle G,\mathcal S_T\psi_T\rangle
   \le
   \langle G,\psi_T\rangle .
\]
Since \(0\le G(x,s)\le Ce^s\),
\[
\begin{aligned}
   \langle G,\psi_T\rangle
   &\le
   C\iint e^s e^{3T/2}\varphi(e^{T/2}x,s+T)\,dx\,ds .
\end{aligned}
\]
Make the change of variables \(x'=e^{T/2}x\), \(s'=s+T\).  Then
\(dx\,ds=e^{-3T/2}dx'\,ds'\), so
\[
   \langle G,\psi_T\rangle
   \le
   Ce^{-T}\iint e^{s'}\varphi(x',s')\,dx'\,ds'
   \xrightarrow[T\to\infty]{}0 .
\]
Hence \(\langle G,\varphi\rangle=0\) for every nonnegative test function
\(\varphi\).  Therefore \(G=0\) distributionally, so also almost everywhere.
Since \(G=e^sF\), it follows that \(F\equiv0\).
\end{proof}

\begin{lemma}[Minimal-scale blow-downs are strongly compact]
\label{lem:first-bad-strong-compactness}
Every minimal-scale blow-down \(W_n\) either has a retained affine/parasitic
descendant in \(\Caff\), or has, after passing to a subsequence, a limit
\(\bar W\) such that
\[
   W_n\to\bar W
   \qquad\hbox{strongly in }L^3_{\loc}(\mathbb R^3\times\mathbb R).
\]
The limit satisfies \eqref{eq:first-bad-barycenter-equation}.
\end{lemma}

\begin{proof}
By \Cref{lem:first-bad-young-variance}, either the affine/parasitic alternative
has occurred, or the Young-measure variance \(\mathcal V\) satisfies the
hypotheses of
\Cref{lem:bounded-variance-liouville}.  Hence \(\mathcal V\equiv0\).  Thus
\[
   \widetilde\nu_{x,s}=\delta_{\bar W(x,s)}
   \qquad\hbox{for a.e. }(x,s).
\]
Uniform \(L^\infty_{\loc}\) boundedness then upgrades convergence in measure to
strong convergence in \(L^3_{\loc}\).  The limit equation has already been
proved in \Cref{lem:first-bad-young-variance}.
\end{proof}

\begin{theorem}[Minimal mesoscopic oscillation scale reduction]
\label{thm:first-bad-mesoscopic-reduction}
Every minimal mesoscopic oscillation scale generated by a realized recurrent
critical-tail extraction after concentration-set removal and affine normalization
has a descendant in
\[
   \Cactive\cup\Caff\cup\Ccohcrit .
\]
\end{theorem}

\begin{proof}
Let \(z_n,r_n\) be a minimal mesoscopic oscillation scale.  If, before passing to the
minimal-scale blow-down, some unit covariant observer in the selected residual
window carries a positive constant-in-time oscillation threshold, then the
concentration alternative is regenerated and the descendant lies in \(\Cactive\).  It remains to
consider the case in which unit concentration has been removed at all fixed positive
thresholds.

By construction, this theorem is invoked only after
\Cref{lem:translation-defect-to-bad-radius,lem:h2-failure-first-bad-scale}
have already produced a genuine spatial mesoscopic scale
\(1\ll r_n\ll R_n\); pure time-translation defects were routed earlier to
regenerated concentration or to the affine/parasitic stratum.

Form the minimal-scale blow-down \(W_n\) of
\Cref{def:first-bad-spatial-blowdown}.  By
\Cref{lem:first-bad-strong-compactness}, either the affine/parasitic stratum
\(\Caff\) is reached, or, after passing to a subsequence,
\[
   W_n\to W
   \qquad\hbox{strongly in }L^3_{\loc},
\]
and \((W,Q)\) satisfies the coherent critical-tail transport equation
\[
   \partial_sW+\frac12x\cdot\nabla W+\frac12W+\nabla Q=0,
   \qquad \nabla\cdot W=0 .
\]
The minimal-scale lower bound and strong convergence give
\[
   \oscsp\bigl(W;B_1(0)\times(-1,1)\bigr)
   \ge \varepsilon_0 .
\]
Thus the limit is not spatially affine/parasitic on the unit critical-tail
window.  If this non-affine lower bound were lost in a further descendant, the
descendant would be assigned to \(\Caff\) by
\Cref{thm:affine-normalization-dichotomy}.  Otherwise the strongly realized
nonzero limit is a coherent critical-tail profile on annular cylinders.
Consider its compact joint log/time translation hull inside the critical-tail
compactification.  If some hull descendant regenerates local concentration or
collapses to the affine/parasitic quotient, then the branch is assigned to
\(\Cactive\) or \(\Caff\), respectively.  Otherwise, by compactness of the hull
and descendant heredity
\Cref{lem:diagonal-lifting-descendants,thm:descendant-heredity}, the hull
contains a sequence-realized recurrent support profile with the same nonzero
spatial quotient oscillation lower bound.  This recurrent support profile is a
coherent critical-tail descendant, hence belongs to \(\Ccohcrit\).  Therefore every minimal mesoscopic oscillation scale is assigned to
\(\Cactive\), \(\Caff\), or \(\Ccohcrit\), as claimed.
\end{proof}

\begin{remark}[Critical-tail reduction theorem]
\label{rem:remaining-hidden-scale-input}
\Cref{thm:first-bad-mesoscopic-reduction} is the critical-tail reduction
theorem proved above.  It is a local multiscale rigidity statement: oscillation
first appearing at an intermediate scale
\(1\ll r_n\ll R_n\) must either regenerate a concentrating observer, collapse into
the affine/parasitic lower stratum, or produce a coherent critical tail.  The
log-diffuse alternative has been reduced above to the same hidden-scale compactness
statement through tight log-window normalized support profiles.
\end{remark}

\begin{theorem}[No hidden-scale variance from minimal mesoscopic scale reduction]
\label{thm:no-hidden-scale-variance-realized-defects}
Every tight macroscopic
annular blow-down generated by a realized recurrent critical-tail extraction
after concentration-set removal and affine normalization has no hidden-scale variance
on compact macroscopic annular cylinders.
\end{theorem}

\begin{proof}
Suppose that hidden-scale variance fails.  By
\Cref{lem:h2-failure-first-bad-scale}, either a unit concentrating observer is
regenerated, the affine stratum is reached, or a minimal mesoscopic oscillation scale
exists.  The first alternative contradicts concentration-set removal.  The affine
alternative is routed by \Cref{thm:oscillatory-entry-normalization} to
\(\Caff\), hence it does not produce a retained normalized residual
representative.  In the
third alternative,
\Cref{thm:first-bad-mesoscopic-reduction} gives a descendant in
\(\Cactive\cup\Caff\cup\Ccohcrit\).  The concentration case again contradicts
concentration-set removal; the affine case is routed by
\Cref{thm:oscillatory-entry-normalization} to \(\Caff\), hence it is not a
retained normalized representative; and
the coherent critical-tail case is closed by
\Cref{cor:coherent-critical-tail-branch-closure}.  Thus all
alternatives are impossible, and hidden-scale variance cannot fail.
\end{proof}

\begin{definition}[Macroscopic critical-tail blow-down map]
\label{def:macroscopic-critical-tail-blowdown}
Let \((U_n,\Pi_n)\subset\mathfrak X_M\), let
\(z_n=(y_n,\sigma_n)\) be covariant observer centers, and let
\(R_n\to\infty\) be a selected macroscopic scale.  Put
\[
   \widetilde U_n:=\calR_{z_n}U_n,\qquad
   \widetilde\Pi_n:=\calR_{z_n}\Pi_n .
\]
The ungauged spatial/logarithmic blow-down is
\[
   V_n(x,s):=\widetilde U_n(R_nx,s),
   \qquad
   \widehat Q_n^{\rm act}(x,s):=R_n^{-1}
      \bigl(\widetilde\Pi_n(R_nx,s)-a_n(s)\bigr),
\]
where \(a_n(s)\) is the time-dependent pressure gauge.  If an affine pressure
component
\[
   L_n(y,s)=b_n(s)\cdot y+a_n(s)
\]
is removed in the affine quotient, the corresponding quotient pressure is
\[
   Q_n^\#(x,s):=R_n^{-1}
      \bigl(\widetilde\Pi_n(R_nx,s)-L_n(R_nx,s)\bigr).
\]
Velocity oscillations are measured after a spatially homogeneous velocity gauge
\[
   W_n(x,s):=V_n(x,s)-c_n(s),
   \qquad |c_n(s)|\le M .
\]
The map is not the Navier--Stokes parabolic scaling: centered time is not
rescaled.  The scaling identities are
\[
\begin{array}{c|c}
\text{term} & \text{blow-down size}\\ \hline
\partial_s V_n & 1\\
\Delta V_n & R_n^{-2}\\
(V_n\cdot\nabla)V_n & R_n^{-1}\\
\frac12 y\cdot\nabla_y U_n & \frac12 x\cdot\nabla_xV_n\\
\nabla_y\widetilde\Pi_n & \nabla_x\widehat Q_n^{\rm act} .
\end{array}
\]
Thus \(V_n\) satisfies the exact distributional equation
\begin{equation}\label{eq:critical-tail-exact-blowdown-equation}
   \partial_sV_n
   -R_n^{-2}\Delta V_n
   +\frac12x\cdot\nabla V_n+\frac12V_n
   +R_n^{-1}(V_n\cdot\nabla)V_n
   +\nabla\widehat Q_n^{\rm act}
   =0 .
\end{equation}
If one passes to the affine-quotient representative
\(Q_n^\#=\widehat Q_n^{\rm act}-b_n(s)\cdot x\), then
\[
   \nabla\widehat Q_n^{\rm act}=\nabla Q_n^\#+b_n(s).
\]
Thus the exact equation is not written with \(Q_n^\#\) alone unless the affine
force \(b_n(s)\) has already been routed to \(\Caff\) or removed by the
canonical affine normalization.  The field \(W_n\) is used for oscillation and
Young-measure compactness, not as an ungauged Navier--Stokes solution unless the
time regularity of \(c_n\) has separately been obtained.
\end{definition}

\begin{lemma}[Critical-tail homogeneous gauge bridge]
\label{lem:critical-tail-homogeneous-gauge-bridge}
Let \(V_n\) be ungauged macroscopic critical-tail blow-downs on fixed annular
cylinders, and let
\[
   W_n(x,s):=V_n(x,s)-c_n(s)
\]
be the affine-normalized representatives with \(|c_n(s)|\le M\) on compact
time intervals.  After passing to a subsequence, exactly one of the following
alternatives occurs:
\begin{enumerate}[label=\textup{(\roman*)}]
   \item the homogeneous gauge sequence has a noncompact quotient defect, and by
   \Cref{lem:homogeneous-gauge-compactness} the corresponding descendant is
   assigned to \(\Caff\);
   \item \(c_n\to c\) strongly in \(L^3_{\loc}(\mathbb R;\mathbb R^3)\).  If
   \(\mathcal Y^V_{x,s}\) and \(\mathcal Y^W_{x,s}\) are the Young measures
   generated by \(V_n\) and \(W_n\), respectively, then
   \[
      \mathcal Y^V_{x,s}
      =
      (\xi\mapsto \xi+c(s))_\#\mathcal Y^W_{x,s},
      \qquad
      \mathcal Y^W_{x,s}
      =
      (\xi\mapsto \xi-c(s))_\#\mathcal Y^V_{x,s}.
   \]
   In particular, Dirac/non-Dirac status, hidden-scale variance, and spatial
   quotient oscillation pass between the ungauged and affine-normalized
   representatives modulo the affine quotient.
\end{enumerate}
\end{lemma}

\begin{proof}
Apply \Cref{lem:homogeneous-gauge-compactness} to the spatially homogeneous
gauge sequence \(c_n\).  This gives the dichotomy between compactness in
\(L^3_{\loc}\) and routing to \(\Caff\).  In the compact case,
\Cref{lem:variance-homogeneous-shift} gives exactly the deterministic
translation relation between the Young measures generated by \(V_n\) and
\(W_n\), together with preservation of variance inequalities.  Since the
spatial quotient oscillation functionals are defined modulo spatially homogeneous
velocity gauges, they are unchanged by adding or subtracting \(c(s)\).
\end{proof}

\begin{theorem}[Critical-tail compactification extraction]
\label{thm:critical-tail-compactification-extraction}
Let \((U_n,\Pi_n)\subset\mathfrak X_M\) be a bounded normalized terminal
residual sequence, \(\|U_n\|_{L^\infty}\le M\).  Assume that, after removal of
all affine-normalized concentrating covariant observers, there remains a nonzero
diffuse exterior oscillation defect.  Then, after passing to a subsequence,
exactly one of the following alternatives occurs:
\begin{enumerate}[label=\textup{(\roman*)},leftmargin=2.2em]
\item a retained affine-pressure or homogeneous-gauge descendant is generated
and the branch is routed to \(\Caff\);
\item a pressure/noncompact descendant is generated and routed to the
corresponding pressure/noncompact branch;
\item after alternatives \textup{(i)}--\textup{(ii)} have been routed away, the
retained critical-tail data generate a nonzero element
\[
   \mathcal T=(\bar W,\mathcal Y,\Lambda;
   \widehat{\mathscr P},\mathscr D,
   \mathfrak q_{\rm aff},\mathfrak v_{\rm vir},\mathfrak s)\in\Tcrit .
\]
This retained alternative is read in priority order: if \(\Lambda\ne0\), it is
the log-diffuse full state, whose annular coordinates are inactive until
log-window selection; if \(\Lambda=0\), it is an annular retained state whose
pressure and viscous defect coordinates are uniformly bounded on the fixed
countable test family.
\end{enumerate}
In the annular retained subcase of alternative \textup{(iii)}, the ungauged
barycentric velocity satisfies the linear critical-tail
transport law
\[
   \partial_s\bar V+\frac12 x\cdot\nabla\bar V+\frac12\bar V+\nabla Q^{\rm act}=0,
   \qquad
   \nabla\cdot\bar V=0
\]
in distributions on macroscopic annular cylinders.  The affine-normalized
barycenter \(\bar W=\bar V-c(s)\) is the corresponding quotient representative;
its equation is understood modulo spatially homogeneous velocity modes and
affine pressure forces, where \(Q\) is obtained by first passing to the limit in
the actual pressure gradients modulo time gauges and only afterwards projecting
to the non-affine affine-quotient class.  The
compactification stores the affine-normalized Young/log state
\((\mathcal Y,\Lambda)\) generated by \(W_n\), together with the pressure,
viscous, affine-quotient, virial, and support coordinates in the displayed
tuple.  The ungauged Young measure
\(\mathcal Y^V\) generated by the exact PDE variables \(V_n\) satisfies the
balance identity of \Cref{lem:critical-tail-young-log-transport}; outside
\(\Caff\), \Cref{lem:critical-tail-homogeneous-gauge-bridge} identifies
\(\mathcal Y^V\) and \(\mathcal Y\) by deterministic translation, so their
Dirac/non-Dirac and hidden-scale alternatives agree.  In particular the
log-diffuse subcase is exactly the case \(\Lambda\ne0\), and the annular Young
subcase satisfies the induced log-dilation balance while nonlinear velocity
observables retain the pressure and viscous defect pairings explicitly recorded
in \eqref{eq:young-transport-identity}.
\end{theorem}

\begin{proof}
Let \(\nu_n^{\rm sp}\) be the spatial quotient oscillation defect measure
after concentration-set removal, as in
\Cref{lem:spatial-quotient-defect-cKN-control}.  By the mixed-diffuse
hypothesis there are escaping
observer windows \(E_n\) and \(m_0>0\) such that
\[
   \nu_n^{\rm sp}(E_n)\ge m_0 .
\]
Since all positive-threshold constant-in-time oscillatory observers have been
removed and
\[
   \oscsp\bigl(U;\mathcal Q_1^{\rm cov}(z)\bigr)
   \le
   \actthree\bigl(U;\mathcal Q_1^{\rm cov}(z)\bigr),
\]
the atom size of \(\nu_n^{\rm sp}\) on \(E_n\) tends to zero.  Applying
\Cref{lem:finite-atomic-observer-window-selection} to
\(\gamma_n=\nu_n^{\rm sp}\), after passing to a subsequence we may replace
\(E_n\) by an escaping finite net-cell subset, still denoted \(E_n\), such
that
\[
   m_0\le \nu_n^{\rm sp}(E_n)\le 2m_0 .
\]
Normalize
\[
   \lambda_n:=
   \frac{\nu_n^{\rm sp}\lfloor E_n}{\nu_n^{\rm sp}(E_n)} .
\]
By \Cref{lem:spatial-quotient-defect-cKN-control} and the proof of
\Cref{thm:diffuse-defect-compactness}, after passing to a subsequence these
probability measures converge to a boundary diffuse-defect profile.  Because all
positive-threshold concentrating observers have been removed in the stronger
\(\actthree\)-sense, the limit gives no mass to ordinary unit concentrating observers
in \(Z\).

We now compactify the critical boundary and split according to annular
tightness.  Apply Prokhorov compactness first in the observer compactification
and then in the logarithmic radial compactification.

First suppose the normalized spatial quotient defect is not tight in any fixed
multiplicative annulus.  Then its logarithmic pushforward has nonzero mass at
the logarithmic boundary.  We define the critical-tail state with
\(\Lambda\ne0\); this is the log-diffuse critical-tail alternative inside
\(\Tcrit\).  In this coarse log-boundary representative the annular Young,
pressure, and viscous slots are not invoked to construct the branch; the
annular coordinates enter only after passing to normalized log-window support
profiles.  The subsequent log-window selection is handled by
\Cref{lem:log-window-descendant-selection,lem:renormalized-log-window-heredity}.
This gives the log-diffuse subcase of alternative \textup{(iii)}.

It remains to consider the annular-tight case.  Then select macroscopic scales
\(R_n\to\infty\), observers \(z_n=(y_n,\sigma_n)\), affine velocity gauges
\(c_n(s)\), and pressure gauges.  On fixed annular cylinders use the exact
blow-down map of \Cref{def:macroscopic-critical-tail-blowdown}; in this case no
mass is assigned to the logarithmic boundary, so \(\Lambda=0\):
\[
   V_n(x,s)=\widetilde U_n(R_nx,s),\qquad
   W_n(x,s)=V_n(x,s)-c_n(s),
\]
with \(\widetilde U_n=\calR_{z_n}U_n\).
Before comparing the affine-normalized compactness variable \(W_n\) with the
exact PDE variable \(V_n\), apply
\Cref{lem:critical-tail-homogeneous-gauge-bridge} to the homogeneous gauges
\(c_n\).  If the gauge sequence has a noncompact homogeneous time defect, then
the corresponding descendant is assigned to \(\Caff\), so this branch is routed
away before entering the retained critical-tail compactification; this is
alternative \textup{(i)} of the present theorem, and there is nothing further to
prove in that branch.  Hence we work in the complementary case, after passing to
a subsequence, in which
\[
   c_n\to c
   \qquad\text{strongly in }L^3_{\loc}(\mathbb R;\mathbb R^3).
\]
By \Cref{lem:bounded-velocity-gauges-blowdown}, the affine-normalized gauges
may be chosen with \(|c_n(s)|\le M\) on compact time intervals.  Hence
\[
   |W_n|\le 2M
\]
on every selected macroscopic cylinder; since also \(|V_n|\le M\),
Banach--Alaoglu gives weak-* compactness of both \(V_n\) and \(W_n\) in
\(L^\infty_{\rm loc}\).  The fundamental theorem of Young measures gives, after
passing to a further subsequence, ungauged barycenter/Young data
\((\bar V,\mathcal Y^V)\) for \(V_n\) and affine-normalized barycenter/Young
data \((\bar W,\mathcal Y)\) for \(W_n\).  By
\Cref{lem:critical-tail-homogeneous-gauge-bridge},
\(\mathcal Y^V\) and \(\mathcal Y\) differ only by deterministic translation.
If the affine-normalized sequence is strongly compact in \(L^3_{\rm loc}\), then
\(\mathcal Y_{x,s}=\delta_{\bar W(x,s)}\); otherwise \(\mathcal Y\) records the
unresolved oscillation defect.
Apply \Cref{thm:critical-tail-compactification-compactness} to the selected
critical-tail family.  If its affine-pressure or pressure/noncompact alternative
occurs, the branch is precisely alternative \textup{(i)} or \textup{(ii)} of the
present theorem, and we stop.  Otherwise we are in the retained
bounded-coordinate branch.  In this branch the pressure-defect
coordinates \(\widehat{\mathscr P}_{n,F}[\psi]\), the viscous-defect coordinates
\(\mathscr D_{n,F}[\psi]\), and, on the realized recurrent scale-critical
virial-tight subhull of \Cref{def:scale-critical-virial-tight-subhull}, the
virial coordinates are uniformly bounded on the fixed countable dense families
entering the compactification.  Diagonal extraction therefore gives the full
tuple in \(\Tcrit\).  Since the normalized defect mass is one, the retained
limit satisfies \(\mathcal T\ne0\), giving alternative \textup{(iii)}.

It remains, in the annular retained subcase of alternative \textup{(iii)}, to
identify the equation satisfied by the barycenter.  The ungauged
fields \(V_n\) satisfy the exact scaled equation
\eqref{eq:critical-tail-exact-blowdown-equation}.  Passing from \(V_n\) to
\(W_n=V_n-c_n(s)\) is used only to measure oscillation.  We do not differentiate
the homogeneous gauge \(c_n(s)\).  Instead, the exact equation is passed to the
limit for \(V_n\) and only afterwards projected to the affine/parasitic
quotient.  The noncompact homogeneous-gauge alternative has already been routed
to \(\Caff\) by \Cref{lem:critical-tail-homogeneous-gauge-bridge}; hence in the
present branch \(\bar W=\bar V-c(s)\) differs from \(\bar V\) only by a
deterministic spatially homogeneous mode.
For every compactly supported test function, the viscous and nonlinear
contributions are bounded respectively by
\[
   CR_n^{-2}\|V_n\|_{L^\infty},
   \qquad
   CR_n^{-1}\|V_n\|_{L^\infty}^2 ,
\]
and hence both tend to zero.  The pressure gauges give distributional
compactness of the remaining gradients on annuli.  Passing to the weak-*
barycentric limit of \(V_n\) yields the coherent transport equation for
\(\bar V\) with the actual pressure representative \(Q^{\rm act}\); projecting
only after this step gives the affine-quotient transport equation for the
non-affine barycenter \(\bar W\), modulo spatially homogeneous velocity modes
and affine pressure forces.  No derivative of the limiting gauge \(c(s)\) is
used unless it is absorbed into that affine pressure quotient.  Applying
\Cref{lem:critical-tail-young-log-transport} gives the exact Young/log balance
for the ungauged Young measure \(\mathcal Y^V\).  By
\Cref{lem:critical-tail-homogeneous-gauge-bridge}, the affine-normalized Young
measure \(\mathcal Y\) stored in \(\Tcrit\) differs from \(\mathcal Y^V\) only
by deterministic translation, so the branch decomposition sees the same
Dirac/non-Dirac and hidden-scale alternatives.  The log-scale component
satisfies the induced log-dilation balance, while nonlinear velocity
observables retain the pressure and viscous defect pairings displayed in
\eqref{eq:young-transport-identity}.  This construction is local on
macroscopic annuli and uses no global critical norm, CKN budget, or energy
estimate.
\end{proof}

\begin{definition}[Young/defect critical-tail branch]
\label{def:young-defect-critical-tail-branch}
A tight critical-tail profile belongs to \(\Cyoung\) if \(\Lambda=0\) and at
least one of the following holds:
\begin{enumerate}[label=\textup{(\roman*)}]
   \item the affine-normalized Young measure \(\mathcal Y\) is non-Dirac on a set
   of positive macroscopic measure;
   \item for some \((F,\psi)\in\mathfrak F_{\rm crit}\),
   \[
      \widehat{\mathscr P}_F[\psi]\ne0;
   \]
   \item for some positive convex test pair
   \((F,\psi)\in\mathfrak F_{\rm crit}^{\rm conv}\),
   \[
      \mathscr D_F[\psi]>0.
   \]
\end{enumerate}
Thus \(\Cyoung\) records every tight critical-tail obstruction to coherence not
already routed to regenerated concentration, affine/parasitic collapse, or a
pressure/noncompact branch.  The viscous subcase is stated using positive convex
tests because only those coordinates have a sign.  Once all positive convex
coordinates vanish, the associated positive matrix-valued viscous defect measure
vanishes, and the remaining signed coordinates vanish by polarization.
\end{definition}

\begin{lemma}[No pure viscous defect in the tight Dirac critical-tail branch]
\label{lem:no-pure-viscous-defect-tight-dirac}
Let \(\mathcal T\in\Tcrit\) be tight, with
\[
   \Lambda=0,
   \qquad
   \mathcal Y=\delta_{\bar W},
\]
and suppose the realized profile is outside \(\Cactive\), outside \(\Caff\),
outside the pressure/noncompact branch, and has no hidden-scale variance on
compact annular cylinders.  Then for every
\((F,\psi)\in\mathfrak F_{\rm crit}^{\rm conv}\),
\[
   \mathscr D_F[\psi]=0 .
\]
\end{lemma}

\begin{proof}
Work on a compact annular cylinder and in the non-affine, non-pressure/noncompact
branch.  The Dirac Young measure and
\Cref{lem:hidden-variance-young-obstruction,lem:critical-tail-homogeneous-gauge-bridge}
give strong \(L^3_{\loc}\) convergence of the affine-normalized fields, and of
the ungauged fields up to the deterministic homogeneous translation.  By the
general affine-pressure coefficient dichotomy
\Cref{lem:affine-pressure-coefficient-dichotomy-general}, the actual pressure is
normalized so that
\[
   \widehat Q_n^{\rm act}-a_n(s)\to0
   \qquad\text{strongly in }L^{3/2}_{\loc},
\]
and the time-only gauge contributes zero to the local-energy pressure flux.
In the present branch, the affine-pressure coefficient has either been routed to
\(\Caff\) or tends to zero, and the non-affine pressure part tends to zero
strongly after time-only gauges.  Hence the ungauged barycenter equation is
represented with an \(L^{3/2}_{\loc}\) pressure representative; in fact, in the
non-affine quotient one may take \(Q=0\).

Apply the blow-down local energy inequality
\Cref{lem:local-energy-blow-down} with nonnegative cutoffs supported in the
annular cylinder.  Passing to the limit uses the strong velocity convergence,
the pressure-flux convergence just described, and the vanishing factors in the
viscous diffusion and convective flux.  The limiting ungauged coherent transport
profile satisfies the renormalized inviscid energy identity from
\Cref{lem:coherent-barycenter-renormalization}.  Subtracting the limiting energy
identity from the limiting local energy inequality leaves only the nonnegative
trace of the scaled viscous defect measure.  Therefore that trace is zero on
every compact annular cylinder.

The positive convex coordinates \(\mathscr D_F[\psi]\) are contractions of this
positive matrix-valued defect measure with the nonnegative matrix
\(\nabla^2_\xi F\) and the nonnegative cutoff \(\psi\).  Since the trace measure
is zero, every such contraction is zero.  This proves the claim without using a
Poincar\'e lower bound from gradient energy to velocity oscillation.
\end{proof}

\begin{lemma}[No pressure defect after viscous-defect removal]
\label{lem:no-pressure-defect-tight-dirac}
Under the hypotheses of
\Cref{lem:no-pure-viscous-defect-tight-dirac}, after affine-pressure routing and
pressure/noncompact routing one has
\[
   \widehat{\mathscr P}_F[\psi]=0
   \qquad\text{for every }(F,\psi)\in\mathfrak F_{\rm crit}.
\]
\end{lemma}

\begin{proof}
Use the actual/time-gauge representative supplied by
\Cref{lem:affine-pressure-coefficient-dichotomy-general} and the non-affine
estimate of \Cref{lem:scaled-pressure-compactness-affine}.  On the fixed annular
cylinder write
\[
   \widehat Q_n^{\rm act}
   =q_n^{\rm na}+b_n(s)\cdot x+a_n(s),
   \qquad
   \|q_n^{\rm na}\|_{L^p(K)}\le C_{p,K,M}R_n^{-1}
\]
for every finite \(p\).  The time-only gauge has zero spatial gradient and
therefore drops out of the pressure-defect pairing.  The affine spatial piece is
kept in gradient form:
\[
   \iint \psi\,\nabla_\xi F(V_n)\cdot\nabla(b_n(s)\cdot x)
   =
   \iint \psi\,\nabla_\xi F(V_n)\cdot b_n(s).
\]
If \(b_n\) has a retained nonzero component, the branch is routed to \(\Caff\).
In the present non-affine branch, \(b_n\to0\) locally in \(L^{3/2}_s\), and the
last display tends to zero because \(\nabla_\xi F(V_n)\) is bounded and
\(\psi\) has compact support.

It remains to treat the non-affine pressure.  Integrating only
\(q_n^{\rm na}\) by parts gives
\[
\begin{aligned}
\iint \psi\,\nabla_\xi F(V_n)\cdot\nabla q_n^{\rm na}
&=-\iint q_n^{\rm na}\,
   \nabla\psi\cdot\nabla_\xi F(V_n)  \\
&\quad
  -\iint q_n^{\rm na}\,
   \psi\,\nabla_\xi^2F(V_n):\nabla V_n .
\end{aligned}
\]
The first term tends to zero by
\(\|q_n^{\rm na}\|_{L^{3/2}(K)}=O(R_n^{-1})\) and the uniform velocity bound.
For the second term, Cauchy--Schwarz and the \(L^2\) case of
\eqref{eq:scaled-pressure-nonaffine-Lp} give
\[
\left|
\iint q_n^{\rm na}\psi\,
\nabla_\xi^2F(V_n):\nabla V_n
\right|
\le
C\Bigl(R_n^{-2}\iint_{\supp\psi}|\nabla V_n|^2\Bigr)^{1/2}.
\]
The factor on the right is the trace viscous defect on the support of \(\psi\),
which vanishes by \Cref{lem:no-pure-viscous-defect-tight-dirac}.  Therefore all
pressure-defect coordinates vanish.
\end{proof}

\begin{lemma}[Critical-tail defect-coordinate closure]
\label{lem:critical-tail-defect-coordinate-routing}
\label{lem:critical-tail-defect-coordinate-closure}
Let \(\mathcal T\in\Tcrit\) be tight and Dirac in the affine-normalized Young
measure.  After regenerated concentration, affine/parasitic,
pressure/noncompact, and hidden-scale branches have been removed, all pressure
and viscous defect coordinates vanish:
\[
   \widehat{\mathscr P}=0,
   \qquad
   \mathscr D=0.
\]
\end{lemma}

\begin{proof}
The non-Dirac Young-measure case is exactly the hidden-scale variance obstruction
by \Cref{lem:hidden-variance-young-obstruction}.  A genuinely retained positive
\(L^3\)-oscillation defect is routed to the minimal mesoscopic branch by
\Cref{thm:first-bad-mesoscopic-reduction}.  In the remaining tight Dirac branch,
pure viscous defects vanish by
\Cref{lem:no-pure-viscous-defect-tight-dirac}.  Once the positive convex viscous
coordinates vanish, all signed viscous coordinates vanish by the polarization
statement in \Cref{def:critical-pressure-gradient-convergence}.  The pressure
coordinates then vanish by \Cref{lem:no-pressure-defect-tight-dirac}.  Thus no
pressure or viscous critical-tail defect coordinate survives in the retained
residual class.
\end{proof}

\begin{corollary}[Coherent extraction under annular tightness and strong compactness]
\label{cor:coherent-critical-tail-extraction}
If the critical-tail extraction is in alternative \textup{(iii)} of
\Cref{thm:critical-tail-compactification-extraction}, producing a retained
defect \(\mathcal T\in\Tcrit\), and if this defect satisfies
\(\Lambda=0\) and
\(\mathcal Y_{x,s}=\delta_{W(x,s)}\) for a.e. \((x,s)\), with
\(\widehat{\mathscr P}=0\), \(\mathscr D=0\), and no retained affine or
pressure/noncompact defect, then there exist
macroscopic scales, covariant observers, affine gauges, and pressure gauges for
which the ungauged blown-down fields converge, up to the deterministic
homogeneous gauge bridge, to an ungauged representative \((V,Q^{\rm act})\)
satisfying the literal coherent transport equation on punctured annular
cylinders.  The affine-normalized fields converge strongly in \(L^3_{\rm loc}\)
and weak-* in \(L^\infty_{\rm loc}\) to the corresponding nonzero quotient
coherent critical tail \((W,Q)\in\Ccohcrit\), with \(W=V-c(s)\).
\end{corollary}

\begin{proof}
This is the coherent subcase of the retained bounded-coordinate alternative in
\Cref{thm:critical-tail-compactification-extraction}.  The condition
\(\Lambda=0\) rules out loss of mass across logarithmic scale, and the
Dirac-valued Young measure gives strong \(L^3_{\rm loc}\) convergence by
uniform \(L^\infty\) boundedness.  The vanishing of
\(\widehat{\mathscr P}\), \(\mathscr D\), and retained affine/pressure defects is
precisely the extra coordinate condition in \(\Tcoh\), so no residual
pressure/viscous coordinate remains outside the coherent transport equation.
The literal equation is carried by the ungauged representative; the displayed
affine-normalized tail is the quotient representative described in
\Cref{def:critical-diffuse-tail-strata}.  Nontriviality follows from the
normalized affine oscillation mass.
\end{proof}

\begin{theorem}[Coherent critical-tail classification]
\label{thm:coherent-critical-tail-classification}
Let \((W,Q)\in\Ccohcrit\) be a coherent critical tail, with \(W\) understood as
the affine-normalized quotient representative of an ungauged coherent tail.  The
classification below uses only this quotient representative, its non-affine
spatial oscillation, and its log/time translation hull.  In logarithmic radial
variables
\[
   x=e^\rho\omega,\qquad Z(\rho,\omega,s)=e^\rho W(e^\rho\omega,s),
\]
its compact log-translation/time-translation hull belongs to exactly one of
\[
   \mathcal C_{\rm affine},\qquad
   \Chomcrit,\qquad
   \Clogper,\qquad
   \Capcrit .
\]
Here \(\Chomcrit\) is the homogeneous \((-1)\)-tail class,
\(\Clogper\) is the nonzero log-periodic tail class, and \(\Capcrit\) is the
remaining aperiodic minimal log-translation hull class.
\end{theorem}

\begin{proof}
Let
\[
   \mathcal H(Z):=
   \overline{\{Z(\rho+a,\omega,s+b):a,b\in\R\}}
\]
be the compact logarithmic translation hull.  If the affine-normalized
oscillation vanishes on every compact log-cylinder, then the tail is
affine/parasitic.  Otherwise, if every element of \(\mathcal H(Z)\) is independent of
\(\rho\), then \(W\) is \((-1)\)-homogeneous and belongs to \(\Chomcrit\).  If
the log-translation orbit has a nonzero period, then \(W\) is discretely
self-similar at infinity and belongs to \(\Clogper\).  If neither alternative
holds, compact-flow decomposition gives a minimal compact invariant subsystem
with no nonzero log period; this is the aperiodic log-translation class
\(\Capcrit\).  These alternatives exhaust compact log-translation hulls.
\end{proof}

\begin{theorem}[Realized critical-tail rigidity by stratification]
\label{thm:realized-critical-tail-rigidity}
Let \(\mathcal T\in\Tcrit\) be a nonzero realized
critical-tail profile generated by a bounded normalized residual Navier--Stokes
terminal sequence after terminal concentration-set removal.  Then the ordered
stratification selects one of the following priority profile alternatives:
\begin{enumerate}[label=\textup{(\alph*)}]
   \item \(\mathcal T\) regenerates affine-normalized local concentration;
   \item \(\mathcal T\) is affine/parasitic;
   \item \(\mathcal T\) is log-diffuse;
   \item \(\mathcal T\) is a Young/defect critical tail;
   \item \(\mathcal T\) is a coherent homogeneous critical tail;
   \item \(\mathcal T\) is a coherent log-periodic critical tail;
   \item \(\mathcal T\) is a coherent aperiodic critical tail.
\end{enumerate}
In the ordered refined decomposition, the log-diffuse, Young, homogeneous,
log-periodic, and aperiodic critical-tail alternatives are lower structured strata.
Thus every nonzero realized critical-tail profile which remains after these
lower strata have been removed either regenerates local concentration or collapses
to the affine/parasitic lower stratum.
\end{theorem}

\begin{proof}
If the covariant observer orbit closure of \(\mathcal T\) intersects a
constant-in-time concentration set \(\{U:\actthree(U;Q_1^{\rm cov})\ge
\eta\}\) for some \(\eta>0\), then alternative \emph{(a)} holds.  This case is
removed first.

If every spatial quotient oscillation observable vanishes on the profile, then
\Cref{thm:affine-normalization-dichotomy,def:affine-constant-lower-stratum}
place the profile in the affine/parasitic lower stratum.  This is alternative
\emph{(b)}.  It remains to consider the case in which the profile has nonzero
non-affine diffuse content and does not regenerate retained unit-scale velocity
concentration.

By \Cref{thm:critical-tail-compactification-extraction}, the residual normalized
diffuse mass either first produces an affine-pressure or pressure/noncompact
alternative, or else is represented by a full \(\Tcrit\)-state.  The
pre-\(\Tcrit\) alternatives are already assigned to their corresponding lower
strata in the priority order; since \(\mathcal T\in\Tcrit\) is the retained
object under consideration, we work in alternative
\textup{(iii)} of that extraction theorem.  Thus the
boundary mass has only two mechanisms: mass escaping every bounded logarithmic
annulus, recorded by \(\Lambda\), and annular blow-down mass, recorded by
\((\bar W,\mathcal Y)\), together with the scalar pressure/viscous defect
coordinates recorded in \(\Tcrit\).  If \(\Lambda\ne0\), the profile is log-diffuse, giving
\emph{(c)}.  If \(\Lambda=0\) but \(\mathcal Y_{x,s}\) is not a Dirac mass on a
set of positive macroscopic measure, the profile is a Young/defect critical tail, giving
\emph{(d)}.  If \(\Lambda=0\) and \(\mathcal Y\) is Dirac but some coordinate
\(\widehat{\mathscr P}_F[\psi]\) or \(\mathscr D_F[\psi]\) is nonzero, then
\Cref{def:young-defect-critical-tail-branch,lem:critical-tail-defect-coordinate-closure}
place the profile in the Young/defect critical-tail alternative before closure;
after the already ordered concentration, affine, pressure/noncompact,
hidden-scale, and minimal mesoscopic branches have been removed, those
coordinates vanish and this alternative is empty.

The remaining case is coherent in the full compactification topology:
\(\Lambda=0\),
\(\mathcal Y_{x,s}=\delta_{\bar W(x,s)}\) a.e.,
\(\widehat{\mathscr P}=0\), \(\mathscr D=0\), and no retained affine or
pressure/noncompact defect.  Then
\Cref{cor:coherent-critical-tail-extraction} gives a coherent critical tail
\((\bar W,Q)\in\Ccohcrit\).  Applying
\Cref{thm:coherent-critical-tail-classification} gives exactly one of the
homogeneous, log-periodic, or aperiodic coherent alternatives, namely
\emph{(e)}, \emph{(f)}, or \emph{(g)}, unless the affine/parasitic case already
occurred.  The alternatives are mutually exclusive after the priority order
above: regenerated concentration is removed first, then affine collapse, then
log-diffuse mass, then Young/defect obstruction, and finally the coherent
log-translation classification.  The final assertion follows from this ordered
classification.
\end{proof}

\subsection{Complete closure of lower alternatives}
\label{sec:complete-branch-closure}

The critical-tail classification above is a classification theorem rather than
an emptiness theorem.  The complete residual decomposition is closed by
local lower-stratum exclusions together with the minimal mesoscopic-scale
theorem \MesoscopicReductionTheorem.  First,
\Cref{thm:oscillatory-entry-normalization} routes affine/parasitic profiles to
the quotient stratum \(\Caff\).  Next,
\Cref{cor:coherent-critical-tail-branch-closure} closes the coherent critical-tail
row by bounded-origin realization and log-hull minimality, independently of the
log-diffuse and hidden-scale arguments.  The minimal mesoscopic-scale theorem
then gives hidden-scale compactness through
\Cref{thm:no-hidden-scale-variance-realized-defects}.  The non-Dirac subcase of
the Young alternative is reduced to that compactness statement by the
Fr\'echet--Kolmogorov/Young-measure equivalence in
\Cref{lem:hidden-variance-young-obstruction}; its pressure/viscous subcases are
closed by \Cref{lem:critical-tail-defect-coordinate-closure}.  The log-diffuse alternative is
reduced by log-window support-profile selection together with the same
hidden-scale theorem and the already closed coherent row; see
\Cref{thm:log-diffuse-from-hidden-scale,cor:log-diffuse-branch-exclusion}.  These are local
statements about realized residual descendants; none is a global
\(L^\infty_tL^3_x\), finite-energy-tail, or global CKN-budget estimate.

\begin{theorem}[Oscillatory entry normalization and affine routing]
\label{thm:oscillatory-entry-normalization}
Let \((U,\Pi)\) be a retained residual profile after the earlier lower classes
recorded in the companion setup paper have been removed.  Then every
pre-quotient retained residual candidate satisfies exactly one of the following:
\begin{enumerate}[label=\textup{(\roman*)}]
   \item after canonical affine/parasitic normalization and an admissible affine
   symmetry \(S_c\), the profile retains positive constant-in-time non-affine
   oscillation:
   \[
      \Phi(S_cU)>0;
   \]
   \item all spatial oscillations \(\oscsp(U;Q)\) vanish on compact cylinders,
   and the profile is assigned to the affine/parasitic quotient stratum
   \(\Caff\).
\end{enumerate}
Since \(\calR^\#(\calS;I,J)\) is defined after this affine/parasitic routing,
the retained normalized residual decomposition has
\[
   \Caff\cap\calR^\#(\calS;I,J)=\varnothing .
\]
\end{theorem}

\begin{proof}
The dichotomy is \Cref{thm:affine-normalization-dichotomy}.  If all
spatial oscillations \(\oscsp(U;Q)\) vanish on compact cylinders, then the profile
is a spatially homogeneous quotient mode and is routed to \(\Caff\).  If this
alternative does not occur, the parasitic-free representative has
\(\Phi(S_cU)>0\) for an admissible affine symmetry \(S_c\).  Thus every
pre-quotient candidate is either retained with positive non-affine oscillation
or assigned to \(\Caff\), and the normalized residual class
\(\calR^\#(\calS;I,J)\) contains no representative lying in \(\Caff\).
\end{proof}

\begin{corollary}[Young alternative exclusion from minimal mesoscopic scale reduction]
\label{cor:young-branch-exclusion}
\[
   \Cyoung=\varnothing .
\]
\end{corollary}

\begin{proof}
Let a retained residual profile generate a nonzero Young/defect critical-tail alternative.
By definition of \(\Cyoung\), its logarithmic component is tight:
\(\Lambda=0\), and one of the subcases in
\Cref{def:young-defect-critical-tail-branch} occurs.  In the non-Dirac Young
measure subcase, let \(W_n\) be the corresponding annular blow-down and
\(\mathcal Y_{x,s}\) its Young measure.  By
\Cref{thm:first-bad-mesoscopic-reduction,thm:no-hidden-scale-variance-realized-defects},
\(W_n\) has no hidden-scale variance on compact annular cylinders.  Hence
\Cref{lem:hidden-variance-young-obstruction} gives
\[
   \mathcal Y_{x,s}=\delta_{\bar W(x,s)}
   \qquad\hbox{for a.e. }(x,s),
\]
contradicting that subcase.  In the pressure/viscous defect-coordinate subcases,
\Cref{lem:no-pure-viscous-defect-tight-dirac,lem:no-pressure-defect-tight-dirac,lem:critical-tail-defect-coordinate-closure}
show that, after affine, pressure/noncompact, regenerated-concentration, and
hidden-scale branches are removed, no nonzero defect coordinate remains.
Therefore
\(\Cyoung=\varnothing\).
\end{proof}

\begin{theorem}[Critical-tail rigidity from minimal mesoscopic scale reduction]
\label{thm:remaining-rigidity-exact-form}
Every nonzero recurrent critical-tail defect generated by a realized residual
terminal sequence belongs to
\[
   \Cactive\cup\Caff\cup\Ccohcrit .
\]
Equivalently, after regenerated concentration, affine/parasitic collapse, and
coherent critical tails have been removed, no recurrent critical-tail defect
remains.
\end{theorem}

\begin{proof}
Let \(\mathcal D\in\Tcrit\) be a nonzero recurrent critical-tail defect.  By
\Cref{thm:realized-critical-tail-rigidity}, the ordered priority alternative is
one of: regenerated concentration, affine/parasitic collapse, log-diffuse defect,
Young/defect critical tail, or coherent critical tail.  The first two alternatives are
precisely \(\Cactive\) and \(\Caff\).  The log-diffuse alternative is empty by
\Cref{cor:log-diffuse-branch-exclusion}.  The Young alternative is empty by
\Cref{cor:young-branch-exclusion}.  The only remaining nonempty critical-tail
alternative is therefore the coherent alternative \(\Ccohcrit\).  This proves the
inclusion.  Once \(\Cactive\), \(\Caff\), and \(\Ccohcrit\) are removed, the
inclusion leaves no recurrent critical-tail defect.
\end{proof}

\begin{lemma}[Bounded velocity gauges for realized blow-downs]
\label{lem:bounded-velocity-gauges-blowdown}
Let \(f:E\to\mathbb R^3\) be measurable on a finite positive-measure set
\(E\), and assume \(|f|\le M\) a.e.  Then the minimization problem
\[
   \inf_{c\in\mathbb R^3}\int_E |f-c|^3
\]
has a minimizer in \(\overline{B_M(0)}\).  Consequently every
affine-normalized blow-down of a sequence satisfying
\(\|U_n\|_{L^\infty}\le M\) may be represented, without increasing any local
oscillation functional, by velocity gauges \(c_n\) with \(|c_n|\le M\).  If the
normalization is performed time-slice by time-slice, the same conclusion holds
for a measurable choice \(c_n(s)\), with \(|c_n(s)|\le M\) for a.e. \(s\).
\end{lemma}

\begin{proof}
The functional is continuous and coercive in \(c\), so a minimizer exists.  Let
\(\pi:\mathbb R^3\to\overline{B_M(0)}\) be the nearest-point projection.  Since
\(\overline{B_M(0)}\) is closed and convex, for every \(v\in\overline{B_M(0)}\)
and every \(c\in\mathbb R^3\),
\[
   |v-\pi(c)|\le |v-c|.
\]
Taking \(v=f(y)\) and integrating gives
\[
   \int_E |f-\pi(c)|^3\le \int_E |f-c|^3 .
\]
Thus any minimizing sequence may be projected into \(\overline{B_M(0)}\), and a
minimizer can be chosen there.  Applying this to the local velocity fields used
to define the affine-normalized observer oscillations gives the asserted
bounded representatives.  The time-slice version is the same argument on each
slice, with measurability obtained by the usual measurable-selection theorem
for a continuous coercive finite-dimensional minimization problem.  The
projection step is pointwise in \(s\), so the bound \(|c_n(s)|\le M\) is
preserved.
\end{proof}

\begin{lemma}[Domain realization dichotomy for critical-tail windows]
\label{lem:critical-tail-domain-realization-dichotomy}
Every critical-tail window used in the bounded-origin realization theorem is of
exactly one of the following two admissible types.
\begin{enumerate}[label=\textup{(\roman*)}]
   \item It is a window of a bounded ancient descendant
   \((U,\Pi)\in\mathfrak X_M\).  Then every fixed macroscopic blow-down cylinder
   is contained in the full domain \(\mathbb R^3\times\mathbb R\).
   \item It is a first-generation local window satisfying the local-window
   condition of \Cref{def:first-generation-local-window-condition} on every fixed
   blow-down cylinder used in the extraction.
\end{enumerate}
If neither condition holds, the window is not used in bounded-origin
realization.  It is assigned by
\Cref{lem:compact-window-expanding-region-protocol} to the exterior/diffuse,
pressure/noncompact, or critical-tail boundary branch before bounded-origin
realization is invoked.
\end{lemma}

\begin{proof}
Such windows are obtained after an ancient bounded profile has already
been extracted, so their domains are global in centered variables.  For a
first-generation window, the local setup estimates are available precisely on
the physical image of the local Type~I cylinder.  If every fixed blow-down
cylinder remains in that image, the local-window condition gives the required
compactness and pressure atlas.  If some fixed blow-down cylinder leaves the
local window, the sequence has escaped the first-generation setup domain; by
\Cref{lem:escaping-first-generation-alternatives}, it must be assigned through the
expanding-region protocol rather than treated as a bounded-origin compactness
input.
\end{proof}

\begin{theorem}[Bounded-origin realization of coherent critical tails]
\label{thm:bounded-origin-realization-coherent-tails}
Let a coherent critical tail be generated by a bounded normalized residual
Navier--Stokes terminal sequence with:
\begin{enumerate}[label=\textup{(\roman*)}]
   \item locally bounded velocity representatives;
   \item pressure gauges bounded modulo time and affine functions on each fixed
   time window;
   \item convergence on every fixed annulus/log-window;
   \item the normalized support-profile property realized from the original blow-up sequence;
   \item the domain alternative in
   \Cref{lem:critical-tail-domain-realization-dichotomy}.
\end{enumerate}
Then it is realized, after subsequence extraction and admissible gauges, as a
weak-\(*\) \(L^\infty_{\loc}\) limit on full macroscopic compact cylinders,
including cylinders meeting the macroscopic origin.  In particular, every
coherent critical tail has a representative which is locally bounded near
\(x=0\) on compact macroscopic time intervals, and whose restrictions to all
tail windows reproduce the coherent tail.
\end{theorem}

\begin{remark}[Use of the bounded-origin representative]
\label{rem:bounded-origin-used-only-as-extension}
The bounded-origin representative is used only as a locally bounded extension
of the annular coherent tail.  The coherent transport equation and pressure
normalization are required only on punctured annular cylinders
\(\{0<r_0<|x|<r_1<\infty\}\times I\).  None of the homogeneous,
log-periodic, or aperiodic exclusions below uses a distributional equation
through \(x=0\).
\end{remark}

\begin{proof}
Let \((W,Q)\in\Ccohcrit\) be a coherent critical tail generated by a bounded
terminal sequence \((U_n,\Pi_n)\subset\mathfrak X_M\), with
\(\|U_n\|_{L^\infty}\le M\).  By coherent extraction, after passing to a
subsequence there are macroscopic scales \(R_n\to\infty\), covariant observer
centers \((y_n,\sigma_n)\), affine velocity gauges, and pressure gauges such
that on every compact annular cylinder
\[
   A_{r_0,r_1,I}:=\{r_0<|x|<r_1\}\times I,
   \qquad 0<r_0<r_1<\infty,
\]
the fields
\[
   \widetilde U_n:=\calR_{(y_n,\sigma_n)}U_n,
   \qquad
   W_n(x,s):=\widetilde U_n(R_nx,s)-c_n(s)
\]
converge to \(W\) strongly in \(L^3(A_{r_0,r_1,I})\) and weak-\(*\) in
\(L^\infty(A_{r_0,r_1,I})\).
By \Cref{lem:critical-tail-domain-realization-dichotomy}, the fixed
macroscopic cylinders used below are contained in the domain either because the
sequence is an already extracted ancient profile, or because the first-generation local-window
condition holds.  If this domain condition failed, the window would already
belong to an exterior/diffuse or noncompact branch and would not be an input to
bounded-origin realization.

By \Cref{lem:bounded-velocity-gauges-blowdown}, the affine-normalized blow-down
may be represented by gauges satisfying
\[
   |c_n(s)|\le M
\]
for a.e. \(s\) on every compact time interval, without increasing the
oscillation functional.  Hence on every full macroscopic cylinder
\[
   K_{R,I}:=B_R(0)\times I,
\]
one has
\[
   |W_n(x,s)|
   \le |\widetilde U_n(R_nx,s)|+|c_n(s)|
   \le 2M .
\]
Thus Banach--Alaoglu gives a weak-\(*\) \(L^\infty_{\rm loc}\) limit on full
cylinders, including cylinders meeting \(x=0\).

The domain realization dichotomy
\Cref{lem:critical-tail-domain-realization-dichotomy} is the hypothesis
which ensures that \(W_n\) is defined on \(K_{R,I}\) for all sufficiently large
\(n\).

Choose an exhaustion of \(\mathbb R^3\times\mathbb R\) by full compact
macroscopic cylinders \(K_j=B_j(0)\times[-j,j]\).  Banach--Alaoglu gives, by a
diagonal extraction that does not change notation, a field
\[
   W^{\rm full}\in L^\infty_{\loc}(\mathbb R^3\times\mathbb R),
\]
such that
\[
   W_n\rightharpoonup^\ast W^{\rm full}
   \qquad\hbox{in }L^\infty(K_j)
\]
for every \(j\).  The subsequence still has the annular strong
\(L^3_{\loc}\) convergence to the coherent tail \(W\).  Therefore, for every
annular cylinder \(A_{r_0,r_1,I}\),
\[
   W^{\rm full}=W
   \qquad\hbox{a.e. on }A_{r_0,r_1,I}.
\]
Since the annuli exhaust
\((\mathbb R^3\setminus\{0\})\times\mathbb R\), the annular coherent tail is the
restriction of the full weak-\(*\) limit \(W^{\rm full}\) away from \(x=0\).

We take \(W^{\rm full}\) as the full-cylinder representative of the coherent
tail.  It is locally bounded on every compact macroscopic cylinder, including
cylinders meeting \(x=0\), and it realizes the same annular tail because it
agrees with \(W\) on all punctured compact cylinders.  Pressure gauges are not
needed for the boundedness conclusion; if one records the pressure as well, the
local pressure-gauge compactness in the terminal profile class gives
distributional pressure-gradient limits on the same full cylinders after the
usual time-dependent gauges.  The construction uses only the uniform local
boundedness inherited from the Type~I terminal profile class and finite-dimensional
affine normalization; it uses no global critical norm or global tail budget.
\end{proof}

\begin{lemma}[Failure of amplitude-normalized bounded-origin heredity is routed]
\label{lem:amplitude-normalized-heredity-failure-routing}
Let \((W,Q)\in\Ccohcrit\) be generated by a bounded terminal sequence and
realized by \Cref{thm:bounded-origin-realization-coherent-tails}.  Write
\[
   Z(\rho,\omega,s)=e^\rho W(e^\rho\omega,s).
\]
Consider amplitude-normalized translated windows
\[
   Z(\rho+a_j,\omega,s+b_j).
\]
If, after passing to a subsequence, these translated windows do not admit a
sequence-realized bounded-origin representative stable under the
amplitude-normalized translation, then the ordered critical-tail protocol
assigns a sequence-realized support profile to one of the already stratified
lower alternatives
\[
   \Cactive,\qquad \Caff,\qquad \Clogdiff,\qquad \Cyoung.
\]
\end{lemma}

\begin{proof}
The translated windows are still generated from the original blow-up sequence,
so every realized subsequential profile extracted from them is a realized
critical-tail profile.  Apply the ordered critical-tail extraction theorem
\Cref{thm:realized-critical-tail-rigidity} to such a realized translated
profile.  Its alternatives are regenerated concentration, affine/parasitic
collapse, log-diffuse escape, Young/defect critical-tail defect, or coherent
critical-tail realization.  The last alternative is exactly the successful
bounded-origin heredity case in which the translated amplitude-normalized
windows admit a bounded-origin representative stable under the normalized
translation.  Therefore failure of successful bounded-origin heredity leaves
only the first four alternatives, namely
\(\Cactive\), \(\Caff\), \(\Clogdiff\), and \(\Cyoung\).  The corresponding
profiles are sequence-realized by the support-transfer statements already built
into the critical-tail compactification, so the translated family is formally
routed to one of these lower strata.
\end{proof}

\begin{lemma}[Realized amplitude-normalized log-hull heredity]
\label{lem:realized-amplitude-log-hull-heredity}
Let \((W,Q)\in\Ccohcrit\) be generated by a bounded terminal sequence and be
realized by \Cref{thm:bounded-origin-realization-coherent-tails}.  Write
\[
   Z(\rho,\omega,s)=e^\rho W(e^\rho\omega,s).
\]
Let \(Z_*\) be a nonzero element of the amplitude-normalized log/time hull,
obtained as a local limit
\[
   Z_*(\rho,\omega,s)
   =
   \lim_{j\to\infty}Z(\rho+a_j,\omega,s+b_j).
\]
Then one of the following alternatives occurs:
\begin{enumerate}[label=\textup{(\roman*)}]
   \item the translated windows are realized by admissible critical-tail
   extractions from the original blow-up sequence, with bounded-origin
   representatives stable under the amplitude-normalized translation, and
   \[
      Z_*(\rho,\cdot,\cdot)\to0
      \qquad\hbox{as }\rho\to-\infty
   \]
   in the local logarithmic topology;
   \item the failure of such realized bounded-origin heredity is routed by
   \Cref{lem:amplitude-normalized-heredity-failure-routing} to one of the
   already stratified lower alternatives
   \[
      \Cactive,\qquad \Caff,\qquad \Clogdiff,\qquad \Cyoung.
   \]
\end{enumerate}
\end{lemma}

\begin{proof}
A translated log profile \(Z(\rho+a_j,\omega,s+b_j)\) represents
\[
   e^{a_j}W(e^{a_j}x,s+b_j),
\]
not merely \(W(e^{a_j}x,s+b_j)\).  Thus it is not admitted into the coherent
critical-tail closure solely by changing the macroscopic scale in the
no-amplitude blow-down.  It is admitted only when the amplitude-normalized
profile is itself realized by the normalized critical-tail extraction from the
original terminal sequence.

Assume first that this realized amplitude-normalized extraction exists along
the translated windows.  The bounded-origin realization theorem applied to
these realized representatives gives \(L^\infty_{\loc}\) bounds on every full
compact macroscopic cylinder, including cylinders meeting \(x=0\), with
constants stable under the local limit.  Hence \(Z_*\) corresponds to a
velocity
\[
   W_*(r\omega,s)=r^{-1}Z_*(\log r,\omega,s)
\]
which is locally bounded near \(r=0\).

For \(r=e^\rho\downarrow0\), local boundedness gives
\[
   |Z_*(\rho,\omega,s)|
   =
   r\,|W_*(r\omega,s)|
   \le C_K r
\]
on every compact angular-time set \(K\).  Thus
\(Z_*(\rho,\cdot,\cdot)\to0\) locally as \(\rho\to-\infty\), proving
alternative \textup{(i)}.

It remains to identify the complementary case.  If the amplitude-normalized
translated windows cannot be realized with the same bounded-origin control,
then \Cref{lem:amplitude-normalized-heredity-failure-routing} applies and gives
alternative \textup{(ii)}.
\end{proof}

\begin{lemma}[Homogeneous tails are incompatible with bounded-origin retention]
\label{lem:homogeneous-tail-bounded-origin-closure}
Let \((W,Q)\in\Chomcrit\) be a coherent critical tail obtained from the
bounded-origin realization theorem.  Assume it carries the retained non-affine
annular activity required in the definition of the coherent tail stratum.  Then
no such \(W\) exists.
\end{lemma}

\begin{proof}
In the homogeneous critical-tail class,
\[
   W(r\omega,s)=r^{-1}A(\omega,s).
\]
The bounded-origin realization gives a locally bounded representative on
\(\{0<r<r_0\}\) for each compact angular-time set.  Hence
\[
   |A(\omega,s)|=r|W(r\omega,s)|\le C_K r
\]
on every compact angular-time set \(K\).  Letting \(r\downarrow0\) gives
\(A=0\), and therefore \(W\equiv0\) on punctured compact cylinders.  This is
incompatible with the retained non-affine annular activity that defines a
nonzero coherent critical tail.
\end{proof}

\begin{lemma}[Log-periodic tails are incompatible with bounded-origin retention]
\label{lem:logperiodic-tail-bounded-origin-closure}
Let \((W,Q)\in\Clogper\) be a coherent critical tail obtained from the
bounded-origin realization theorem.  Assume it carries retained non-affine
annular activity.  Then no such \(W\) exists.
\end{lemma}

\begin{proof}
Write
\[
   Z(\rho,\omega,s)=e^\rho W(e^\rho\omega,s).
\]
Bounded-origin realization gives \(Z(\rho,\cdot,\cdot)\to0\) in the local
logarithmic topology as \(\rho\to-\infty\).  If the log-period is
\(L\ne0\), then
\[
   Z(\rho,\omega,s)=Z(\rho-kL,\omega,s)
   \qquad(k\in\mathbb Z).
\]
For \(L>0\) take \(k\to+\infty\), and for \(L<0\) take \(k\to-\infty\).  In
both cases \(\rho-kL\to-\infty\), so the right-hand side converges to zero.
Thus \(Z\equiv0\), hence \(W\equiv0\) on punctured annular cylinders, contrary
to retained non-affine annular activity.
\end{proof}

\begin{lemma}[Aperiodic coherent tails are incompatible with bounded-origin retention]
\label{lem:aperiodic-tail-bounded-origin-closure}
Let \((W,Q)\in\Capcrit\) be a coherent critical tail obtained from the
bounded-origin realization theorem.  Assume it carries retained non-affine
annular activity.  Then no such \(W\) exists.
\end{lemma}

\begin{proof}
Let \(Z(\rho,\omega,s)=e^\rho W(e^\rho\omega,s)\), and let \(\mathcal M\) be
the nonzero minimal aperiodic logarithmic subsystem in the definition of
\(\Capcrit\), taken inside the realized amplitude-normalized hull.  Choose
\(Z_*\in\mathcal M\).  By \Cref{lem:realized-amplitude-log-hull-heredity},
either a lower stratum has already occurred, or \(Z_*\) has a bounded-origin
representative and hence
\[
   Z_*(\rho,\cdot,\cdot)\to0
   \qquad\hbox{as }\rho\to-\infty .
\]
By the ordered definition of the coherent aperiodic stratum, the object under
consideration is already inside the realized bounded-origin coherent hull.
Hence the routing alternative of
\Cref{lem:amplitude-normalized-heredity-failure-routing} is not admissible for
an object genuinely lying in \(\Capcrit\): if it occurred, the construction
would have been assigned earlier to \(\Cactive\), \(\Caff\), \(\Clogdiff\), or
\(\Cyoung\), and the profile would not have entered the coherent aperiodic row.
Therefore, for a profile genuinely lying in \(\Capcrit\), only alternative
\textup{(i)} of \Cref{lem:realized-amplitude-log-hull-heredity} remains.  In
that case \(0\in\mathcal H(Z_*)\).
Minimality gives \(\mathcal H(Z_*)=\mathcal M\).  Since the zero profile is a
closed invariant subsystem, \(\mathcal M=\{0\}\).  This contradicts the
retained non-affine annular activity that makes the aperiodic coherent tail a
nonzero residual tail.
\end{proof}

\begin{corollary}[Coherent critical-tail alternative closure]
\label{cor:coherent-critical-tail-branch-closure}
One has
\[
   \Chomcrit=\varnothing,\qquad
   \Clogper=\varnothing,\qquad
   \Capcrit=\varnothing .
\]
\end{corollary}

\begin{proof}
\Cref{lem:homogeneous-tail-bounded-origin-closure} excludes
\(\Chomcrit\), \Cref{lem:logperiodic-tail-bounded-origin-closure} excludes
\(\Clogper\), and \Cref{lem:aperiodic-tail-bounded-origin-closure} excludes
\(\Capcrit\).  Each contradiction uses the same retained non-affine annular
activity that makes the coherent tail a nonzero realized residual object, so no
coherent critical-tail lower stratum remains.
\end{proof}

\begin{corollary}[Critical-tail exclusion, complete form]
\label{cor:critical-tail-exclusion-complete}
No nonzero realized critical-tail profile remains after terminal
concentration-set removal and affine/parasitic lower-stratum removal.
\end{corollary}

\begin{proof}
Let a nonzero realized critical-tail profile remain after terminal
concentration-set removal and affine/parasitic lower-stratum removal.  The
extraction first either routes to an affine-pressure or pressure/noncompact
branch, or else produces a full
\(\mathcal T\in\Tcrit\) in alternative \textup{(iii)} of
\Cref{thm:critical-tail-compactification-extraction}.  The pre-\(\Tcrit\)
branches are assigned by affine routing and the pressure/noncompact terminal
protocol to lower strata, so they cannot remain in the retained residual class.
In the
remaining case, by \Cref{thm:realized-critical-tail-rigidity}, the ordered
priority alternative is one of: regenerated concentration, affine/parasitic
collapse, log-diffuse defect, Young/defect critical tail, coherent homogeneous
tail, coherent log-periodic tail, or coherent aperiodic tail.

Regenerated concentration contradicts terminal concentration-set removal.  The
affine/parasitic alternative is routed by
\Cref{thm:oscillatory-entry-normalization} to \(\Caff\), hence it is not a
retained normalized residual representative.  The log-diffuse alternative is
impossible by \Cref{cor:log-diffuse-branch-exclusion}, and the Young
alternative is impossible by \Cref{cor:young-branch-exclusion}.  The three coherent alternatives are impossible by
\Cref{cor:coherent-critical-tail-branch-closure}.  All alternatives are
therefore impossible, so no such \(\mathcal T\) remains.
\end{proof}

\begin{theorem}[No mixed compact--diffuse residual profile]
\label{thm:no-mixed-compact-diffuse}
Let \((U,\Pi)\) be a retained terminal profile in the refined residual class,
after the affine normalization of \Cref{def:affine-constant-lower-stratum}.
Then \(U\) cannot have a diffuse exterior oscillation remainder coexisting with
its compact retained core.  Equivalently, the paired terminal configurations
\[
   \hbox{one compact concentration core}+\mathcal R_{\rm diff},
\]
and their finite, infinite, or continuum concentration-set analogues do not occur in
the refined residual class.
\end{theorem}

\begin{proof}
Assume that such a diffuse exterior remainder exists.  By the paired
concentration-set decomposition \Cref{thm:terminal-exhaustion-main}, remove
neighborhoods of the concentration set at every positive constant-in-time oscillation
threshold.  The mixed assumption says that a nonzero diffuse oscillation
residual remains.  Apply \Cref{thm:diffuse-defect-trichotomy}.

If regenerated concentration occurs, then a positive oscillatory observer
survives after all positive-threshold concentration sets were removed, contradicting
the definition of the diffuse residual.  If affine/parasitic collapse occurs,
then \Cref{def:affine-constant-lower-stratum,thm:descendant-heredity} places
the relevant descendant hull in an already closed lower stratum.  This cannot
occur inside \(\Cgen\).

The only remaining alternative is a critical diffuse tail.  By
\Cref{thm:critical-tail-compactification-extraction}, it either routes first to
an affine-pressure or pressure/noncompact branch, or else generates a nonzero
retained \(\mathcal T\in\Tcrit\).  The routed
branches have already been assigned to lower terminal strata and cannot remain
inside the refined residual class; the retained \(\Tcrit\) case contradicts
\Cref{cor:critical-tail-exclusion-complete}.  Hence no mixed compact--diffuse residual
profile exists.  If the mixed configuration were witnessed by an escaping
retained recentering, \Cref{lem:escaping-retained-recentering-routing} would
first route that recentering to the appropriate terminal branch; the
contradiction occurs only after the corresponding branch closure theorem is
applied.
\end{proof}

\begin{lemma}[Escaping retained recenterings are routed by the terminal protocol]
\label{lem:escaping-retained-recentering-routing}
An escaping retained recentering coexisting with a compact retained core is not
an immediate contradiction.  By
\Cref{lem:retained-recentering-alternatives}, such a recentering is assigned to
one of the following: a locally compact retained descendant, a diffuse exterior
defect, a pressure/noncompact defect, or a critical-tail defect.  The locally
compact case is recorded as a separated retained successor or a
concentration-chain successor.  The diffuse, pressure/noncompact, and
critical-tail branches are then handled by the corresponding terminal closure
arguments.
\end{lemma}

\begin{proof}
This is the routing content of
\Cref{lem:retained-recentering-alternatives}.  Alternatives \textup{(i)} and
\textup{(ii)} there produce locally compact descendants, which are precisely the
separated-family or chain successors used in the terminal protocol.  Alternative
\textup{(iii)} routes the escaping recentering to the diffuse,
pressure/noncompact, or critical-tail terminal branch.  No contradiction is
claimed at the classification stage; the contradiction appears only after the
corresponding branch closure theorem is applied.
\end{proof}

\begin{definition}[Finite separated family of retained velocity concentration profiles]
\label{def:finite-separated-profile-family}
A finite separated family of retained velocity concentration profiles is a finite
family of recentering sequences.  For the original terminal extraction these
are the parabolic recenterings of \Cref{def:parabolic-recentering-sequence}; for
descendants of a centered profile they are the covariant tail recenterings of
\Cref{def:spacetime-covariant-action}.  We write them uniformly as
\(g_{1,n},\dots,g_{J,n}\), \(J\ge2\), such that:
\begin{enumerate}[label=(\roman*)]
   \item the recentering sequences are mutually separated;
   \item each recentered sequence converges locally smoothly for velocity and
   locally in \(L^{3/2}\) for pressure modulo gauges to a terminal profile
   \((U_j,\Pi_j)\);
   \item each limiting profile has retained local velocity concentration:
\[
   \Vel_1(U_j;0)\ge\eta_*
   \qquad (j=1,\dots,J).
\]
\end{enumerate}
This is a finite local family, not a statement that all concentration profiles
in the whole space have been counted.  The separated-profile proof below treats
the possible continuation of escaping recenterings by local concentration
alternatives, not by a global counting estimate.
\end{definition}

\section{Separated concentration profiles, indecomposability, and sequence-\texorpdfstring{\(L^3\)}{L3}}
\label{sec:separated-profiles-atomic}

We now close the terminal alternatives carrying retained local velocity concentration.
The proof uses no global critical-mass budget.  Instead, finite separated
concentration families are treated
through local recentering and compactness.  Iterating retained tail limits from
such a finite family has only local profile outcomes: an indecomposable
retained velocity concentration profile, a mixed compact--diffuse or compact--noncompact
successor, or a chain of retained concentrating descendants.  The mixed
compact--diffuse successors are excluded by \Cref{thm:no-mixed-compact-diffuse},
while the compact--noncompact successors are routed by
\Cref{lem:escaping-retained-recentering-routing}; infinite chains of retained concentrating descendants and finite
loops are excluded by the compact recurrence theorem.

\begin{definition}[Tail limits and terminal indecomposability]
\label{def:atomic-terminal-profile}
Let \((U,\Pi)\) be a terminal profile with retained local velocity concentration.  A
retained concentrating tail limit is a limit of covariant tail recenterings
\[
   \mathscr T_{(x_n,\tau_n)}(U,\Pi),
   \qquad |x_n|+|\tau_n|^{1/2}\to\infty,
\]
   in the local smooth velocity topology and the local \(L^{3/2}\) pressure
   topology modulo gauges, with positive retained local velocity concentration
   at the origin, equivalently with retained local velocity concentration at
   some rational level \(\eta\in\mathbb Q_{>0}\).  An ordinary fixed spatial
   shift is not used in this definition; by
   \Cref{rem:raw-translates,def:spacetime-covariant-action}, it would introduce an
   unbounded spurious drift for escaping centers.  Diffuse exterior concentration
   means that there are \(\tau_n\to-\infty\), \(R_n\to\infty\), and \(\eta>0\),
   such that
  \[
     \int_{|y|>R_n}|U(y,\tau_n)|^3\,dy\ge\eta,
  \]
   while every unit-scale exterior covariant recentering sequence based in the
   parabolic exterior windows
   \[
      \{|y|>R_n\}\times(\tau_n-1,\tau_n+1)
   \]
   has velocity concentration below the velocity threshold, unless that
   recentering sequence has loss of local compactness.  If a pressure-inclusive
   CKN lower bound remains without velocity concentration, it is assigned by
   \Cref{lem:no-pressure-only-retained-profile} to pressure loss, the
   affine/parasitic quotient, or neighboring velocity concentration.
   Thus the word ``diffuse'' is always parabolic: localization is tested by
   unit parabolic cylinders, not only by spatial balls on the distinguished time
   slices.  An escaping non-compact recentering sequence is a
   covariant tail recentering sequence for which the local compactness,
suitability, or pressure control needed to extract a terminal profile fails.  A
retained velocity concentration profile is terminally indecomposable if it has
no covariant tail recentering limit carrying positive retained local velocity
concentration at any rational level \(\eta\in\mathbb Q_{>0}\), no diffuse
exterior concentration, and no escaping non-compact recentering sequence.
\end{definition}

\begin{lemma}[Compact state-space time-slice concentration]
\label{lem:timeslice-to-ckn-activity}
Let \(\calK\subset\mathfrak X_M^\#\) be compact in the normalized local
profile topology, with pressure representatives taken from the fixed setup
pressure atlas and with affine/parasitic modes already removed by the
canonical representative or assigned to \(\Caff\).  For every \(\eps_0>0\)
there are
\[
   \delta_{\rm sl}=\delta_{\rm sl}(\calK,\eps_0)\in(0,1),
   \qquad
   c_{\rm sl}=c_{\rm sl}(\calK,\eps_0)>0,
\]
such that every \((U,\Pi)\in\calK\) satisfying
\[
   \int_{B_1(0)}|U(y,0)|^3\,dy\ge\eps_0>0,
\]
also satisfies, after the standard local pressure gauge,
\[
   \iint_{B_1(0)\times(-\delta_{\rm sl}^2,0)} |U|^3\,dy\,d\tau
   \ge c_{\rm sl},
   \qquad\hbox{and hence}\qquad
   \Vel_1(U;0)\ge c_{\rm sl} .
\]
The same constants apply to any covariant tail recentering family whose
normalized pressure-atlas closure is contained in the same compact set
\(\calK\).
\end{lemma}

\begin{proof}
For \((U,\Pi)\in\calK\) define
\[
   F(U,\tau):=\int_{B_1}|U(y,\tau)|^3\,dy .
\]
The topology of \(\mathfrak X_M^\#\) gives local \(C^0\) convergence of the
velocity.  Hence \((U,\tau)\mapsto F(U,\tau)\) is continuous on
\(\calK\times[-1,1]\).  Since this set is compact, the continuity is uniform.
Choose \(\delta_{\rm sl}\in(0,1)\) so small that, for every
\((U,\Pi)\in\calK\),
\[
   |F(U,\tau)-F(U,0)|\le \frac12\eps_0
   \qquad\hbox{whenever } -\delta_{\rm sl}^2<\tau<0 .
\]
If \(F(U,0)\ge\eps_0\), then \(F(U,\tau)\ge\eps_0/2\) throughout that time
interval.  Therefore
\[
   \iint_{B_1\times(-\delta_{\rm sl}^2,0)}|U|^3\,dy\,d\tau
   \ge \frac12\eps_0\delta_{\rm sl}^2 .
\]
Set \(c_{\rm sl}:=\eps_0\delta_{\rm sl}^2/2\).  The pressure contribution is
not used to certify nontriviality; the pressure atlas is used only to define
the compact normalized state space in which the uniform time modulus is valid.

The dependence on \(\calK\), rather than only on \(M\), is essential.  A
spatially homogeneous parasitic family \(U_N(y,\tau)=b_N(\tau)\), paired with
the affine pressure
\(-(\partial_\tau b_N+\frac12 b_N)\cdot y\), can have a uniform velocity
\(L^\infty\) bound and fixed time-slice mass while losing every uniform
backward time modulus.  Such a family is not compact in the normalized
pressure-atlas topology unless the homogeneous modes have a convergent modulus,
and nonzero homogeneous modes are assigned to the affine/parasitic stratum
\(\Caff\).
\end{proof}

\begin{lemma}[Quantitative exterior slice-to-defect alternative]
\label{lem:exterior-slice-to-diffuse-defect}
Let \((U,\Pi)\in\mathfrak X_M^\#\) be represented after canonical
affine/parasitic normalization.  Let \(\tau_n\to-\infty\), \(R_n\to\infty\), and
\[
   E_n:=\{|y|>R_n\}
\]
satisfy
\[
   \int_{E_n}|U(y,\tau_n)|^3\,dy\ge m_0>0 .
\]
If the exterior covariant recenterings in the windows
\[
   \{|y|>R_n-2\}\times(\tau_n-1,\tau_n+1)
\]
do not have compact closure in the normalized local topology with the setup
pressure atlas fixed, then alternative \textup{(i)} below occurs.  Otherwise,
after passing to a subsequence, let \(\calK_{\rm ext}\subset\mathfrak X_M^\#\)
be a compact set containing all these recentered profiles.  Fix \(\alpha>0\),
and let \(c_{\rm sl}(\calK_{\rm ext},\alpha)>0\) be the parabolic lower bound
given by \Cref{lem:timeslice-to-ckn-activity} for the compact state class
\(\calK_{\rm ext}\) and unit ball time-slice mass \(\alpha\).  Define
\[
   m_n^{\rm vel}
   :=
   \sup_{\substack{|x|>R_n-2\\ |s-\tau_n|\le1}}
   \Vel_1\!\left(\mathscr T_{(x,s)}U;0\right).
\]
If \(m_n^{\rm vel}<c_{\rm sl}(\calK_{\rm ext},\alpha)\), then every unit ball
with center in the exterior collar satisfies
\[
   \int_{B_1(x)}|U(y,\tau_n)|^3\,dy<\alpha
   \qquad (|x|>R_n-1).
\]
Consequently there are measurable exterior subsets \(F_n\subset E_n\), escaping
every compact set, such that
\[
   m_0\le\int_{F_n}|U(y,\tau_n)|^3\,dy\le2m_0,
\]
and, after restricting \(F_n\) to a bounded-overlap unit cover, the normalized
exterior time-slice measures have no unit-observer atom larger than
\(C_{\rm ov}\alpha/m_0\), where \(C_{\rm ov}\) is the overlap constant.

Moreover, at least one of the following alternatives occurs after passing to a
subsequence:
\begin{enumerate}[label=\textup{(\roman*)}]
   \item an exterior covariant recentering has loss of local compactness or
   pressure-gauge compactness;
   \item a pressure-inclusive exterior unit cylinder has a lower bound not
   carried by velocity, and is assigned by
   \Cref{lem:no-pressure-only-retained-profile};
   \item the exterior mass is critical only after the logarithmic or
   normalized-tail compactification and enters the corresponding critical-tail
   alternative;
   \item the remaining exterior mass is a diffuse exterior concentration in the
   sense of \Cref{def:atomic-terminal-profile}.
\end{enumerate}
\end{lemma}

\begin{proof}
If the exterior recentering family has no compact closure in the normalized
pressure-atlas topology, this is precisely the local compactness or pressure
gauge failure alternative \textup{(i)}.  We therefore work under the compact
closure assumption and use the compact set \(\calK_{\rm ext}\) from the
statement.

Since \(|U(\cdot,\tau_n)|^3dy\) is a nonatomic measure absolutely continuous
with respect to Lebesgue measure, the lower bound on \(E_n\) allows us to
choose a measurable subset \(F_n\subset E_n\) with
\[
   m_0\le M_n:=\int_{F_n}|U(y,\tau_n)|^3\,dy\le2m_0 .
\]
For instance, take an exterior truncation with mass larger than \(m_0\) and,
if necessary, cut it by a level set of a smooth radial coordinate to reduce the
mass below \(2m_0\).  Because \(F_n\subset E_n\) and \(R_n\to\infty\), these
sets escape every compact set.

Let \(\{B_1(x_{n,j})\}_j\) be a measurable unit-ball cover of \(F_n\), chosen
with centers \(|x_{n,j}|>R_n-1\) and with overlap at most \(C_{\rm ov}\) after
enlarging the balls by a fixed harmless factor.  Suppose that for some \(j\),
\[
   \int_{B_1(x_{n,j})}|U(y,\tau_n)|^3\,dy\ge\alpha .
\]
Apply \Cref{lem:timeslice-to-ckn-activity} to the covariantly recentered
solution \(\mathscr T_{(x_{n,j},\tau_n)}(U,\Pi)\), which belongs to
\(\calK_{\rm ext}\).  At local time \(0\) the ball \(B_1(0)\) is the original
ball \(B_1(x_{n,j})\).  Hence
\[
   \Vel_1\!\left(\mathscr T_{(x_{n,j},\tau_n)}U;0\right)
   \ge c_{\rm sl}(\calK_{\rm ext},\alpha),
\]
so the parabolic velocity part of the raw atlas defect measure on that unit
cylinder has mass at least \(c_{\rm sl}(\calK_{\rm ext},\alpha)\).  This is the only
``thickening'' used here: it applies to a unit ball with a quantitative
time-slice lower bound, not to an arbitrarily diffuse exterior slice as a
whole.  The displayed lower bound contradicts
\(m_n^{\rm vel}<c_{\rm sl}(\calK_{\rm ext},\alpha)\).  This proves the stated
unit-ball time-slice bound.

Define discrete exterior slice measures on the observer centers by
\[
   \lambda_n^{\rm sl}
   :=
   \frac{1}{M_n}
   \sum_j
   \left(\int_{B_1(x_{n,j})\cap F_n}|U(y,\tau_n)|^3\,dy\right)
   \delta_{(x_{n,j},\tau_n)} .
\]
Their total masses are bounded above and below by constants depending only on
the overlap convention, and each unit observer ball has mass at most
\(C_{\rm ov}\alpha/m_0\).  Since \(R_n\to\infty\) and
\(\tau_n\to-\infty\), these normalized measures escape every compact subset of
the covariant observer space.  Passing to the observer compactification gives a
nonzero boundary measure unless one of the local extraction inputs fails along
some exterior observer sequence.

If local compactness or pressure-gauge compactness fails along such a sequence,
we are in alternative \textup{(i)}.  If compactness holds but a lower bound in a
unit exterior cylinder is pressure-only, then
\Cref{lem:no-pressure-only-retained-profile} gives alternative \textup{(ii)}.
If the mass is not seen at unit scale but becomes critical only after the
logarithmic or normalized-tail compactification, it is the already
listed critical-tail alternative \textup{(iii)}.  In the complementary case the
escaping exterior slice measure persists, all unit exterior velocity observers are
below the velocity threshold fixed by \(c_{\rm sl}(\calK_{\rm ext},\alpha)\), pressure-only
observers have been assigned to the pressure alternatives, and no noncompact
recentering remains.  The result is the
parabolic diffuse exterior concentration alternative in
\Cref{def:atomic-terminal-profile}.  The argument is local on unit covariant
cylinders and uses no global \(L^3\) or CKN budget.
\end{proof}

\begin{lemma}[One retained tail limit gives two separated concentration profiles]
\label{lem:descendant-to-separated-family}
Let \(U\) be a retained velocity concentration terminal profile.  If \(U\) has a
retained concentrating tail limit, then the corresponding terminal extraction contains a
finite separated family of retained velocity concentration profiles.
\end{lemma}

\begin{proof}
The origin recentering sequence carries retained velocity concentration:
\[
   \Vel_1(U;0)\ge\eta_*.
\]
Let \(W\) be a retained concentrating tail limit.  By definition there are
covariant relative tail recentering parameters
\(\zeta_n=(x_n,\tau_n)\), with
\(|x_n|+|\tau_n|^{1/2}\to\infty\), such that
\(\mathscr T_{\zeta_n}(U,\Pi)\to(W,\Pi_W)\) in the terminal topology and
\[
   \Vel_1(W;0)\ge\eta_*.
\]
By strong local convergence of the velocity,
\[
   \Vel_1(\mathscr T_{\zeta_n}U;0)
   \ge\frac12\eta_*
\]
for all sufficiently large \(n\).  Thus, for large \(n\), the origin
recentering sequence and the covariant tail sequence \(\zeta_n\) are separated
and have retained local velocity concentration.
They form a finite separated two-profile concentration family.
\end{proof}

\begin{proposition}[No finite separated family and local exclusions imply indecomposability]
\label{prop:no-multi-implies-atomic}
After finite separated-family, mixed compact--diffuse, and noncompact escaping
alternatives have been excluded, every single retained velocity concentration terminal
profile arising from a refined residual candidate is terminally indecomposable.
\end{proposition}

\begin{proof}[Proof deferred]
The proof is given immediately after
\Cref{thm:no-finite-separated-profile-family}.  This proposition is not used in
the proof of that theorem.
\end{proof}

\begin{lemma}[Indecomposable sequence-\(L^3\) completeness]
\label{lem:atomic-sequence-L3}
Let \((U,\Pi)\in\mathfrak X_M^\#\) be represented after canonical
affine/parasitic normalization.  Suppose \(U\) has retained local velocity
concentration and is terminally indecomposable in the sense of
\Cref{def:atomic-terminal-profile}.  Then there exists a sequence
\(\tau_k\to-\infty\) such that
\[
   \sup_k\|U(\cdot,\tau_k)\|_{L^3(\R^3)}<\infty .
\]
\end{lemma}

\begin{proof}
We prove the contrapositive.  Define
\[
   A(\tau):=\int_{\R^3}|U(y,\tau)|^3\,dy\in[0,\infty].
\]
If the conclusion fails, then
\[
   \liminf_{\tau\to-\infty}A(\tau)=\infty.
\]
Choose \(\tau_n\to-\infty\) such that \(A(\tau_n)\to\infty\), allowing the
value \(+\infty\).  Since \(U\) is bounded by some \(M\), for each finite
\(R\),
\[
   \int_{|y|\le R}|U(y,\tau_n)|^3\,dy\le M^3|B_R|.
\]
Choose \(R_n\to\infty\) so slowly that, after discarding finitely many \(n\),
\[
   \int_{|y|>R_n}|U(y,\tau_n)|^3\,dy\ge1.
\]
If \(A(\tau_n)=+\infty\), take \(R_n=n\); otherwise choose \(R_n\) with
\(M^3|B_{R_n}|\le A(\tau_n)/2\).

Define the parabolic exterior velocity concentration level
\[
   m_n^{\rm vel}
   :=
   \sup_{\substack{|x|>R_n-2\\ |s-\tau_n|\le1}}
   \Vel_1\!\left(
      \mathscr T_{(x,s)}U;0
   \right).
\]
This is the quantity which can be passed to a limiting profile by strong local
\(L^3\) convergence.  Pressure-inclusive CKN mass in the same exterior windows
is handled separately by \Cref{lem:no-pressure-only-retained-profile}.

If some unit exterior covariant recentering sequence in the windows
\[
   \{|y|>R_n-2\}\times(\tau_n-1,\tau_n+1)
\]
has loss of local compactness or pressure-gauge compactness, terminal
indecomposability is contradicted by the noncompact defect alternative.  Hence
we may assume that the normalized pressure-atlas closure of all such exterior
recentering profiles is compact; denote one compact closure by
\(\calK_{\rm ext}\subset\mathfrak X_M^\#\).

If \(\limsup_{n\to\infty}m_n^{\rm vel}=\gamma>0\), then, passing to a
subsequence, choose \((x_n,s_n)\), with
\[
   |x_n|>R_n-2,\qquad |s_n-\tau_n|\le1,
\]
such that
\[
   \Vel_1\!\left(\mathscr T_{(x_n,s_n)}U;0\right)
   \ge\frac{\gamma}{2}.
\]
Since \(R_n\to\infty\) and \(\tau_n\to-\infty\), these are escaping covariant
tail recenterings.  If they have loss of local compactness, terminal
indecomposability is contradicted.  Otherwise local compactness gives a
subsequential terminal limit \((W,\Pi_W)\), and strong local \(L^3\) convergence gives
\[
   \Vel_1(W;0)\ge\frac{\gamma}{4}>0 .
\]
Choose a rational \(0<\eta<\gamma/4\).  Thus \(W\) is a covariant tail limit
with retained local velocity concentration at level \(\eta\), contradicting the
all-positive-threshold definition of terminal indecomposability.  Therefore
\[
   m_n^{\rm vel}\to0 .
\]

Fix a small \(\alpha_0>0\), and let
\[
   c_{\rm sl}(\calK_{\rm ext},\alpha_0)>0
\]
be the compact-state constant from \Cref{lem:timeslice-to-ckn-activity}, as
used in \Cref{lem:exterior-slice-to-diffuse-defect}.  Since
\(m_n^{\rm vel}\to0\), for all large \(n\),
\[
   m_n^{\rm vel}<c_{\rm sl}(\calK_{\rm ext},\alpha_0).
\]
If the
velocity concentration is small but the pressure-inclusive CKN mass is not,
\Cref{lem:no-pressure-only-retained-profile} assigns the recentering either to
pressure compactness failure, to the affine/parasitic pressure quotient, or to
a neighboring unit cylinder with positive velocity concentration.  Each case is
already excluded by terminal indecomposability or by lower-stratum removal.  If the
exterior mass is critical only after log-window normalization, the resulting
normalized support profile is a critical-tail defect and is one of the already excluded
terminal alternatives.  If no such recentering loses compactness and no
critical-tail normalization is selected, then the exterior time-slice lower bound
\[
   \int_{|y|>R_n}|U(y,\tau_n)|^3\,dy\ge1
\]
together with the parabolic velocity smallness
\(m_n^{\rm vel}<c_{\rm sl}(\calK_{\rm ext},\alpha_0)\)
is now
handled quantitatively by \Cref{lem:exterior-slice-to-diffuse-defect}.  Apply
that lemma with \(m_0=1\), \(\alpha=\alpha_0\), and the displayed inequality.
If the lemma produces local compactness loss, pressure-gauge loss, pressure-only
mass, or a critical-tail normalization, the branch is one of the already covered
noncompact, pressure, affine/parasitic, or critical-tail alternatives.  In the
remaining case it produces diffuse exterior concentration in the sense of
\Cref{def:atomic-terminal-profile}.  This again contradicts terminal
indecomposability.  Both parabolic alternatives are impossible, so the desired
sequence exists.
\end{proof}

\begin{proposition}[Stability of the bounded mild formulation]
\label{prop:bounded-mild-stability}
Let \((U_m,\Pi_m)\) be normalized parasitic-free centered profiles whose
physical pullbacks are bounded mild Navier--Stokes solutions on every finite
terminal interval.  Assume \(U_m\to U\) locally smoothly in centered variables
and that the \(U_m\) are uniformly bounded on compact centered time intervals.
Then the physical pullback of \(U\) is also bounded mild on every finite
terminal interval.  Let
\[
   u(x,t):=(-t)^{-1/2}
   U\!\left(\frac{x}{\sqrt{-t}},-\log(-t)\right),
   \qquad t<0,
\]
be the physical pullback.  For each fixed \(T<0\), define the finite terminal
shift
\[
   u^T(x,s):=u(x,T+s)
   =
   (-(T+s))^{-1/2}
   U\!\left(\frac{x}{\sqrt{-(T+s)}},-\log (-(T+s))\right),
   \qquad s<0,
\]
Then, after translating the interval so that its terminal time is \(0\), one
has, for all \(-\infty<s<t<0\),
\begin{equation}\label{eq:finite-shift-duhamel}
   u^T(t)=e^{(t-s)\Delta}u^T(s)
   -\int_s^t e^{(t-\sigma)\Delta}
      \mathbb P\nabla\cdot(u^T\otimes u^T)(\sigma)\,d\sigma
\end{equation}
in the bounded mild duality sense used in the setup paper.  The same conclusion is
preserved by centered time translations, covariant translations, retained tail
limits, concentrating descendants, recurrent-core limits, and affine normalization
outside \(\Caff\).
\end{proposition}

\begin{proof}
For each \(m\), the normalized setup representative satisfies the
Duhamel formula on every finite physical interval.  Fix \(T<0\) and
\(-\infty<s<t<0\).  Pull the centered convergence back to the compact physical
time interval \([s,t]\) after the finite terminal shift.  The velocities
converge locally smoothly and are uniformly bounded; hence
\[
   u_m^T\otimes u_m^T\to u^T\otimes u^T
\]
locally in every finite \(L^q\), and in particular in distributions.  Testing
the Duhamel identity for \(u_m^T\) against a compactly supported smooth vector
field and using the heat-kernel/Leray-projection duality gives convergence of
the linear heat term and of the bilinear Duhamel term.  The bounded mild
duality formulation used in the setup paper then extends the identity from
compactly supported tests to the standard bounded mild class, yielding
\eqref{eq:finite-shift-duhamel} for \(u^T\).

Centered time translations and covariant translations are exact symmetries of
\eqref{eq:centered-main}; in physical variables they are compositions of
translations, scalings, and finite time shifts, which preserve the Duhamel
identity.  Retained tail limits, concentrating descendants, and recurrent-core limits
are locally smooth limits of such normalized representatives, so the preceding
limit argument applies.  The affine normalization either removes a spatially
homogeneous parasitic component or assigns the profile to \(\Caff\); outside
\(\Caff\) no homogeneous parasitic forcing remains, and the same Duhamel
identity holds for the canonical representative.
\end{proof}

\begin{theorem}[Parasitic-free mildness inheritance for normalized profiles]
\label{thm:mildness-inheritance-main}
Let \((U,\Pi)\) be a retained terminal profile, retained covariant tail
descendant, concentrating descendant limit, or recurrent-core limit represented in the
normalized residual profile class generated by the Seregin extraction in the
setup paper.
Then exactly one of the following alternatives applies.
\begin{enumerate}[label=\textup{(\roman*)}]
   \item \((U,\Pi)\) belongs to the affine/parasitic lower stratum of
   \Cref{def:affine-constant-lower-stratum};
   \item after the canonical affine normalization already built into the
   residual representative, the physical pullback
\[
   u(x,t)=(-t)^{-1/2}
   U\!\left(\frac{x}{\sqrt{-t}},-\log(-t)\right),
   \qquad t<0,
\]
is parasitic-free and has the following property: for every \(T<0\), the
finite-terminal shift
\[
   u^T(x,s):=u(x,T+s),\qquad s<0,
\]
is a bounded mild ancient Navier--Stokes solution in the Albritton--Barker
sense on \(\mathbb R^3\times(-\infty,0)\).
\end{enumerate}
\end{theorem}

\begin{proof}
Affine-constant centered profiles are precisely the homogeneous parasitic
profiles removed by \Cref{def:affine-constant-lower-stratum}; they are assigned
to alternative (i).  We therefore consider a representative outside that lower
stratum.

The setup paper constructs normalized Seregin limits as centered bounded mild
ancient representatives, after removal of the spatially homogeneous parasitic
mode.  The stability statement
\Cref{prop:bounded-mild-stability} shows that the limiting operations used
in this paper preserve the bounded mild formulation: exact centered symmetries
preserve the Duhamel identity, local limits preserve it by bounded mild duality,
and affine normalization either removes the homogeneous parasitic mode or places
the profile in \(\Caff\).  Thus retained terminal profiles, covariant tail
descendants, concentrating descendant limits, and recurrent-core limits remain
parasitic-free mild unless they have fallen into the affine/parasitic lower
stratum.

Finally fix \(T<0\).  On \(\R^3\times(-\infty,T]\), the physical pullback
\[
   u(x,t)=(-t)^{-1/2}
   U\!\left(\frac{x}{\sqrt{-t}},-\log(-t)\right)
\]
is smooth and bounded by \((-T)^{-1/2}\|U\|_{L^\infty}\), and it solves the
physical Navier--Stokes equations by the direct self-similar change of
variables.  By the inherited parasitic-free mild formulation, for
\(-\infty<s<t\le T\),
\[
   u(t)=e^{(t-s)\Delta}u(s)
   -\int_s^t e^{(t-\sigma)\Delta}
      \mathbb P\nabla\cdot(u\otimes u)(\sigma)\,d\sigma .
\]
After translating the finite terminal time \(T\) to \(0\), this is the mild
ancient formulation for \(u^T\) on every finite forward interval
\((-\infty,0)\).  The bound
\(\|u^T\|_{L^\infty(\mathbb R^3\times(-\infty,0))}
\le (-T)^{-1/2}\|U\|_{L^\infty}\) is finite.  Thus \(u^T\) is a genuine
bounded mild ancient solution in the sense required by the
Albritton--Barker endpoint theorem.
\end{proof}

\begin{theorem}[Albritton--Barker endpoint Liouville theorem]
\label{thm:AB-main}
Let \(u\in L^\infty(\mathbb R^3\times(-\infty,0))\) be a bounded mild ancient
solution of the three-dimensional Navier--Stokes equations.  If there exists
\(t_k\to-\infty\) such that
\[
   \sup_k\|u(\cdot,t_k)\|_{L^3(\R^3)}<\infty,
\]
then \(u\equiv0\).
\end{theorem}

\begin{proof}
This is the endpoint ancient Liouville theorem in the sequence-\(L^3\) form
proved by Albritton and Barker \cite{AlbrittonBarker2019}.
\end{proof}

\begin{lemma}[No retained indecomposable concentration profile]
\label{lem:no-atomic-active}
No retained terminally indecomposable concentration profile exists.
\end{lemma}

\begin{proof}
Let \(U\) have retained local velocity concentration and be terminally indecomposable.  By
\Cref{lem:atomic-sequence-L3}, choose \(\tau_k\to-\infty\) such that
\[
   \sup_k\|U(\cdot,\tau_k)\|_{L^3}<\infty .
\]
Let \(u\) be the physical pullback and set \(t_k=-e^{-\tau_k}\).  Since
\(\tau_k\to-\infty\), the sequence satisfies \(t_k\to-\infty\), and
critical scaling gives
\[
   \|u(\cdot,t_k)\|_{L^3(\R^3)}
   =
   \|U(\cdot,\tau_k)\|_{L^3(\R^3)}.
\]
By \Cref{thm:mildness-inheritance-main}, either \(U\) is in the affine/parasitic
lower stratum, which is not part of the normalized residual class, or else for
each \(T<0\)
\[
   u^T(x,s):=u(x,T+s),\qquad s<0,
\]
is a bounded mild ancient solution.  Choose \(T<0\) and discard finitely many
indices so that \(t_k<T\).  Then \(s_k:=t_k-T\to-\infty\), and
\[
   \|u^T(\cdot,s_k)\|_{L^3(\R^3)}
   =
   \|u(\cdot,t_k)\|_{L^3(\R^3)}
   \le \sup_j\|U(\cdot,\tau_j)\|_{L^3(\R^3)}<\infty .
   \]
\Cref{thm:AB-main} gives \(u^T\equiv0\).  Since \(T<0\) was arbitrary, the
physical pullback vanishes on \(\mathbb R^3\times(-\infty,0)\), and therefore
\(U\equiv0\).  This contradicts retained local velocity concentration.
\end{proof}

\begin{lemma}[Diagonal lifting of retained concentrating tail limits]
\label{lem:diagonal-lifting-descendants}
Let \(A=\{\mathfrak s_1,\dots,\mathfrak s_J\}\), \(J\ge2\), be a finite
retained velocity concentration separated family for a refined residual candidate, represented by
separated recentering sequences \(g_{j,n}=(y_{j,n},\tau_{j,n})\).  Let
\((U_j,\Pi_j)\) be the terminal profile obtained in the \(j\)-th recentering.  If
\((U_j,\Pi_j)\) has a retained concentrating tail limit \((W,\Pi_W)\), then
there is a combined parabolic-covariant descendant recentering sequence
\(\mathfrak d=(d_n)\), in the sense of
\Cref{def:combined-raw-covariant-recentering}, for the original terminal
sequence such that:
\begin{enumerate}[label=(\roman*)]
   \item \(\mathfrak d\) is asymptotically separated from \(g_j\);
   \item the recentered original sequence at \(\mathfrak d\) converges locally
   to \((W,\Pi_W)\), after admissible pressure gauges;
   \item \(\mathfrak d\) has retained local velocity concentration.
\end{enumerate}
\end{lemma}

\begin{proof}
By definition of retained concentrating tail limit, there are relative
covariant tail recenterings \(\zeta_m=(\xi_m,\sigma_m)\), with
\[
   |\xi_m|+|\sigma_m|^{1/2}\to\infty,
\]
such that
\[
   \mathscr T_{\zeta_m}(U_j,\Pi_j)\to(W,\Pi_W)
\]
locally in the terminal topology, and
\[
   \Vel_1(W;0)\ge\eta_* .
\]
Fix \(m\).  Since \(V_n^{g_j}\to U_j\) locally smoothly and
\((P_n^{g_j})\to\Pi_j\) locally in \(L^{3/2}\) modulo gauges, the convergence
also holds on the fixed covariant tube
\[
   \{(y+e^{s/2}\xi_m,s+\sigma_m):(y,s)\in Q_m(0,0)\}.
\]
Therefore one can choose \(n(m)\) so large that the covariantly recentered
original fields
\[
   V_{n(m)}(y+y_{j,n(m)}+e^{s/2}\xi_m,\,
             s+\tau_{j,n(m)}+\sigma_m)
\]
and the corresponding pressure gauges are within \(m^{-1}\) of the
\(\mathscr T_{\zeta_m}\)-recentered fields of \((U_j,\Pi_j)\) on
\(Q_m(0,0)\).

Set
\[
   d_m:=(g_{j,n(m)};\xi_m,\sigma_m).
\]
Relabeling the subsequence \(n(m)\) as the terminal index gives a combined
parabolic-covariant descendant recentering sequence \(\mathfrak d=(d_m)\) for the
original terminal sequence.  Since the covariant parabolic distance from \(d_m\) to
\(g_{j,n(m)}\) is
\(|\xi_m|+|\sigma_m|^{1/2}\), the sequence \(\mathfrak d\) is separated from the
\(j\)-th recentering sequence.  The diagonal choice gives local convergence of
the original sequence at \(\mathfrak d\) to \((W,\Pi_W)\), after the admissible
pressure gauges used above.  The diagonal convergence gives strong local
\(L^3\)-convergence of the velocity, so
\[
   \liminf_{m\to\infty}
      \Vel_1\bigl(V_{n(m)}^{\mathfrak d};0\bigr)
   \ge
   \Vel_1(W;0)
   \ge\eta_* .
\]
Thus \(\mathfrak d\) has retained local velocity concentration.  The pressure
gauges are inherited through the same diagonal construction, but
pressure-inclusive \(\calE_1\) is not used to certify retained nonzero
concentration.
\end{proof}

\begin{theorem}[Heredity of retained descendants]
\label{thm:descendant-heredity}
Let \((U,\Pi)\) be a normalized retained residual profile in
\(\mathfrak X_M\), and let \((W,\Pi_W)\) be a retained covariant tail
descendant of \((U,\Pi)\).  Then \((W,\Pi_W)\in\mathfrak X_M\), it carries
retained local velocity concentration at the descendant threshold fixed in
\Cref{def:threshold-hierarchy}, and it has the same parasitic-free normalization
as a terminal profile generated directly by the Seregin extraction in the
setup paper.
Consequently exactly one of the following ordered outcomes occurs:
\begin{enumerate}[label=\textup{(\roman*)}]
   \item \((W,\Pi_W)\) belongs to one of the previously closed lower strata,
   in which case that descendant is already excluded;
   \item \((W,\Pi_W)\) belongs to the normalized residual terminal profile class,
   so the concentrating descendant and terminal-indivisibility arguments may be iterated
   from \(W\).
\end{enumerate}
\end{theorem}

\begin{proof}
The descendant is obtained as a limit of covariant tail recenterings
\(\mathscr T_{\zeta_n}(U,\Pi)\).  By
\Cref{lem:state-space-local-compactness}, the recentered sequence is locally
precompact in \(\mathfrak X_M\), and the limit remains a bounded smooth
locally suitable centered solution with the same \(L^\infty\) bound and
compatible pressure gauges.  Retained local velocity concentration is part of the
definition of retained descendant; when the descendant is produced through a
limit passage, the velocity semicontinuity in \Cref{lem:gauge-compatibility}
fixes the threshold loss once and
for all through \Cref{def:threshold-hierarchy}.  The affine/parasitic
normalization is canonical for the residual representative and is preserved by
\(T_a\), \(\Theta_s\), and their composition
\(\mathscr T_{\zeta_n}\); if a limit is affine-parasitic, it is assigned to the
affine-constant lower stratum of
\Cref{def:affine-constant-lower-stratum}.  Otherwise it is represented in the
same normalized residual profile class.  The ordered residual decomposition then
places the descendant either in a previously closed lower class or in the
generic residual profile class.  This is a local profile reduction and does
not use a global critical norm or global tail budget.
\end{proof}

\begin{lemma}[Escaping recenterings in a separated family]
\label{lem:finite-successor-alternatives}
Let \(A=\{\mathfrak s_1,\dots,\mathfrak s_J\}\), \(J\ge2\), be a finite
retained velocity concentration separated family for a refined residual candidate.  For each component
\(\mathfrak s_j\), the following ordered alternatives exhaust the possible
successor profiles:
\begin{enumerate}[label=(\roman*)]
   \item the profile \((U_j,\Pi_j)\) is terminally indecomposable;
   \item \((U_j,\Pi_j)\) has diffuse exterior concentration or an escaping
   recentering sequence with loss of local compactness;
   \item \((U_j,\Pi_j)\) has a retained concentrating tail limit which lifts to a
   recentering sequence of the original terminal sequence with retained local
   velocity concentration.
\end{enumerate}
\end{lemma}

\begin{proof}
This is the definition of terminal indecomposability plus the diagonal lifting
lemma, with no disjointness assertion.  If \((U_j,\Pi_j)\) is terminally
indecomposable, (i) holds.  If it is not terminally indecomposable, at least one
of the three excluded alternatives in \Cref{def:atomic-terminal-profile} occurs:
a retained concentrating tail limit, diffuse exterior concentration, or an
escaping recentering sequence with loss of local compactness.  The diffuse
exterior and local-compactness-failure cases give (ii).  In the retained
concentrating tail-limit case, \Cref{lem:diagonal-lifting-descendants} lifts the
limit to a recentering sequence of the original terminal sequence with retained
local velocity concentration, and \Cref{thm:descendant-heredity} places the descendant
either in a closed lower stratum or back in the normalized residual profile class,
giving (iii) in the iterative residual case.  If several successor profiles
coexist, they are all recorded; the next lemma removes the nonconcentrating ones.
\end{proof}

\begin{lemma}[No diffuse or non-compact escaping recentering from a retained family]
\label{lem:no-inactive-successor}
In the setting of \Cref{lem:finite-successor-alternatives}, alternative
\emph{(ii)} cannot occur.
\end{lemma}

\begin{proof}
Suppose that \((U_j,\Pi_j)\) has diffuse exterior concentration or an escaping
recentering sequence with loss of local compactness.  By the same diagonal
procedure used in \Cref{lem:diagonal-lifting-descendants}, this tail profile is
realized by a terminal extraction of the original refined residual candidate,
after passing to a subsequence and choosing the local pressure gauges supplied
by the imported setup results.  The fact that the diffuse or
noncompact profile is not the sole source of the original compact concentration
does not exclude coexistence with the compact core \((U_j,\Pi_j)\).  In the
diffuse case, \Cref{thm:no-mixed-compact-diffuse} rules out exactly this
coexistence in the refined residual class.  In the escaping noncompact case,
\Cref{lem:escaping-retained-recentering-routing} shows that the profile is not
an immediate contradiction but is routed to a locally compact retained
descendant, a diffuse defect, a pressure/noncompact defect, or a critical-tail
defect.  The locally compact routed case is recorded as a successor, so it does
not belong to alternative \emph{(ii)}.  The diffuse, pressure/noncompact, and
critical-tail routed cases are closed by the corresponding terminal branch
arguments.  Therefore alternative \emph{(ii)} cannot occur for a refined
residual family with retained local velocity concentration.
\end{proof}

\begin{definition}[Closed retained velocity concentration successor relation]
\label{def:closed-active-successor-relation}
For a fixed threshold \(\eta>0\), let
\[
   \mathfrak A_\eta:=\mathfrak A_\eta^\#
   =
   \{U\in\mathfrak X_M^\#:\Phi(U)\ge\eta\}
\]
be the compact non-affine oscillation set of
\Cref{cor:compact-normalized-activity-set}.  We write
\[
   U\,\mathcal R_\eta\,V
\]
if \(U,V\in\mathfrak A_\eta\) and \(V\) is a retained concentrating covariant tail
descendant of \(U\).
\end{definition}

\begin{lemma}[Closedness of the retained velocity concentration successor relation]
\label{lem:closed-active-successor-relation}
The relation
\[
   \mathcal R_\eta\subset\mathfrak A_\eta\times\mathfrak A_\eta
\]
is closed in the local \(\mathfrak X_M\) topology.
\end{lemma}

\begin{proof}
Let \(U_m\to U\), \(V_m\to V\) in \(\mathfrak X_M^\#\), with
\(U_m\,\mathcal R_\eta\,V_m\).  For each \(m\), choose a retained covariant tail
observer sequence \(\zeta_{m,\ell}\) for \(U_m\) whose recentered fields converge
to \(V_m\) on \(Q_m(0,0)\) to accuracy \(m^{-1}\), after admissible pressure
gauges.  A diagonal choice in \((m,\ell)\), using
\Cref{lem:state-space-local-compactness,lem:gauge-compatibility}, gives a
single escaping covariant tail sequence for \(U\) converging locally to \(V\).
Closedness of \(\mathfrak A_\eta\) follows from
\Cref{cor:compact-normalized-activity-set}.  The retained descendant threshold
is preserved by strong local \(L^3\) convergence for the velocity part, or by
closedness of the non-affine oscillation functional \(\Phi\) when the successor
relation is recorded in \(\mathfrak A_\eta\).  No pressure-only lower bound is
used to keep the edge.
Hence \(U\,\mathcal R_\eta\,V\).
\end{proof}

\begin{theorem}[Recurrence for retained velocity concentration chains]
\label{thm:active-path-space-recurrence}
Assume there exists an infinite chain of retained concentrating descendants
\[
   U_0\,\mathcal R_\eta\,U_1\,\mathcal R_\eta\,U_2\,
   \mathcal R_\eta\cdots .
\]
Then the path space
\[
   \mathscr P_\eta
   :=
   \left\{
      (W_j)_{j\ge0}\in\mathfrak A_\eta^{\mathbb N}:
      W_j\,\mathcal R_\eta\,W_{j+1}\ \hbox{for all }j
   \right\}
\]
is nonempty and compact, is invariant under the left shift
\[
   S(W_0,W_1,W_2,\ldots)=(W_1,W_2,\ldots),
\]
and carries an \(S\)-invariant Borel probability measure \(\mathbb P\).  Its
projected stationary profile measure \(\nu=(\pi_0)_\#\mathbb P\) is supported
entirely on \(\mathfrak A_\eta\).
\end{theorem}

\begin{proof}
The set \(\mathfrak A_\eta=\mathfrak A_\eta^\#\) is compact by
\Cref{cor:compact-normalized-activity-set}.  Hence
\(\mathfrak A_\eta^{\mathbb N}\) is compact and metrizable.  By
\Cref{lem:closed-active-successor-relation},
\[
   \mathscr P_\eta
   =
   \bigcap_{j=0}^\infty
   \{(W_k)_{k\ge0}:(W_j,W_{j+1})\in\mathcal R_\eta\}
\]
is closed, hence compact.  The given infinite chain is an element of
\(\mathscr P_\eta\), so the path space is nonempty, and it is plainly invariant
under the left shift.

For the path \(p=(U_0,U_1,\ldots)\), define empirical measures
\[
   \mathbb P_N:=\frac1N\sum_{k=0}^{N-1}\delta_{S^kp}.
\]
Weak-* compactness of probability measures on compact \(\mathscr P_\eta\) gives
a convergent subsequence \(\mathbb P_{N_\ell}\rightharpoonup\mathbb P\).  For
each continuous \(F\) on \(\mathscr P_\eta\),
\[
   \int F\circ S\,d\mathbb P_N-\int F\,d\mathbb P_N
   =
   \frac{F(S^Np)-F(p)}{N}\to0.
\]
Thus \(\mathbb P\) is shift-invariant.  Since every coordinate of every path in
\(\mathscr P_\eta\) belongs to \(\mathfrak A_\eta\), the projected measure
\(\nu=(\pi_0)_\#\mathbb P\) is supported in \(\mathfrak A_\eta\).
\end{proof}

\begin{lemma}[Concatenation of retained concentrating descendants]
\label{lem:active-descendant-concatenation}
Let
\[
   U_0\,\mathcal R_\eta\,U_1\,\mathcal R_\eta\cdots
   \mathcal R_\eta\,U_N
\]
be a finite chain of retained concentrating descendants.  Then \(U_N\) is obtained from \(U_0\) by a
single diagonal covariant recentering sequence, after admissible pressure
gauges, and it remains in \(\mathfrak A_\eta\).  The pressure gauges in the
intermediate limits can be chosen compatibly with the final diagonal limit.
\end{lemma}

\begin{proof}
For \(N=1\) this is the definition of \(\mathcal R_\eta\).  Assume it holds up
to \(N-1\).  The relation \(U_{N-1}\mathcal R_\eta U_N\) gives escaping
covariant observers for \(U_{N-1}\) which converge locally to \(U_N\).  Compose
these observers with the diagonal observers which realize \(U_{N-1}\) from
\(U_0\).  The semidirect covariance identity in
\Cref{lem:semidirect-covariance} shows that the composition is again a
covariant observer sequence for \(U_0\).  A diagonal choice on \(Q_m(0,0)\),
together with the pressure-gauge compatibility in
\Cref{lem:gauge-compatibility}, gives local convergence to \(U_N\).  Since each
link lies in \(\mathfrak A_\eta\), the endpoint does too.
\end{proof}

\begin{lemma}[Recurrent path produces a compact retained velocity concentration core]
\label{lem:recurrent-path-compact-core}
Let \(\mathbb P\) be a shift-invariant probability measure on the retained velocity concentration path
space \(\mathscr P_\eta\).  For \(\mathbb P\)-a.e. recurrent path
\((W_j)_{j\ge0}\), the orbit closure of its coordinates, together with all
covariant tail limits generated by the returning blocks, contains a nonempty
compact invariant set \(\calM_\eta\subset\mathfrak A_\eta\).  Every element of
\(\calM_\eta\) is a retained descendant of the original profile realized from
the original blow-up sequence.
\end{lemma}

\begin{proof}
Poincar\'e recurrence for the left shift gives a full-measure set of paths which
return arbitrarily close to themselves in the product topology.  Fix such a
path.  For each returning block, use
\Cref{lem:active-descendant-concatenation} to compose the finitely many concentrating
successor observers into one covariant recentering of the initial profile.  The
endpoints of the returning blocks remain in the compact set
\(\mathfrak A_\eta\).  Taking the closure of all such endpoints and their
covariant tail limits inside the compact normalized profile class
\(\mathfrak X_M^\#\) gives a nonempty compact set.  Closedness of
\(\mathcal R_\eta\) and compactness of \(\mathfrak A_\eta\) keep the closure
inside \(\mathfrak A_\eta\).  Intersecting over all closed invariant subsets of
this compact set and applying Zorn's lemma gives a compact invariant
minimal core \(\calM_\eta\).  The diagonal construction in the concatenation
lemma preserves realization from the original sequence and retained velocity concentration,
so every element of \(\calM_\eta\) is realized from the original blow-up sequence.
\end{proof}

\begin{lemma}[Retention and gauges on the recurrent concentration core]
\label{lem:active-core-retention-gauges}
The compact core \(\calM_\eta\) of
\Cref{lem:recurrent-path-compact-core} carries the same retained oscillation
threshold:
\[
   \Phi(U)\ge\eta\qquad (U,\Pi)\in\calM_\eta,
\]
and its pressure gauges are compatible with the invariant averaging used in
\Cref{thm:recurrent-tail-core-rigidity}.
\end{lemma}

\begin{proof}
The inequality \(\Phi\ge\eta\) is closed under the local topology by
\Cref{cor:compact-normalized-activity-set}, and all approximating profiles lie in
\(\mathfrak A_\eta\).  Hence every element of the compact core lies in
\(\mathfrak A_\eta\).  Gauge compatibility follows from
\Cref{lem:gauge-compatibility}: on each compact cylinder the pressure gauges are
chosen by a fixed normalization, and diagonal limits preserve the normalized
local \(L^{3/2}\) pressure topology.  Therefore the invariant measures on
\(\calM_\eta\) used in the recurrent-core rigidity theorem average genuine
normalized profiles with retained oscillation, not untracked pressure representatives.
\end{proof}

\begin{theorem}[No infinite retained velocity concentration chain]
\label{thm:compact-active-descendant}
No infinite chain of retained concentrating descendants can occur in the refined residual
class.
\end{theorem}

\begin{proof}
Suppose an infinite chain of retained concentrating descendants exists.  If
\(\limsup_j\Phi(U_j)=0\), local compactness gives a subsequential limit
\(U_{j_k}\to U_\infty\) with \(\Phi(U_\infty)=0\).  The diagonal lifting used in
\Cref{lem:diagonal-lifting-descendants} shows that this limit is still a
retained tail descendant of the original residual profile.  By
\Cref{thm:affine-normalization-dichotomy}, \(U_\infty\) belongs to the
affine/parasitic lower stratum, contradicting descendant heredity in the
refined residual class.  Hence \(\limsup_j\Phi(U_j)>0\).  Choose
\(\eta>0\) and an infinite ordered subchain contained in \(\mathfrak A_\eta\).
By transitivity of retained tail descendants through
\Cref{lem:diagonal-lifting-descendants}, this ordered subchain is an
\(\mathcal R_\eta\)-chain.

\Cref{thm:active-path-space-recurrence} gives a compact shift-invariant path
space and a stationary path measure supported on \(\mathfrak A_\eta\).
\Cref{lem:active-descendant-concatenation,lem:recurrent-path-compact-core,lem:active-core-retention-gauges}
turn a recurrent path into a compact tail-minimal recurrent core for
\eqref{eq:centered-main} whose profiles remain in \(\mathfrak A_\eta\) with
compatible pressure gauges.  The recurrent-core rigidity theorem
\Cref{thm:recurrent-tail-core-rigidity} forces this core to be affine-constant
after the canonical affine normalization.  But affine/parasitic profiles have
\(\Phi=0\), whereas \(\mathfrak A_\eta=\{\Phi\ge\eta\}\).  This contradiction
rules out the infinite chain of retained concentrating descendants.
\end{proof}

\begin{lemma}[Finite separated families produce indecomposable profiles or concentration chains]
\label{lem:finite-graph-cycle}
Let \(A=\{\mathfrak s_1,\dots,\mathfrak s_J\}\), \(J\ge2\), be a finite
retained velocity concentration separated family for a refined residual candidate.  Then either one of
the terminal profiles \((U_j,\Pi_j)\) is terminally indecomposable, or the
family generates an infinite chain of retained concentrating descendants.  A finite recurrent loop
is counted as an infinite chain by periodic repetition.
\end{lemma}

\begin{proof}
Assume no \((U_j,\Pi_j)\) is terminally indecomposable.  By
\Cref{lem:finite-successor-alternatives,lem:no-inactive-successor}, each
\((U_j,\Pi_j)\) has a retained concentrating tail limit which lifts to a
recentering sequence of the original terminal sequence with retained local
velocity concentration.

Starting from \(A\), enlarge the separated family whenever a lifted retained
tail descendant is separated from every profile already selected.  If this
enlargement can be carried out indefinitely, the successive lifted descendants
already form an infinite chain of retained concentrating descendants.  Otherwise the procedure
stops after finitely many enlargements at a finite separated family \(A'\)
which is maximal for this descendant-lifting procedure.

Form a finite directed graph whose vertices are the retained profiles in
\(A'\).  Draw an edge \(i\to j\) if a retained tail descendant of \(U_i\),
lifted by \Cref{lem:diagonal-lifting-descendants}, is equivalent to \(U_j\).
There is no escape edge by the construction of \(A'\): an escape edge would be
exactly another separated retained unit concentrating cylinder in the original
sequence, so the enlargement procedure would not have stopped.  Every
decomposable vertex has an outgoing edge by the preceding paragraph.

A finite directed graph in which every vertex has an outgoing edge contains a
directed cycle.  Therefore, if no terminally indecomposable profile appears, the
directed recentering relation contains a finite loop
\[
   \mathfrak s_{j_0}\to\mathfrak s_{j_1}\to\cdots
   \to\mathfrak s_{j_q}=\mathfrak s_{j_0}.
\]
Repeating the loop periodically gives an infinite chain of retained concentrating descendants.  The
nonconcentrating successor alternatives at each vertex are excluded by
\Cref{lem:no-inactive-successor}, while coexistence of a compact core with a
diffuse remainder is handled by \Cref{thm:no-mixed-compact-diffuse} and an
escaping noncompact recentering is routed by
\Cref{lem:escaping-retained-recentering-routing}.  This proof is graph/recurrence based and
does not use a global counting or mass budget.
\end{proof}

\begin{lemma}[Finite recurrent loops of retained tail limits are impossible]
\label{lem:no-finite-recurrent-cycle}
No refined residual terminal extraction contains a finite recurrent loop of retained
concentrating tail limits.
\end{lemma}

\begin{proof}
Let
\[
   \mathfrak s_{j_0}\to\mathfrak s_{j_1}\to\cdots
   \to\mathfrak s_{j_q}=\mathfrak s_{j_0}
\]
be such a loop, and write \((U_\ell,\Pi_\ell)\) for the corresponding retained
concentration terminal profiles.  Each arrow means that \(U_{\ell+1}\) is a
retained concentrating tail limit of \(U_\ell\).  Thus, for each \(\ell\), there are escaping
tail recentering sequences whose recentered fields converge locally to \(U_{\ell+1}\).

Let \(\calK\) be the closure, in the local smooth velocity topology and local
\(L^{3/2}\) pressure topology modulo gauges, of all profiles obtained from the
finite set \(\{U_0,\dots,U_{q-1}\}\) by time translations and covariant
terminal translations.  The boundedness inherited from the Type I Seregin
extraction theorem and the local pressure compactness from
the imported setup results give compactness of \(\calK\) on every fixed terminal
cylinder; a diagonal argument gives compactness of \(\calK\).  The set is
invariant under the same time and covariant translations by construction.  The
recurrent-loop relations imply that each \(U_\ell\)
returns to \(\calK\) through tail limits, so the closed invariant subset of
\(\calK\) generated by the loop is nonempty, compact, invariant, and contains
the retained velocity concentration profiles.

Repeating the loop periodically gives an infinite chain of retained concentrating descendants.
This contradicts \Cref{thm:compact-active-descendant}.  Therefore no finite
recurrent loop of retained concentrating tail limits exists.
\end{proof}

\begin{theorem}[Finite separated retained velocity concentration profile families are impossible]
\label{thm:no-finite-separated-profile-family}
No refined residual candidate admits a finite separated family of retained
velocity concentration profiles.
\end{theorem}

\begin{proof}
Suppose, for contradiction, that a refined residual candidate admits a finite
retained velocity concentration separated family \(A=\{\mathfrak s_1,\dots,\mathfrak s_J\}\),
\(J\ge2\).  By \Cref{lem:finite-graph-cycle}, either some component of the
family is terminally indecomposable or the family generates an infinite chain of
retained concentrating descendants.  The first alternative contradicts
\Cref{lem:no-atomic-active}.  The second contradicts
\Cref{thm:compact-active-descendant}.  Hence the finite retained velocity concentration
separated family cannot exist.
\end{proof}

\begin{proof}[Proof of \Cref{prop:no-multi-implies-atomic}]
Let \(U\) be a single retained velocity concentration terminal profile after the finite
separated-family, mixed compact--diffuse, and noncompact escaping alternatives
have been excluded.  If \(U\) had a retained concentrating tail limit,
\Cref{lem:descendant-to-separated-family} would produce a finite two-profile
separated family with retained velocity concentration, contradicting
\Cref{thm:no-finite-separated-profile-family}.  If \(U\) had diffuse exterior
concentration or an escaping non-compact recentering sequence, the corresponding
terminal extraction would be a mixed profile configuration consisting of the
retained compact core \(U\) plus a diffuse or noncompact tail remainder.  This
is excluded in the diffuse case by \Cref{thm:no-mixed-compact-diffuse}, while
the escaping noncompact case is routed by
\Cref{lem:escaping-retained-recentering-routing}.  Hence no retained
concentrating tail limit, diffuse exterior concentration, or escaping
non-compact recentering sequence exists, and \(U\) is terminally
indecomposable.
\end{proof}

The terminal endpoint argument uses the following acyclic order:
\[
\begin{array}{c}
\text{concentration-set extraction}\\
\Downarrow\\
\text{diffuse, pressure, noncompact, and critical-tail closures}\\
\Downarrow\\
\text{sequence-}L^3\text{ for terminally indecomposable atomic profiles}\\
\Downarrow\\
\text{no atomic active profile by bounded-mild endpoint Liouville}\\
\Downarrow\\
\text{no finite separated family and no infinite chain}\\
\Downarrow\\
\text{terminal stratification assembly.}
\end{array}
\]
In particular, \Cref{lem:atomic-sequence-L3} does not use
\Cref{thm:terminal-stratification}; it uses only terminal indecomposability and
the already covered exterior, diffuse, pressure/noncompact, and critical-tail
alternatives.

\begin{longtable}{@{}L{0.27\textwidth}L{0.43\textwidth}L{0.22\textwidth}@{}}
\caption{Acyclic dependency table for the terminal concentration argument.}
\label{tab:terminal-dependency-table}\\
\hline
\textbf{Result} & \textbf{May use} & \textbf{Must not use}\\
\hline
\endfirsthead
\hline
\textbf{Result} & \textbf{May use} & \textbf{Must not use}\\
\hline
\endhead
\hline
\endfoot
Concentration-set extraction &
Local Seregin compactness, pressure atlas, and the defect measures \(\mu_n\). &
Terminal assembly theorem.\\
Diffuse-defect compactness &
Observer compactification, concentration-neighborhood removal, and amenable
averaging. &
Atomic Liouville contradiction.\\
Critical-tail closure &
Mesoscopic reduction, affine quotient, bounded-origin realization, and coherent
tail classification. &
Terminal assembly theorem.\\
Mixed compact--diffuse exclusion &
Concentration-set extraction, diffuse-defect compactness, and critical-tail
closure. &
Finite separated-family theorem.\\
Sequence-\(L^3\) lemma &
Terminal indecomposability and absence of retained velocity, diffuse,
noncompact, pressure-only, and critical-tail descendants. &
Albritton--Barker endpoint theorem.\\
Atomic contradiction &
Sequence-\(L^3\), finite-shift bounded mildness, and the Albritton--Barker
endpoint theorem. &
Terminal assembly theorem.\\
Finite separated-family exclusion &
No atomic profile and no infinite retained velocity concentration chain. &
Terminal assembly theorem.\\
Terminal stratification assembly theorem &
All previous rows. &
Itself.\\
\hline
\end{longtable}

\begin{theorem}[Terminal stratification assembly theorem]
\label{thm:terminal-stratification}
Let \((V_n,P_n)\) be a terminal residual sequence realized from the original blow-up sequence and satisfying
the imported setup admissibility hypotheses, the pressure atlas hypotheses, and
a retained unit velocity concentration lower bound
\[
   \liminf_{n\to\infty}\Vel_1(V_n;z_n)\ge\eta_*>0 .
\]
Assume all lower classes already excluded by the setup paper have been removed.  The compact
subthreshold regular residual alternative \textup{(b)} of
\Cref{lem:compact-window-expanding-region-protocol} is excluded before the
terminal stratification by \Cref{lem:compact-subthreshold-regular-residual}
(CKN \(\varepsilon\)-regularity); hence it does not appear as a terminal
residual alternative below.  Then the paired
terminal extraction of \Cref{thm:terminal-exhaustion-main} gives the following
ordered alternatives:
\begin{enumerate}[label=\textup{(\arabic*)},leftmargin=2.2em]
   \item a compact concentrating descendant occurs;
   \item a finite separated family of retained velocity concentration profiles occurs;
   \item an infinite chain of retained concentrating descendants occurs;
   \item a diffuse exterior defect occurs after concentration-neighborhood removal;
   \item a noncompact or pressure-gauge defect occurs;
   \item a critical-tail defect occurs;
   \item affine/parasitic routing occurs into \(\Caff\);
   \item an indecomposable atomic concentration profile occurs and satisfies a
   sequence-\(L^3\) bound.
\end{enumerate}
Alternative \textup{(1)} is a reduction alternative: by descendant heredity it is
iterated inside the same normalized profile class until one reaches
\textup{(2)}, \textup{(3)}, a mixed diffuse/noncompact/critical-tail branch, affine routing, or
the atomic case \textup{(8)}.  Alternatives \textup{(2)}--\textup{(7)} are
excluded by the preceding branch closures and local terminal arguments.
Alternative \textup{(8)} is excluded by the endpoint bounded-mild ancient
Liouville theorem after finite terminal shift.  Hence no terminal residual
profile satisfying the retained unit velocity concentration lower bound remains.
\end{theorem}

\begin{proof}
This is the assembly of the local results and explicit profile-class arguments in
\Cref{sec:terminal-concentration,sec:separated-profiles-atomic}.  The local
concentration decomposition is the paired concentration-set/remainder decomposition
\Cref{thm:terminal-exhaustion-main}.  Closed or continuum active loci are reduced
to the single-profile, finite-separated, or infinite-chain cases by
\Cref{lem:closed-active-locus-single-finite-chain}, so no additional continuum
active-locus branch is hidden in the compact concentration set.  This theorem does not introduce an
additional compactness principle; it only combines the preceding extraction,
reduction, and closure results in the order recorded in
\Cref{tab:terminal-dependency-table}.  The pure radiative alternatives enter
only as sole-source exclusions, and are used only in that form through
\Cref{thm:terminal-inactive-exclusion}.  Coexistence of a compact concentration core
with a diffuse tail remainder is ruled out by
\Cref{thm:no-mixed-compact-diffuse}, while escaping noncompact recenterings are
routed by \Cref{lem:escaping-retained-recentering-routing}.  Pressure-only or
pressure/noncompact defects are assigned by
\Cref{lem:no-pressure-only-retained-profile}.  Critical tails are closed by
\Cref{thm:realized-critical-tail-rigidity,cor:critical-tail-exclusion-complete}.  Infinite concentration chains and finite
recurrent loops are ruled out by the retained-chain recurrence theorem
\Cref{thm:compact-active-descendant}, while finite separated concentration families are ruled out by
\Cref{thm:no-finite-separated-profile-family}.  Therefore any retained
concentration terminal profile that survives the ordered lower-class removals is
single and has no retained concentrating descendant or diffuse/noncompact tail
remainder.  By \Cref{prop:no-multi-implies-atomic} it is terminally
indecomposable, and then
\Cref{lem:atomic-sequence-L3,thm:mildness-inheritance-main,lem:no-atomic-active}
give the sequence-\(L^3\), parasitic-free finite-shift mildness, and endpoint
contradiction.  All estimates used in this assembly are localized to fixed
terminal cylinders, except for the final sequence-\(L^3\) conclusion, which is
derived by contradiction from terminal indecomposability and is not assumed as a
global bound.
\end{proof}

\begin{remark}[No global budget in the terminal assembly theorem]
\label{rem:no-global-budget-terminal-theorem}
The terminal stratification assembly theorem assumes no global \(L^3\), global CKN
budget, finite energy bound on expanding regions, or global counting measure
for concentration profiles.  Compactness and recurrence are local, normalized, or
realized from the original blow-up sequence.
\end{remark}

\begin{corollary}[Rotational residual closure]
\label{cor:R2-closure}
The ordered rotational relative-equilibrium residual class is empty:
\[
   \Crot(\mathcal S;I,J)=\varnothing .
\]
\end{corollary}

\begin{proof}
By \Cref{prop:R2-reduction-to-terminal}, every R2 residual either belongs to a
previously closed lower class or produces a terminal concentration profile
realized from the original blow-up sequence with retained unit velocity concentration.
The lower classes are empty by the imported setup results, and the terminal
concentration case is empty by
\Cref{thm:terminal-stratification}.  Hence no R2 residual remains.
\end{proof}

\begin{corollary}[Stationary-hull residual closure]
\label{cor:R3-closure}
For every admissible stationary phase space \(\mathfrak B\),
\[
   \Cstat(\mathcal S;I,J;\mathfrak B)=\varnothing .
\]
\end{corollary}

\begin{proof}
By
\Cref{prop:R3-stationary-hull-reduction,prop:R3S-R3-terminal-reduction}, every R3
stationary-hull residual either lies in a lower class already excluded by the setup paper or produces
a stationary terminal concentration profile realized from the original blow-up
sequence.  The lower classes are empty by the imported setup results.  The stationary terminal
profile satisfies the retained unit velocity concentration hypotheses of
\Cref{thm:terminal-stratification}, so it is reduced to the atomic alternative
only after all separated, diffuse, noncompact, critical-tail, and affine
alternatives have been removed.  In that atomic case
\Cref{lem:atomic-sequence-L3,thm:mildness-inheritance-main,lem:no-atomic-active}
give the sequence-\(L^3\) bound, finite-shift bounded mildness, and the endpoint
Albritton--Barker contradiction.  Hence no stationary terminal concentration profile remains.
This proves the deferred stationary concentration and degenerate stationary-family
occurrence statements
\Cref{thm:R3S-no-active-stationary-carrier,cor:R3S-no-active-degenerate-family,thm:R3-no-attainable-degenerate-family},
and therefore the R3 class is empty.
\end{proof}

\section{Closure of the generic terminal residual}
\label{sec:R4-closure}

\begin{definition}[Generic residual profile]
\label{def:R4-profile-main}
A generic residual profile is a smooth bounded normalized Seregin ancient solution of
\eqref{eq:centered-main} which carries persistent compact velocity
\(L^3\) concentration and
belongs to none of the previously closed alternatives, none of
\(\Cax,\Crot,\Cstat\), and none of the structured or tight
classes removed earlier, including the ordered critical-tail lower strata
\(\Clogdiff,\Cyoung,\Chomcrit,\Clogper,\Capcrit\).
\end{definition}

\begin{theorem}[Generic terminal exhaustion]
\label{thm:generic-terminal-exhaustion}
Let \((U,\Pi)\in\Cgen\) be represented after canonical affine/parasitic
normalization and after removal of the lower alternatives already excluded by the setup paper.  Then a
terminal extraction of \((U,\Pi)\), with the priority order used throughout the
terminal section, selects one of the following priority alternatives:
\begin{enumerate}[label=\textup{(\arabic*)},leftmargin=2.2em]
   \item a retained compact active descendant;
   \item a finite separated family of retained velocity profiles;
   \item an infinite chain of retained active descendants;
   \item a diffuse exterior defect after active-neighborhood removal;
   \item a noncompact or pressure-gauge defect;
   \item a critical-tail defect;
   \item a single terminally indecomposable retained velocity profile.
\end{enumerate}
This is a concentration-compactness exhaustion of the realized terminal
sequence, not a convention in the definition of \(\Cgen\).
\end{theorem}

\begin{proof}
The generic representative is first put in the affine-normalized quotient.  If
all non-affine oscillation disappears, then
\Cref{thm:affine-normalization-dichotomy} places the profile in the
affine/parasitic lower stratum, contradicting the assumption that we are in
\(\Cgen\).  Thus a positive non-affine oscillation threshold survives, while the
setup results supply retained compact velocity concentration.

Apply the paired terminal concentration-set/remainder decomposition
\Cref{thm:terminal-exhaustion-main} together with the compact-window to
expanding-region protocol
\Cref{lem:compact-window-expanding-region-protocol}.  The retained velocity
mass forces at least one selected active or residual alternative by
\Cref{cor:persistent-forces-active}.  If active neighborhoods persist, the
priority order first records a compact retained descendant; iterating the same
realized-descendant construction either produces a finite separated retained
family, an infinite active chain, or terminates at a single retained profile
with no further retained active descendant.  If active neighborhoods do not
exhaust the residual measure on the selected terminal regions, the protocol
selects a diffuse exterior defect, a loss of local compactness or pressure-gauge
compactness, or a scale-critical tail window.  These alternatives are
\textup{(4)}--\textup{(6)}.  If none of the separated, chained, diffuse,
noncompact, or critical-tail remainders occurs, then the remaining single
profile is terminally indecomposable by \Cref{prop:no-multi-implies-atomic}.
The priority order makes the selected alternative unique.
\end{proof}

\begin{theorem}[Closure of the generic residual class by local concentration analysis]
\label{thm:R4-final-closure}
One has
\[
   \Cgen=\varnothing .
\]
\end{theorem}

\begin{proof}
Suppose, for contradiction, that \(V\in\Cgen\).  Apply
\Cref{thm:generic-terminal-exhaustion}.  A retained compact active descendant is
not a terminal stopping case: by descendant heredity it is reinserted into the
same ordered terminal extraction until it gives a finite separated family, an
infinite active chain, a diffuse/noncompact/critical-tail remainder, or a single
indecomposable retained profile.  Finite separated families are excluded by
\Cref{thm:no-finite-separated-profile-family}, and infinite active chains by
\Cref{thm:compact-active-descendant}.  Diffuse remainders are excluded by
\Cref{thm:terminal-inactive-exclusion,thm:no-mixed-compact-diffuse}.  Escaping
noncompact remainders are routed by
\Cref{lem:retained-recentering-alternatives,lem:escaping-retained-recentering-routing}
to the diffuse, pressure/noncompact, or critical-tail branches; those routed
branches are closed by the corresponding theorems cited here.  Critical
tails are excluded by
\Cref{thm:realized-critical-tail-rigidity,cor:critical-tail-exclusion-complete}.
The pressure-only case is assigned by
\Cref{lem:no-pressure-only-retained-profile} into the pressure/noncompact or
affine/parasitic lower stratum, both already removed.  In the remaining
single-profile case, \Cref{prop:no-multi-implies-atomic} gives terminal
indecomposability and \Cref{lem:no-atomic-active} gives the endpoint
contradiction.  Every selected alternative in the generic terminal exhaustion is
impossible, so \(\Cgen=\varnothing\).
\end{proof}

\section{Final assembly}
\label{sec:final-assembly}

The class closures above form the following complete dependency chain.  The rotational
relative-equilibrium arrow is the localized Coriolis-flux argument; the terminal
concentration-set compactness, no finite separated-family, and indecomposability
arrows are proved in this paper from local estimates:
\[
\begin{aligned}
\Cax
   &\xrightarrow{\;\Gamma\hbox{-Liouville}\;}
   \varnothing,\\
\Crot
   &\xrightarrow{\;\hbox{local flux reduction + terminal assembly}\;}
   \varnothing,\\
\Cstat
   &\xrightarrow{\;\hbox{stationary-hull reduction + terminal assembly}\;}
   \varnothing,\\
\Caff
   &\xrightarrow{\;\hbox{affine/parasitic routing}\;}
   \hbox{no retained normalized representative},\\
\Clogdiff
   &\xrightarrow{\;\hbox{log-window support profiles + hidden-scale closure}\;}
   \varnothing,\\
\Cyoung
   &\xrightarrow{\;\hbox{hidden-scale variance exclusion}\;}
   \varnothing,\\
\Chomcrit\cup\Clogper\cup\Capcrit
   &\xrightarrow{\;\hbox{bounded-origin realization + log-hull minimality}\;}
   \varnothing,\\
\Cgen
   &\xrightarrow{\;\hbox{critical-tail closure + terminal concentration analysis}\;}
   \varnothing .
\end{aligned}
\]
The positive local velocity concentration inherited by every normalized Seregin limit,
after affine normalization through the non-affine oscillation functional
\(\Phi\), is the common contradiction target in the complete ordered
decomposition.

Equivalently, the final exclusion is the composition of eight ordered steps:
\begin{enumerate}[label=\textup{(\arabic*)}]
   \item \emph{Entry step:} The setup paper supplies a positive local
   velocity concentration sequence; the endpoint weak Serrin local Type~I bounds
   \Cref{thm:AB-weak-serrin-typeI-estimate} gives the terminal local
   energy/no-escape estimate in
   \Cref{thm:local-typeI-terminal-energy-estimates}.  Thus the sequence is
   admissible in the sense of \Cref{def:admissible-local-typeI-sequence}, and
   \Cref{hyp:base-seregin-hypotheses} turns it into a normalized Seregin profile
   with retained compact velocity concentration.
   \item \emph{Stratification step:} the ordered residual decomposition
   places every remaining retained profile in the complete list of alternatives
   \[
      \Cax\cup\Crot\cup\Cstat\cup\Caff
      \cup\Clogdiff\cup\Cyoung\cup\Chomcrit\cup\Clogper\cup\Capcrit
      \cup\Cgen .
   \]
   \item \emph{Closed-strata step:} the small, stationary \(L^3\),
   tight, closed structured/decay, axisymmetric, rotational, and
   stationary-hull strata are excluded by the results cited above.
   \item \emph{Affine-routing step:} oscillatory entry normalization routes every
   affine/parasitic representative to \(\Caff\), so no retained normalized
   representative lies in that quotient row.
   \item \emph{Coherent critical-tail step:} the coherent rows
   \(\Chomcrit,\Clogper,\Capcrit\) are excluded by bounded-origin realization
   and log-hull minimality, independently of the later hidden-scale/log-diffuse
   arguments.
   \item \emph{Hidden-scale step:} minimal mesoscopic reduction yields
   \Cref{thm:no-hidden-scale-variance-realized-defects}.  This excludes the
   Young branch \(\Cyoung\) by \Cref{cor:young-branch-exclusion} and, together
   with log-window support-profile selection and the already closed coherent
   rows, excludes \(\Clogdiff\) by \Cref{cor:log-diffuse-branch-exclusion}.
   \item \emph{Residual terminal step:} the generic residual is reduced
   by the paired concentration-set decomposition, diffuse-defect compactness,
   diffuse-defect trichotomy, critical-tail compactification, realized
   critical-tail rigidity, descendant heredity, active-chain exclusion, finite
   separated-family exclusion, and terminal indecomposability.
   \item \emph{Endpoint Liouville step:} a terminally indecomposable
   retained profile is either routed to the affine/parasitic lower stratum
   or has a parasitic-free bounded mild physical pullback with a backward
   \(L^3\)-bounded sequence.  In the latter case it is zero by
   Albritton--Barker.
\end{enumerate}

\begin{proof}[Proof of \Cref{thm:refined-residual-closure}, revisited]
Assume, for contradiction, that a retained admissible residual profile exists.
Apply the setup-paper export interface to the corresponding remaining
generated state.  This produces a fixed residual object carrying the same
centered profile, retained compact velocity mass, pressure-gauge compatibility,
and terminal ancestry as the setup construction.  The refined decomposition is
then used as a proof-carrying routing datum for this same object: the selected
alternative is not allowed to change the setup export, and every later
descendant or quotient profile is linked back to that export by the realization
inheritance theorem.  By the refined ordered decomposition, after extraction
the residual object belongs to one of
\[
   \Cax,\Crot,\Cstat,\Caff,\Clogdiff,\Cyoung,
   \Chomcrit,\Clogper,\Capcrit,\Cgen .
\]
The axisymmetric class is empty by \Cref{thm:R1-exclusion}.  The rotational
class is reduced to the terminal stratification assembly theorem by \Cref{prop:R2-reduction-to-terminal} and
is empty by \Cref{cor:R2-closure}.  The stationary-hull class is reduced to
the terminal stratification assembly theorem by
\Cref{prop:R3-stationary-hull-reduction,prop:R3S-R3-terminal-reduction} and is
empty by \Cref{cor:R3-closure}.  The affine/parasitic quotient row
\(\Caff\) contains no retained normalized representative by the routing theorem
\Cref{thm:oscillatory-entry-normalization}.  The
log-diffuse and Young branches are empty by
\Cref{cor:log-diffuse-branch-exclusion,cor:young-branch-exclusion}.  The
coherent critical-tail classes are empty by
\Cref{cor:coherent-critical-tail-branch-closure}.  The generic terminal class
is empty by \Cref{thm:terminal-stratification,thm:R4-final-closure}.

Every refined residual alternative is impossible, contradicting the ordered
decomposition.  Hence
\(\calR^\#(\calS;I,J)=\varnothing\), which is the residual hypothesis of the
setup paper.  In particular, the proof constructs the setup residual-closure
hypothesis from the ordered decomposition and the individual branch closures;
it does not assume an independent contradiction for the residual export.
\end{proof}

\begin{theorem}[Paper IV residual closure]
\label{thm:paperIV-residual-closure}
For every normalized Seregin collection \(\calS\) satisfying the imported setup
admissibility hypotheses,
\[
   \calR^\#(\calS;I,J)=\varnothing .
\]
Equivalently, Paper~IV proves exactly the missing setup hypothesis
that this residual class is empty.
\end{theorem}

\begin{proof}
This is the conclusion of \Cref{thm:refined-residual-closure}, proved in the
preceding final assembly.  The statement is repeated here to separate the output
of Paper~IV from the final consequence imported from the setup paper.
\end{proof}

\begin{corollary}[Verification of the setup residual hypothesis]
\label{cor:setup-residual-hypothesis-proof}
The residual-class hypothesis
\SetupLabel{p1:hyp:no-remainder} holds.
\end{corollary}

\begin{proof}
By \Cref{thm:imported-setup-results}, the setup residual hypothesis is
the assertion that the retained residual class
\(\calR^\#(\calS;I,J)\) is empty after the lower setup classes are removed.
This is \Cref{thm:paperIV-residual-closure}.
\end{proof}

\begin{corollary}[Local pointwise Type-I exclusion from the two-paper chain]
\label{cor:two-paper-local-typeI-exclusion}
By the setup paper's conditional final assembly
\SetupLabel{p1:thm:final-assembly} and
\Cref{cor:setup-residual-hypothesis-proof}, no local pointwise Type~I
singularity exists under the stated admissibility hypotheses.
\end{corollary}

\begin{proof}
This is the conditional final assembly imported in
\Cref{thm:imported-setup-results}, with the residual hypothesis verified by
\Cref{cor:setup-residual-hypothesis-proof}.
\end{proof}

\begin{remark}[Inherited inputs and new arguments]
\label{rem:imported-vs-proved}
The paper inherits the base setup, including local concentration, the local Type I
Type II dichotomy, the Type I Seregin extraction theorem, and the initial
structural decomposition, from the setup paper summarized in
\Cref{thm:imported-setup-results}.  All remaining case-closure steps are proved in this
paper using local concentration and terminal dynamics:
\begin{itemize}
\item terminal concentration-set extraction;
\item compactness of descendant/tail families and pressure controls;
\item recurrence for chains of retained velocity concentration profiles;
\item finite separated-family exclusion;
\item oscillatory-entry normalization;
\item hidden-scale compactness reduction for critical tails;
\item bounded-origin realization and critical-tail alternative exclusion;
\item and the compactness and alternative exclusions culminating in
 \Cref{thm:R4-final-closure}.
\end{itemize}
\end{remark}

\begin{remark}[No hidden global Liouville or critical-norm input]
\label{rem:no-hidden-global-liouville}
The residual closure does not invoke any global critical-norm bound on the
original residual class, any global Coriolis--Leray Liouville theorem for
bounded centered profiles, or any global Liouville theorem for bounded ancient
solutions of the centered equation.  The only external Liouville input used
after residual closure is the Albritton--Barker endpoint Liouville theorem
\Cref{thm:AB-main}, applied in exactly two places.  First, the local
weak-Serrin Type~I estimate
\Cref{thm:AB-weak-serrin-typeI-estimate,thm:local-typeI-terminal-energy-estimates}
upgrades the local pointwise Type~I envelope to terminal velocity
no-escape; this is an entry implication, optional inside Paper~IV in the sense
of \Cref{sec:local-typeI-implication}.  Second, in
\Cref{lem:no-atomic-active}, the endpoint ancient Liouville theorem is applied
to a parasitic-free bounded mild physical pullback that carries a backward
\(L^3\)-bounded sequence supplied by indecomposable sequence-\(L^3\)
completeness \Cref{lem:atomic-sequence-L3}.  No other global Liouville
theorem and no global critical-norm estimate is used in any of the closure
steps for \(\Cax,\Crot,\Cstat,\Caff,\Clogdiff,\Cyoung,\Chomcrit,\Clogper,
\Capcrit,\Cgen\).  In particular, the rotational closure
\Cref{cor:R2-closure} proceeds through the localized Coriolis-flux
identity and terminal concentration-set decomposition, and the variance/critical-tail
closures proceed through the minimal mesoscopic-scale reduction theorem
\Cref{thm:first-bad-mesoscopic-reduction} together with bounded-origin
realization \Cref{thm:bounded-origin-realization-coherent-tails}.
\end{remark}

\appendix

\section{Setup Paper Label Dictionary}
\label{app:setup-label-dictionary}

The imported setup results displayed in \Cref{thm:imported-setup-results} are
cited by their companion-paper theorem numbers.  This table records the same
dictionary in one place so that Paper~IV can be read using the public theorem
numbering and the imported statements displayed here.
\begingroup
\scriptsize
\begin{longtable}{@{}L{0.60\textwidth}L{0.34\textwidth}@{}}
\SetupLabel{p0:thm:singular-positive} & Positive local concentration at a
Type~I singular point.\\
\SetupLabel{p0:prop:terminal-concentration} & Terminal concentration theorem.\\
\SetupLabel{p0:thm:local-weak-serrin-typeI} & Local weak-Serrin/Type~I
implication.\\
\SetupLabel{p0:cor:paper0-local-typeI-admissible} & Admissibility of
local Type~I sequences.\\
\SetupLabel{p0:lem:local-seregin-extraction} & Local Seregin extraction.\\
\SetupLabel{p1:thm:extraction} & Normalized Seregin extraction theorem.\\
\SetupLabel{p1:prop:ancient-limit} & Ancient suitable limit.\\
\SetupLabel{p1:prop:nonzero} & Nonzero retained velocity limit.\\
\SetupLabel{p1:lem:pressure} & Pressure compactness.\\
\SetupLabel{p1:lem:gauge-compatibility} & Pressure gauge compatibility.\\
\SetupLabel{p1:lem:pressure-atlas} & Pressure atlas.\\
\SetupLabel{p1:lem:mild-representation} & Mild ancient representation.\\
\SetupLabel{p2:lem:mild-limit} & Mildness under limits.\\
\SetupLabel{p2:lem:physical-mild} & Finite-shift physical mildness.\\
\SetupLabel{p1:thm:classical-classes} & Small and stationary lower classes.\\
\SetupLabel{p1:thm:known-structure-decay-liouville} & Known structured/decay
classes.\\
\SetupLabel{p2:thm:no-tight-normalized-collection} & No uniformly tight
normalized collection.\\
\SetupLabel{p2:thm:tight-liouville} & Tight Liouville theorem.\\
\SetupLabel{p2:thm:AB-endpoint} & Bounded mild ancient endpoint theorem.\\
\SetupLabel{p1:hyp:no-remainder} & Residual hypothesis proved here.\\
\SetupLabel{p1:thm:final-assembly} & Conditional final assembly in the setup
paper.\\
\end{longtable}
\endgroup
Paper~IV uses these results only through the imported theorem statements and
the companion-paper numbering displayed in this appendix.

\section{Generator calculus on compact recurrent cores}
\label{app:generator-calculus}

This appendix collects the operator-theoretic facts used implicitly in the
amenable-averaging and pressure-observable arguments of the body of the
paper.  The setting is a compact metrizable invariant minimal core
\((M,\mu,\Phi)\) of the centered translation hull, where \(\mu\) is a
\(G\)-invariant Borel probability measure (existence follows from
compactness together with the amenable-averaging protocol of
\Cref{prop:amenable-averaging-realized-diffuse-defects}; uniqueness is not
used).  The acting group is
\(G:=\mathbb R^3\rtimes\mathbb R\) (covariant translations and parabolic
time shift).

\subsection{Hilbert space and unitary representation}
Let \(H:=L^2(M,\mu;\mathbb C)\).  The Koopman representation
\(U_g f:=f\circ\Phi_g^{-1}\) (\(g\in G\)) is a strongly continuous unitary
representation of \(G\) on \(H\), by \(G\)-invariance of \(\mu\) and
the dominated convergence theorem.  In the pressure argument \(\Pi_0\) denotes
the orthogonal projection onto the subspace invariant under the three spatial
covariant translations \(X_1,X_2,X_3\).  This is the zero spatial-frequency
projection used in \Cref{lem:recurrent-core-pressure-gauges}.  It is not the
projection onto the full spacetime-invariant subspace; time translations are
used only to average total time derivatives.  The arguments below require only
this spatial projection, not ergodicity or uniqueness of \(\mu\).

For each one-parameter subgroup \(\{e^{tX_k}\}_{t\in\mathbb R}\) of \(G\)
(\(k=0,1,2,3\): time translation \(X_0\) and three covariant spatial
translations \(X_1,X_2,X_3\)), Stone's theorem gives a self-adjoint
generator \(A_k\) on a dense domain \(\mathcal D(A_k)\subset H\) with
\(U_{e^{tX_k}}=e^{itA_k}\).  The pressure inverse uses only the spatial
generators.  Define the non-negative self-adjoint operator
\(\mathcal L_{\rm sp}\) by the closed quadratic form
\[
   \mathfrak q_{\rm sp}(f)
   :=
   \sum_{k=1}^3\|A_k f\|_{L^2(M,\mu)}^2,
   \qquad
   f\in\bigcap_{k=1}^3\mathcal D(A_k).
\]
Thus \(\mathcal L_{\rm sp}\) is the Friedrichs generator associated with the formal
operator
\[
   \mathcal L_{\rm sp}=\sum_{k=1}^3 A_k^2,
\]
and \(\mathcal L_{\rm sp}\Pi_0=\Pi_0\mathcal L_{\rm sp}=0\).  All pressure
inverses below are understood as resolvent limits of
\((\varepsilon+\mathcal L_{\rm sp})^{-1}(I-\Pi_0)\).
Equivalently, with the skew-adjoint infinitesimal generators
\[
   D_k:=iA_k ,
\]
the non-negative spatial generator is
\[
   \mathcal L_{\rm sp}=-\sum_{k=1}^3D_k^2
\]
on the common cylinder core.  This is the operator denoted by
\(-\sum_{k=1}^3D_k^2\) in the recurrent-core pressure argument.

\subsection{Pressure observables and bounded continuous core}
The pressure observables used in
\Cref{def:diffuse-defect-observables}\,(iv)--(v) (compatible-gauge averaged
pressure differences \(\Pi(\cdot)-a(\cdot)\) integrated against compactly
supported test functions, plus the affine quotient observables
\([\Pi]\in D_1\)) define bounded measurable functions
\(b_{ij}^{(\varphi,a)}\colon M\to\mathbb R\) for each test function
\(\varphi\in C_c^\infty\) and each gauge selection \(a\) compatible with
\Cref{def:pressure-gauge-affine-quotient}.  These observables generate a
unital \(C^*\)-subalgebra
\(\mathcal B\subset C(M)\) which is dense in \(L^2(M,\mu)\).  The
restriction of the unitary representation \(U_g\) to \(\mathcal B\) is the
Koopman action by composition.
Affine pressure modes have already been separated into the
affine/parasitic stratum \(\Caff\) before this calculus is invoked; the
observables below are therefore taken modulo the affine quotient fixed in
\Cref{def:pressure-gauge-affine-quotient}.

\subsection{Regularized inverse and distributional pressure identity}
For each \(\varepsilon>0\) and each spatial generator \(A_k\),
\(k=1,2,3\), define the
regularized distributional pressure functional by
\[
   \left\langle q_{k,\varepsilon}^{(\varphi,a)},F\right\rangle
   :=
   \left\langle
      (I-\Pi_0)D_iD_j b_{ij}^{(\varphi,a)},
      (\varepsilon+\mathcal L_{\rm sp})^{-1}D_k^*F
   \right\rangle
\]
for cylinder test observables \(F\in\mathcal B\).  Here \(D_k=iA_k\) is the
skew-adjoint generator, \(D_k^*=-D_k\), and \(D_iD_jb_{ij}\) is interpreted in
the weak generator sense on the cylinder core.  The operator
\((\varepsilon+\mathcal L_{\rm sp})^{-1}\) is bounded on \(H\) with
norm \(\le\varepsilon^{-1}\) and commutes with \(\Pi_0\); each \(D_k\) is
closed on \(\mathcal D(A_k)\), and
\((\varepsilon+\mathcal L_{\rm sp})^{-1}D_k^*F\) lies in the form domain.  Therefore the
displayed pairing is well-defined for every \(\varepsilon>0\), \(\varphi\),
\(a\), and \(F\in\mathcal B\).  If the weak source
\((I-\Pi_0)D_iD_jb_{ij}^{(\varphi,a)}\) is represented by an \(L^2\) function,
then this functional is represented by the corresponding \(H\)-element; the
arguments in the paper need only the weak pairing.
In the formal notation used in the body this is
\[
   q_{k,\varepsilon}^{(\varphi,a)}
   =
   D_k(\varepsilon+\mathcal L_{\rm sp})^{-1}
   (I-\Pi_0)D_iD_jb_{ij}^{(\varphi,a)}
\]
understood through the preceding weak pairing.

\begin{lemma}[Distributional pressure identity in the generator calculus]
\label{lem:generator-pressure-identity}
For every test function \(\varphi\in C_c^\infty\) and every compatible
gauge \(a\), the family
\(\{q_{k,\varepsilon}^{(\varphi,a)}\}_{\varepsilon>0}\) has a distributional
limit along the cylinder core as \(\varepsilon\downarrow0\).  The limit is
independent of the chosen gauge \(a\) in the affine quotient of
\Cref{def:pressure-gauge-affine-quotient} and represents the centered
distributional pressure identity
\(\Pi(\cdot)-a(\cdot)=R_iR_j(\chi u_i u_j)+h\) of
\Cref{lem:local-pressure-decomposition} averaged against \(\mu\).
\end{lemma}

\begin{proof}
The spectral theorem applied to \(\mathcal L_{\rm sp}\) gives
\[
   D_k(\varepsilon+\mathcal L_{\rm sp})^{-1}(I-\Pi_0)
   \xrightarrow{\;\varepsilon\downarrow0\;}
   D_k\mathcal L_{\rm sp}^{-1}(I-\Pi_0)
\]
in the weak resolvent sense on the form domain, after projecting away the
invariant kernel.  For \(F\in\mathcal B\), the test
\((\varepsilon+\mathcal L_{\rm sp})^{-1}D_k^*F\) converges in the form topology to
\(\mathcal L_{\rm sp}^{-1}D_k^*F\) modulo \(\Pi_0H\).  Hence the displayed weak pairing
with \((I-\Pi_0)D_iD_jb_{ij}^{(\varphi,a)}\) converges and defines the
distributional limit \(q_k^{(\varphi,a)}\).  Gauge independence in the affine
quotient is the statement that two compatible
gauges \(a,a'\) differ by an affine pressure mode whose distributional
contribution to \(D_iD_j b_{ij}\) vanishes (because affine functions are
killed by two derivatives).  The identification with the centered
distributional pressure identity is the statement that, in a fixed local
covariant chart, the operators \(D_k\) on \(M\) restrict to the
distributional spatial derivatives along the chart and
\(\mathcal L_{\rm sp}^{-1}(I-\Pi_0)\) restricts to the inverse Laplacian on the
centered ancient solution modulo the spatially invariant pressure kernel; this
restriction is the
content of the local pressure atlas
(\Cref{thm:imported-setup-results}\,(3)).
\end{proof}

\begin{remark}[Gauge independence in the affine quotient]
\label{rem:generator-gauge-independence}
The limit \(q_k^{(\varphi,a)}\) depends on the gauge \(a\) only through
its affine quotient \([a]\in D_1\) of
\Cref{def:pressure-gauge-affine-quotient}.  If two compatible gauges
\(a,a'\) give the same affine quotient, then \(a-a'\) is the time-only
contribution allowed by the gauge atlas, and the pressure identity in
\Cref{lem:generator-pressure-identity} differs by a quantity in the
kernel of \(D_iD_j\), which is killed in the regularized formula and
therefore does not contribute to the limit.  This is the precise
statement underlying the use of pressure observables in the
amenable-averaging arguments of
\Cref{thm:diffuse-defect-compactification-construction} and
\Cref{thm:diffuse-defect-compactness}: averaging is performed on the
gauge-quotient pressure observables, and the regularized inverse
\(q_{k,\varepsilon}\) provides a regularized weak representative for each fixed
\(\varepsilon\) before the limit \(\varepsilon\downarrow0\) gives the
gauge-independent pressure identity.
\end{remark}

\begin{thebibliography}{99}

\bibitem{AlbrittonBarker2019}
D.~Albritton and T.~Barker,
\emph{On local Type I singularities of the Navier--Stokes equations and
Liouville theorems},
J. Math. Fluid Mech. 21 (2019), Article 43.

\bibitem{CKN1982}
L.~Caffarelli, R.~Kohn, and L.~Nirenberg,
\emph{Partial regularity of suitable weak solutions of the Navier--Stokes
equations},
Comm. Pure Appl. Math. 35 (1982), 771--831.

\bibitem{ESS2003}
L.~Escauriaza, G.~Seregin, and V.~Sverak,
\emph{\(L_{3,\infty}\)-solutions of Navier--Stokes equations and backward
uniqueness},
Russian Math. Surveys 58 (2003), 211--250.

\bibitem{KNSS2009}
G.~Koch, N.~Nadirashvili, G.~Seregin, and V.~Sverak,
\emph{Liouville theorems for the Navier--Stokes equations and applications},
Acta Math. 203 (2009), 83--105.

\bibitem{LRZ2022}
Z.~Lei, J.~Ren, and X.~Zhang,
\emph{A Liouville theorem for the axially symmetric Navier--Stokes equations},
J. Differential Equations 313 (2022), 213--238.

\bibitem{LZZ2017}
Z.~Lei, Q.~Zhang, and N.~Zhao,
\emph{Improved Liouville theorems for axially symmetric Navier--Stokes
equations},
J. Math. Fluid Mech. 19 (2017), 273--285.

\bibitem{NRS1996}
J.~Necas, M.~Ruzicka, and V.~Sverak,
\emph{On Leray's self-similar solutions of the Navier--Stokes equations},
Acta Math. 176 (1996), 283--294.

\bibitem{Seregin2012}
G.~Seregin,
\emph{A certain necessary condition of potential blow up for Navier--Stokes
equations},
Comm. Math. Phys. 312 (2012), 833--845.

\end{thebibliography}

\end{document}
